\documentclass[11pt,reqno]{amsart}
% REVIEW SOURCE: clean, with no revision colors. All other manuscript prose is unchanged.
\usepackage[T1]{fontenc}
\usepackage[utf8]{inputenc}
\usepackage{lmodern}
\usepackage[margin=1in]{geometry}
\usepackage{amsmath,amssymb,amsthm,mathtools,mathrsfs,microtype}
\usepackage{xcolor,aliascnt}
\definecolor{RED}{rgb}{1,0,0}
\usepackage[shortlabels]{enumitem}
\usepackage{tikz}
\usetikzlibrary{arrows.meta,positioning}
\usepackage[colorlinks=true,linkcolor=blue,citecolor=magenta,urlcolor=magenta]{hyperref}
\usepackage[nameinlink,capitalise,noabbrev]{cleveref}
\allowdisplaybreaks
\numberwithin{equation}{section}
\setlist{listparindent=\parindent,parsep=0pt}
\setlist[enumerate]{topsep=4pt,itemsep=3pt}
\theoremstyle{plain}
\newtheorem{thm}{Theorem}[section]
\newaliascnt{prop}{thm}\newtheorem{prop}[prop]{Proposition}\aliascntresetthe{prop}
\newaliascnt{lemma}{thm}\newtheorem{lemma}[lemma]{Lemma}\aliascntresetthe{lemma}
\newaliascnt{cor}{thm}\newtheorem{cor}[cor]{Corollary}\aliascntresetthe{cor}
\theoremstyle{remark}
\newaliascnt{remark}{thm}\newtheorem{remark}[remark]{Remark}\aliascntresetthe{remark}
\crefname{thm}{theorem}{theorems}\Crefname{thm}{Theorem}{Theorems}
\crefname{prop}{proposition}{propositions}\Crefname{prop}{Proposition}{Propositions}
\crefname{lemma}{lemma}{lemmas}\Crefname{lemma}{Lemma}{Lemmas}
\crefname{cor}{corollary}{corollaries}\Crefname{cor}{Corollary}{Corollaries}
\crefname{remark}{remark}{remarks}\Crefname{remark}{Remark}{Remarks}
\DeclarePairedDelimiter{\norm}{\lVert}{\rVert}
\DeclarePairedDelimiter{\abs}{\lvert}{\rvert}
\renewcommand{\d}{\mathsf{d}}\newcommand{\s}{\mathsf{s}}
\newcommand{\T}{\mathbb T}\newcommand{\R}{\mathbb R}\newcommand{\Z}{\mathbb Z}
\newcommand{\E}{\mathbb E}\newcommand{\Q}{\mathbb Q}\newcommand{\Pb}{\mathbb P}
\newcommand{\K}{\mathsf K}\newcommand{\g}{\mathsf g}\newcommand{\Fr}{\mathsf F}
\newcommand{\Hr}{\mathsf H}\newcommand{\dd}{\,\mathrm d}
\DeclareMathOperator{\diver}{div}
\DeclareMathOperator{\Ent}{Ent}\DeclareMathOperator{\TV}{TV}\DeclareMathOperator{\Var}{Var}
\newcommand{\parheading}[1]{\par\medskip\noindent\textbf{#1}\ }
\title[Sharp local chaos for diffusive log/Riesz dynamics]{Sharp local propagation of chaos for diffusive log/Riesz gas dynamics in the Hilbert--Schmidt regime}
\author[M. G. Delgadino]{Matias G. Delgadino}
\address{Matias Delgadino, Department of Mathematics, The University of Texas at Austin, PMA 8.100, 2515 Speedway, Austin, TX 78712, USA}
\email{matias.delgadino@math.utexas.edu}
\urladdr{https://math.utexas.edu/directory/matias-delgadino}
\author[R. Gvalani]{Rishabh Gvalani}
\address{Rishabh Gvalani, School of Mathematics, University of Edinburgh, James Clerk Maxwell Building, Peter Guthrie Tait Road, Edinburgh EH9 3FD, United Kingdom}
\email{rishabh.gvalani@ed.ac.uk}
\urladdr{https://webhomes.maths.ed.ac.uk/~rgvalani/}
\author[M. Rosenzweig]{Matthew Rosenzweig}
\address{Matthew Rosenzweig, Department of Mathematical Sciences, Carnegie Mellon University, 7127 Wean Hall, 5000 Forbes Avenue, Pittsburgh, PA 15213, USA}
\email{mrosenz2@andrew.cmu.edu}

\date{September 8, 2026}
\subjclass[2020]{60K35, 35Q84, 82C31, 60H10}
\keywords{propagation of chaos, logarithmic and Riesz interactions, modulated Gibbs measures, Hilbert--Schmidt kernels, time-dependent integrability}
\hypersetup{pdftitle={Sharp local propagation of chaos for diffusive log/Riesz gas dynamics in the Hilbert--Schmidt regime},pdfauthor={Matias G. Delgadino, Rishabh Gvalani, Matthew Rosenzweig}}

\newcommand{\Xspace}{\mathcal X}
\newcommand{\Aop}{\mathcal A}
\newcommand{\Id}{\mathrm{Id}}

% Public working-draft presentation (2026-10-28).
\usepackage{draftwatermark}
\SetWatermarkText{DRAFT}
\SetWatermarkAngle{45}
\SetWatermarkScale{0.8}
\SetWatermarkColor[gray]{0.92}

\begin{document}
\raggedbottom
\begin{abstract}

{We prove sharp bounds on the correlations of diffusive particle systems with repulsive logarithmic and Riesz interactions in $\R^\d$, $\d\ge2$, in the Hilbert--Schmidt regime $0\le\s<\d/2$.} Starting from a modulated Gibbs measure at any fixed inverse temperature greater than half that of the dynamics, the relative entropy and chi-square divergence of the $k$-particle marginal from the mean-field product measure are bounded by $Ck^2/N^2$, and its total variation distance by $Ck/N$. We obtain finite-time estimates under one-body regularity and Poincar\'e bounds and a sufficient bound on the time integral of $\norm{\nabla u_t}_{L^\infty}$, where $u_t$ is the mean-field continuity velocity. For small logarithmic perturbations of a confined thermal equilibrium measure, the estimates are uniform in time at every fixed positive inverse temperature. The confinement may be anisotropic and need not be globally convex. The rates are optimal, as shown by the stationary Gibbs subclass. Self-similar coordinates give corresponding unconfined results for dilated thermal equilibrium measures and sufficiently diffuse Gaussian perturbations. The proof combines partition-function estimates with a decreasing-exponent estimate relative to a moving Gibbs measure and weighted estimates for the first marginal. We also construct a four-dimensional whole-space example with a stationary mean-field measure and independent initial data. The full particle measure has bounded relative entropy but a chi-square divergence that tends to infinity with the particle number at a fixed positive time, despite zero integrability expenditure. Finally, for the four-dimensional periodic logarithmic gas, we construct a smooth trajectory along which the full relative entropy remains bounded but the terminal full chi-square divergence tends to infinity on a fixed interval of time. Thus, even when the particle system starts from the modulated Gibbs measure associated with the initial mean-field measure and the mean-field trajectory is smooth, an $N$-uniform full-particle chi-square bound need not hold on an arbitrary fixed time interval.
\end{abstract}

\maketitle
\begingroup

\section{Introduction and main results}\label{sec:intro}

The purpose of this paper is to quantify the correlations generated by the overdamped dynamics of a logarithmic or Riesz gas. The mean-field equation describes the limiting one-particle density. We ask how accurately its product law approximates a labeled group of particles, with an estimate that keeps track of both the size of the group and the elapsed time. A sharp dependence on the group size determines how many particles remain approximately independent. Its persistence in time depends on the evolution of the mean-field measure and on the correlations present initially. We establish quantitative estimates under conditions on both, and construct two examples showing why those conditions matter when the estimate includes the full particle law.


The particle dynamics combine singular pairwise repulsion, an external force, and independent Brownian noise. Fix $\d\ge2$, $\beta>0$, and $0\le\s<\d/2$, and set
\begin{equation}\label{eq:kernel}
 \g(x)=\begin{cases}-\log|x|,&\s=0,\\ |x|^{-\s}/\s,&\s>0.\end{cases}
\end{equation}
Logarithmic and Riesz interactions provide basic models of long-range particle systems. The Coulomb cases describe electrostatic repulsion, while the logarithmic interaction also appears in random matrix theory and two-dimensional vortex dynamics. Varying the Riesz exponent permits one to study how the strength of the singularity affects correlations and the mean-field limit. We refer to \cite{SerfatyBook,RSdiffusive,RSselfsimilar} for further background and references.



For $N\ge1$, write $[N]:=\{1,\ldots,N\}$ and $X_N:=(x_1,\ldots,x_N)$ for a configuration. The mean-field Hamiltonian is
\begin{equation}\label{eq:intro-hamiltonian}
 \mathcal H_N(X_N)=\sum_{i=1}^N V(x_i)+\frac1N\sum_{i<j}\g(x_i-x_j).
\end{equation}
We consider its overdamped Langevin dynamics,
\begin{equation}\label{eq:sde}
 \dd x_i=-\left(\nabla V(x_i)+\frac1N\sum_{j\ne i}\nabla\g(x_i-x_j)\right)\dd t
            +\sqrt{2/\beta}\,\dd B_i,\qquad 1\le i\le N,
\end{equation}
where $(B_i)_{i\in[N]}$ are independent Brownian motions. We write $X_N(t):=(x_1(t),\ldots,x_N(t))$ for the full random configuration. The inverse temperature $\beta$ is independent of $N$. In this scaling, the mean-field interaction and Brownian diffusion both contribute to the limiting equation; this is the diffusive temperature regime of \cite{DGRpartition}. The class includes the planar Coulomb gas and the three-dimensional Coulomb gas. More generally, the restriction $\s<\d/2$ is the local square-integrability condition on the interaction potential. Square integrability of the force is not assumed.



The corresponding mean-field equation is
\begin{equation}\label{eq:mf}
 \partial_t\mu=\beta^{-1}\Delta\mu+
    \diver\bigl(\mu(\nabla V+\nabla\g*\mu)\bigr).
\end{equation}
It replaces the empirical interaction force in \eqref{eq:sde} by the force generated by the one-particle density. Qualitative propagation of chaos asks whether each fixed labeled marginal converges to the corresponding product of mean-field measures as $N\to\infty$. Quantitative estimates must also describe the residual dependence at finite $N$ and the time intervals on which it remains small; see, for example, \cite{Lacker}. Our main results concern confined whole-space dynamics, and a self-similar change of variables gives explicit unconfined classes. The sharp local estimates also carry over, mutatis mutandis, to the periodic setting, where the absence of spatial tails simplifies the proof; see \cref{rem:periodic}. The hypotheses on $V$ and the solution of \eqref{eq:mf} are stated with each result.



\parheading{Initial preparation.}
For the sharp local chaos estimates, we choose an initial measure that incorporates the singular pair repulsion around a prescribed probability density. {This permits nonstationary mean-field trajectories while preserving quantitative control of the initial correlations.} The measure is a modulated Gibbs measure; at a matching thermal equilibrium it reduces to the canonical Gibbs ensemble. See the work of Serfaty and the third author \cite{RSlsi} for its role in the modulated-free-energy method, and \cite{SerfatyBook} for a broader account of modulated Gibbs measures and the modulated-energy framework. Its normalization and equilibrium correlations are studied in \cite{DGRpartition,DGRchaos}.



For a probability density $\mu$, denote its interaction potential and energy, the modulated energy, and the modulated Gibbs measure, respectively, by
\begin{equation}\label{eq:background}
 h_\mu=\g*\mu,\qquad I_\mu=\frac12\int h_\mu\dd\mu,
\end{equation}
\begin{equation}\label{eq:modulated}
 N\Fr_N(X_N,\mu)=\frac1N\sum_{i<j}\g(x_i-x_j)-\sum_i h_\mu(x_i)+NI_\mu,
\end{equation}
\begin{equation}\label{eq:gibbs}
 \dd\Q_{N,b}(\mu)=G_{N,b}(\mu)\dd\mu^{\otimes N},\qquad
 G_{N,b}(\mu)=\K_{N,b}(\mu)^{-1}e^{-bN\Fr_N(X_N,\mu)},\qquad b\ge0.
\end{equation}
The normalizing constant $\K_{N,b}(\mu)$ is the modulated partition function. The subtracted one-body and constant terms in \eqref{eq:modulated} remove the mean-field contribution of the prescribed background. The normalizations are those of \cite{DGRpartition,DGRchaos}.



For an initial inverse temperature $\beta_0\ge0$, consider the preparation
\begin{equation}\label{eq:preparation}
 P_{N,0}=\Q_{N,\beta_0}(\mu_0).
\end{equation}
{The parameter $\beta_0$ controls only the initial correlations; both the particle and mean-field dynamics are at inverse temperature $\beta$.} The choice $\beta_0=\beta$ is allowed but is not required. The first marginal of \eqref{eq:preparation} need not equal $\mu_0$, and \cref{prop:initial-mark} supplies the required initial one-particle estimate. The actual law at positive time need not be $\Q_{N,\beta}(\mu_t)$: the evolving mean-field measure determines a reference Gibbs law, not the solution of the particle forward equation. The proof estimates the difference between these laws and then its effect on the marginals. \Cref{prop:iid-counterexample} explains why the initial correlations matter when the conclusion includes $k=N$; that example starts from independent and identically distributed (iid) samples.


The initial configuration is independent of the Brownian motions. It is almost surely pairwise distinct because \eqref{eq:preparation} is absolutely continuous with respect to a product density. Under the bounded-Hessian assumption on $V$ used below, \cref{prop:particle-wellposedness} gives a pathwise unique global strong solution of \eqref{eq:sde}, with no collisions on finite time intervals. We write $P_{N,t}$ for its law and $P_{N,t}^{(k)}$ for its labeled $k$-particle marginal.

\subsection{Local, mesoscopic, and full-particle information}\label{sec:intro-scales}

We measure the difference between $P_{N,t}^{(k)}$ and $\mu_t^{\otimes k}$ in relative entropy, chi-square divergence, and total variation. For probability measures $P\ll Q$, these are defined by
\begin{equation}\label{eq:intro-entropy-definition}
 \Hr(P\mid Q):=\int\log\!\left(\frac{\dd P}{\dd Q}\right)\dd P,
\end{equation}
\begin{equation}\label{eq:intro-chi-definition}
 \chi^2(P\mid Q):=\int\left|\frac{\dd P}{\dd Q}-1\right|^2\dd Q,
\end{equation}
\begin{equation}\label{eq:intro-TV-definition}
 \TV(P,Q):=\sup_{A}|P(A)-Q(A)|
 =\frac12\int\left|\frac{\dd P}{\dd Q}-1\right|\dd Q,
\end{equation}
where the supremum is over measurable sets. When $P\not\ll Q$, the first two quantities are $+\infty$, while the supremum still defines total variation. Relative entropy is unnormalized. A bound on the full relative entropy does not by itself give the sharp marginal rate. Indeed, entropy subadditivity and exchangeability give
\begin{equation}\label{eq:intro-entropy-subadditivity}
 \Hr(P_{N,t}^{(k)}\mid\mu_t^{\otimes k})
 \le\frac{2k}{N}\Hr(P_{N,t}\mid\mu_t^{\otimes N}),\qquad 1\le k\le N.
\end{equation}
Thus, an order-one full entropy bound yields only an order-$k/N$ estimate through \eqref{eq:intro-entropy-subadditivity}. The estimates below have the sharper order $k^2/N^2$ in relative entropy and chi-square divergence, and order $k/N$ in total variation.



For fixed $k$, these are local size-of-chaos estimates. They also imply convergence for every mesoscopic sequence $k=k_N\to\infty$ with $k_N/N\to0$. At $k=N$, the same inequalities have a different meaning: they bound the full density in $L^2(\mu_t^{\otimes N})$, without asserting convergence of the full law to a product. Our sharp local estimates hold for all $1\le k\le N$ under their stated hypotheses.


The full-particle conclusion makes the role of the initial law and the time restriction particularly visible. \Cref{prop:iid-counterexample} shows that independent initialization can lead to an unbounded full chi-square at a fixed later time even when the mean-field measure is stationary. {\Cref{thm:periodic-counterexample} preserves the Gibbs initialization but shows that smoothness of a nonstationary trajectory does not suffice on an arbitrary finite interval.} Both examples have bounded full relative entropy. They concern the stronger full-particle estimate, not a failure of fixed-marginal propagation of chaos.

\subsection{Finite-time control along a prescribed trajectory}

{We first give a finite-time criterion in terms of the prescribed mean-field solution and the static partition and transport-commutator estimates proved in \cite{DGRpartition}.} Set
\begin{equation}\label{eq:finite-drift}
 B_t=-\nabla V-\nabla\g*\mu_t
\end{equation}
and
\begin{equation}\label{eq:velocity}
 u_t=-\nabla V-\nabla\g*\mu_t-\beta^{-1}\nabla\log\mu_t.
\end{equation}
Thus, $B_t$ is the drift of the nonlinear diffusion, and \eqref{eq:mf} is also the continuity equation with velocity $u_t$.



\begin{thm}[Finite-time sharp local estimate under an integrated condition]\label{thm:finite}
\begingroup
Fix $\beta>0$ and $\beta_0>\beta/2$. Let $V\in C_{\mathrm{loc}}^3(\R^\d)$ satisfy
\begin{equation}\label{eq:intro-V-regularity}
 \sup_{x\in\R^\d}\frac{|\nabla V(x)|}{1+|x|}+\norm{\nabla^2V}_{L^\infty(\R^\d)}+\norm{\nabla^3V}_{L^\infty(\R^\d)}<\infty.
\end{equation}
Let $\mu_t$ be a smooth positive solution of \eqref{eq:mf} on $[0,T]$.
\endgroup Assume the following bounds, with constants independent of $t\in[0,T]$.
\begin{enumerate}[label=\textup{(\roman*)},leftmargin=2.4em,itemsep=7pt]
\item \emph{Density and logarithmic Hessian.} There are $M,a>0$ and $L<\infty$ such that
\begin{equation}\label{eq:intro-density-bound}
 0<\mu_t(x)\le M e^{-a|x|^2},\qquad x\in\R^\d,
\end{equation}
and, writing $U_t=-\log\mu_t$,
\begin{equation}\label{eq:intro-hessian-bound}
 -L\Id\le\nabla^2U_t(x)\le L\Id,\qquad x\in\R^\d.
\end{equation}
\item \emph{One-body Poincar\'e inequality.} There is $C_{\mathrm P}<\infty$ such that
\begin{equation}\label{eq:intro-poincare-bound}
 \int\left|\varphi-\int\varphi\dd\mu_t\right|^2\dd\mu_t
 \le C_{\mathrm P}\int|\nabla\varphi|^2\dd\mu_t,
 \qquad \varphi\in C_c^\infty(\R^\d).
\end{equation}
\item \emph{Drift derivatives.} There is $C_B<\infty$ such that
\begin{equation}\label{eq:intro-drift-bound}
 \norm{\nabla B_t}_{L^\infty}+\norm{\nabla^2B_t}_{L^\infty}\le C_B.
\end{equation}
\item \emph{Static estimates.} The estimates of \cref{prop:static-inputs} hold uniformly along $\mu_t$. Let $\kappa>0$ be a common commutator coefficient in \eqref{eq:static-commutator}.
\end{enumerate}
If the \emph{integrability expenditure along the mean-field trajectory} satisfies
\begin{equation}\label{eq:expenditure}
 A_\mu(T):=\frac\beta\kappa\int_0^T\norm{\nabla u_t}_{L^\infty}\dd t
 <\log\left(2\min\left\{1,\frac{\beta_0}{\beta}\right\}\right),
\end{equation}
then the solution of \eqref{eq:sde} with preparation \eqref{eq:preparation} satisfies, for every $N\ge2$ and $1\le k\le N$,
\begin{align}
 \sup_{0\le t\le T}\left\{\chi^2(P_{N,t}^{(k)}\mid\mu_t^{\otimes k})+
              \Hr(P_{N,t}^{(k)}\mid\mu_t^{\otimes k})\right\}
       &\le C_T\frac{k^2}{N^2},\label{eq:finite-rates}\\
 \sup_{0\le t\le T}\TV(P_{N,t}^{(k)},\mu_t^{\otimes k})
       &\le C_T\frac{k}{N}.\label{eq:intro-finite-TV}
\end{align}
The constant depends only on the preceding one-body bounds, the static constants, $\beta_0$, the model parameters, and the positive gap in \eqref{eq:expenditure}.
\end{thm}



Assumption \textup{(iv)} follows from the Gaussian upper bound in \textup{(i)} by \cref{prop:static-inputs}; we state it separately to identify the commutator coefficient entering \eqref{eq:expenditure}. We have not sought optimal regularity or decay assumptions on the mean-field solution. Extensions under weaker one-body hypotheses are left to the interested reader.



The right-hand side of \eqref{eq:expenditure} is positive precisely because $\beta_0>\beta/2$. When $\beta_0\ge\beta$, it reduces to $\log2$. This sufficient condition is not claimed sharp. Its origin is the loss of integrability in the decreasing-exponent estimate of \cref{sec:decreasing}, together with the comparison between the initial Gibbs measure and the moving reference. The theorem does not require that $\mu_t$ remain close to an equilibrium. We next give a class of initial measures for which the required bounds hold on $[0,\infty)$.


\subsection{Uniform-in-time control near a confined thermal equilibrium measure}
We say that the confinement force is \emph{asymptotically linear with a positive-definite limiting matrix} if there exist $A=A^{\mathsf T}>0$ and $b\in\R^\d$ such that
\begin{equation}\label{eq:confinement-class}
 \widetilde V(x):=V(x)-\tfrac12x\cdot Ax-b\cdot x
\end{equation}
has a smooth gradient, all derivatives of $\nabla\widetilde V$ are bounded, and
\begin{equation}\label{eq:confinement-decay}
 \lim_{R\to\infty}\sup_{|x|\ge R}|\nabla^{j+1}\widetilde V(x)|=0
 \qquad\text{for every }j\ge0.
\end{equation}
Writing $\lambda_A$ for the smallest eigenvalue of $A$, these assumptions imply
\begin{align}
 V(x)&=\frac12x\cdot Ax+b\cdot x+o(|x|),\label{eq:confinement-potential-asymptotic}\\
 \nabla^2V(x)&\longrightarrow A,\label{eq:confinement-hessian-asymptotic}\\
 \nabla^2V(x)&\ge\frac{\lambda_A}{2}\Id\qquad\text{for }|x|\ge R_V
 \label{eq:exterior-curvature}
\end{align}
for some $R_V<\infty$. The equations are invariant under a common spatial translation, and adding a constant to $V$ changes neither equation. {We may therefore take $b=0$ and $V(0)=0$ in the proofs and continue to write $V$ for the resulting confinement.} The class includes anisotropic positive quadratic confinements, arbitrary smooth compactly supported perturbations of them, and unbounded sublinear perturbations such as a multiple of $\log(1+|x|^2)$. No global convexity assumption is imposed. The finite-time theorem uses less: the asymptotically linear-force condition is used only to verify uniform-in-time logarithmic stability by comparison with the anisotropic Ornstein--Uhlenbeck semigroup.

For a confinement satisfying \eqref{eq:confinement-class}--\eqref{eq:confinement-decay}, the thermal probability density is the unique solution
\begin{equation}\label{eq:thermal}
 \pi(x)=Z_\pi^{-1}\exp\left[-\beta\left(V(x)+(\g*\pi)(x)\right)\right].
\end{equation}
Its existence, smoothness, Gaussian tail bounds, and one-body Poincar\'e inequality are recorded with proofs in \cref{lem:thermal}. At the matching inverse temperature, $\Q_{N,\beta}(\pi)$ is the usual canonical Gibbs ensemble.

\parheading{The logarithmic perturbation class.}
A logarithmic perturbation is normalized as
\begin{equation}\label{eq:log-tilt}
 \mu_0=\frac{\pi e^{z_0}}{\int\pi e^{z_0}}.
\end{equation}
The simplest class covered below is
\begin{equation}\label{eq:explicit-class}
 z_0(x)=\ell\cdot x+\varphi(x),\qquad
 \ell\in\R^\d,\quad\varphi\in C_c^\infty(\R^\d).
\end{equation}
In particular, we allow unbounded logarithmic relative densities and do not assume two-sided boundedness of $\mu_0/\pi$.

We also use the completion of this class. For an integer $m\ge3$, let $\Xspace^m$ be the completion of smooth affine functions plus $C_c^\infty$, modulo constants, in the norm
\begin{equation}\label{eq:X-norm}
 \norm z_{\Xspace^m}=\sum_{j=1}^m\norm{\nabla^jz}_\infty.
\end{equation}
A representative is fixed, when needed, by $\int z\dd\pi=0$. The completion and its use in the one-body stability argument are discussed in \cref{sec:backgrounds}.

\begin{thm}[Uniform-in-time sharp local chaos near a confined thermal equilibrium measure]\label{thm:uniform}
Let $V$ satisfy \eqref{eq:confinement-class}--\eqref{eq:confinement-decay}, and fix $\d,\s,\beta$, an initial inverse temperature $\beta_0>\beta/2$, and an integer $m\ge3$. There are $\delta>0$ and $C<\infty$ with the following property.


Let $z_0$ be of the form \eqref{eq:explicit-class}, and assume
\begin{equation}\label{eq:intro-small-perturbation}
 \norm{z_0}_{\Xspace^m}\le\delta.
\end{equation}
Define $\mu_0$ by \eqref{eq:log-tilt}. Then \eqref{eq:mf} has a global smooth positive solution, and the solution of \eqref{eq:sde} with preparation \eqref{eq:preparation} satisfies, for every $N\ge2$ and $1\le k\le N$,
\begin{align}
 \sup_{t\ge0}\left\{\chi^2(P_{N,t}^{(k)}\mid\mu_t^{\otimes k})+
              \Hr(P_{N,t}^{(k)}\mid\mu_t^{\otimes k})\right\}
       &\le C\frac{k^2}{N^2},\label{eq:main-rates}\\
 \sup_{t\ge0}\TV(P_{N,t}^{(k)},\mu_t^{\otimes k})
       &\le C\frac{k}{N}.\label{eq:intro-uniform-TV}
\end{align}
The constants depend only on $\d,\s,\beta,\beta_0,m$, the matrix $A$, and finitely many uniform bounds and decay moduli for $\nabla\widetilde V$ through an order determined by $m$. They are independent of $N,k,t$ and of higher derivatives of $z_0$. The same conclusion holds for smooth data in the completion for which the higher norms needed for classical regularity are finite.
\end{thm}



The dynamical inverse temperature $\beta$ is arbitrary and fixed, and the initial inverse temperature may be any fixed $\beta_0>\beta/2$; the allowed size $\delta$ depends on both. Positivity of the centered interaction operator is used to prove stability of the mean-field solution and of the homogeneous equation for the first-marginal error. The initial first-marginal estimate instead follows from a partition bound at a larger inverse temperature; see \cref{prop:initial-mark}.



\begin{cor}[Short-time sharp local chaos]\label{cor:local}
Fix $\beta>0$ and $\beta_0>\beta/2$, and let the confinement force satisfy \eqref{eq:confinement-class}--\eqref{eq:confinement-decay}. For every $z_0$ of the form \eqref{eq:explicit-class}, without a smallness assumption, there is $T_*>0$ such that the conclusions of \cref{thm:finite} hold on $[0,T_*]$.
\end{cor}



\begin{remark}[Sharpness through equilibrium]\label{rem:intro-sharpness}
The admissible class contains the matched choice $\beta_0=\beta$. For the thermal equilibrium measure \eqref{eq:thermal}, the corresponding preparation is
\begin{equation}\label{eq:stationary-preparation}
 \dd\Q_{N,\beta}(\pi)\ \propto\
 \exp\!\left[-\beta\left(\sum_{i=1}^NV(x_i)+\frac1N\sum_{i<j}\g(x_i-x_j)\right)\right]\dd X_N.
\end{equation}
Both this canonical Gibbs measure and the mean-field equilibrium measure are stationary. The two-sided equilibrium estimates of \cite[Theorem~2.1]{DGRchaos} therefore establish optimality of the present rates. In \cref{prop:stationary-sharpness}, we include a proof of the required stationary lower bounds using the estimates already available in this paper and the static fluctuation formula of \cite{DGRpartition}, in an effort to make the sharpness argument self-contained. Together with \cref{thm:uniform}, these bounds show that, for $2\le k\le N$, the chi-square divergence and relative entropy have order $k^2/N^2$, and total variation has order $k/N$. Consequently, the rates in \eqref{eq:finite-rates} and \eqref{eq:main-rates} cannot be improved uniformly over the stated classes. The lower bounds start at $k=2$; the upper estimates also hold for $k=1$.



\end{remark}



\begin{remark}[Unconfined consequences]\label{rem:intro-unconfined}
\Cref{thm:unconfined} gives two explicit classes through the self-similar transformation \eqref{eq:selfsimilar}. For the logarithmic interaction, the initial measures are dilates of small logarithmic perturbations of the thermal equilibrium measure for $V(x)=|x|^2/4$. For positive Riesz exponents, they are sufficiently diffuse dilates of small Gaussian perturbations. The lower bound on the dilation scale is part of the Riesz assertion. The full statements and their proofs are given in \cref{sec:selfsimilar}; \cref{rem:periodic} records the corresponding periodic result.



\end{remark}


\subsection{The initial condition and the full-particle chi-square}\label{sec:intro-iid-counterexample}

The first counterexample concerns the initial law of the $N$-particle system. Independent samples have zero chi-square divergence from the product reference at time zero. Nevertheless, for a suitable confined system, the interaction generates fluctuations that make the full chi-square unbounded in $N$ at a fixed later time of order one with respect to $N$. The mean-field measure in this example is stationary, so its continuity velocity and integrability expenditure vanish identically.


\begin{prop}[Stationary counterexample from independent initial data]\label{prop:iid-counterexample}
Let $\d=4$ and $0\le\s<1$. There are families of inverse temperatures $\beta_\delta>0$, confinements $V_\delta\in C^\infty_{\mathrm{loc}}(\R^4)$, smooth positive probability densities $\mu_\delta$, and times $T_\delta>0$, indexed by $0<\delta<1/2$, with the following properties.
\begin{enumerate}[label=\textup{(\roman*)},leftmargin=2.4em,itemsep=7pt]
\item The confinement satisfies \eqref{eq:confinement-class}--\eqref{eq:confinement-decay} with $A=(4\beta_\delta)^{-1}\Id$ and $b=0$. More explicitly, set
\begin{equation}\label{iid:eq:force-class}
 \widetilde V_\delta(x):=V_\delta(x)-\frac{|x|^2}{8\beta_\delta},\qquad
 \lim_{R\to\infty}\sup_{|x|\ge R}|\nabla^{j+1}\widetilde V_\delta(x)|=0
 \quad(j\ge0),
\end{equation}
and every derivative of $\nabla\widetilde V_\delta$ is bounded. The density $\mu_\delta$ is the thermal equilibrium measure for $V_\delta$ and $\beta_\delta$. Its mean-field trajectory satisfies
\begin{equation}\label{iid:eq:stationary-trajectory}
 \mu_t=\mu_\delta,\qquad u_t=0\quad(t\ge0),\qquad
 A_{\mu_\delta}(T)=0\quad(T<\infty).
\end{equation}
\item For every $N\ge2$, consider \eqref{eq:sde} with confinement $V_\delta$, inverse temperature $\beta_\delta$, and initial law
\begin{equation}\label{iid:eq:initial-law}
 P_{N,0}=\mu_\delta^{\otimes N},
\end{equation}
independent of the Brownian motions. The system has a pathwise unique global noncolliding strong solution. Its law satisfies
\begin{equation}\label{iid:eq:entropy-main}
 \sup_{N\ge2}\sup_{t\ge0}\Hr(P_{N,t}\mid\mu_\delta^{\otimes N})<\infty.
\end{equation}
\item For every fixed $N\ge2$ and $t<\infty$, the full chi-square is finite, and it is zero at $t=0$. For every sufficiently small fixed $\delta$,
\begin{equation}\label{iid:eq:chi-main}
 \lim_{N\to\infty}\chi^2(P_{N,T_\delta}\mid\mu_\delta^{\otimes N})=+\infty.
\end{equation}
\end{enumerate}
All parameters, the confinement, the initial density, and the terminal time are fixed before $N\to\infty$. The bound in \eqref{iid:eq:entropy-main} may depend on $\s$ and $\delta$.
\end{prop}


The construction and proof are given in Appendix~\ref{app:iid-counterexample}. We have chosen a dimension and range of Riesz exponents that simplify the construction, without seeking the greatest generality. Obtaining such a counterexample throughout the dimension and exponent range of \cref{thm:finite} is left to the interested reader. Stationarity here concerns $\mu_\delta$, not the particle law initialized by $\mu_\delta^{\otimes N}$. The stationary trajectory satisfies the intrinsic one-body and static hypotheses of \cref{thm:finite}, and it is the zero logarithmic perturbation of its own thermal equilibrium measure. {Thus, replacing the preparation required in \cref{thm:finite,thm:uniform} by iid initialization while keeping their full-particle conclusion is impossible under these trajectory conditions alone.} This conclusion does not result from setting $\beta_0=0$ in \eqref{eq:expenditure}, and it does not determine the optimal range of initial Gibbs temperatures. The confinement is asymptotically isotropic but not exactly quadratic.


The divergence in \eqref{iid:eq:chi-main} does not preclude local or mesoscopic chaos. In fact, \cref{iid:rem:local} gives a marginal entropy bound of order $k/N$, uniform for all time, and hence convergence whenever $k=o(N)$. The sharp local iid rates remain undecided by this example. The next result concerns the time restriction: even the matched Gibbs initial law does not give unrestricted full-particle chi-square control along every smooth mean-field trajectory.

\subsection{Failure of unrestricted full-particle chi-square control}\label{sec:intro-counterexample}

Within this subsection, we recycle $\g$ for the mean-zero spectral logarithm on the unit torus $\T^4$, normalized by
\begin{equation}\label{eq:periodic-ce-kernel}
 \widehat\g(n)=\frac{1}{2\pi^2|n|^4},\qquad n\in\Z^4\setminus\{0\}.
\end{equation}
It has local singularity $-\log|x|$. Consider the periodic versions of \eqref{eq:sde} and \eqref{eq:mf}, with $V=0$ and preparation $P_{N,0}=\Q_{N,\beta}(\mu_0)$ at the same inverse temperature as the dynamics.



\begin{thm}[Failure of unrestricted full-particle chi-square control]\label{thm:periodic-counterexample}
For every sufficiently large fixed $\beta$, there is $\mu_0\in C^\infty(\T^4)$ with $\inf\mu_0>0$ such that the corresponding mean-field trajectory is smooth and strictly positive on $[0,5]$. The laws $P_{N,t}$ of the corresponding particle system satisfy
\begin{equation}\label{eq:periodic-ce-failure}
 \lim_{N\to\infty}\chi^2(P_{N,5}\mid\mu_5^{\otimes N})=+\infty,
\end{equation}
while each finite-$N$ divergence is finite, the initial full chi-square divergences are uniformly bounded in $N$, and
\begin{equation}\label{eq:periodic-ce-entropy}
 \sup_N\sup_{0\le t\le5}\Hr(P_{N,t}\mid\mu_t^{\otimes N})<\infty.
\end{equation}
Moreover, there are a real-valued $\psi\in C^\infty(\T^4)$ and $v>2\Var_{\mu_5}(\psi)>0$ such that, under $P_{N,5}$,
\begin{equation}\label{eq:intro-periodic-clt}
 \frac1{\sqrt N}\sum_{i=1}^N
 \left(\psi(x_i)-\int\psi\dd\mu_5\right)
 \underset{N\to\infty}{\longrightarrow}\mathcal N(0,v)
 \quad\text{in law}.
\end{equation}
Consequently, an estimate of the form
\begin{equation}\label{eq:periodic-ce-unrestricted-target}
 \chi^2(P_{N,t}^{(k)}\mid\mu_t^{\otimes k})
 +\Hr(P_{N,t}^{(k)}\mid\mu_t^{\otimes k})
 \le C_T\frac{k^2}{N^2},\qquad 1\le k\le N,
\end{equation}
cannot hold for every smooth positive periodic log/Riesz trajectory throughout the Hilbert--Schmidt class.
\end{thm}



The mean-field measure and inverse temperature are fixed before $N$ tends to infinity. Unlike \cref{prop:iid-counterexample}, this example starts from a modulated Gibbs measure whose inverse temperature $\beta_0=\beta$ agrees with that of the dynamics. Its mechanism is a collective translation. Diffusion continually injects fluctuations into a collective translation mode that is invisible to the relative-coordinate interaction, while the product reference continues to resolve the much narrower one-particle scale. The accumulated Brownian variance is too large for a uniform relative second moment. The proof is given in \cref{sec:periodic-counterexample-proof}. Together, the two examples show that neither stationarity of the mean-field measure with independent initialization nor Gibbs initialization along an arbitrary smooth trajectory suffices for the full-particle bound. The bounded relative entropies in \eqref{iid:eq:entropy-main} and \eqref{eq:periodic-ce-entropy} do not control these relative second moments.



We carry out this Gibbs-initialized collective-translation construction in the periodic setting to simplify the presentation. Preliminary calculations suggest an analogous construction for the confined whole-space systems of \cref{thm:uniform}, using a collective Ornstein--Uhlenbeck translation mode in the isotropic case. We leave this extension to the interested reader. The stationary example in \cref{prop:iid-counterexample} uses independent initial data and a different fluctuation mechanism.


\subsection{Relation to previous work}\label{sec:previous-work}
Sharp local entropy estimates for mean-field diffusions were developed by Lacker \cite{Lacker} through a hierarchy of marginal relative entropies. Lacker and Le Flem \cite{LackerLeFlem} obtained uniform-in-time estimates under one-body logarithmic Sobolev inequalities and quantitative interaction and temperature conditions. These results distinguish the local entropy scale from the weaker estimate obtained by subadditivity. Hess-Childs and Rowan \cite[Theorem 1.4 and Remarks 1.6 and 1.8]{HCR} prove $L^2$ expansions for bounded interaction forces, including the sharp leading marginal estimate and higher-order corrections on a quantitative growing-marginal range. Their stated tensorized initialization can be relaxed when the initial laws satisfy the corresponding expansion bounds. Thus, correlated initialization is not itself a distinction from that method.


For singular interactions, relative-entropy and modulated-energy methods provide mean-field estimates without requiring a Lipschitz pair force; see \cite{JabinWang2018,BJW,BJWPKS,NguyenRosenzweigSerfaty,RSdiffusive,RSlsi,CdCRSRepulsive,CdCRSAttractive}. These are global approaches: relative entropy and modulated free energy are evaluated on the full $N$-particle measure, while modulated energy compares the empirical measure with the mean-field measure. For more details on these approaches, we refer the reader to the third author's recent survey \cite{RosenzweigJEDP}.

\begingroup
Wang \cite[Theorems 1--2 and Remark 2]{WangSharpLocal} treats periodic forces with a divergence-free $W^{-1,\infty}$ part and a bounded remainder. His entropy theorem gives the sharp $k^2/N^2$ dependence under a smallness condition, with a uniform-in-time conclusion under additional decay assumptions. His short-time weighted $L^2$ theorem gives the sharp $N^{-2}$ order for fixed $k$, with exponential rather than quadratic dependence on $k$. These force hypotheses differ from those of the repulsive whole-space gradient systems considered here.
\par\endgroup



Duerinckx and Jabin \cite[Corollary 1.3(i) and Remark 1.2]{DJBrownian} obtain short-time order-$N^{-1}$ estimates for every fixed marginal in a weighted density norm, starting from independent particles, and analyze further correlation corrections. Their overdamped result assumes square integrability of the force and a lower bound on its divergence. The square-integrability requirement alone restricts a Riesz force to $\s<\d/2-1$ at the collision singularity, whereas our estimates reach the full Hilbert--Schmidt range $\s<\d/2$ for the potential. In particular, they include the planar logarithm and three-dimensional Coulomb potential. The comparisons therefore concern the interaction assumptions, marginal-size dependence, initial laws, and time regime, rather than the fixed-$k$ power of $N$ alone. In their subsequent work \cite[Theorem~1.1 and Appendix~A]{DJFluctuations}, Duerinckx and Jabin derive the linearized Vlasov equation from signed infinitesimal perturbations of a stationary kinetic Gibbs equilibrium. That result concerns linear response around equilibrium, whereas the present theorems estimate the complete marginal probability densities along a nonlinear mean-field trajectory.



We also mention that, for equilibrium Gibbs measures, Lacker \cite{LackerGibbs} obtains sharp marginal Fisher-information and corresponding relative-entropy bounds under suitable functional inequalities, using a perturbative force-level hierarchy. Ren and Wang \cite{RenWangGibbs} derive sharp relative-entropy estimates for all marginal sizes from a hierarchy of conditional entropies, under functional-inequality, transport, and force-integrability hypotheses.

The estimates on singular Gibbs measures used below come from the authors' two companion works. The partition-function paper \cite[Theorems 2.2 and 2.8]{DGRpartition} supplies the partition and transport-commutator estimates, and also treats diffusive mean-field limits and fluctuations. See also the concurrent work of Wang and Zhao \cite{WZ26} for overlapping partition-function, commutator, and fluctuation results, and the earlier high-temperature partition bounds and limiting formulas of Duerinckx and Jabin \cite[Theorem~2.1 and Proposition~2.2]{DJFluctuations}. The static fluctuation formula and finite-particle adjoint identity \cite[Propositions 2.11 and 8.2]{DGRpartition} enter the periodic counterexample. The iid counterexample uses the same adjoint identity and commutator estimates, but its initial covariance is the ordinary iid covariance. Here, these inputs are combined with estimates for the evolution of complete marginal densities and with two obstructions to unrestricted full-particle chi-square control.



\begingroup
The authors' equilibrium paper \cite[Lemma~5.2 and Lemma~6.1]{DGRchaos} also develops the screened inverse and orthogonal product-space decomposition used here; we give proofs in \cref{lem:screened-inverse,lem:product-projection}. Its quantitative marginal-density expansion identifies the first Bogolyubov correction and, together with the associated divergence-limit theorem, gives the leading constants for fixed and growing sublinear marginals under the separate hypotheses of \cite[Theorems~2.3 and~2.4]{DGRchaos}. Neither its cavity expansion nor these asymptotic results are needed below. We instead use a weighted elliptic argument for the initial first marginal, since the same weighted estimates enter its propagation, and give the stationary lower-bound deduction in \cref{prop:stationary-sharpness}. The Hilbert--Schmidt threshold obstruction in \cite[Theorem~2.5]{DGRchaos} concerns the collision singularity. Both dynamical counterexamples here lie below that threshold: they concern amplification of initial iid fluctuations or accumulation of Brownian fluctuations, rather than loss of square integrability of the potential.
\par\endgroup



\begingroup
The attractive logarithmic interaction is treated in our companion work \cite{DGRAttractiveII}. In dimension two, for example, we prove there the corresponding sharp local chaos estimates throughout the full subcritical range $0<\beta<4$.
\par\endgroup



\begin{remark}[Kinetic dynamics]\label{rem:intro-kinetic}
Extending the sharp local estimates of this paper to underdamped Langevin dynamics is an interesting problem. {Feng and Wang \cite{FW26} recently introduced a kinetic version of the modulated Gibbs formalism popularized by Serfaty and the third author \cite{RSlsi}.} Together with the weighted marginal estimates of Bresch, Jabin, and Soler \cite{BJS}, this yields a short-time mean-field limit for singular kinetic systems, including the three-dimensional Vlasov--Poisson--Fokker--Planck equation. Obtaining the sharp dependence on the number of observed particles in relative entropy and chi-square divergence requires additional estimates.
\end{remark}


\subsection{Proof strategy}\label{sec:proof-strategy}

The proof of \cref{thm:finite} reduces to two estimates: a bounded full-density second moment and an order-$N^{-1}$ first-marginal error in the weighted $L^2$ norm. Indeed, the product-space projection of \cref{lem:product-projection} gives
\begin{equation}\label{eq:intro-projection-bound}
 \chi^2(P_{N,t}^{(k)}\mid\mu_t^{\otimes k})
 \le k\chi^2(P_{N,t}^{(1)}\mid\mu_t)
 +\frac{k(k-1)}{N(N-1)}\chi^2(P_{N,t}\mid\mu_t^{\otimes N}),\qquad 1\le k\le N.
\end{equation}
The full moment is also used inside the first-marginal argument, where it controls the centered two-particle source through \eqref{eq:pair-projection}. \Cref{fig:proof} records these dependencies.



Set
\begin{equation}\label{eq:three-densities}
 \Gamma_{N,t}=\Q_{N,\beta}(\mu_t),\quad
 G_{N,t}=\frac{\dd\Gamma_{N,t}}{\dd\mu_t^{\otimes N}},\quad
 q_{N,t}=\frac{\dd P_{N,t}}{\dd\Gamma_{N,t}},\quad
 f_{N,t}=q_{N,t}G_{N,t}.
\end{equation}
{The reference $\Gamma_{N,t}$ incorporates the singular pair repulsion.} Its failure to solve the particle forward equation is the centered transport commutator computed in \eqref{eq:defect-coupling}, whose exponential moment under $\Gamma_{N,t}$ is controlled by \cite{DGRpartition}. The Donsker--Varadhan lemma converts that estimate into a relative-entropy term in the equation for the norm of $q_{N,t}$. Choosing a decreasing exponent in the $L^{p(t)}(\Gamma_{N,t})$ estimate cancels this term. As long as $p(t)>2$, H\"older's inequality and static moments of $G_{N,t}$ give the full $L^2(\mu_t^{\otimes N})$ bound. The initial norm is estimated by comparing $\Q_{N,\beta_0}(\mu_0)$ and $\Q_{N,\beta}(\mu_0)$ in \cref{lem:initial-temperature}; it equals one in the matched case. This comparison and the subsequent loss of integrability give the temperature-dependent sufficient condition \eqref{eq:expenditure}.




\begin{figure}[htbp]
\centering

\begin{tikzpicture}[box/.style={draw,rounded corners=2pt,align=center,text width=9.4cm,inner sep=6pt,font=\small},arr/.style={-{Stealth[length=2mm]},semithick},node distance=6mm]
\node[box] (static) {Static partition and transport-commutator moments\\for the instantaneous modulated Gibbs measure};
\node[box,below=of static] (q) {Decreasing-exponent estimate for $q_{N,t}$\\and full product-space second moment of $f_{N,t}$};
\node[box,below=of q] (mark) {Canonical pair projection of order $N^{-1}$\\$+$\\weighted first-marginal elliptic and parabolic estimates};
\node[box,below=of mark] (all) {First-marginal density error of order $N^{-1}$\\$+$\\full product-space second moment\\$\Downarrow$\\sharp $k$-marginal bounds through $k=N$\\under the sufficient bound on integrability expenditure or near equilibrium};
\draw[arr] (static)--(q);\draw[arr] (q)--(mark);\draw[arr] (mark)--(all);
\end{tikzpicture}
\caption{Proof of the sharp local estimates through the full-particle scale. The full moment, supplied under the sufficient bound on integrability expenditure or near equilibrium, controls the pair source inside the first-marginal estimate and is used again in the final product-space projection.}
\label{fig:proof}
\end{figure}



We estimate the initial first marginal directly from the modulated Gibbs measure. Differentiation of its density and weighted elliptic duality convert the static second-moment estimate into an order-$N^{-1}$ first-marginal bound; see \cref{prop:initial-mark}. We then propagate this estimate using only the exact first-marginal equation \eqref{eq:first-mark-equation}. The full second moment controls its centered two-particle source through \eqref{eq:pair-projection}, so no evolution equation for a higher marginal is required. Testing this source against the backward mean-field propagator produces a difference of two test-function gradients. {The weighted commutator estimate of \cref{lem:weighted-commutator} combines gradient and Hessian bounds while leaving a factor of the background density along the segment joining the two points.} For $|x-y|>1$, the force bound and the mass of the background suffice. Backward smoothing gives a time singularity that is integrable throughout $\s<\d/2$. This proves the first-marginal estimate needed in \eqref{eq:intro-projection-bound}, without square integrability of the force or a positive lower bound on the whole-space density.



The hierarchy argument of Duerinckx and Jabin \cite[Sections~3--4]{DJBrownian} estimates linear correlations of all orders with time-dependent Gibbs-type weights. Their later equilibrium-fluctuation argument \cite[Sections~3--4]{DJFluctuations} already uses global density bounds and orthogonal product-space projections, followed by the first-marginal equation. This reduction is also central here. The additional task is to control a Gibbs reference moving along a nonlinear mean-field trajectory and to obtain a quantitative first-marginal error. The decreasing-exponent estimate controls the reference defect, and the weighted commutator estimate controls the singular source in that equation.



The moving Gibbs formalism is part of the modulated-free-energy method \cite{BJW,RSlsi}. Varying-exponent differentiation is classical in hypercontractivity and smoothing \cite{Gross,CarlenLoss}; the decreasing exponent used here controls the reference defect without a many-particle logarithmic Sobolev inequality. Previous works also use time-dependent norms; for example, in \cite{BJS}, the time dependence is in an energy weight, while the Lebesgue exponent remains fixed.



For \cref{thm:uniform}, two stability arguments complement the finite-time estimates. {We estimate the homogeneous first-marginal equation in the quadratic form associated with the mean-field linearization, with the finite-particle factor $(N-1)/N$ included in the interaction term; see \eqref{eq:screened-energy}.} This form combines the weighted $L^2$ norm of the signed density error with its interaction energy. Near equilibrium, contraction on fixed time intervals, combined with the source estimate above, prevents accumulation of the first-marginal error. Separately, the mean-field stability argument transfers the $L^2(\pi)$ spectral gap of the linearized equation to the logarithmic derivative norm by comparison with the anisotropic Ornstein--Uhlenbeck semigroup. This gives the uniform background bounds and integrable velocity used in the proof. The first argument is given in \cref{prop:uniform-mark}; the mean-field estimates are stated in \cref{sec:backgrounds} and proved in Appendix~\ref{app:background}. The unconfined results follow through the self-similar transformation of Serfaty and the third author \cite{RSselfsimilar}, with an additional decreasing-coupling argument for $\s>0$.



The counterexamples compare a terminal linear statistic under the actual particle law and under the product reference. A Gaussian limit with variance greater than twice the product variance has infinite chi-square divergence from the product Gaussian limit. Data processing and weak lower semicontinuity then exclude a bounded full-particle chi-square divergence. The two constructions obtain this variance separation in different ways. For independent initialization, the backward linearized evolution amplifies the initial product fluctuation, while the Brownian contribution is lower order; see \eqref{iid:eq:covariance} and \cref{iid:prop:variance-crossing}. For the periodic Gibbs preparation, the Brownian contribution alone exceeds twice the product variance; see \cref{ce:prop:noiseobstruction}. The terminal statistic is, up to a vanishing error, the sum of an initial fluctuation and an asymptotically independent Gaussian Brownian fluctuation. The Brownian variance already forces infinite chi-square divergence of every subsequential limit from the product Gaussian, irrespective of the initial limiting distribution. The static central limit theorem for the initial modulated Gibbs measure then identifies the initial fluctuation as Gaussian and gives the explicit terminal Gaussian limit in \eqref{eq:intro-periodic-clt}.


\subsection{Organization of the paper}\label{sec:organization}

\Cref{sec:static,sec:decreasing} give the static inputs, the moving-reference defect, and the decreasing-exponent estimate. \Cref{sec:weighted,sec:marginals} prove the weighted first-marginal estimates and the passage to all marginals. \Cref{sec:confined-proofs} verifies the confined background hypotheses, justifies approximation, and completes the confined proofs. \Cref{sec:selfsimilar} proves the unconfined results. The proof of \cref{thm:periodic-counterexample} is in \cref{sec:periodic-counterexample-proof}. The required one-body stability and finite-particle well-posedness arguments are in Appendices~\ref{app:background} and~\ref{app:particle-wp}, respectively. The proof of \cref{prop:iid-counterexample}, including the weaker marginal entropy estimate, is in Appendix~\ref{app:iid-counterexample}.


\subsection{Statement on AI use}





This working manuscript was generated with the assistance of generative AI tools, including OpenAI's ChatGPT and Codex and Anthropic's Claude, for literature identification, the development and testing of mathematical arguments, exposition, and consistency checks. The authors have not yet reviewed and independently verified all AI-generated content. They are in the process of formalizing the results in Lean and revising the natural-language presentation. The manuscript is therefore being shared as a draft, and its statements, proofs, references, and exposition remain subject to correction. Any completed formal verification applies only to the statements and versions explicitly identified in the accompanying verification records.


\par\endgroup

\begingroup

\clearpage
\section{Static inputs and the moving-reference defect}\label{sec:static}

This section supplies the exponential estimate for the moving-reference defect used in \cref{sec:decreasing}. We first recall the partition and commutator estimates, including their uniformity under heat truncation. We then compute the defect and estimate its exponential moment. The screened inverse recorded below will also enter the initial one-mark estimate in \cref{prop:initial-mark}.

To include the self-similar applications in the same calculation, we allow a prescribed $C_{\mathrm{loc}}^1$ nonincreasing interaction coefficient $c(t)\in[0,c_*]$. The particle force and mean-field force are both multiplied by $c(t)$, and the reference is $\Q_{N,\beta c(t)}(\mu_t)$. The constant case $c=1$ gives the confined theorems; decreasing coefficients enter \cref{thm:unconfined}. Set
\begin{equation}\label{eq:energy-short}
 E_{N,t}=N\Fr_N(X_N,\mu_t),\qquad C_N[v]=\frac N2\mathsf A_1(v,X_N,\mu_t),
\end{equation}
where
\begin{equation}\label{eq:A1}
 \mathsf A_1(v,X_N,\mu)=\iint_{x\ne y}(v(x)-v(y))\cdot\nabla\g(x-y)
 \dd\left(\frac1N\sum_i\delta_{x_i}-\mu\right)^{\otimes2}(x,y).
\end{equation}
Thus, $C_N$ denotes the transport commutator, with precisely half the normalization of $N\mathsf A_1$.

\subsection{Static and heat-truncation estimates}
All backgrounds in the applications satisfy a common bound
\begin{equation}\label{eq:gaussian-envelope}
 0<\mu(x)\le M e^{-a|x|^2},\qquad a>0.
\end{equation}
No positive lower bound is asserted. The static results require substantially less; \eqref{eq:gaussian-envelope} conveniently supplies every logarithmic moment used here.

We collect four inputs and indicate their later roles. \Cref{prop:static-inputs} gives the partition and commutator estimates used in \cref{sec:decreasing}. \Cref{lem:heat-potential-estimates} records the potential estimates that make those static bounds uniform under heat truncation. \Cref{lem:heat-force-estimates} supplies the cutoff-uniform force estimates used in \cref{lem:weighted-commutator,lem:fixed-background-approximation,lem:approximation}. Finally, \cref{lem:screened-inverse} gives the inverse bound used in the initial one-mark estimate \cref{prop:initial-mark}.

\begin{prop}[Partition and commutator estimates]\label{prop:static-inputs}
For the kernels \eqref{eq:kernel}, every fixed $B<\infty$ admits constants, uniform over \eqref{eq:gaussian-envelope}, such that
\begin{equation}\label{eq:static-partition}
 e^{bI_\mu}\le\K_{N,b}(\mu)\le C_B,\quad 0\le b\le B,\qquad
 |I_\mu|+\norm{\widetilde\g}_{L^2(\mu^{\otimes2})}+\norm{h_\mu}_{L^2(\mu)}\le C.
\end{equation}
Here, $\widetilde\g$ is the centered kernel
\begin{equation}\label{eq:static-centered-kernel}
 \widetilde\g(x,y)=\g(x-y)-h_\mu(x)-h_\mu(y)+2I_\mu.
\end{equation}
In particular,
\begin{equation}\label{eq:static-powers}
 \int G_{N,b}^{p}\dd\mu^{\otimes N}=\frac{\K_{N,pb}(\mu)}{\K_{N,b}(\mu)^p}\le C_{p,B},
       \qquad 0\le b\le B,
\end{equation}
for every fixed $1\le p<\infty$.

For fixed $\beta,c_*>0$, there are $\kappa,K>0$ such that
\begin{equation}\label{eq:static-commutator}
 \E_{\Q_{N,\beta c}(\mu)}[e^{\sigma\kappa c C_N[v]}]\le K,
 \qquad 0\le c\le c_*,\quad\norm{\nabla v}_{L^\infty}\le1,\quad\sigma=\pm1.
\end{equation}
These estimates hold uniformly for the heat truncations defined in \eqref{eq:heat-cutoff}.
\end{prop}
\begin{proof}
The partition estimate is \cite[Theorem 2.2]{DGRpartition}, and the commutator estimate is \cite[Theorem 2.8 and Sections 7.3--7.4]{DGRpartition}. Appendix A of that reference makes their background dependence explicit. The bound \eqref{eq:gaussian-envelope} controls the logarithmic energy and all fixed polynomial moments; the local singularity is square integrable exactly when $2\s<\d$. The lower bound in \eqref{eq:static-partition} is Jensen's inequality, since the product expectation of $E_N$ is $-I_\mu$. In the whole-space logarithmic case $I_\mu$ need not be nonnegative: the useful uniform lower bound is $e^{-B|I_\mu|}$, not one. Uniformity for $0\le b\le B$ also follows by H\"older from the estimate at $b=B$, and \eqref{eq:static-powers} follows directly from the definition of $G_{N,b}$.

It remains to explain how the local numerator estimate of \cite[Lemma~7.2]{DGRpartition} yields the coefficient proportional to $c$ in \eqref{eq:static-commutator}. With $C_N=N\mathsf A_1/2$, the cited estimate has the form
\begin{equation}\label{eq:local-numerator}
 \E_{\mu^{\otimes N}}[e^{\sigma a_0 b C_N[v]-bE_N}]\le C,
                  \qquad 0<b\le b_0,
\end{equation}
where $a_0,b_0>0$ are common to the background class and $\norm{\nabla v}_{L^\infty}\le1$. For $b=\beta c>0$, set $b_1=\min\{b,b_0\}$. Cauchy--Schwarz gives
\begin{equation}\label{eq:commutator-coupling-interpolation}
 \E_{\mu^{\otimes N}}[e^{\sigma a_0b_1C_N[v]/2-bE_N}]
 \le \E_{\mu^{\otimes N}}[e^{\sigma a_0b_1C_N[v]-b_1E_N}]^{1/2}
       \K_{N,2b-b_1}(\mu)^{1/2}\le C.
\end{equation}
The constants are uniform for $0<b\le\beta c_*$, since $0\le2b-b_1\le2\beta c_*$. Dividing by $\K_{N,b}(\mu)$ and using its lower bound in \eqref{eq:static-partition} gives the corresponding Gibbs moment. Finally,
\begin{equation}\label{eq:commutator-coefficient-choice}
 \kappa\le\frac{a_0}{2}\min\left\{\beta,\frac{b_0}{c_*}\right\}
 \quad\Longrightarrow\quad
 \kappa c\le\frac{a_0b_1}{2}.
\end{equation}
H\"older's inequality gives \eqref{eq:static-commutator} at the smaller coefficient. At $c=0$, its integrand is one. The cutoff-uniformity follows from \cref{lem:heat-potential-estimates} and the corresponding proofs in the cited reference.
\end{proof}

For the heat truncation, set
\begin{equation}\label{eq:heat-normalization}
 p_t(x):=(4\pi t)^{-\d/2}e^{-|x|^2/(4t)},\qquad
 a_{\d,\s}:=\frac{(4\pi)^{\d/2}}{2^{\s+1}\Gamma(1+\s/2)}.
\end{equation}
For $\s>0$, define
\begin{equation}\label{eq:heat-cutoff}
 \g_\varepsilon(x):=a_{\d,\s}\int_\varepsilon^\infty
             t^{(\d-\s)/2-1}p_t(x)\dd t,
             \qquad \varepsilon>0.
\end{equation}
For $\s=0$, replace $p_t(x)$ by $p_t(x)-p_t(e)$, where $e$ is a fixed unit vector. This fixes $\g_\varepsilon(e)=0$ and the limit $-\log|x|$. We set $\g_0:=\g$. The parameter $\varepsilon$ is a heat time: the corresponding length scale in \cite[Section~4.1]{DGRpartition} is $\sqrt\varepsilon$. In the logarithmic case our normalization differs from that reference by an additive constant, which cancels from the normalized Gibbs density. Each auxiliary particle equation and its modulated Gibbs reference use the same truncated kernel.

For $\eta>0$ and $\varepsilon\ge0$, set
\begin{equation}\label{eq:heat-short-range}
 \mathsf f_{\eta,\varepsilon}(x):=
 a_{\d,\s}\int_\varepsilon^{\max\{\varepsilon,\eta^2\}}
                   t^{(\d-\s)/2-1}p_t(x)\dd t.
\end{equation}

\begin{lemma}[Heat-truncation estimates for the potential]\label{lem:heat-potential-estimates}
Let $\mu$ satisfy \eqref{eq:gaussian-envelope}, let $0\le\varepsilon\le1$, and let $0<\eta\le1$. Then the following estimates hold with constants independent of $\varepsilon$ and $\eta$.
\begin{enumerate}[(i)]
\item The short-range part satisfies
\begin{equation}\label{eq:heat-short-mass}
 \norm{\mathsf f_{\eta,\varepsilon}}_{L^1}
 \le C\eta^{\d-\s},\qquad
 \norm{\mathsf f_{\eta,\varepsilon}*\mu}_{L^\infty}
 \le CM\eta^{\d-\s},
\end{equation}
and, for every Lipschitz vector field $v$,
\begin{equation}\label{eq:heat-short-transport}
 |(v(x)-v(y))\cdot\nabla\mathsf f_{\eta,\varepsilon}(x-y)|
 \le C\norm{\nabla v}_{L^\infty}\mathsf f_{\sqrt2\eta,2\varepsilon}(x-y).
\end{equation}
\item The remaining potential $\g_{\max\{\varepsilon,\eta^2\}}$ has Fourier multiplier
\begin{equation}\label{eq:heat-long-fourier}
 0\le\widehat{\g_{\max\{\varepsilon,\eta^2\}}}(\xi)
 \le C|\xi|^{\s-\d},\qquad \xi\ne0.
\end{equation}
If $q$ is obtained from this potential by a smooth Fourier cutoff bounded by one and vanishing on $|\xi|\le K_0$, where $K_0\ge1$, then
\begin{equation}\label{eq:heat-high-frequency}
 \norm q_{L^2(\R^\d)}^2
 \le C\int_{|\xi|>K_0}|\xi|^{2\s-2\d}\dd\xi
 \le CK_0^{2\s-\d},
\end{equation}
\begin{equation}\label{eq:heat-row-operator}
 \sup_x\int|q(x-y)|^2\dd\mu(y)\le M\norm q_2^2,\qquad
 \norm{q*\!(\mu\,\cdot)}_{L^2(\mu)\to L^2(\mu)}
 \le\norm{q(x-y)}_{L^2(\mu^{\otimes2})}\le\sqrt M\norm q_2.
\end{equation}
Centering $q(x-y)$ preserves these bounds up to a numerical factor. Moreover,
\begin{equation}\label{eq:heat-long-diagonal}
 \left|\g_{\max\{\varepsilon,\eta^2\}}(0)\right|
 \le\begin{cases}C\eta^{-\s},&\s>0,\\ C(1+|\log\eta|),&\s=0,\end{cases}
\end{equation}
and, in the logarithmic case,
\begin{equation}\label{eq:heat-log-infrared}
 t^{\d/2-1}|p_t(z)-p_t(e)|\le C(1+|z|^2)t^{-2},
 \qquad t\ge1.
\end{equation}
\item The centered interaction operator associated with $\g_\varepsilon$ is nonnegative on $L^2_0(\mu)$.
\end{enumerate}
\end{lemma}
\begin{proof}
Integrating the heat kernel in space gives the first estimate in \eqref{eq:heat-short-mass}; the second follows from $\norm\mu_\infty\le M$. For \eqref{eq:heat-short-transport}, use $|v(x)-v(y)|\le\norm{\nabla v}_{L^\infty}|x-y|$ and the Gaussian bound
$|z||\nabla p_t(z)|\le Cp_{2t}(z)$, followed by the change of heat-time variable $r=2t$.

Taking Fourier transforms in the heat representation gives \eqref{eq:heat-long-fourier}. Plancherel's theorem then gives \eqref{eq:heat-high-frequency}. The density bound in \eqref{eq:gaussian-envelope} yields the first estimate in \eqref{eq:heat-row-operator}, while the second is the Hilbert--Schmidt bound for the integral operator with kernel $q(x-y)$. Subtracting either one-variable average, and then the total average, changes these estimates only by a numerical factor.

Evaluating the heat representation at the origin gives \eqref{eq:heat-long-diagonal}. For the logarithm, the subtraction at $e$ makes the large-time integral convergent; the mean-value theorem applied to the Gaussian exponential gives \eqref{eq:heat-log-infrared}. Hence, changing the lower heat cutoff does not affect the large-time cancellation, and the moments supplied by \eqref{eq:gaussian-envelope} control the resulting one-body terms.

Finally, the Fourier multiplier in \eqref{eq:heat-long-fourier} is nonnegative. The additive constant in the logarithmic representation does not contribute to the quadratic form on zero-mass functions. Centering the kernel likewise leaves that quadratic form unchanged, which proves the last assertion.
\end{proof}

The estimates in \cref{lem:heat-potential-estimates} are the short-range, high-frequency, row, and diagonal bounds used in the cutoff-uniform partition and commutator arguments of \cite[Sections~4, 5.1, and~7]{DGRpartition}. The partition argument also uses nonnegativity of the centered quadratic form, whereas the commutator argument retains the deterministic factor $N$ in the diagonal bound.

\begin{lemma}[Heat-truncation estimates for the force]\label{lem:heat-force-estimates}
For $0\le\varepsilon\le1$,
\begin{equation}\label{eq:force-cutoff}
 |\nabla\g_\varepsilon(x)|\le C|x|^{-\s-1},\qquad
 \sup_{0\le\varepsilon\le1}
 \left(\norm{\nabla\g_\varepsilon}_{L^1(B_1)}
       +\norm{\nabla\g_\varepsilon}_{L^\infty(B_1^c)}\right)\le C.
\end{equation}
As $\varepsilon\downarrow0$, the force converges in $L^1(B_1)$ and with all derivatives locally away from the origin. Moreover,
\begin{equation}\label{eq:heat-weighted-force}
 \sup_{0\le\varepsilon\le1}\int_{|z|\le1}|z|^{2\alpha}
               |\nabla\g_\varepsilon(z)|^2\dd z
 \le C\int_0^1r^{\d-1+2\alpha-2\s-2}\dd r<\infty,
 \qquad\alpha>\s+1-\d/2.
\end{equation}
\end{lemma}
\begin{proof}
Differentiating \eqref{eq:heat-cutoff} and using the scaling of $\nabla p_t$ gives the pointwise estimate in \eqref{eq:force-cutoff}. Its local $L^1$ and exterior $L^\infty$ consequences follow because $\s+1<\d$. The difference $\nabla\g-\nabla\g_\varepsilon$ is the integral over $0<t<\varepsilon$; dominated convergence gives convergence in $L^1(B_1)$, while differentiation under the integral gives convergence of every derivative on compact sets separated from the origin. Finally, the pointwise bound in \eqref{eq:force-cutoff} and radial integration give \eqref{eq:heat-weighted-force}.
\end{proof}

\begin{lemma}[Screened inverse]\label{lem:screened-inverse}
Let $T_\mu$ be the integral operator on $L^2_0(\mu)$ with kernel $\widetilde\g$. Suppose that the bounds in \eqref{eq:static-partition} hold up to coupling $2b$, where $b\ge0$. Then
\begin{equation}\label{eq:screened-inverse}
 \left\|\left(\Id+b\frac{N-1}{N}T_\mu\right)^{-1}\right\|_{L^2_0(\mu)\to L^2_0(\mu)}
 \le\frac{2N}{N+1}\le2.
\end{equation}
The conclusion is uniform under the heat truncation and does not require positive definiteness of the interaction.
\end{lemma}
\begin{proof}
We give a short proof of the estimate in \cite[Lemma~5.2]{DGRchaos} using the partition bound already stated here. If $b=0$ or $N=1$, the inverse is the identity. Suppose that $b>0$ and $N\ge2$. Fix a bounded real $\varphi\in L^2_0(\mu)$ and set $\nu=(1+\eta\varphi)\mu$, where $|\eta|\norm{\varphi}_\infty<1/2$. The bounds in \eqref{eq:static-partition} imply that $(x,y)\mapsto\g(x-y)$ belongs to $L^2(\mu^{\otimes2})$, so all the following integrals are finite. By \eqref{eq:modulated},
\begin{equation}\label{eq:screened-tilted-energy}
 \E_{\nu^{\otimes M}}[M\Fr_M(X_M,\mu)]
 =\frac{M\eta^2}{2}\langle\varphi,T_\mu\varphi\rangle_{L^2(\mu)}-I_\nu.
\end{equation}
Changing the product reference from $\mu^{\otimes M}$ to $\nu^{\otimes M}$ and applying Jensen's inequality gives
\begin{equation}\label{eq:screened-jensen}
 \log\K_{M,2b}(\mu)
 \ge -M\bigl(\Hr(\nu\mid\mu)+b\eta^2\langle\varphi,T_\mu\varphi\rangle_{L^2(\mu)}\bigr)+2bI_\nu.
\end{equation}
For fixed $\eta$, divide by $M$ and let $M\to\infty$, using the uniform partition bound. Then
\begin{equation}\label{eq:screened-free-energy}
 \Hr((1+\eta\varphi)\mu\mid\mu)+b\eta^2\langle\varphi,T_\mu\varphi\rangle_{L^2(\mu)}\ge0.
\end{equation}
Since $\varphi$ is bounded and centered, the entropy is $\eta^2\norm{\varphi}_{L^2(\mu)}^2/2+o(\eta^2)$. Dividing by $\eta^2$ and then letting $\eta\to0$ yields
\begin{equation}\label{eq:screened-form-bound}
 \langle\varphi,T_\mu\varphi\rangle_{L^2(\mu)}
 \ge-\frac{1}{2b}\norm{\varphi}_{L^2(\mu)}^2.
\end{equation}
Bounded centered functions are dense in $L^2_0(\mu)$, and $T_\mu$ is bounded and self-adjoint. Thus, \eqref{eq:screened-form-bound} holds throughout that space, and
\begin{equation}\label{eq:screened-coercivity}
 \Id+b\frac{N-1}{N}T_\mu\ge\frac{N+1}{2N}\Id.
\end{equation}
The spectral theorem proves \eqref{eq:screened-inverse}. The argument applies at each heat cutoff with the same final constant. For the repulsive kernels, \cref{lem:heat-potential-estimates}\textup{(iii)} gives $T_\mu\ge0$ and hence the stronger inverse bound $1$.
\end{proof}


\subsection{The exact defect, including a varying coupling}
Set $\Gamma_{N,t}=\Q_{N,\beta c(t)}(\mu_t)$, and let $L_{N,t}$ be the particle generator with interaction coefficient $c(t)$. The continuity velocity is
\begin{equation}\label{eq:velocity-coupling}
 u_t=-\nabla V-c(t)\nabla h_t-\beta^{-1}\nabla\log\mu_t,
              \qquad\partial_t\mu_t+\diver(\mu_tu_t)=0.
\end{equation}
The score identity is
\begin{equation}\label{eq:score}
 b_i^N=u_t(x_i)+\beta^{-1}\nabla_i\log\Gamma_{N,t}.
\end{equation}
Consequently, the failure of $\Gamma_{N,t}$ to solve the particle forward equation is
\begin{equation}\label{eq:defect}
 \partial_t\Gamma_{N,t}-L_{N,t}^*\Gamma_{N,t}=R_{N,t}\Gamma_{N,t},
\end{equation}
\begin{equation}\label{eq:defect-coupling}
 R_{N,t}=-\beta\left(c(t)C_N[u_t]+c'(t)E_{N,t}
                  -\E_{\Gamma_{N,t}}[c(t)C_N[u_t]+c'(t)E_{N,t}]\right).
\end{equation}
Indeed, the continuity equation and integration by parts give
\begin{equation}\label{eq:static-material-potential}
 (\partial_t+u_t(x)\cdot\nabla)h_t(x)
 =\int(u_t(x)-u_t(y))\cdot\nabla\g(x-y)\dd\mu_t(y),
\end{equation}
and
\begin{equation}\label{eq:static-material-background-energy}
 \frac{\dd}{\dd t}I_{\mu_t}
 =\frac12\iint(u_t(x)-u_t(y))\cdot\nabla\g(x-y)
                    \dd\mu_t(x)\dd\mu_t(y).
\end{equation}
Applying these identities to the pair, one-body, and constant terms in \eqref{eq:modulated}, with their original coefficients, yields
\begin{equation}\label{eq:static-material-particle-energy}
 \left(\partial_t+\sum_{i=1}^N u_t(x_i)\cdot\nabla_i\right)E_{N,t}
 =C_N[u_t].
\end{equation}
The product reference satisfies the corresponding continuity equation. Integrating the resulting equation for the unnormalized Gibbs density therefore gives
\begin{equation}\label{eq:static-normalizer-derivative}
 \frac{\dd}{\dd t}\log\K_{N,\beta c(t)}(\mu_t)
 =-\beta\E_{\Gamma_{N,t}}[c(t)C_N[u_t]+c'(t)E_{N,t}].
\end{equation}
Dividing the density $e^{-\beta c(t)E_{N,t}}\mu_t^{\otimes N}$ by its integral $\K_{N,\beta c(t)}(\mu_t)$ gives $\Gamma_{N,t}$. Using \eqref{eq:score}, its logarithmic source is
\begin{equation}
 \frac{\partial_t\Gamma_{N,t}-L_{N,t}^*\Gamma_{N,t}}{\Gamma_{N,t}} =-\beta\bigl(c(t)C_N[u_t]+c'(t)E_{N,t}\bigr) -\frac{\dd}{\dd t}\log\K_{N,\beta c(t)}(\mu_t) =R_{N,t}. \label{eq:static-normalized-defect}
\end{equation}
This proves \eqref{eq:defect-coupling}. These calculations hold at every smooth heat cutoff, where Gaussian tails justify the integrations by parts; \cref{sec:completion} gives the passage to the singular particle laws.

For $c=1$, Jensen's inequality with both signs in \eqref{eq:static-commutator} gives
\begin{equation}\label{eq:scalar-defect}
 \log\E_{\Gamma_{N,t}}[e^{-R_{N,t}/\lambda_0(t)}]\le C_0,
              \qquad\lambda_0(t)=\frac\beta\kappa\norm{\nabla u_t}_{L^\infty}.
\end{equation}
Indeed, if $\norm{\nabla u_t}_{L^\infty}>0$, \eqref{eq:defect-coupling} and Jensen's inequality give
\begin{align}
 \log\E_{\Gamma_{N,t}}[e^{-R_{N,t}/\lambda_0(t)}]
 &=\log\E_{\Gamma_{N,t}}[e^{\kappa C_N[u_t]/\norm{\nabla u_t}_{L^\infty}}]
       -\frac{\kappa}{\norm{\nabla u_t}_{L^\infty}}\E_{\Gamma_{N,t}}[C_N[u_t]]\notag\\
 &\le\log\E_{\Gamma_{N,t}}[e^{\kappa C_N[u_t]/\norm{\nabla u_t}_{L^\infty}}]
       +\log\E_{\Gamma_{N,t}}[e^{-\kappa C_N[u_t]/\norm{\nabla u_t}_{L^\infty}}]
 \le2\log K.
 \label{eq:static-defect-exponential-calculation}
\end{align}
If $\norm{\nabla u_t}_{L^\infty}=0$, the velocity is spatially constant, so $C_N[u_t]=R_{N,t}=0$.

For nonincreasing coupling an additional one-sided energy moment suffices. For every fixed $b>0$,
\begin{equation}\label{eq:energy-exponential}
 \sup_{\substack{N\ge2,\ 0\le c\le c_*\\ \mu\ \text{satisfies }\eqref{eq:gaussian-envelope}}}
 \E_{\Q_{N,\beta c}(\mu)}
       [e^{-b(E_N-\E_{\Q_{N,\beta c}(\mu)}[E_N])}]\le C_b<\infty.
\end{equation}
To estimate the mean subtracted here, the centered degree-two decomposition gives
\begin{equation}\label{eq:static-energy-decomposition}
 E_N+I_\mu=\frac1N\sum_{i<j}\widetilde\g(x_i,x_j)
                -\frac1N\sum_{i=1}^N\bigl(h_\mu(x_i)-2I_\mu\bigr).
\end{equation}
Orthogonality under $\mu^{\otimes N}$ then yields
\begin{equation}\label{eq:static-energy-L2}
 \norm{E_N}_{L^2(\mu^{\otimes N})}^2
 =I_\mu^2+\frac{N-1}{2N}\norm{\widetilde\g}_{L^2(\mu^{\otimes2})}^2
         +\frac1N\norm{h_\mu-2I_\mu}_{L^2(\mu)}^2\le C.
\end{equation}
Consequently, Cauchy--Schwarz and \eqref{eq:static-powers} give
\begin{equation}\label{eq:static-energy-mean}
 \sup_{\substack{N\ge2,\ 0\le c\le c_*\\ \mu\ \text{satisfies }\eqref{eq:gaussian-envelope}}}
 \left|\E_{\Q_{N,\beta c}(\mu)}[E_N]\right|
 \le\sup_{N,c,\mu}\norm{G_{N,\beta c}}_{L^2(\mu^{\otimes N})}
                      \norm{E_N}_{L^2(\mu^{\otimes N})}\le C.
\end{equation}
The second supremum is over the same $N,c,\mu$. The exact identity
\begin{equation}\label{eq:static-centered-energy-moment}
 \E_{\Q_{N,\beta c}(\mu)}
 [e^{-b(E_N-\E_{\Q_{N,\beta c}(\mu)}[E_N])}]
 =e^{b\E_{\Q_{N,\beta c}(\mu)}[E_N]}
   \frac{\K_{N,\beta c+b}(\mu)}{\K_{N,\beta c}(\mu)}
\end{equation}
and \eqref{eq:static-partition} prove \eqref{eq:energy-exponential}.

Set
\begin{equation}\label{eq:lambda-coupling}
 \lambda_0(t):=\frac\beta\kappa\norm{\nabla u_t}_{L^\infty}+\frac\beta b|c'(t)|.
\end{equation}
For $\lambda_0(t)>0$, H\"older's inequality, \eqref{eq:static-commutator}, and \eqref{eq:energy-exponential} give
\begin{equation}\label{eq:static-combined-defect}
 \log\E_{\Gamma_{N,t}}[e^{-R_{N,t}/\lambda_0(t)}]
 \le \frac{2\beta\norm{\nabla u_t}_{L^\infty}\log K}{\kappa\lambda_0(t)}
       +\frac{\beta|c'(t)|\log C_b}{b\lambda_0(t)}
 \le\max\{2\log K,\log C_b\}.
\end{equation}
Terms with zero coefficient are omitted; if $\lambda_0(t)=0$, then $R_{N,t}=0$. The sign $c'\le0$ is important: the required energy exponential has the repulsive, negative sign. Its coefficient $b$ can be arbitrarily large and fixed, because the static partition estimate holds at every fixed inverse temperature. Thus, decreasing the coupling creates no uncontrolled positive-energy exponential.

\par\endgroup

\begingroup

\section{Full-particle integrability with a decreasing exponent}\label{sec:decreasing}

The aim of this section is to bound the full density in $L^2(\mu_t^{\otimes N})$. We first estimate the density relative to the moving Gibbs reference, allowing the exponent to decrease. We then compare the initial Gibbs preparation with that reference and pass to the product reference using the static density moments. The three steps are \cref{prop:decreasing,lem:initial-temperature,cor:full-moment}. The whole-space integrations by parts at a fixed heat cutoff are justified in \cref{lem:approximation}.

\begin{prop}[A varying-power estimate for a moving reference]\label{prop:decreasing}
Let $P_t$ solve the diffusion forward equation and let a positive moving probability density $\Gamma_t$ satisfy \eqref{eq:defect}. Suppose
\begin{equation}\label{eq:abstract-exponential}
 \log\E_{\Gamma_t}[e^{-R_t/\lambda_0(t)}]\le C_0,
\end{equation}
where $\lambda_0$ is nonnegative and locally integrable and $R_t=0$ when $\lambda_0(t)=0$. Work either with smooth positive data on a compact manifold or with the smooth whole-space approximations of \cref{lem:approximation}. For finite $p_0>1$, define
\begin{equation}\label{eq:exponent-ode}
 A(t)=\int_0^t\lambda_0(r)\dd r,\qquad
 1-\frac1{p(t)}=\left(1-\frac1{p_0}\right)e^{-A(t)}.
\end{equation}
Then $q_t=\dd P_t/\dd\Gamma_t$ satisfies
\begin{equation}\label{eq:varying-bound}
 \log\norm{q_t}_{L^{p(t)}(\Gamma_t)}
 \le\log\norm{q_0}_{L^{p_0}(\Gamma_0)}
       +C_0\left(1-\frac1{p_0}\right)(1-e^{-A(t)}).
\end{equation}
\end{prop}
\begin{proof}
Write $b$ for the diffusion drift, retain the diffusion coefficient $\beta^{-1}$, and suppress the time variable when no confusion can arise. Substituting $P=q\Gamma$ into the forward equation gives
\begin{equation}\label{eq:q-general}
 \partial_tq=\beta^{-1}\Delta q+
     (2\beta^{-1}\nabla\log\Gamma-b)\cdot\nabla q-Rq.
\end{equation}
For fixed $p>1$, multiply by $pq^{p-1}\Gamma$ and differentiate the moving measure. Integration by parts cancels the transport terms. The defect contributes $-pR+R$, so
\begin{equation}\label{eq:fixed-power}
 \frac\dd{\dd t}\int q^p\dd\Gamma
 =-\frac{4(p-1)}{\beta p}\int|\nabla(q^{p/2})|^2\dd\Gamma
       -(p-1)\int Rq^p\dd\Gamma.
\end{equation}
Set $M_p=\int q^p\dd\Gamma$ and introduce the normalized power tilt
\begin{equation}\label{eq:power-tilt}
 \dd\nu_p=\frac{q^p}{M_p}\dd\Gamma
              =\frac{q^{p-1}}{M_p}\dd P.
\end{equation}
The measure $\nu_p$ identifies the contribution from the varying exponent as a relative entropy. Indeed, differentiating with $p=p(t)$ gives
\begin{equation}
 \frac\dd{\dd t}\log\norm q_{L^{p(t)}(\Gamma_t)} ={}-\frac{4(p-1)}{\beta p^2} \frac{\int|\nabla(q^{p/2})|^2\dd\Gamma}{M_p} +\frac{p-1}{p}\E_{\nu_p}[-R] +\frac{p'}{p^2}\Hr(\nu_p\mid\Gamma). \label{eq:log-power}
\end{equation}
The term from differentiating the exponent is $p'p^{-2}(p\E_{\nu_p}[\log q]-\log M_p)$, and $p\E_{\nu_p}[\log q]-\log M_p=\Hr(\nu_p\mid\Gamma)$ by \eqref{eq:power-tilt}. For $\lambda_0(t)>0$, apply the Donsker--Varadhan lemma, as in \cite[Section~7.1]{DGRpartition}, to $-R/\lambda_0$. Equivalently, use nonnegativity of relative entropy with respect to the probability measure whose density relative to $\Gamma$ is proportional to $e^{-R/\lambda_0}$. Together with \eqref{eq:abstract-exponential}, this gives
\begin{equation}\label{eq:entropy-duality}
 \E_{\nu_p}[-R]\le\lambda_0(t)(\Hr(\nu_p\mid\Gamma)+C_0).
\end{equation}
On the other hand, differentiating \eqref{eq:exponent-ode} gives
\begin{equation}\label{eq:pprime}
 p'=-\lambda_0(t)p(p-1).
\end{equation}
Thus, after applying \eqref{eq:entropy-duality} in \eqref{eq:log-power}, the two coefficients of $\Hr(\nu_p\mid\Gamma)$ cancel. Discarding the nonpositive Dirichlet term, we obtain
\begin{equation}\label{eq:cancelled}
 \frac\dd{\dd t}\log\norm q_{L^{p(t)}(\Gamma_t)}
 \le C_0\lambda_0(t)\left(1-\frac1{p(t)}\right)
 =-C_0\frac\dd{\dd t}\left(1-\frac1{p(t)}\right).
\end{equation}
Integration proves the assertion. When $\lambda_0=0$, both the defect and exponent derivative vanish. The same computation holds almost everywhere for an absolutely continuous exponent.
\end{proof}

\parheading{The cost of discarding diffusion.}
At a fixed exponent $p>1$, the entropy contribution can instead be absorbed into diffusion when
\begin{equation}\label{eq:fixed-exponent-absorption}
 \lambda_0\Ent_\Gamma(q^p)
 \le \frac{4}{\beta p}\int|\nabla(q^{p/2})|^2\dd\Gamma.
\end{equation}
A logarithmic Sobolev inequality with a sufficiently small constant would provide this estimate, but the static density moments do not supply it. Positive definiteness of the interaction operator is not convexity of the particle Hamiltonian: the tangential Hessian of a repulsive Riesz potential is negative and singular. The varying exponent cancels the entropy directly. Its price is the monotone loss of integrability, with $p(t)\downarrow1$ if $A(t)\to\infty$.

\begin{lemma}[Comparison of Gibbs preparations]\label{lem:initial-temperature}
Let $b_0,b_1\ge0$ and $p>1$. Then
\begin{equation}\label{eq:initial-temperature-moment}
 \int\left(\frac{\dd\Q_{N,b_0}(\mu)}{\dd\Q_{N,b_1}(\mu)}\right)^p
 \dd\Q_{N,b_1}(\mu)
 =\frac{\K_{N,b_1}(\mu)^{p-1}
 \K_{N,pb_0-(p-1)b_1}(\mu)}{\K_{N,b_0}(\mu)^p},
\end{equation}
whenever the right-hand side is finite. The same identity holds for the heat truncations. If the static estimates of \cref{prop:static-inputs} hold and
\begin{equation}\label{eq:initial-temperature-positive-coupling}
 pb_0-(p-1)b_1\ge0,
\end{equation}
then the left-hand side of \eqref{eq:initial-temperature-moment} is bounded uniformly in $N$ and in the heat cutoff. For fixed $p$, the bound is uniform when $b_0$ and $b_1$ range over a compact interval and \eqref{eq:initial-temperature-positive-coupling} holds, with a constant depending only on that interval, $p$, and the static background class.
\end{lemma}
\begin{proof}
Substitution of the two Gibbs densities gives \eqref{eq:initial-temperature-moment} directly. Under \eqref{eq:initial-temperature-positive-coupling}, every partition function on its right-hand side has a nonnegative fixed coupling. The upper bounds in \eqref{eq:static-partition}, together with the Jensen lower bound for $\K_{N,b_0}(\mu)$, prove the uniform estimate. The cutoff statement is identical.
\end{proof}

\begin{cor}[Full-particle second moment]\label{cor:full-moment}
Assume the static inputs hold uniformly for a family of consistent approximations and the scalar exponential estimate holds with a common $\lambda_0$. Suppose that, for some finite $p_0>2$,
\begin{equation}\label{eq:initial-relative-moment}
 \sup_{N,\varepsilon}\norm{q_{N,0}^{\varepsilon}}_{L^{p_0}(\Gamma_{N,0}^{\varepsilon})}<\infty,
 \qquad
 \left(1-\frac1{p_0}\right)e^{-A(T)}>\frac12,
 \qquad A(T)=\int_0^T\lambda_0(t)\dd t.
\end{equation}
Then
\begin{equation}\label{eq:full-endpoint-moment}
 \sup_{N,\varepsilon,\,0\le t\le T}
 \norm{f_{N,t}^{\varepsilon}}_{L^2((\mu_t^{\varepsilon})^{\otimes N})}^2<\infty.
\end{equation}
The assertion is time-uniform when \eqref{eq:initial-relative-moment} holds with $T=\infty$.
\end{cor}
\begin{proof}
We suppress the approximation parameter and set $p_*=p(T)>2$, with $p(t)$ defined by \eqref{eq:exponent-ode}. Then $p(t)\ge p_*$ on $[0,T]$. Since $\Gamma_t$ is a probability measure, \cref{prop:decreasing} and the $L^{p(t)}(\Gamma_t)$ to $L^{p_*}(\Gamma_t)$ inclusion bound $\norm{q_t}_{L^{p_*}(\Gamma_t)}$ uniformly over the family. Using $f_{N,t}=q_{N,t}G_{N,t}$ and H\"older with conjugate exponents $p_*/2$ and $p_*/(p_*-2)$, we obtain
\begin{equation}
 \int f_{N,t}^2\dd\mu_t^{\otimes N} =\int q_{N,t}^2G_{N,t}\dd\Gamma_{N,t} \le\norm{q_{N,t}}_{L^{p_*}(\Gamma_{N,t})}^2 \left(\int G_{N,t}^{\frac{2(p_*-1)}{p_*-2}}\dd\mu_t^{\otimes N}\right)^{\frac{p_*-2}{p_*}}. \label{eq:full-holder}
\end{equation}
The last factor is uniformly bounded by \eqref{eq:static-powers}; its exponent is finite because $p_*>2$. This proves \eqref{eq:full-endpoint-moment}.

In the matched case $P_{N,0}=\Gamma_{N,0}$, one may choose $p_0$ after fixing any $A(T)<\log2$. For example, taking $p_0=8$ and $A(T)\le\log(7/6)$ gives $p(t)\ge4$ and $\log\norm{q_t}_{L^4(\Gamma_t)}\le C_0/8$. Therefore,
\begin{equation}\label{eq:explicit-full}
 \int f_{N,t}^2\dd\mu_t^{\otimes N}
       \le e^{C_0/4}\left(\int G_{N,t}^3\dd\mu_t^{\otimes N}\right)^{1/2}.
\end{equation}
The constants do not depend on time when the corresponding assumptions hold on the whole half-line.
\end{proof}

For the constant-coupling dynamics, take $b_0=\beta_0$ and $b_1=\beta$ in \cref{lem:initial-temperature}. Writing $r=1-1/p_0$, the following strict inequalities are sufficient for \eqref{eq:initial-relative-moment}:
\begin{equation}\label{eq:temperature-exponent-choice}
 \frac12e^{A(T)}<r<\min\left\{1,\frac{\beta_0}{\beta}\right\}.
\end{equation}
Such an $r$ exists exactly under the sufficient condition in \eqref{eq:expenditure}. In the matched case, the resulting threshold is $\log2$; it records the need to leave an exponent above two in \eqref{eq:full-holder}. Neither this value nor its different-temperature counterpart is claimed sharp. One cannot restart the calculation with relative density one at a later time without replacing the actual particle law by the instantaneous Gibbs law. The periodic counterexample in \cref{thm:periodic-counterexample} shows that some restriction is genuinely needed if the target is a uniform full-particle second moment, but it does not identify the optimal hypothesis or a sharp numerical threshold.

\par\endgroup

\begingroup

\section{Weighted two-point, elliptic, and parabolic estimates}\label{sec:weighted}

The first-marginal argument in \cref{sec:marginals} uses three estimates. The two-point commutator is bounded by a mixed norm involving both the gradient and Hessian of the test function. Weighted elliptic duality controls that norm in the initial preparation argument, while backward smoothing controls it along the mean-field trajectory. The latter two estimates follow from the one-body Bochner calculus and nonautonomous Kolmogorov theory; we give the proofs in the weighted form needed here.

\subsection{A weighted two-point commutator}
For a symmetric $\eta\in L^2(\mu^{\otimes2})$ and $0\le\varepsilon\le1$, define
\begin{equation}\label{eq:two-point-form}
 \langle S_\eta^\varepsilon,\psi\rangle
 :=-\frac12\iint\eta(x,y)(\nabla\psi(x)-\nabla\psi(y))\cdot\nabla\g_\varepsilon(x-y)
          \dd\mu(x)\dd\mu(y),\qquad S_\eta:=S_\eta^0.
\end{equation}
For $\varepsilon>0$, symmetry of $\eta$ and oddness of $\nabla\g_\varepsilon$ identify this functional with the distribution
\begin{equation}\label{eq:two-point-divergence}
 S_\eta^\varepsilon
 =\diver_x\left(\mu(x)\int\nabla\g_\varepsilon(x-y)\eta(x,y)\dd\mu(y)\right).
\end{equation}
We do not assume that $\int\eta(x,y)\dd\mu(y)=0$ for $\mu$-almost every $x$; only symmetry of $\eta$ is used below.

\begin{lemma}[Weighted two-point commutator estimate]\label{lem:weighted-commutator}
Let $\mu=e^{-U}$ be a probability density satisfying
\begin{equation}\label{eq:weighted-density-hypotheses}
 \norm\mu_\infty\le M,\qquad \nabla^2U\ge-K\Id,\qquad K\ge0.
\end{equation}
Fix $R>0$ and $0\le\alpha\le1$, and assume
\begin{align}
 J_{\alpha,R}^2&:=\sup_{0\le\varepsilon\le1}\int_{|z|\le R}|z|^{2\alpha}
                              |\nabla\g_\varepsilon(z)|^2\dd z<\infty,
 \label{eq:local-force-norms}\\
 B_R&:=\sup_{0\le\varepsilon\le1}\mathop{\mathrm{ess\,sup}}_{|z|>R}
                                     |\nabla\g_\varepsilon(z)|<\infty.
 \label{eq:far-force-norm}
\end{align}
Then, for every $0\le\varepsilon\le1$,
\begin{align}
 |\langle S_\eta^\varepsilon,\psi\rangle|
 \le\norm\eta_{L^2(\mu^{\otimes2})}\Big(&C_\alpha\sqrt M\,e^{\alpha KR^2/16}J_{\alpha,R}
       \norm{\nabla\psi}_{L^2(\mu)}^{1-\alpha}
       \norm{\nabla^2\psi}_{L^2(\mu)}^\alpha\notag\\
       &+B_R\norm{\nabla\psi}_{L^2(\mu)}\Big),
       \qquad \forall\psi\in C_c^\infty(\R^\d).
 \label{eq:weighted-commutator}
\end{align}
Consequently, $S_\eta^\varepsilon$ extends uniquely to a bounded linear functional on the closure of $C_c^\infty(\R^\d)$ in $H^2(\mu)$, with the same estimate and constants independent of $\varepsilon$. Here, $H^2(\mu)$ carries the norm obtained by summing the weighted $L^2$ norms of derivatives of orders zero, one, and two.
\end{lemma}
\begin{proof}
The lower Hessian bound on $U$ gives, for $w=(1-\theta)x+\theta y$,
\begin{equation}\label{eq:segment-weight}
 \mu(x)\mu(y)\le
 e^{K\theta(1-\theta)|x-y|^2/2}\mu(w)\mu(x)^\theta\mu(y)^{1-\theta}
 \le M e^{K|x-y|^2/8}\mu(w).
\end{equation}
We estimate the gradient difference at each displacement $z$. Splitting its square into the two gradient terms and bounding the unused density by $M$, we obtain
\begin{equation}\label{eq:segment-crude}
 \int\mu(x)\mu(x-z)|\nabla\psi(x)-\nabla\psi(x-z)|^2\dd x
 \le4M\norm{\nabla\psi}_{L^2(\mu)}^2.
\end{equation}
Alternatively, integrate the Hessian along the segment from $x$ to $x-z$, apply Jensen's inequality in $\theta$, and use \eqref{eq:segment-weight}. The change of variables $x\mapsto x-\theta z$ then gives
\begin{align}
 &\int\mu(x)\mu(x-z)|\nabla\psi(x)-\nabla\psi(x-z)|^2\dd x\notag\\
 &\quad\le |z|^2\int_0^1\int\mu(x)\mu(x-z)
                    |\nabla^2\psi(x-\theta z)|^2\dd x\dd\theta
 \le M e^{K|z|^2/8}|z|^2\norm{\nabla^2\psi}_{L^2(\mu)}^2.
 \label{eq:segment-hessian}
\end{align}
The left-hand sides of \eqref{eq:segment-crude} and \eqref{eq:segment-hessian} are the same. Raising their right-hand sides to the powers $1-\alpha$ and $\alpha$, respectively, and integrating against $|\nabla\g_\varepsilon(z)|^2$ on $|z|\le R$, yields
\begin{align}
 &\iint_{|x-y|\le R}|\nabla\psi(x)-\nabla\psi(y)|^2
                  |\nabla\g_\varepsilon(x-y)|^2\dd\mu(x)\dd\mu(y)\notag\\
 &\quad\le4^{1-\alpha}M e^{\alpha KR^2/8}J_{\alpha,R}^2
       \norm{\nabla\psi}_{L^2(\mu)}^{2(1-\alpha)}
       \norm{\nabla^2\psi}_{L^2(\mu)}^{2\alpha}.
 \label{eq:near-interpolation}
\end{align}
On the complementary region $|x-y|>R$, we use $|\nabla\g_\varepsilon(x-y)|\le B_R$ and $\mu(\R^\d)=1$:
\begin{align}
 &\iint_{|x-y|>R}|\nabla\psi(x)-\nabla\psi(y)|^2
                 |\nabla\g_\varepsilon(x-y)|^2\dd\mu(x)\dd\mu(y)\notag\\
 &\quad\le 2B_R^2\iint\bigl(|\nabla\psi(x)|^2+|\nabla\psi(y)|^2\bigr)
                         \dd\mu(x)\dd\mu(y)
 =4B_R^2\norm{\nabla\psi}_{L^2(\mu)}^2.
 \label{eq:weighted-far-integral}
\end{align}
Cauchy--Schwarz in \eqref{eq:two-point-form}, followed by \eqref{eq:near-interpolation} and \eqref{eq:weighted-far-integral}, proves \eqref{eq:weighted-commutator}. For $\alpha=0$ or $\alpha=1$, use \eqref{eq:segment-crude} or \eqref{eq:segment-hessian} directly. The extension follows by density in the indicated completion.
\end{proof}

By \cref{lem:heat-force-estimates}, the log/Riesz kernels \eqref{eq:kernel} verify the force assumptions with $R=1$ whenever
\begin{equation}\label{eq:alpha-choice}
 \max\{0,\s+1-\d/2\}<\alpha<1.
\end{equation}
This interval is nonempty throughout $0\le\s<\d/2$. Choosing $\alpha<1$ will make the time singularity in the application to \cref{prop:finite-mark} integrable.

\subsection{Weighted elliptic duality}
Let
\begin{equation}\label{eq:weighted-elliptic-operator}
 \Aop_\mu:=-\mu^{-1}\diver(\mu\nabla)
          =-\Delta+\nabla U\cdot\nabla
\end{equation}
denote the nonnegative operator associated with the closed form $\int|\nabla\psi|^2\dd\mu$ on $L^2(\mu)$.

\begin{lemma}[One-body elliptic estimate]\label{lem:weighted-elliptic}
Suppose $\mu=e^{-U}$ is a smooth probability density satisfying
\begin{equation}\label{eq:elliptic-density-hypotheses}
 \norm{\nabla^2U}_{L^\infty}<\infty,\qquad
 \nabla^2U\ge-K\Id,\qquad K\ge0.
\end{equation}
Assume also the Poincar\'e inequality
\begin{equation}\label{eq:one-body-PI}
 \int|\psi-\mu(\psi)|^2\dd\mu\le C_P\int|\nabla\psi|^2\dd\mu.
\end{equation}
For each $\zeta\in L^2_0(\mu)$, the centered solution $\Aop_\mu\psi=\zeta$ satisfies
\begin{equation}\label{eq:weighted-elliptic-estimate}
 \norm{\nabla\psi}_{L^2(\mu)}\le\sqrt{C_P}\norm\zeta_{L^2(\mu)},
 \qquad
 \norm{\nabla^2\psi}_{L^2(\mu)}\le\sqrt{1+KC_P}\norm\zeta_{L^2(\mu)}.
\end{equation}
\end{lemma}
\begin{proof}
The Poincar\'e inequality makes $\int\nabla\psi\cdot\nabla\varphi\dd\mu$ a Hilbert-space inner product on the centered form domain, and $\varphi\mapsto\int\zeta\varphi\dd\mu$ is continuous for the resulting norm. The Riesz--Fr\'echet representation theorem (see, e.g., \cite[Theorem~5.5]{Brezis}) therefore gives the centered weak solution. Testing its equation with $\psi$ proves the first estimate in \eqref{eq:weighted-elliptic-estimate}.

The integrated Bochner identity for the weighted diffusion operator is
\begin{equation}\label{eq:weighted-bochner}
 \int|\Aop_\mu\psi|^2\dd\mu
 =\int|\nabla^2\psi|^2\dd\mu+
   \int\nabla^2U(\nabla\psi,\nabla\psi)\dd\mu.
\end{equation}
This is the standard $\Gamma_2$ identity; see \cite[Chapters~1 and~4]{BGL}. We justify its extension to the weak solution. Local elliptic regularity gives $\psi\in H^2_{\mathrm{loc}}$. Choose $\chi\in C_c^\infty$ equal to one on $B_1$ and zero outside $B_2$, and set $\chi_R(x):=\chi(x/R)$. Then
\begin{equation}\label{eq:elliptic-cutoff-identity}
 \Aop_\mu(\chi_R\psi)
 =\chi_R\zeta-2\nabla\chi_R\cdot\nabla\psi
   +\psi\Aop_\mu\chi_R.
\end{equation}
The bounded Hessian of $U$ gives at most linear growth of $\nabla U$. Thus, $\Aop_\mu\chi_R$ is uniformly bounded for $R\ge1$ and supported in $B_{2R}\setminus B_R$. Since $\psi,\nabla\psi\in L^2(\mu)$, we obtain
\begin{equation}\label{eq:elliptic-graph-cutoff}
 \norm{\chi_R\psi-\psi}_{L^2(\mu)}
 +\norm{\Aop_\mu(\chi_R\psi)-\zeta}_{L^2(\mu)}
 \longrightarrow0\qquad(R\to\infty).
\end{equation}
For each fixed $R$, mollify the compactly supported $H^2$ function $\chi_R\psi$, extended by zero to $\R^\d$. On a compact neighborhood of its support, the weight is bounded above and below and $\nabla U$ is bounded. Convergence of these mollifications in unweighted $H^2$ therefore implies convergence in the graph norm in \eqref{eq:elliptic-graph-cutoff}. Choosing $R=n$ and a sufficiently small mollification scale gives $\psi_n\in C_c^\infty(\R^\d)$ such that
\begin{equation}\label{eq:elliptic-graph-mollification}
 \norm{\psi_n-\psi}_{L^2(\mu)}
 +\norm{\Aop_\mu\psi_n-\zeta}_{L^2(\mu)}\longrightarrow0.
\end{equation}
The Dirichlet-form identity gives convergence of $\nabla\psi_n$ in $L^2(\mu)$. Applying \eqref{eq:weighted-bochner} to $\psi_n-\psi_m$ gives
\begin{equation}\label{eq:elliptic-hessian-Cauchy}
 \norm{\nabla^2(\psi_n-\psi_m)}_{L^2(\mu)}^2
 \le\norm{\Aop_\mu(\psi_n-\psi_m)}_{L^2(\mu)}^2
       +K\norm{\nabla(\psi_n-\psi_m)}_{L^2(\mu)}^2\longrightarrow0.
\end{equation}
Local positivity of $\mu$ identifies the limit as the distributional Hessian of $\psi$. The upper bound on $\nabla^2U$ then permits passage to the limit in \eqref{eq:weighted-bochner}. Using its lower bound and the first estimate, we conclude that
\begin{equation}
 \norm{\nabla^2\psi}_{L^2(\mu)}^2
 \le \norm\zeta_{L^2(\mu)}^2+K\norm{\nabla\psi}_{L^2(\mu)}^2
 \le(1+KC_P)\norm\zeta_{L^2(\mu)}^2.
\end{equation}
\end{proof}

Combining the two lemmas gives the estimate used in \cref{prop:initial-mark}:
\begin{align}
 \sup_{0\le\varepsilon\le1}|\langle S_\eta^\varepsilon,\Aop_\mu^{-1}\zeta\rangle|
 &\le\Big(C_\alpha\sqrt M e^{\alpha KR^2/16}J_{\alpha,R}
          C_P^{(1-\alpha)/2}(1+KC_P)^{\alpha/2}
           +B_R\sqrt{C_P}\Big)\notag\\
 &\qquad\times\norm\eta_{L^2(\mu^{\otimes2})}\norm\zeta_{L^2(\mu)},
 \qquad \zeta\in L^2_0(\mu).
 \label{eq:weighted-elliptic-source}
\end{align}

\subsection{Backward smoothing in the moving weight}
Let $\mu_t$ be a smooth positive probability density satisfying
\begin{equation}\label{eq:backward-generator-density}
 \mathcal L_t:=\beta^{-1}\Delta+B_t\cdot\nabla,\qquad
 \partial_t\mu_t=\mathcal L_t^*\mu_t.
\end{equation}
Under the spatial Lipschitz and linear-growth bounds below, the stochastic differential equation with drift $B_t$ and diffusion coefficient $\sqrt{2/\beta}$ has a unique nonexplosive strong solution from each point. This is the standard existence and uniqueness theorem (see, e.g., \cite[Theorem~5.2.1]{Oksendal}). Its transition operators are
\begin{equation}\label{eq:backward-propagator-definition}
 P_{r,t}\psi(x):=\E[\psi(x_t^{r,x})],\qquad
 \dd x_v^{r,x}=B_v(x_v^{r,x})\dd v+\sqrt{2/\beta}\dd W_v,
 \quad x_r^{r,x}=x.
\end{equation}
The Markov property, also for time-dependent coefficients, follows by restarting the solution with independent Brownian increments; see \cite[Section~7.1]{Oksendal}. Thus, $P_{r,t}=P_{r,v}P_{v,t}$ for $r\le v\le t$, and uniqueness of the forward equation gives $\mu_rP_{r,t}=\mu_t$. In \cref{prop:finite-mark}, the drift is $B_t=-\nabla V-c(t)\nabla h_t$.

\begin{lemma}[Weighted backward derivative estimates]\label{lem:backward-smoothing}
Assume that
\begin{equation}\label{eq:backward-regularity-growth}
 B\in C([0,T];C_{\mathrm{loc}}^2(\R^\d;\R^\d)),\qquad
 \sup_{t\in[0,T]}\sup_{x\in\R^\d}\frac{|B_t(x)|}{1+|x|}<\infty,
\end{equation}
where continuity into $C_{\mathrm{loc}}^2$ means locally uniform continuity of the spatial derivatives through order two. Suppose also that
\begin{equation}\label{eq:backward-drift-hypotheses}
 \sup_{t\in[0,T]}\norm{\nabla B_t}_{L^\infty}<\infty,\qquad
 \sup_{t\in[0,T]}\norm{\nabla^2B_t}_{L^\infty}<\infty.
\end{equation}
For $0\le r<t\le T$ and $\psi\in L^2(\mu_t)$,
\begin{align}
 \norm{P_{r,t}\psi}_{L^2(\mu_r)}&\le\norm\psi_{L^2(\mu_t)},
 \label{eq:backward-smoothing}\\
 \norm{\nabla P_{r,t}\psi}_{L^2(\mu_r)}&\le C_T(t-r)^{-1/2}\norm\psi_{L^2(\mu_t)},
 \label{eq:backward-gradient}\\
 \norm{\nabla^2P_{r,t}\psi}_{L^2(\mu_r)}&\le C_T(t-r)^{-1}\norm\psi_{L^2(\mu_t)}.
 \label{eq:backward-hessian}
\end{align}
The constants depend only on $T,\beta,\d$ and the two bounds in \eqref{eq:backward-drift-hypotheses}.
\end{lemma}
\begin{proof}
Jensen's inequality and $\mu_rP_{r,t}=\mu_t$ give
\begin{equation}\label{eq:backward-contraction-calculation}
 \int|P_{r,t}\psi|^2\dd\mu_r
 \le\int P_{r,t}(|\psi|^2)\dd\mu_r
 =\int|\psi|^2\dd\mu_t.
\end{equation}
For the derivatives, we use the Bernstein--Bochner calculation for the nonautonomous propagator; see \cite{KLL,AngiuliLorenzi,Lunardi} for the evolution-family and derivative estimates, and \cite{BGL} for the diffusion calculus.

First let $\psi\in C_c^\infty(\R^\d)$ and set $F_r:=P_{r,t}\psi$. Its backward equation and \eqref{eq:backward-generator-density} give
\begin{equation}\label{eq:backward-differentiation-rule}
 \partial_rF_r=-\mathcal L_rF_r,\qquad
 \frac{\dd}{\dd r}\int a_r\dd\mu_r
 =\int(\partial_ra_r+\mathcal L_ra_r)\dd\mu_r.
\end{equation}
The second identity is initially applied with a smooth spatial cutoff. Set
\begin{equation}\label{eq:backward-energy-definitions}
 E(r):=\int F_r^2\dd\mu_r,\qquad
 G(r):=\int|\nabla F_r|^2\dd\mu_r,\qquad
 H(r):=\int|\nabla^2F_r|^2\dd\mu_r.
\end{equation}
The product rule for $\mathcal L_r$, applied to $F_r^2$, $|\nabla F_r|^2$, and $|\nabla^2F_r|^2$, yields
\begin{equation}\label{eq:backward-EG}
 E'(r)=\frac2\beta G(r),
\end{equation}
\begin{equation}
 G'(r) =\frac2\beta H(r) -2\sum_{i,j=1}^{\d}\int\partial_iF_r\,\partial_iB_r^j\,\partial_jF_r\dd\mu_r \ge\frac2\beta H(r)-CG(r), \label{eq:backward-G-calculation}
\end{equation}
and
\begin{align}
 H'(r)
 &=\frac2\beta\int|\nabla^3F_r|^2\dd\mu_r
  -4\sum_{i,j,\ell=1}^{\d}\int
       \partial_{ij}F_r\,\partial_iB_r^\ell\,\partial_{j\ell}F_r\dd\mu_r\notag\\
 &\quad-2\sum_{i,j,\ell=1}^{\d}\int
       \partial_{ij}F_r\,\partial_{ij}B_r^\ell\,\partial_\ell F_r\dd\mu_r\notag\\
 &\ge-C\bigl(H(r)+G(r)\bigr).
 \label{eq:backward-H-calculation}
\end{align}
The first two derivatives of $B_r$ bound the drift terms, and the nonnegative third-derivative integral in \eqref{eq:backward-H-calculation} has been discarded.

We justify the integrated equality in \eqref{eq:backward-EG} and the integrated inequalities in \eqref{eq:backward-G-calculation} and \eqref{eq:backward-H-calculation} without a decay assumption on $\mu_r$. For smooth compactly supported $\psi$, the equations for the first two spatial derivatives of the stochastic flow give bounded $F_r$, $\nabla F_r$, and $\nabla^2F_r$ on $0\le r\le t$. Apply It\^o's formula along the diffusion stopped on balls. Nonexplosion and these bounds permit passage to the limit in the terminal evaluations and drift terms; Fatou's lemma applies to the nonnegative derivative integrals.

The computation may first be made after spatial mollification of $B$. To pass to the regularity in \eqref{eq:backward-regularity-growth}, fix $r<t$, start each mollified diffusion with law $\mu_r$, and use its own law at each intermediate time in the preceding calculation. The first two drift-derivative bounds are preserved by mollification. The stochastic flows and their first two derivatives converge after localization, while the evaluations of $\psi$ and its derivatives remain bounded uniformly in the mollification parameter. Dominated convergence therefore gives the stated integrated relations with the original laws $\mu_v$, $r\le v\le t$.

Since \eqref{eq:backward-G-calculation} implies $G'(r)+CG(r)\ge0$, the function $e^{Cr}G(r)$ is nondecreasing. Hence $G(r)\le e^{C(v-r)}G(v)$ for $r\le v\le t$. Setting $h:=t-r$ and integrating \eqref{eq:backward-EG}, we obtain
\begin{align}
 \int_r^tG(v)\dd v&\le\frac\beta2E(t),
 \label{eq:backward-gradient-integral}\\
 G(r)&\le\frac{e^{CT}}h\int_r^tG(v)\dd v\le C_Th^{-1}E(t).
 \label{eq:backward-gradient-average}
\end{align}
Apply the latter estimate also at $r+h/2$. Integrating \eqref{eq:backward-G-calculation} over $[r,r+h/2]$ gives
\begin{equation}
 \int_r^{r+h/2}H(v)\dd v \le\frac\beta2\left(G(r+h/2)-G(r) +C\int_r^{r+h/2}G(v)\dd v\right) \le C_Th^{-1}E(t). \label{eq:backward-hessian-average}
\end{equation}
By \eqref{eq:backward-H-calculation}, for $r\le v\le r+h/2$,
\begin{equation}\label{eq:backward-hessian-propagation}
 H(r)\le e^{C(v-r)}H(v)+C\int_r^v e^{C(a-r)}G(a)\dd a.
\end{equation}
Averaging in $v$ and using \eqref{eq:backward-gradient-integral} and \eqref{eq:backward-hessian-average}, we conclude that
\begin{equation}\label{eq:backward-hessian-final}
 H(r)\le\frac{C_T}h\int_r^{r+h/2}H(v)\dd v
                +C_T\int_r^tG(v)\dd v
       \le C_Th^{-2}E(t),
\end{equation}
where $h\le T$ was used in the last inequality. Taking square roots proves the derivative estimates for $\psi\in C_c^\infty$.

For a general $\psi\in L^2(\mu_t)$, choose $\psi_n\in C_c^\infty(\R^\d)$ with $\psi_n\to\psi$ in $L^2(\mu_t)$. Applying \eqref{eq:backward-smoothing}--\eqref{eq:backward-hessian} to $\psi_n-\psi_m$ gives, for $j=0,1,2$,
\begin{equation}\label{eq:backward-test-approximation}
 \norm{\nabla^jP_{r,t}(\psi_n-\psi_m)}_{L^2(\mu_r)}
 \le C_T(t-r)^{-j/2}\norm{\psi_n-\psi_m}_{L^2(\mu_t)}\longrightarrow0.
\end{equation}
For $j=0$, the limit is $P_{r,t}\psi$ by \eqref{eq:backward-contraction-calculation}. Positivity and local smoothness of $\mu_r$ identify the other two limits as its weak derivatives. Passing to the limit proves the lemma.
\end{proof}

\par\endgroup

\begingroup

\section{One-mark control and the global all-marginal enhancement}\label{sec:marginals}

The full second moment from \cref{sec:decreasing} controls the two-particle projection that appears as a source in the first-marginal equation \eqref{eq:first-mark-equation}. We first estimate the initial marginal by weighted elliptic duality, then propagate this estimate by the backward bounds of \cref{sec:weighted}. For the uniform-in-time result, we combine a homogeneous energy estimate with the finite-time source estimate on successive intervals. Finally, \cref{lem:product-projection} transfers the first-marginal estimate and the full second moment to every $1\le k\le N$.

We first work with the heat-truncated interaction. In this section, we suppress the cutoff from $\g$, the corresponding background, and the particle densities; the constants in the estimates are independent of the cutoff. The interaction coefficient is the function $c$ introduced in \cref{sec:static}. On a finite interval $[0,T]$, the common background assumptions are \eqref{eq:gaussian-envelope}, the static estimates of \cref{prop:static-inputs}, and
\begin{equation}\label{eq:mark-background-bounds}
 \sup_{0\le t\le T}\left(\norm{\nabla^2U_t}_{L^\infty}
  +\norm{\nabla B_t}_{L^\infty}+\norm{\nabla^2B_t}_{L^\infty}\right)<\infty,
 \qquad U_t:=-\log\mu_t.
\end{equation}
Here, $B_t=-\nabla V-c(t)\nabla h_t$ is the nonlinear diffusion drift, as in \eqref{eq:finite-drift} for $c\equiv1$, and $\partial_t\mu_t=\mathcal L_t^*\mu_t$ with the generator \eqref{eq:backward-generator-density}. The Poincar\'e inequality for $\mu_t$ is assumed where the elliptic and uniform-in-time estimates require it. The concrete verification of these hypotheses is in Section~\ref{sec:backgrounds} and Appendix~\ref{app:background}; the passage to the singular particle law is in \cref{sec:completion}.

We will use the following consequence of \cref{lem:heat-force-estimates}. If $\norm\mu_{L^\infty}\le M$, then
\begin{equation}\label{eq:force-schur-row}
 \sup_x\int|\nabla\g(x-y)|\dd\mu(y)
 \le M\norm{\nabla\g}_{L^1(B_1)}
       +\norm{\nabla\g}_{L^\infty(B_1^c)}\le C_M.
\end{equation}
Schur's test applied to the kernel $|\nabla\g(x-y)|$ gives
\begin{equation}\label{eq:force-schur}
 \norm{\nabla\g*(\mu a)}_{L^2(\mu)}\le C_M\norm a_{L^2(\mu)},\qquad
 \norm{\nabla h_\mu}_\infty\le C_M.
\end{equation}
Only local integrability of $\nabla\g$, supplied by $\s+1<\d$, is used here; square integrability of the force is not required.

\subsection{The initial marginal}
The next estimate uses the screened inverse of \cref{lem:screened-inverse} and does not require positive definiteness of the interaction.

\begin{prop}[Weighted initial one-mark estimate]\label{prop:initial-mark}
Let $\mu_0=e^{-U_0}$ satisfy the assumptions of \cref{lem:weighted-elliptic} and the Gaussian upper bound \eqref{eq:gaussian-envelope}. For every fixed $b\ge0$, the initial law $\Q_{N,b}(\mu_0)$ satisfies
\begin{equation}\label{eq:initial-mark}
 \norm{G_{N,b}^{(1)}-1}_{L^2(\mu_0)}\le\frac C N.
\end{equation}
The constant depends on $b$, the one-body Poincar\'e and Hessian bounds, the constants in \eqref{eq:gaussian-envelope}, the static second moment, and the inverse bound in \cref{lem:screened-inverse}. It is uniform when $b$ ranges over a fixed compact interval.
\end{prop}
\begin{proof}
For $b=0$, the initial law is the product law and there is nothing to prove. For $b>0$, suppress the initial time and the preparation parameter in the densities, and set
\begin{equation}\label{eq:initial-projections}
 a:=G_N^{(1)}-1,\qquad
 \eta(x,y):=G_N^{(2)}(x,y)-1-a(x)-a(y),\qquad
 M_0:=\sup_{N\ge2}\int G_N^2\dd\mu^{\otimes N}.
\end{equation}
The moment $M_0$ is finite by \eqref{eq:static-powers}. Both $a$ and $\eta$ have zero average in each of their coordinates. The orthogonal decomposition in \cref{lem:product-projection} gives
\begin{align}
 \norm\eta_{L^2(\mu^{\otimes2})}
 &\le\left(\frac{2(M_0-1)}{N(N-1)}\right)^{1/2}
 \le\frac{2\sqrt{M_0-1}}N,\label{eq:static-pair-projection}\\
 \norm a_{L^2(\mu)}&\le\sqrt{M_0-1}.\label{eq:static-one-projection}
\end{align}

Differentiating the Gibbs density before removing the cutoff, and then integrating the other $N-1$ coordinates, gives
\begin{equation}\label{eq:initial-score}
 \nabla a=\frac bN(1+a)\nabla h
 -b\frac{N-1}{N}\left(\nabla\g*(\mu a)
               +\int\nabla\g(x-y)\eta(x,y)\dd\mu(y)\right).
\end{equation}
In fact, the two-particle term in the differentiated marginal is the integral of $\nabla\g(x-y)$ against $\mu(y)[1+a(x)+a(y)+\eta(x,y)]$. Its first two terms are $(1+a(x))\nabla h(x)$, which leave precisely the coefficient $b/N$ in \eqref{eq:initial-score}.

Since $a$ is centered, $\nabla T_\mu a=\nabla\g*(\mu a)$. Hence the centered function
\begin{equation}\label{eq:initial-screened-variable}
 z:=\left(\Id+b\frac{N-1}{N}T_\mu\right)a
\end{equation}
satisfies the distributional equation
\begin{equation}\label{eq:screened-equation}
 -\diver(\mu\nabla z)=-\frac bN\diver(\mu(1+a)\nabla h)
                      +b\frac{N-1}{N}S_\eta.
\end{equation}
For $\zeta\in L^2_0(\mu)$, let $\psi$ be the centered solution of $\Aop_\mu\psi=\zeta$ from \cref{lem:weighted-elliptic}. Pairing \eqref{eq:screened-equation} with $\psi$ yields
\begin{equation}\label{eq:initial-duality-identity}
 \int z\zeta\dd\mu
 =\frac bN\int(1+a)\nabla h\cdot\nabla\psi\dd\mu
       +b\frac{N-1}{N}\langle S_\eta,\psi\rangle.
\end{equation}
The first term is controlled by \eqref{eq:force-schur}. For the second, apply \cref{lem:weighted-commutator} with $\alpha=1$. Thus,
\begin{align}
 \left|\int z\zeta\dd\mu\right|
 &\le \frac C N(1+\norm a_{L^2(\mu)})\norm{\nabla\psi}_{L^2(\mu)}
 +C\norm\eta_{L^2(\mu^{\otimes2})}
       \left(\norm{\nabla^2\psi}_{L^2(\mu)}+\norm{\nabla\psi}_{L^2(\mu)}\right)
 \notag\\
 &\le C\left(\frac{1+\norm a_{L^2(\mu)}}N
                      +\norm\eta_{L^2(\mu^{\otimes2})}\right)
                         \norm\zeta_{L^2(\mu)}.\label{eq:initial-duality}
\end{align}
The second inequality is \cref{lem:weighted-elliptic}. Its graph-core approximation justifies \eqref{eq:initial-duality-identity}: the left side is continuous in the operator graph norm, and the two sources are continuous in the weighted gradient and Hessian norms. Taking the supremum over centered $\zeta$ of unit norm, using \eqref{eq:static-pair-projection}--\eqref{eq:static-one-projection}, and applying \cref{lem:screened-inverse} gives \eqref{eq:initial-mark}.
\end{proof}

\subsection{The exact first-marginal equation}
Write $f_{N,t}^{(k)}:=\frac{\dd P_{N,t}^{(k)}}{\dd\mu_t^{\otimes k}}$, and set
\begin{align}
 a_{N,t}&:=f_{N,t}^{(1)}-1,\qquad
 w_{N,t}:=\mu_ta_{N,t},\qquad\theta_N:=\frac{N-1}{N},\notag\\
 \eta_{N,t}(x,y)&:=f_{N,t}^{(2)}(x,y)-1-a_{N,t}(x)-a_{N,t}(y).
 \label{eq:dynamic-projections}
\end{align}
The normalization and marginal identities imply
\begin{equation}\label{eq:dynamic-projection-centering}
 \int a_{N,t}\dd\mu_t=0,\qquad
 \int\eta_{N,t}(x,y)\dd\mu_t(y)=\int\eta_{N,t}(x,y)\dd\mu_t(x)=0
 \quad\text{a.e.}
\end{equation}
If the full moment is at most $M_*$, \cref{lem:product-projection} gives
\begin{equation}\label{eq:pair-projection}
 \norm{\eta_{N,t}}_{L^2(\mu_t^{\otimes2})}\le\frac{2\sqrt{M_*-1}}N.
\end{equation}

The two-particle law has Lebesgue density
\begin{equation}
 \label{eq:two-marginal-expansion} \frac{\dd P_{N,t}^{(2)}}{\dd x\dd y}(x,y) ={}\mu_t(x)\mu_t(y)+w_{N,t}(x)\mu_t(y)+\mu_t(x)w_{N,t}(y) +\mu_t(x)\mu_t(y)\eta_{N,t}(x,y).
\end{equation}
Since the one-body part of the particle drift is $B_t+c(t)\nabla h_t$, the first equation of the marginal hierarchy is
\begin{equation}
 \label{eq:first-marginal-hierarchy} \partial_t(\mu_t f_{N,t}^{(1)}) ={}\mathcal L_t^*(\mu_t f_{N,t}^{(1)}) +c(t)\diver\left(\theta_N\mu_t(x)\int\nabla\g(x-y)f_{N,t}^{(2)}(x,y)\dd\mu_t(y) -\mu_t f_{N,t}^{(1)}\nabla h_t\right),
\end{equation}
where $\mathcal L_t^*=\beta^{-1}\Delta-\diver(B_t\,\cdot)$. Substituting \eqref{eq:two-marginal-expansion} into \eqref{eq:first-marginal-hierarchy} and subtracting $\partial_t\mu_t=\mathcal L_t^*\mu_t$ gives
\begin{equation}
 \partial_tw_{N,t}={}\mathcal L_t^*w_{N,t} +\theta_Nc(t)\diver\bigl(\mu_t\nabla\g*w_{N,t}\bigr) -\frac{c(t)}N\diver(w_{N,t}\nabla h_t) -\frac{c(t)}N\diver(\mu_t\nabla h_t) +\theta_Nc(t)S_{\eta_{N,t}}. \label{eq:first-mark-equation}
\end{equation}
The terms proportional to $\nabla h_t$ in \eqref{eq:first-marginal-hierarchy} have the combined coefficient $c(t)(\theta_N-1)=-c(t)/N$. They therefore contribute $-\frac{c(t)}N\diver((\mu_t+w_{N,t})\nabla h_t)$ to \eqref{eq:first-mark-equation}. The part linear in $w_{N,t}$ remains in the homogeneous equation; the last two terms of \eqref{eq:first-mark-equation} are the prescribed sources.

\begin{prop}[Finite-time weighted one-mark estimate]\label{prop:finite-mark}
Suppose the bounds for $\mu_t$ and $B_t$ at the beginning of this section hold on $[0,T]$, and assume
\begin{align}
 \sup_{N\ge2}\sup_{0\le t\le T}
       \norm{f_{N,t}}_{L^2(\mu_t^{\otimes N})}^2&\le M_*<\infty,
              \label{eq:finite-mark-assumptions}\\
 \norm{a_{N,0}}_{L^2(\mu_0)}&\le\frac{C_0}N.
              \label{eq:finite-mark-initial}
\end{align}
Then
\begin{equation}\label{eq:finite-mark}
 \sup_{0\le t\le T}\norm{a_{N,t}}_{L^2(\mu_t)}\le\frac{C_T}N.
\end{equation}
More generally, for an interval length $L>0$, whenever the same bounds for $\mu_t$ and $B_t$ hold on $[r_0,r_0+L]$, the linear equation \eqref{eq:first-mark-equation} with a prescribed measurable symmetric two-particle source satisfying \eqref{eq:pair-projection} satisfies
\begin{equation}\label{eq:finite-mark-shifted}
 \sup_{r_0\le t\le r_0+L}\norm{a_{N,t}}_{L^2(\mu_t)}
 \le C_L\left(\norm{a_{N,r_0}}_{L^2(\mu_{r_0})}+\frac1N\right).
\end{equation}
The constant $C_L$ depends on the interval length $L$, the common bounds for $\mu_t$ and $B_t$, $\beta,c_*$, and $M_*$. The constant $C_T$ has the same dependencies with $L=T$ and also depends on $C_0$. In particular, the constant in \eqref{eq:finite-mark-shifted} is independent of $r_0$ when those bounds are uniform in time. The zero-initial-data response to the last two terms of \eqref{eq:first-mark-equation} is therefore bounded by $C_L/N$.
\end{prop}
\begin{proof}
Fix $t\in(0,T]$ and $\psi\in L^2(\mu_t)$. Testing with the backward propagator $P_{r,t}\psi$ from \cref{lem:backward-smoothing} gives
\begin{align}
 \int a_{N,t}\psi\dd\mu_t
 ={}&\int a_{N,0}P_{0,t}\psi\dd\mu_0\notag\\
 &-\int_0^t c(r)\int\left(\theta_N\nabla\g*(\mu_ra_{N,r})
                 -\frac{1+a_{N,r}}N\nabla h_r\right)
                   \cdot\nabla P_{r,t}\psi\dd\mu_r\dd r\notag\\
 &+\theta_N\int_0^tc(r)\langle S_{\eta_{N,r}},P_{r,t}\psi\rangle\dd r.
 \label{eq:first-mark-duhamel}
\end{align}
The contraction in \eqref{eq:backward-smoothing} bounds the first term by
$\norm{a_{N,0}}_{L^2(\mu_0)}\norm\psi_{L^2(\mu_t)}$.
For the flux term, \eqref{eq:force-schur} and \eqref{eq:backward-gradient} give
\begin{equation}\label{eq:first-mark-flux-bound}
 \left|\int\!\left(\theta_N\nabla\g*(\mu_ra_{N,r})-\frac{1+a_{N,r}}N\nabla h_r\right)\!\cdot\nabla P_{r,t}\psi\dd\mu_r\right|
 \le C_T\frac{\norm{a_{N,r}}_{L^2(\mu_r)}+N^{-1}}{\sqrt{t-r}}\norm\psi_{L^2(\mu_t)}.
\end{equation}
Choose $\alpha$ as in \eqref{eq:alpha-choice}, so in particular $\alpha<1$.
Combining \cref{lem:weighted-commutator,lem:backward-smoothing} with \eqref{eq:pair-projection} bounds the two-particle integrand by
\begin{equation}\label{eq:source-smoothing}
 \frac{C_T}N\left((t-r)^{-(1+\alpha)/2}+(t-r)^{-1/2}\right)\norm\psi_{L^2(\mu_t)}.
\end{equation}
Indeed, the mixed term in the commutator estimate is bounded by
$(t-r)^{-(1-\alpha)/2-\alpha}\norm\psi_{L^2(\mu_t)}$.
Both time singularities in \eqref{eq:source-smoothing} are integrable. After integrating the source terms and taking the supremum over $\norm\psi_{L^2(\mu_t)}\le1$, we obtain
\begin{equation}\label{eq:weighted-volterra}
 y(t)\le\frac{C_T}N+C_T\int_0^t(t-r)^{-1/2}y(r)\dd r,
             \qquad y(t):=\norm{a_{N,t}}_{L^2(\mu_t)}.
\end{equation}
The constant includes the initial bound $C_0/N$. Substituting this inequality into its integral term and using
\begin{equation}\label{eq:volterra-convolution}
 \int_s^t(t-r)^{-1/2}(r-s)^{-1/2}\dd r=\pi,
 \qquad 0\le s<t,
\end{equation}
gives
\begin{equation}\label{eq:volterra-gronwall}
 y(t)\le\frac{C_T(1+2C_T\sqrt T)}N
               +\pi C_T^2\int_0^t y(s)\dd s.
\end{equation}
Gronwall's inequality proves \eqref{eq:finite-mark}. Keeping the initial norm rather than replacing it by $C_0/N$, and translating the interval, proves \eqref{eq:finite-mark-shifted}.

The same argument constructs the auxiliary linear solution for a prescribed source. At fixed cutoff, the right side of \eqref{eq:first-mark-duhamel} defines a bounded functional of $\psi$ and hence an element of $L^2(\mu_t)$. On an interval of length $\delta$, the part depending linearly on $a$ has operator norm at most $2C_T\sqrt\delta$ in the norm $\sup_t\norm{a_t}_{L^2(\mu_t)}$. Choosing this number below one gives a contraction. Iterating on successive intervals constructs the solution, and the same argument proves uniqueness. The source integral is well defined by \eqref{eq:source-smoothing}; its mass is zero because $P_{r,t}1=1$ and both source distributions annihilate constants. No probability density on the full configuration space is required for this auxiliary solution. The testing of the actual cutoff law is justified in \cref{lem:approximation}.
\end{proof}

\subsection{A time-uniform closure by homogeneous energy and fixed intervals}
We now separate the homogeneous evolution in \eqref{eq:first-mark-equation} from its last two source terms. The quadratic form \eqref{eq:screened-energy}, associated with the positive operator $\Id+\beta\theta_Nc(t)T_{\mu_t}$, gives contraction of the homogeneous evolution over a sufficiently long fixed interval. The estimate \eqref{eq:finite-mark-shifted} bounds the response to the sources on the same interval. Iterating these two estimates will control the first marginal for all times, always retaining its actual value at the beginning of the next interval.

\begin{prop}[Uniform weighted one-mark estimate]\label{prop:uniform-mark}
Assume the common bounds for $\mu_t$ and $B_t$ hold uniformly for $t\ge0$, and assume
\begin{equation}\label{eq:uniform-mark-moment}
 \sup_{N\ge2}\sup_{t\ge0}\norm{f_{N,t}}_{L^2(\mu_t^{\otimes N})}^2\le M_*<\infty.
\end{equation}
Suppose the interaction coefficient satisfies
\begin{equation}\label{eq:uniform-mark-coupling}
 c\in C_{\mathrm{loc}}^1([0,\infty)),\qquad 0\le c(t)\le c_*,\qquad c'(t)\le0,
 \qquad t\ge0.
\end{equation}
Assume that each $\mu_t$ satisfies the Poincar\'e inequality with a common constant $C_P$:
\begin{equation}\label{eq:uniform-mark-poincare}
 \int\left|\varphi-\int\varphi\dd\mu_t\right|^2\dd\mu_t
 \le C_P\int|\nabla\varphi|^2\dd\mu_t,
 \qquad \varphi\in C_c^\infty(\R^\d),\quad t\ge0.
\end{equation}
Suppose also that
\begin{equation}\label{eq:integrable-velocity}
 \int_0^\infty\norm{\nabla u_t}_{L^\infty}\dd t<\infty
\end{equation}
and
\begin{equation}\label{eq:uniform-mark-initial}
 \norm{a_{N,0}}_{L^2(\mu_0)}\le C_0/N.
\end{equation}
Then, for the repulsive kernels \eqref{eq:kernel},
\begin{equation}\label{eq:uniform-mark}
 \sup_{t\ge0}\norm{a_{N,t}}_{L^2(\mu_t)}\le C/N.
\end{equation}
\end{prop}
\begin{proof}
For a signed density $w=\mu_ta$ of mass zero, set
\begin{equation}\label{eq:screened-energy}
 \mathcal Q_{N,t}(w):=\int a^2\dd\mu_t+
              \beta\theta_Nc(t)\iint\g(x-y)w(x)w(y)\dd x\dd y.
\end{equation}
On $L^2_0(\mu_t)$, the centered interaction operator $T_{\mu_t}$ is nonnegative and uniformly bounded, by \cref{lem:heat-potential-estimates}. Consequently,
\begin{equation}\label{eq:screened-comparison}
 \norm a_{L^2(\mu_t)}^2\le\mathcal Q_{N,t}(w)\le C_B\norm a_{L^2(\mu_t)}^2.
\end{equation}
Set
\begin{equation}\label{eq:homogeneous-screened-variable}
 z:=\left(\Id+\beta\theta_Nc(t)T_{\mu_t}\right)a.
\end{equation}
Then $\int z\dd\mu_t=0$ and
\begin{equation}\label{eq:screened-z-bound}
 \norm z_{L^2(\mu_t)}^2-\mathcal Q_{N,t}(w)
 =\beta\theta_Nc(t)\langle a,T_{\mu_t}a\rangle_{L^2(\mu_t)}
   +\beta^2\theta_N^2c(t)^2\norm{T_{\mu_t}a}_{L^2(\mu_t)}^2\ge0.
\end{equation}
This is where positivity of the repulsive interaction is used, in contrast to the screened-inverse estimate for the initial marginal.

Since $B_t=u_t+\beta^{-1}\nabla\log\mu_t$, the homogeneous part of \eqref{eq:first-mark-equation} is
\begin{equation}\label{eq:homogeneous-transport}
 \partial_tw+\diver(wu_t)=\beta^{-1}\diver(\mu_t\nabla z)
                         -\frac{c(t)}N\diver(w\nabla h_t).
\end{equation}
We compute first for smooth data. The continuity equation for $\mu_t$ gives
\begin{equation}\label{eq:homogeneous-local-derivative}
 \frac\dd{\dd t}\int\frac{w^2}{\mu_t}\dd x
 =2\int\frac w{\mu_t}\bigl(\partial_tw+\diver(wu_t)\bigr)\dd x.
\end{equation}
Symmetry of $\g$ and integration by parts likewise give
\begin{align}
 \frac\dd{\dd t}\iint\g(x-y)w(x)w(y)\dd x\dd y
 ={}&2\int(\g*w)\bigl(\partial_tw+\diver(wu_t)\bigr)\dd x\notag\\
 &+\iint(u_t(x)-u_t(y))\cdot\nabla\g(x-y)w(x)w(y)\dd x\dd y.
 \label{eq:homogeneous-interaction-derivative}
\end{align}
In the terms containing $\partial_tw+\diver(wu_t)$, the function $a+\beta\theta_Nc(t)\g*w$ differs from $z$ only by a spatial constant; the signed density in that pairing has mass zero. Substitution of \eqref{eq:homogeneous-transport} therefore yields
\begin{align}
 \frac\dd{\dd t}\mathcal Q_{N,t}(w)
 ={}&-\frac2\beta\int|\nabla z|^2\dd\mu_t\notag\\
 &+\beta\theta_Nc(t)\iint (u_t(x)-u_t(y))\cdot\nabla\g(x-y)w(x)w(y)\dd x\dd y\notag\\
 &+\beta\theta_Nc'(t)\iint\g(x-y)w(x)w(y)\dd x\dd y
 +\frac{2c(t)}N\int a\nabla h_t\cdot\nabla z\dd\mu_t.
 \label{eq:homogeneous-energy-identity}
\end{align}
The term containing $c'$ is nonpositive. For the transport term, \cref{lem:heat-force-estimates} implies
\begin{equation}\label{eq:homogeneous-commutator-row}
 \sup_x\int |x-y|\,|\nabla\g(x-y)|\dd\mu_t(y)
 \le C\left(M\int_{|z|\le1}|z|^{-\s}\dd z+1\right)\le C.
\end{equation}
Schur's test bounds the transport term in absolute value by $C\norm{\nabla u_t}_{L^\infty}\norm a_{L^2(\mu_t)}^2$. The last term satisfies
\begin{equation}\label{eq:homogeneous-young}
 \frac{2c(t)}N\left|\int a\nabla h_t\cdot\nabla z\dd\mu_t\right|
 \le\frac1\beta\norm{\nabla z}_{L^2(\mu_t)}^2
             +\frac D{N^2}\norm a_{L^2(\mu_t)}^2,
\end{equation}
where we used \eqref{eq:force-schur}. Since $z$ is centered, Poincar\'e's inequality and \eqref{eq:screened-z-bound} give
\begin{equation}\label{eq:homogeneous-poincare}
 \norm{\nabla z}_{L^2(\mu_t)}^2
 \ge C_P^{-1}\norm z_{L^2(\mu_t)}^2
 \ge C_P^{-1}\mathcal Q_{N,t}(w).
\end{equation}
Combining these inequalities proves
\begin{equation}\label{eq:homogeneous-contraction}
 \frac\dd{\dd t}\mathcal Q_{N,t}(w)
 \le\left(-\rho+\frac D{N^2}+C\norm{\nabla u_t}_{L^\infty}\right)\mathcal Q_{N,t}(w),
                      \qquad\rho:=(\beta C_P)^{-1}>0.
\end{equation}
At fixed cutoff, spatial localization justifies the calculation for smooth initial signed densities, as in \cref{lem:approximation}. Approximation of a centered $L^2(\mu_r)$ initial value by smooth compactly supported centered functions, followed by the finite-time linear construction in \cref{prop:finite-mark}, extends the integrated form of \eqref{eq:homogeneous-contraction} to such initial values. The bounds in \eqref{eq:screened-comparison} are uniform, so no cutoff-dependent constant enters this inequality.

Choose $N_0$ such that $D/N^2\le\rho/2$ for $N\ge N_0$. Integrating \eqref{eq:homogeneous-contraction} gives, for a homogeneous solution,
\begin{equation}\label{eq:homogeneous-propagator-bound}
 \mathcal Q_{N,t}(w_t)^{1/2}
 \le \exp\left(-\frac\rho4(t-r)+\frac C2\int_r^t\norm{\nabla u_s}_{L^\infty}\dd s\right)
                 \mathcal Q_{N,r}(w_r)^{1/2},\qquad t\ge r.
\end{equation}
By \eqref{eq:integrable-velocity}, we can choose a fixed length $L>0$ such that
\begin{equation}\label{eq:block-length-choice}
 \exp\left(-\frac{\rho L}4+\frac C2\int_0^\infty\norm{\nabla u_t}_{L^\infty}\dd t\right)\le\frac12.
\end{equation}
On $[jL,(j+1)L]$, write
\begin{equation}\label{eq:mark-interval-decomposition}
 w_{N,t}=w_{N,j,t}^{\mathrm{hom}}+w_{N,j,t}^{\mathrm{src}},\qquad
 w_{N,j,jL}^{\mathrm{hom}}=w_{N,jL},\qquad w_{N,j,jL}^{\mathrm{src}}=0.
\end{equation}
Here, $w^{\mathrm{hom}}$ solves \eqref{eq:homogeneous-transport}, whereas $w^{\mathrm{src}}$ solves \eqref{eq:first-mark-equation} with the prescribed $\eta_{N,t}$ and zero initial value. By \eqref{eq:homogeneous-propagator-bound}--\eqref{eq:block-length-choice}, and by \eqref{eq:finite-mark-shifted} together with \eqref{eq:screened-comparison}, respectively,
\begin{align}
 \mathcal Q_{N,(j+1)L}(w_{N,j,(j+1)L}^{\mathrm{hom}})^{1/2}
 &\le\frac12\mathcal Q_{N,jL}(w_{N,jL})^{1/2},
 \label{eq:mark-homogeneous-interval-bound}\\
 \sup_{jL\le t\le(j+1)L}\mathcal Q_{N,t}(w_{N,j,t}^{\mathrm{src}})^{1/2}
 &\le\frac{\sqrt{C_B}C_L}N.
 \label{eq:mark-source-interval-bound}
\end{align}
The triangle inequality for $\mathcal Q_{N,(j+1)L}^{1/2}$ therefore gives
\begin{equation}\label{eq:block-recurrence}
 Y_{j+1}\le\frac12Y_j+\frac{C_L\sqrt{C_B}}N,
                   \qquad Y_j:=\mathcal Q_{N,jL}(w_{N,jL})^{1/2}.
\end{equation}
The initial estimate and \eqref{eq:screened-comparison} imply $Y_0\le\sqrt{C_B}C_0/N$, and hence
\begin{equation}\label{eq:block-summation}
 Y_j\le 2^{-j}Y_0+\frac{C_L\sqrt{C_B}}N\sum_{i=0}^{j-1}2^{-i}
 \le\frac{\sqrt{C_B}(C_0+2C_L)}N.
\end{equation}
Applying \eqref{eq:finite-mark-shifted} within each interval proves \eqref{eq:uniform-mark} for $N\ge N_0$. For $2\le N<N_0$, the full moment and contraction of conditional expectation give directly
\begin{equation}\label{eq:uniform-mark-small-N}
 N\norm{a_{N,t}}_{L^2(\mu_t)}\le N_0\sqrt{M_*-1},\qquad t\ge0.
\end{equation}
This proves the assertion for every $N\ge2$.
\end{proof}

\subsection{Product-space projection}
We record as a lemma the orthogonal projection argument used in \cite{DGRchaos}. Besides giving the final all-marginal estimate, it justifies the two-particle bounds used earlier in this section.

\begin{lemma}[Product-space projection]\label{lem:product-projection}
Let $\mu$ be a probability measure and let $f_N\in L^2(\mu^{\otimes N})$ be a symmetric probability density, with $N\ge2$. Set $P_N:=f_N\mu^{\otimes N}$ and $f_N^{(k)}:=\frac{\dd P_N^{(k)}}{\dd\mu^{\otimes k}}$, so that
\begin{equation}\label{eq:projection-marginal-identities}
 f_N^{(k)}(x_1,\ldots,x_k)
 =\int f_N(x_1,\ldots,x_N)\dd\mu^{\otimes(N-k)}(x_{k+1},\ldots,x_N),
 \qquad 1\le k\le N,
\end{equation}
with $f_N^{(N)}=f_N$. Then, for every $1\le k\le N$,
\begin{equation}\label{eq:projection-transfer}
 \norm{f_N^{(k)}-1}_{L^2(\mu^{\otimes k})}^2
 \le k\norm{f_N^{(1)}-1}_{L^2(\mu)}^2
       +\frac{k(k-1)}{N(N-1)}\norm{f_N-1}_{L^2(\mu^{\otimes N})}^2.
\end{equation}
Moreover, for $\eta(x,y):=f_N^{(2)}(x,y)-f_N^{(1)}(x)-f_N^{(1)}(y)+1$,
\begin{equation}\label{eq:projection-pair-bound}
 \binom N2\norm\eta_{L^2(\mu^{\otimes2})}^2
 \le\norm{f_N-1}_{L^2(\mu^{\otimes N})}^2.
\end{equation}
\end{lemma}
\begin{proof}
We use the Hoeffding decomposition (see, e.g., \cite[Section~1.6]{Lee1990}). Set $f_N^{(0)}:=1$. For $1\le j\le N$, the symmetric Hoeffding kernel is
\begin{equation}\label{eq:hoeffding-definition}
 h_j(x_1,\ldots,x_j):=
 \sum_{A\subset[j]}(-1)^{j-|A|}f_N^{(|A|)}(x_A).
\end{equation}
Here, $x_A$ is the tuple with indices in $A$, in increasing order. The kernels satisfy
\begin{equation}\label{eq:hoeffding-centering}
 \int h_j(x_1,\ldots,x_j)\dd\mu(x_i)=0,
       \qquad 1\le i\le j,
\end{equation}
for almost every choice of the remaining coordinates, and
\begin{equation}\label{eq:hoeffding-expansion}
 f_N-1=\sum_{j=1}^N\ \sum_{\substack{A\subset[N]\\|A|=j}}h_j(x_A).
\end{equation}
For distinct nonempty index sets $A,B\subset[N]$, the orthogonality relation is
\begin{equation}\label{eq:hoeffding-orthogonality}
 \int h_{|A|}(x_A)h_{|B|}(x_B)\dd\mu^{\otimes N}=0.
\end{equation}
Together with the marginal identities \eqref{eq:projection-marginal-identities}, these properties give
\begin{equation}\label{eq:hoeffding-marginal-norms}
 \norm{f_N^{(k)}-1}_{L^2(\mu^{\otimes k})}^2
 =\sum_{j=1}^k\binom kj\norm{h_j}_{L^2(\mu^{\otimes j})}^2,
 \qquad 1\le k\le N.
\end{equation}
For $2\le j\le k\le N$,
\begin{equation}\label{eq:hoeffding-ratio}
 \frac{\binom kj}{\binom Nj}
 =\frac{k(k-1)}{N(N-1)}\prod_{i=2}^{j-1}\frac{k-i}{N-i}
 \le\frac{k(k-1)}{N(N-1)},
\end{equation}
where the product is one when $j=2$. Separating $j=1$, for which $h_1=f_N^{(1)}-1$, and summing the remaining terms proves \eqref{eq:projection-transfer}. When $k=1$, the assertion is an equality. Finally, $h_2=\eta$, so \eqref{eq:projection-pair-bound} follows by retaining the $j=2$ term in \eqref{eq:hoeffding-marginal-norms}.
\end{proof}

If the full second moment is at most $M_*$ and the one-mark error is at most $U/N$, \cref{lem:product-projection} gives
\begin{equation}\label{eq:all-k}
 \chi^2(P_{N,t}^{(k)}\mid\mu_t^{\otimes k})
 \le\frac{kU^2}{N^2}+(M_*-1)\frac{k(k-1)}{N(N-1)}.
\end{equation}
The standard information inequalities (see, e.g., \cite[(16), (139)--(141)]{NishiyamaSason2020}) give
\begin{align}
 \Hr(P\mid\nu)&\le\log(1+\chi^2(P\mid\nu))\le\chi^2(P\mid\nu),
                     \label{eq:entropy-TV}\\
 \TV(P,\nu)&\le\frac12\sqrt{\chi^2(P\mid\nu)}.
                     \label{eq:TV-from-square}
\end{align}
Together with \eqref{eq:all-k}, these inequalities give the claimed relative-entropy and total-variation rates.

\par\endgroup

\begingroup

\section{Confined backgrounds and completion of the proofs}\label{sec:confined-proofs}

We first verify the one-body hypotheses for perturbations of a confined thermal equilibrium measure. We then justify the regularized calculations and pass the estimates of \cref{sec:decreasing,sec:marginals} to the singular particle system. The background stability result is stated in \cref{sec:backgrounds} and proved in \cref{app:background}; the remainder of this section completes the proofs of \cref{thm:uniform,thm:finite,cor:local}.

\subsection{Verification of the confined backgrounds}\label{sec:backgrounds}

The perturbation in \eqref{eq:log-tilt} is measured through derivatives of $\log(\mu_t/\pi)$. This choice controls the logarithmic Hessian of $\mu_t$ throughout $\R^\d$, including the region where $\pi$ is small. Smallness of an unweighted Sobolev norm of $\mu_t-\pi$ would not give the required estimate there.

Recall the space $\Xspace^m$ defined by \eqref{eq:X-norm}. Its elements have representatives of at most linear growth. For each such representative there is a vector $a_z\in\R^\d$ such that
\begin{equation}\label{eq:X-tail-properties}
 \lim_{R\to\infty}\sup_{|x|\ge R}|\nabla z(x)-a_z|=0,
 \qquad
 \lim_{R\to\infty}\sup_{|x|\ge R}|\nabla^jz(x)|=0
 ,\quad(2\le j\le m).
\end{equation}
The highest derivatives are uniformly continuous. These properties follow by approximation with the affine-plus-compact smooth functions used to define $\Xspace^m$. They will be used in \cref{app:linear-semigroup} to construct the linearized evolution on this space.

\begin{prop}[Logarithmic stability at every fixed inverse temperature]\label{prop:background-stability}
Fix $\d,\s,\beta$ and an integer $m\ge3$ as in \cref{thm:uniform}, and let $V$ satisfy \eqref{eq:confinement-class}--\eqref{eq:confinement-decay}. There exist $\delta_0,C,\omega>0$ with the following properties.
\begin{enumerate}[(i)]
\item If $z_0\in\Xspace^m$ and $\norm{z_0}_{\Xspace^m}\le\delta_0$, the initial density \eqref{eq:log-tilt} has a global mean-field solution satisfying
\begin{align}
 \norm{\log(\mu_t/\pi)}_{\Xspace^m}
    &\le C\norm{z_0}_{\Xspace^m}e^{-\omega t},\label{eq:log-decay}\\
 \norm{u_t}_{W^{1,\infty}}
    &\le C\norm{z_0}_{\Xspace^m}e^{-\omega t}.\label{eq:background-velocity-decay}
\end{align}
There are constants $a,M,C_P,C_1>0$, independent of $t\ge0$ and of $z_0$ in the stated ball, such that
\begin{align}
 \mu_t(x)&\le Me^{-a|x|^2},\qquad(t\ge0,\ x\in\R^\d),\label{eq:background-Gaussian-bound}\\
 \operatorname{Var}_{\mu_t}(\psi)&\le C_P\int|\nabla\psi|^2\dd\mu_t,
       \qquad(\psi\in H^1(\mu_t),\ t\ge0),\label{eq:background-Poincare-bound}
\end{align}
\begin{equation}\label{eq:background-derivative-bound}
 \sup_{t\ge0}\left(\norm{\nabla^2\log\mu_t}_{L^\infty}
       +\norm{\nabla B_t}_{L^\infty}+\norm{\nabla^2B_t}_{L^\infty}\right)\le C_1.
\end{equation}
Here, $B_t$ and $u_t$ are defined in \eqref{eq:finite-drift} and \eqref{eq:velocity}. For smooth initial data with all higher $\Xspace^\ell$ norms finite, the solution remains smooth on every finite interval. The constants in the preceding estimates do not depend on those higher norms.

\item Let $\pi^\varepsilon$ be the thermal equilibrium measure for the kernel $\g_\varepsilon$ of \eqref{eq:heat-cutoff}, and let $\mu_t^\varepsilon$ solve the corresponding mean-field equation from
\begin{equation}\label{eq:cutoff-background-preparation}
 \mu_0^\varepsilon=\frac{\pi^\varepsilon e^{z_0}}{\int\pi^\varepsilon e^{z_0}}.
\end{equation}
Set $\pi^0:=\pi$, $\mu_t^0:=\mu_t$, and $u_t^0:=u_t$. There is $\varepsilon_*>0$ such that the estimates in (i) hold with the same constants for $0<\varepsilon\le\varepsilon_*$. For every $T<\infty$,
\begin{equation}\label{eq:background-cutoff-convergence}
 \sup_{0\le t\le T}\left(
 \left\|\log\frac{\mu_t^\varepsilon}{\pi^\varepsilon}
             -\log\frac{\mu_t}{\pi}\right\|_{\Xspace^m}
 +\norm{\mu_t^\varepsilon-\mu_t}_{L^1}
 +\norm{u_t^\varepsilon-u_t}_{W^{1,\infty}}\right)\xrightarrow[\varepsilon\downarrow0]{}0.
\end{equation}
For the smooth data in (i), the densities also converge locally smoothly in space, uniformly on finite time intervals.

\item Let $z_0\in\bigcap_{\ell\ge m}\Xspace^\ell$ be arbitrary and smooth. There exist $T_0>0$ and $\varepsilon_*>0$ such that the singular equation and all regularized equations with $0<\varepsilon\le\varepsilon_*$ have solutions on $[0,T_0]$. The bounds \eqref{eq:background-Gaussian-bound}--\eqref{eq:background-derivative-bound} hold on this interval with constants independent of $t\in[0,T_0]$ and $\varepsilon\in[0,\varepsilon_*]$, and \eqref{eq:background-cutoff-convergence} holds for every $T\le T_0$. The time $T_0$ and these constants may depend on $\norm{z_0}_{\Xspace^m}$. No exponential-decay assertion is made in this case.
\end{enumerate}
\end{prop}

We prove the proposition in \cref{app:background}. Here, we check the integrability expenditure along the mean-field trajectory. Let $\norm{z_0}_{\Xspace^m}\le\delta\le\delta_0$, and choose $\varepsilon_*>0$ as in part (ii). By \eqref{eq:background-velocity-decay},
\begin{equation}\label{eq:small-expenditure}
 \frac\beta\kappa\int_0^\infty\sup_{0\le\varepsilon\le\varepsilon_*}\norm{\nabla u_t^\varepsilon}_{L^\infty}\dd t
              \le\frac{\beta C\delta}{\kappa\omega}.
\end{equation}
Here, $\kappa$ is chosen for the Gaussian upper bound, uniform in $t\ge0$ and $\varepsilon\in[0,\varepsilon_*]$, in \eqref{eq:background-Gaussian-bound}. Since $\beta_0>\beta/2$, the right side can be made smaller than the positive bound in \eqref{eq:expenditure} by decreasing $\delta$. Thus, the same one-body argument verifies both the expenditure condition and the background hypotheses used in \cref{prop:uniform-mark}.

\par\endgroup

\begingroup

\subsection{Approximation and completion of the confined proofs}\label{sec:completion}

For \cref{thm:finite}, the hypotheses concern a prescribed mean-field trajectory. We keep that trajectory fixed by compensating the force truncation with a one-body term that vanishes as the cutoff is removed. For \cref{thm:uniform}, the particle system and mean-field equation instead use the same cutoff, with the backgrounds constructed in \cref{prop:background-stability}. After treating the fixed-trajectory construction, \cref{lem:approximation} justifies nonlinear testing and passage to the singular system for both families, as well as for the self-similar family used in \cref{sec:selfsimilar}.

\begin{lemma}[Regularization with a fixed mean-field trajectory]\label{lem:fixed-background-approximation}
Assume the one-body hypotheses of \cref{thm:finite} on $[0,T]$. For the heat truncation \eqref{eq:heat-cutoff}, define
\begin{equation}\label{eq:fixed-background-confinement}
 V_t^\varepsilon=V+(\g-\g_\varepsilon)*\mu_t.
\end{equation}
Use the auxiliary particle equation
\begin{equation}\label{eq:fixed-background-sde}
 \dd x_i^\varepsilon=-\left(\nabla V_t^\varepsilon(x_i^\varepsilon)
 +\frac1N\sum_{j\ne i}\nabla\g_\varepsilon(x_i^\varepsilon-x_j^\varepsilon)\right)\dd t
 +\sqrt{2/\beta}\,\dd B_i,
 \qquad P_{N,0}^\varepsilon=\Q_{N,\beta_0}^\varepsilon(\mu_0),
\end{equation}
where the superscript on $\Q$ means that $\g$ is replaced by $\g_\varepsilon$ in \eqref{eq:modulated}--\eqref{eq:gibbs}. At fixed $N,\varepsilon$, this equation has a unique strong solution on $[0,T]$. Its exact mean-field trajectory is $\mu_t$; writing $u_t^\varepsilon$ for its continuity velocity, we have
\begin{equation}\label{eq:fixed-background-drift}
 -\nabla V_t^\varepsilon-\nabla\g_\varepsilon*\mu_t=B_t,
 \qquad u_t^\varepsilon=u_t.
\end{equation}
Moreover,
\begin{equation}\label{eq:fixed-background-force-bound}
 \sup_{0\le t\le T}\norm{\nabla V_t^\varepsilon-\nabla V}_{W^{2,\infty}}
 \le C\varepsilon^{(\d-\s-1)/2}\longrightarrow0.
\end{equation}
The moving-reference defect and the first-marginal equation are exactly \eqref{eq:defect-coupling} and \eqref{eq:first-mark-equation}, with $c=1$, kernel $\g_\varepsilon$, and the unchanged background $\mu_t$. The moving reference remains $\Q_{N,\beta}^\varepsilon(\mu_t)$, while the initial laws in \eqref{eq:fixed-background-sde} converge in total variation to \eqref{eq:preparation}.
\end{lemma}

\begin{proof}
\parheading{Density derivatives and the force correction.}
Write $U_t=-\log\mu_t$ and let $L$ be a common bound for $\norm{\nabla^2U_t}_\infty$. At each fixed time, $|\nabla U_t(x)|\le|\nabla U_t(0)|+L|x|$. The Gaussian upper bound therefore makes $\nabla\mu_t=-\mu_t\nabla U_t$ integrable. Spatial cutoffs give $\int\nabla\mu_t\dd x=0$, whence
\begin{equation}\label{eq:fixed-background-score-origin}
 |\nabla U_t(0)|=
 \left|\int\bigl(\nabla U_t(0)-\nabla U_t(x)\bigr)\dd\mu_t(x)\right|
 \le L\int|x|\dd\mu_t(x)\le C.
\end{equation}
The final constant is uniform in $t$. Together with
\begin{equation}\label{eq:fixed-background-density-derivatives}
 \nabla\mu_t=-\mu_t\nabla U_t,
 \qquad \nabla^2\mu_t=\mu_t(\nabla U_t\otimes\nabla U_t-\nabla^2U_t),
\end{equation}
this gives common polynomial--Gaussian bounds for the density and its first two derivatives, in particular a uniform $W^{2,\infty}$ bound.

Differentiating the heat representation gives the $L^1$-valued integral
\begin{equation}\label{eq:fixed-background-force-difference}
 \nabla\g-\nabla\g_\varepsilon
 =a_{\d,\s}\int_0^\varepsilon r^{(\d-\s)/2-1}\nabla p_r\dd r.
\end{equation}
The logarithmic constant subtraction disappears upon differentiation. Since $\norm{\nabla p_r}_1=C_\d r^{-1/2}$,
\begin{equation}\label{eq:fixed-background-force-L1}
 \norm{\nabla\g-\nabla\g_\varepsilon}_{L^1}
 \le C\int_0^\varepsilon r^{(\d-\s-3)/2}\dd r
 \le C\varepsilon^{(\d-\s-1)/2}.
\end{equation}
The exponent is positive: for $\d=2$ we have $\s<1$, while for $\d>2$ we have $\d-\s-1>\d/2-1>0$. Placing zero, one, or two derivatives on $\mu_t$ in the convolution proves \eqref{eq:fixed-background-force-bound}. The potential in \eqref{eq:fixed-background-confinement} is well defined by local integrability and the Gaussian upper bound.

At fixed cutoff, the force and all its positive-order derivatives are bounded. Equation \eqref{eq:fixed-background-force-bound} and the bounded Hessian of $V$ make the drift in \eqref{eq:fixed-background-sde} globally Lipschitz, with at most linear growth, uniformly on $[0,T]$. The existence and uniqueness theorem for stochastic differential equations with globally Lipschitz coefficients gives the asserted strong solution on this interval (see, e.g., \cite[Section~5.2]{Oksendal}). The cancellation \eqref{eq:fixed-background-drift} shows that $\mu_t$ solves its mean-field equation exactly. Thus, the time-dependent correction is confined to the auxiliary system and tends to zero by \eqref{eq:fixed-background-force-bound}. The limiting particle equation is \eqref{eq:sde}.

\parheading{The reference defect and the marginal equation.}
The cancellation in \eqref{eq:fixed-background-drift} allows the computation leading to \eqref{eq:defect-coupling} to be used without changing $\mu_t$ or $u_t$. We record the identities needed for the approximation argument. Let $\Gamma_{N,t}^\varepsilon=\Q_{N,\beta}^\varepsilon(\mu_t)$ and use $E_{N,t}^\varepsilon$ and $C_N^\varepsilon$ for \eqref{eq:energy-short} with the truncated kernel. Differentiating the reference density gives
\begin{equation}\label{eq:fixed-background-score}
 \nabla_i\log\Gamma_{N,t}^\varepsilon
 =\nabla\log\mu_t(x_i)-\beta\left(
 \frac1N\sum_{j\ne i}\nabla\g_\varepsilon(x_i-x_j)
 -\nabla\g_\varepsilon*\mu_t(x_i)\right).
\end{equation}
Consequently, the particle drift is $u_t(x_i)+\beta^{-1}\nabla_i\log\Gamma_{N,t}^\varepsilon$, and its generator satisfies
\begin{equation}\label{eq:fixed-background-generator}
 (L_{N,t}^\varepsilon)^*\Gamma_{N,t}^\varepsilon
 =-\sum_i\diver_i\bigl(u_t(x_i)\Gamma_{N,t}^\varepsilon\bigr).
\end{equation}
The continuity equation $\partial_t\mu_t+\diver(\mu_tu_t)=0$ cancels the product-background derivative. In the remaining energy, material differentiation of the pair term gives
$N^{-1}\sum_{i<j}(u_t(x_i)-u_t(x_j))\cdot\nabla\g_\varepsilon(x_i-x_j)$; the one-body and background terms supply the two cross terms and the background double integral. Thus,
\begin{equation}\label{eq:fixed-background-material-energy}
 \left(\partial_t+\sum_i u_t(x_i)\cdot\nabla_i\right)E_{N,t}^\varepsilon
 =C_N^\varepsilon[u_t].
\end{equation}
Integrating the material derivative of the unnormalized Gibbs density, with spatial cutoffs and the Gaussian tails, gives
\begin{equation}\label{eq:fixed-background-normalizer}
 \frac{\dd}{\dd t}\log\K_{N,\beta}^\varepsilon(\mu_t)
 =-\beta\E_{\Gamma_{N,t}^\varepsilon}[C_N^\varepsilon[u_t]].
\end{equation}
Combining these identities proves
\begin{equation}\label{eq:fixed-background-defect}
 \partial_t\Gamma_{N,t}^\varepsilon-(L_{N,t}^\varepsilon)^*\Gamma_{N,t}^\varepsilon
 =-\beta\left(C_N^\varepsilon[u_t]
 -\E_{\Gamma_{N,t}^\varepsilon}[C_N^\varepsilon[u_t]]\right)\Gamma_{N,t}^\varepsilon.
\end{equation}

For the first marginal, write $h_t^\varepsilon=\g_\varepsilon*\mu_t$ and use the decomposition \eqref{eq:dynamic-projections} for the approximate law. Subtracting the exact mean-field equation from the first particle marginal equation gives
\begin{align}
 \partial_tw_{N,t}^\varepsilon={}&\mathcal L_t^*w_{N,t}^\varepsilon
 +\frac{N-1}{N}\diver\bigl(\mu_t\nabla\g_\varepsilon*w_{N,t}^\varepsilon\bigr)
 -\frac1N\diver\bigl(w_{N,t}^\varepsilon\nabla h_t^\varepsilon\bigr)\notag\\
 &-\frac1N\diver\bigl(\mu_t\nabla h_t^\varepsilon\bigr)
 +\frac{N-1}{N}S^\varepsilon_{\eta_{N,t}^\varepsilon},
 \qquad \mathcal L_t^*=\beta^{-1}\Delta-\diver(B_t\,\cdot).
 \label{eq:fixed-background-first-marginal}
\end{align}
In particular, the backward propagator is unchanged. The two terms with coefficient $-1/N$ arise from $(N-1)/N-1=-1/N$, when the particle interaction is compared with the mean-field force; see \eqref{eq:first-mark-equation}. The fixed-cutoff testing of these identities is justified in \cref{lem:approximation}.

\parheading{Convergence of the initial Gibbs preparations.}
Set
\begin{equation}\label{eq:fixed-background-short-potential}
 H_\varepsilon(x):=a_{\d,\s}\int_0^\varepsilon r^{(\d-\s)/2-1}p_r(x)\dd r\ge0.
\end{equation}
Here, $r$ is the heat parameter in \eqref{eq:heat-cutoff}, and $\norm{H_\varepsilon}_1=C_{\d,\s}\varepsilon^{(\d-\s)/2}$. For $\s>0$, $\g-\g_\varepsilon=H_\varepsilon$. For $\s=0$, the difference is $H_\varepsilon-H_\varepsilon(e)$, and $H_\varepsilon(e)\to0$ by the factor $e^{-1/(4r)}$ in the heat kernel. It follows that $h_{\mu_0}^\varepsilon\to h_{\mu_0}$ uniformly, $I_{\mu_0}^\varepsilon\to I_{\mu_0}$, and the pair potential converges off the collision diagonal.

For $\s>0$, the pair potential is nonnegative and the one-body convolution is bounded above uniformly in $\varepsilon$. The unnormalized Gibbs density is therefore bounded by $C_N\mu_0^{\otimes N}$. For the logarithm and $0<\varepsilon\le1$,
\begin{equation}\label{eq:fixed-background-log-lower}
 \g_\varepsilon(x)\ge-C-\log(1+|x|).
\end{equation}
Indeed, $\g-H_\varepsilon\ge\g-H_1$; the right-hand side is smooth on bounded sets and has this lower bound at infinity, and the normalization constants $H_\varepsilon(e)$ are bounded. The positive part of $h_{\mu_0}^\varepsilon$ is bounded uniformly by local integrability of the logarithm and boundedness of $\mu_0$. The energies $I_{\mu_0}^\varepsilon$ are bounded as well. Since $1+|x-y|\le(1+|x|)(1+|y|)$, a fixed-$N$ majorant for the unnormalized density is
\begin{equation}\label{eq:fixed-background-gibbs-majorant}
 C_N\prod_{i=1}^N(1+|x_i|)^{\beta_0(N-1)/N}e^{-a|x_i|^2}.
\end{equation}
This is integrable. Dominated convergence and the positive lower bound \eqref{eq:static-partition} on the normalizers give total-variation convergence of the normalized initial laws.
\end{proof}

\begin{lemma}[Nonlinear testing and passage to the actual singular law]\label{lem:approximation}
Fix a finite interval $[0,T]$ on which the relevant background bounds hold. Consider one of the following heat-regularized particle systems:
\begin{enumerate}[(i)]
\item the auxiliary equation \eqref{eq:fixed-background-sde}, with $\mu_t^\varepsilon=\mu_t$ and $c\equiv1$;
\item the system with the original confinement, $c\equiv1$, and the cutoff backgrounds of \cref{prop:background-stability};
\item the self-similar system of \cref{sec:selfsimilar}, with its Gaussian backgrounds and decreasing Riesz coupling $c(t)$.
\end{enumerate}
In (ii) and (iii), use the same heat cutoff in the particle equation, mean-field equation, and Gibbs laws. In all three cases, set
\begin{equation}\label{eq:approximation-laws}
 P_{N,0}^\varepsilon:=\Q_{N,\beta_0c(0)}^\varepsilon(\mu_0^\varepsilon),
 \qquad \Gamma_{N,t}^\varepsilon:=\Q_{N,\beta c(t)}^\varepsilon(\mu_t^\varepsilon).
\end{equation}
Then the following conclusions hold.
\begin{enumerate}[(a)]
\item At each fixed $N,\varepsilon>0$, the identities \eqref{eq:fixed-power}, \eqref{eq:log-power}, \eqref{eq:initial-score}, and \eqref{eq:first-mark-equation} are justified on $[0,T]$. The resulting estimates of \cref{sec:decreasing,sec:marginals} have constants independent of the cutoff whenever the one-body and static bounds entering those estimates have constants independent of the cutoff.
\item For each fixed $N$, the particle path laws satisfy
\begin{equation}\label{eq:approximation-path-laws}
 \operatorname{Law}(X_N^\varepsilon)\rightharpoonup\operatorname{Law}(X_N)
 \quad\text{in }\mathcal P\bigl(C([0,T];\R^{\d N})\bigr),
 \qquad \varepsilon\downarrow0,
\end{equation}
where $X_N$ is the corresponding singular particle system with initial law $\Q_{N,\beta_0c(0)}(\mu_0)$. In case (iii), it is the exact self-similar transform of the unconfined particle system. For each $1\le k\le N$ and $t\in[0,T]$,
\begin{equation}\label{eq:approximation-square-transfer}
 \chi^2(P_{N,t}^{(k)}\mid\mu_t^{\otimes k})
 \le\liminf_{\varepsilon\downarrow0}
       \chi^2(P_{N,t}^{\varepsilon,(k)}\mid(\mu_t^\varepsilon)^{\otimes k}).
\end{equation}
In particular, any bound on the right independent of the cutoff passes to the singular law, including the full-particle bound at $k=N$.
\end{enumerate}
\end{lemma}
\begin{proof}
\parheading{Testing at fixed particle number and cutoff.}
Fix $N$ and $\varepsilon>0$. The force and its positive-order derivatives are bounded. The particle drift is globally Lipschitz and has at most linear growth, including the time-dependent one-body correction by \eqref{eq:fixed-background-force-bound}. Since the Hessian of $U_t^\varepsilon=-\log\mu_t^\varepsilon$ is bounded, its gradient has at most linear growth. The reference $\Gamma_{N,t}^\varepsilon$ therefore has Gaussian or polynomial--Gaussian tails, a score of at most linear growth, and bounded logarithmic Hessian, with constants that may depend on $N,\varepsilon,T$.

The coefficient in the relative-density equation \eqref{eq:q-general} is
\begin{equation}\label{eq:cutoff-defect-coefficient}
 R_{N,t}^\varepsilon:=
 \frac{\partial_t\Gamma_{N,t}^\varepsilon-(L_{N,t}^\varepsilon)^*\Gamma_{N,t}^\varepsilon}
      {\Gamma_{N,t}^\varepsilon},
 \qquad
 \sup_{0\le t\le T}\norm{R_{N,t}^\varepsilon}_{W^{1,\infty}(\R^{\d N})}
 \le C_{N,\varepsilon,T}.
\end{equation}
To prove the bound for $c\equiv1$, use \eqref{eq:fixed-background-defect} and the fixed-cutoff bounds on $|z||\nabla\g_\varepsilon(z)|$, $\nabla\g_\varepsilon(z)$, and $|z||\nabla^2\g_\varepsilon(z)|$, together with the bound on $\nabla u_t^\varepsilon$. For decreasing Riesz coupling, \eqref{eq:defect-coupling} also contains $c'E_{N,t}^\varepsilon$. The potential $\g_\varepsilon$ and its first derivative are bounded when $\s>0$, so this contribution and its first spatial derivative satisfy the same fixed-$N,\varepsilon,T$ bound.

For an explicit justification of the relative-density equation, set
\begin{equation}\label{eq:cutoff-initial-couplings}
 b_0=\beta_0c(0),\qquad b_1=\beta c(0),
\end{equation}
and write $E_{N,0}^\varepsilon=N\Fr_N^\varepsilon(\cdot,\mu_0^\varepsilon)$. The initial density relative to the moving reference is
\begin{equation}\label{eq:cutoff-initial-ratio}
 q_{N,0}^\varepsilon
 =\frac{\K_{N,b_1}^\varepsilon(\mu_0^\varepsilon)}
        {\K_{N,b_0}^\varepsilon(\mu_0^\varepsilon)}
   \exp\bigl((b_1-b_0)E_{N,0}^\varepsilon\bigr).
\end{equation}
For $\s>0$, the cutoff potential and $E_{N,0}^\varepsilon$ are bounded. For $\s=0$, the cutoff is smooth at the origin and grows at most logarithmically at infinity, while its force is bounded. Hence, for fixed $N,\varepsilon$, there are $C,m<\infty$ such that
\begin{equation}\label{eq:cutoff-initial-growth}
 C^{-1}(1+|X_N|)^{-m}\le q_{N,0}^\varepsilon(X_N)
 \le C(1+|X_N|)^m,
 \qquad |\nabla q_{N,0}^\varepsilon(X_N)|\le C(1+|X_N|)^m.
\end{equation}

Solve \eqref{eq:q-general} with initial density \eqref{eq:cutoff-initial-ratio} by the Feynman--Kac formula (see, e.g., \cite[Section~8.2]{Oksendal}), evaluating the time-dependent coefficients in reverse time for this initial-value problem. The auxiliary diffusion drift $2\beta^{-1}\nabla\log\Gamma-b^N$ is globally Lipschitz and has at most linear growth. Its polynomial moment bounds and \eqref{eq:cutoff-defect-coefficient} propagate the polynomial upper bound in \eqref{eq:cutoff-initial-growth}. Jensen's inequality applied to the reciprocal gives a polynomial lower bound, and differentiation of the spatial stochastic flow gives a polynomial bound for the first derivative of $q$. All these bounds are on a fixed finite interval and may depend on $N,\varepsilon,T$.

Combined with the Gaussian or polynomial--Gaussian reference tails, they justify the power tests, logarithms, and integrations by parts in \eqref{eq:fixed-power} and \eqref{eq:log-power}. Local parabolic regularity gives the classical equation on compact subsets of $(0,T]\times\R^{\d N}$ (see, e.g., \cite[Chapter~5]{Lunardi}). The product $q\Gamma$ solves the particle forward equation and has mass one, by spatial cutoff and the tail bounds. It starts from $P_{N,0}^\varepsilon$, so uniqueness identifies it with the density of the regularized particle system. The same domination justifies differentiation of the initial Gibbs marginal and the first-marginal equation. These fixed-$N,\varepsilon$ estimates justify the calculations; the cutoff-independent bounds come from \cref{sec:decreasing,sec:marginals}.

\parheading{Convergence of the initial laws.}
For the auxiliary system \eqref{eq:fixed-background-sde}, total-variation convergence follows from \cref{lem:fixed-background-approximation}. For the other two families, the convergence and bounds for the initial mean-field densities give
\begin{equation}\label{eq:approximation-initial-density-convergence}
 \mu_0^\varepsilon\longrightarrow\mu_0\quad\text{locally uniformly},
 \qquad 0<\mu_0^\varepsilon(x)\le M e^{-a|x|^2},
 \qquad 0<\varepsilon\le\varepsilon_*.
\end{equation}
The heat-kernel identity \eqref{eq:fixed-background-short-potential} gives, for each $R<\infty$,
\begin{equation}
 \sup_{x\in B_R}|\g_\varepsilon*\mu_0^\varepsilon(x)-\g*\mu_0(x)| \le \norm{(\g_\varepsilon-\g)*\mu_0^\varepsilon}_{L^\infty} +\sup_{x\in B_R}|\g*(\mu_0^\varepsilon-\mu_0)(x)| \xrightarrow[\varepsilon\downarrow0]{}0. \label{eq:approximation-initial-potential-convergence}
\end{equation}
Indeed, the first term is at most $M\norm{H_\varepsilon}_{L^1}$ for $\s>0$, with the additional constant $H_\varepsilon(e)$ for $\s=0$. For the second term, split the convolution into $|y|\le R'$ and $|y|>R'$. Local uniform convergence and local integrability of $\g$ control the first part; the Gaussian upper bound in \eqref{eq:approximation-initial-density-convergence} makes the second part uniformly small as $R'\to\infty$. The same bounds give $I_{\mu_0^\varepsilon}^\varepsilon\to I_{\mu_0}$.

Let $F_{N,0}^\varepsilon$ denote the unnormalized initial Gibbs density at coupling $b_0=\beta_0c(0)$. Off the collision diagonal,
\begin{equation}\label{eq:approximation-initial-Gibbs-domination}
 F_{N,0}^\varepsilon(X_N)\longrightarrow F_{N,0}(X_N),
 \qquad
 0\le F_{N,0}^\varepsilon(X_N)
 \le C_N\prod_{i=1}^N(1+|x_i|)^{b_0(N-1)/N}e^{-a|x_i|^2}.
\end{equation}
For $\s>0$, the polynomial factor can be omitted. The bound follows as in \eqref{eq:fixed-background-gibbs-majorant} and is integrable at fixed $N$. Dominated convergence and the lower partition bound therefore give
\begin{equation}\label{eq:approximation-initial-TV}
 \TV(P_{N,0}^\varepsilon,P_{N,0})\longrightarrow0,
 \qquad \varepsilon\downarrow0.
\end{equation}

\parheading{Convergence of the particle paths.}
The target singular system has a global noncolliding strong solution by \cref{prop:particle-wellposedness}; for case (iii), use its exact self-similar transform. Couple the initial configurations so that $X_{N,0}^\varepsilon\to X_{N,0}$ in probability, and drive all equations by the same Brownian motions, independent of these initial configurations.

For $R>1$, set
\begin{equation}\label{eq:approximation-compact-sets}
 \mathcal K_R:=\{X_N:|X_N|\le R,\ \min_{i<j}|x_i-x_j|\ge R^{-1}\}.
\end{equation}
The particle drifts converge with their first derivatives uniformly on $[0,T]\times\mathcal K_{2R}$. For \eqref{eq:fixed-background-sde}, this uses \eqref{eq:fixed-background-force-bound} for the one-body correction. Choose a smooth configuration-space cutoff equal to one on a neighborhood of $\mathcal K_R$ and supported in the interior of $\mathcal K_{2R}$. Write $b^{\varepsilon,R}$ and $b^R$ for the extended drifts, and $X_N^{\varepsilon,R}$ and $X_N^R$ for their solutions. Their Lipschitz constants are bounded by a constant $L_R$ independent of small $\varepsilon$, and Gronwall gives
\begin{equation}\label{eq:approximation-path-Gronwall}
 \sup_{t\le T}|X_{N,t}^{\varepsilon,R}-X_{N,t}^R|
 \le e^{L_RT}\left(|X_{N,0}^\varepsilon-X_{N,0}|
          +T\norm{b^{\varepsilon,R}-b^R}_{L^\infty([0,T]\times\R^{\d N})}\right).
\end{equation}
On the event $X_{N,t}\in\mathcal K_R$ for every $t\le T$, the limiting extended path equals the singular path. The positive distance from $\mathcal K_R$ to the boundary of the region where the cutoff equals one and \eqref{eq:approximation-path-Gronwall} identify the approximate extended path with the original approximate path, with probability tending to one as $\varepsilon\downarrow0$. Thus, for every $\eta>0$,
\begin{equation}\label{eq:approximation-path-probability}
 \limsup_{\varepsilon\downarrow0}
 \Pb\left(\sup_{t\le T}|X_{N,t}^\varepsilon-X_{N,t}|>\eta\right)
 \le\Pb\left(X_{N,t}\notin\mathcal K_R\text{ for some }t\le T\right)
 \xrightarrow[R\to\infty]{}0.
\end{equation}
The last limit follows from continuity, nonexplosion, and positive minimum pair separation on $[0,T]$. This proves \eqref{eq:approximation-path-laws} at each fixed $N$.

\parheading{Passage of the second-moment bounds.}
Suppose that $P_j\rightharpoonup P$, $\nu_j\rightharpoonup\nu$, and
$C:=\liminf_j\norm{\mathrm dP_j/\mathrm d\nu_j-1}_{L^2(\nu_j)}<\infty$.
Along a subsequence realizing this lower limit, Cauchy--Schwarz and weak convergence give, for every $\psi\in C_c(\R^{\d k})$,
\begin{equation}\label{eq:L2-lower-semicontinuity}
 \left|\int\psi\dd(P-\nu)\right|
 =\lim_j\left|\int\psi\dd(P_j-\nu_j)\right|
 \le C\left(\int\psi^2\dd\nu\right)^{1/2}.
\end{equation}
The Riesz--Fr\'echet representation theorem and density of $C_c$ in $L^2(\nu)$ give $\mathrm dP/\mathrm d\nu-1\in L^2(\nu)$ with norm at most $C$. Apply this with
\begin{equation}\label{eq:approximation-terminal-measures}
 P_j=P_{N,t}^{\varepsilon_j,(k)},\quad
 \nu_j=(\mu_t^{\varepsilon_j})^{\otimes k},\quad
 P=P_{N,t}^{(k)},\quad \nu=\mu_t^{\otimes k}.
\end{equation}
Path-law convergence gives $P_j\rightharpoonup P$, while the corresponding background convergence gives $\nu_j\rightharpoonup\nu$. This proves \eqref{eq:approximation-square-transfer}. In particular, if an interval $I$ carries the bound
\begin{equation}\label{eq:approximation-uniform-time-transfer}
 \sup_{\varepsilon}\sup_{t\in I}
 \chi^2(P_{N,t}^{\varepsilon,(k)}\mid(\mu_t^\varepsilon)^{\otimes k})\le C_{N,k},
 \qquad\text{then}\qquad
 \sup_{t\in I}\chi^2(P_{N,t}^{(k)}\mid\mu_t^{\otimes k})\le C_{N,k}.
\end{equation}
For $I=[0,\infty)$, perform the path-law convergence on each finite interval and apply the bound at each fixed time before taking the supremum. No uniform-in-time path convergence is used.
\end{proof}

\begin{proof}[Proof of \cref{thm:uniform}]
\Cref{prop:background-stability} supplies the one-body bounds uniformly in time and cutoff. Choose the perturbation radius so that \eqref{eq:small-expenditure} gives \eqref{eq:expenditure} with $T=\infty$, and then choose $p_0>2$ as in \eqref{eq:temperature-exponent-choice}. By \cref{lem:initial-temperature}, with $b_0=\beta_0$ and $b_1=\beta$, the initial relative $L^{p_0}$ norm is bounded. \Cref{cor:full-moment} therefore gives the full second-moment bound on $[0,\infty)$.

Apply \cref{prop:initial-mark} at $b=\beta_0$ and propagate its one-mark estimate by \cref{prop:uniform-mark}. The product-space projection and information inequalities \eqref{eq:all-k}--\eqref{eq:TV-from-square} give the claimed marginal rates, with constants independent of the cutoff. Finally, \cref{lem:approximation} passes these bounds to the singular system. The dependence on the initial perturbation is only through the low norm stated in the theorem.
\end{proof}

\begin{proof}[Proof of \cref{thm:finite,cor:local}]
Use the auxiliary regularization of \cref{lem:fixed-background-approximation}. Its mean-field trajectory and continuity velocity are exactly $\mu_t$ and $u_t$. The one-body constants, weighted backward propagator, and expenditure are therefore independent of the cutoff. The static bounds are uniform by \cref{prop:static-inputs}, and \eqref{eq:fixed-background-defect} supplies the same exponential estimate for the coefficient in the relative-density equation, with
\begin{equation}\label{eq:fixed-background-expenditure}
 \lambda_0(t)=\frac\beta\kappa\norm{\nabla u_t}_{L^\infty},
 \qquad \int_0^T\lambda_0(t)\dd t=A_\mu(T).
\end{equation}
Thus, the exponent equation is the same at every cutoff. Choose $p_0>2$ satisfying \eqref{eq:temperature-exponent-choice}, as permitted by \eqref{eq:expenditure}. \Cref{lem:initial-temperature}, with $b_0=\beta_0$ and $b_1=\beta$, and \cref{cor:full-moment} give the full product-space second-moment bound.

For the first marginal, apply \cref{prop:initial-mark} at $b=\beta_0$ and then \cref{prop:finite-mark}. The latter applies because \eqref{eq:fixed-background-first-marginal} has the original backward propagator and the same two terms with coefficient $-1/N$. Equations \eqref{eq:all-k}--\eqref{eq:TV-from-square} give the marginal rates with a constant independent of $N,k,\varepsilon$, and \cref{lem:approximation} proves \eqref{eq:finite-rates} after passage to the singular system.

For a confinement satisfying \eqref{eq:confinement-class}--\eqref{eq:confinement-decay} and any smooth logarithmic perturbation in \eqref{eq:explicit-class}, the local part of \cref{prop:background-stability} provides the intrinsic one-body bounds on a positive interval. Since the right-hand side of \eqref{eq:expenditure} is positive, shortening the interval gives that condition, so \cref{thm:finite} proves the corollary.
\end{proof}

\begin{prop}[Stationary sharpness]\label{prop:stationary-sharpness}
Let $\beta>0$ and let $\pi$ be a thermal equilibrium measure in the class of \cref{thm:uniform}. There is $c>0$, depending only on the fixed equilibrium $\pi$ and the parameters $\d,\s,\beta$, such that, for every $N\ge2$ and $2\le k\le N$,
\begin{equation}\label{eq:stationary-TV-lower}
 \TV\bigl(\Q_{N,\beta}^{(k)}(\pi),\pi^{\otimes k}\bigr)\ge c\frac{k}{N},
\end{equation}
\begin{equation}\label{eq:stationary-entropy-lower}
 \Hr\bigl(\Q_{N,\beta}^{(k)}(\pi)\mid\pi^{\otimes k}\bigr)\ge2c^2\frac{k^2}{N^2},
\end{equation}
\begin{equation}\label{eq:stationary-chi-lower}
 \chi^2\bigl(\Q_{N,\beta}^{(k)}(\pi)\mid\pi^{\otimes k}\bigr)\ge4c^2\frac{k^2}{N^2}.
\end{equation}
\end{prop}
\begin{proof}
For sufficiently large $N$, these lower bounds are a stationary specialization of \cite[Theorem~2.1]{DGRchaos}. We give a proof using the upper bounds just established and the static fluctuation formula already used in this paper.

Set $\Gamma_N:=\Q_{N,\beta}(\pi)$ and $f_{N,k}:=\dd\Gamma_N^{(k)}/\dd\pi^{\otimes k}$. The stationary case of \cref{thm:uniform} gives
\begin{equation}\label{eq:sharpness-L2-upper}
 \norm{f_{N,k}-1}_{L^2(\pi^{\otimes k})}\le C_0\frac{k}{N},\qquad 1\le k\le N.
\end{equation}
The nonnegative operator $T_\pi$ is nonzero: on any fixed ball, $h_\pi$ is bounded while $\g(x-y)$ diverges at the diagonal, and $\pi$ has a positive lower bound there. Thus, its centered kernel cannot vanish almost everywhere. By density, choose a bounded real $\varphi\in L^2_0(\pi)$ with $T_\pi\varphi\ne0$. The spectral theorem gives
\begin{equation}\label{eq:sharpness-variance-gap}
 v:=\norm{\varphi}_{L^2(\pi)}^2,
 \qquad v_\beta:=\langle\varphi,(\Id+\beta T_\pi)^{-1}\varphi\rangle_{L^2(\pi)}<v.
\end{equation}
The Gaussian upper bound on $\pi$ verifies the logarithmic integrability hypotheses, when needed, of \cite[Proposition~2.11]{DGRpartition}. In the notation
\begin{equation}\label{eq:sharpness-statistic}
 S_N:=\frac1{\sqrt N}\sum_{i=1}^N\varphi(x_i),
\end{equation}
that result gives $S_N\to\mathcal N(0,v_\beta)$ in law under $\Gamma_N$. Boundedness and centering of $\varphi$ give a uniform eighth-moment bound for $S_N$ under $\pi^{\otimes N}$: in its expansion, every index occurring only once integrates to zero, leaving at most four distinct indices and $O(N^4)$ terms. By \eqref{eq:static-powers},
\begin{equation}\label{eq:sharpness-moment-transfer}
 \sup_{N\ge2}\E_{\Gamma_N}[|S_N|^4]
 \le\sup_{N\ge2}\norm{G_{N,\beta}(\pi)}_2
       \bigl(\E_{\pi^{\otimes N}}[|S_N|^8]\bigr)^{1/2}<\infty.
\end{equation}
Hence, the second moments $\E_{\Gamma_N}[S_N^2]$ converge to $v_\beta$. Also, \eqref{eq:sharpness-L2-upper} at $k=1$ gives $\E_{\Gamma_N}[\varphi(x_1)^2]\to v$. Exchangeability yields
\begin{equation}\label{eq:sharpness-pair-covariance}
 (N-1)\E_{\Gamma_N}[\varphi(x_1)\varphi(x_2)]
 =\E_{\Gamma_N}[S_N^2]-\E_{\Gamma_N}[\varphi(x_1)^2]\longrightarrow v_\beta-v<0.
\end{equation}
For $2\le k\le N$, define the centered pair statistic
\begin{equation}\label{eq:sharpness-pair-statistic}
 Z_k:=\frac1k\sum_{\substack{1\le i,j\le k\\i\ne j}}\varphi(x_i)\varphi(x_j)
 =\left(\frac1{\sqrt k}\sum_{i=1}^k\varphi(x_i)\right)^2-\frac1k\sum_{i=1}^k\varphi(x_i)^2.
\end{equation}
Its product expectation is zero, and the same eighth-moment calculation gives
\begin{equation}\label{eq:sharpness-pair-fourth-moment}
 K_\varphi:=\sup_{k\ge2}\E_{\pi^{\otimes k}}[|Z_k|^4]<\infty.
\end{equation}
For all sufficiently large $N$, \eqref{eq:sharpness-pair-covariance} implies, uniformly for $2\le k\le N$,
\begin{equation}\label{eq:sharpness-pair-response}
 \left|\int Z_k(f_{N,k}-1)\dd\pi^{\otimes k}\right|
 =(k-1)\left|\E_{\Gamma_N}[\varphi(x_1)\varphi(x_2)]\right|
 \ge\frac{v-v_\beta}{4}\frac{k}{N}.
\end{equation}
Let $Z_{k,R}:=\max\{-R,\min\{Z_k,R\}\}$. Since $\norm{Z_k-Z_{k,R}}_2\le K_\varphi^{1/2}/R$, \eqref{eq:sharpness-L2-upper} gives
\begin{equation}\label{eq:sharpness-clipping-error}
 \left|\int(Z_k-Z_{k,R})(f_{N,k}-1)\dd\pi^{\otimes k}\right|
 \le\frac{C_0K_\varphi^{1/2}}{R}\frac{k}{N}.
\end{equation}
Choose $R\ge8C_0K_\varphi^{1/2}/(v-v_\beta)$, independently of $N,k$. By total-variation duality,
\begin{equation}\label{eq:sharpness-TV-duality}
 2R\,\TV(\Gamma_N^{(k)},\pi^{\otimes k})
 \ge\left|\int Z_{k,R}(f_{N,k}-1)\dd\pi^{\otimes k}\right|
 \ge\frac{v-v_\beta}{8}\frac{k}{N}.
\end{equation}
This proves \eqref{eq:stationary-TV-lower} for all sufficiently large $N$.

For each remaining $N\ge2$, we claim that $\Gamma_N^{(2)}\ne\pi^{\otimes2}$. Otherwise, its one- and two-body expectations would give $\E_{\Gamma_N}[N\Fr_N(X_N,\pi)]=-I_\pi$, and \eqref{eq:gibbs} would imply
\begin{equation}\label{eq:sharpness-finite-N-entropy}
 \Hr(\Gamma_N\mid\pi^{\otimes N})=\beta I_\pi-\log\K_{N,\beta}(\pi)\le0.
\end{equation}
Thus, $\Gamma_N=\pi^{\otimes N}$, which is impossible: $N\Fr_N$ is not constant, since it diverges at a pair collision with the other coordinates separated. All the expectations used here are finite by \eqref{eq:static-partition} and \eqref{eq:static-powers}. By projection, $\TV(\Gamma_N^{(k)},\pi^{\otimes k})\ge\TV(\Gamma_N^{(2)},\pi^{\otimes2})>0$ for every $k\ge2$. Decreasing $c$ over the finite set of remaining particle numbers proves \eqref{eq:stationary-TV-lower} for every $N\ge2$. Finally, Pinsker's inequality and Cauchy--Schwarz give $\Hr\ge2\TV^2$ and $\chi^2\ge4\TV^2$, respectively, and prove the remaining bounds.
\end{proof}

\begin{remark}[Periodic geometry]\label{rem:periodic}
Fix $\beta>0$ and $\beta_0>\beta/2$. On the unit-volume torus, take $V=0$ and the mean-zero spectral log/Riesz kernel. With initial law $\Q_{N,\beta_0}(\mu_0)$, the same terminal estimates hold for every smooth strictly positive mean-field measure on each interval satisfying
\begin{equation}\label{eq:periodic-expenditure}
 \frac\beta\kappa\int_0^T\norm{\nabla u_t}_{L^\infty}\dd t
 <\log\left(2\min\left\{1,\frac{\beta_0}{\beta}\right\}\right),
\end{equation}
and hence on an initial interval for arbitrary such data. They hold uniformly for all time when $\norm{\mu_0-1}_{H^\ell}$ is sufficiently small, for any fixed integer $\ell>\d/2+2$.

The static inputs follow from \cite[Remark~2.1]{DGRpartition}. On each finite interval, smoothness and strict positivity give upper and positive lower density bounds, hence a one-body Poincar\'e constant. The weighted elliptic and parabolic estimates apply on the compact manifold. To obtain a background satisfying the expenditure condition for all time, write $w=\mu-1$. Its equation is
\begin{equation}\label{eq:periodic-background}
 \partial_tw=(\beta^{-1}\Delta+\Delta\g*)w+
                     \diver(w\nabla\g*w).
\end{equation}
The nonzero Fourier multipliers of its linear semigroup are
$\exp[-4\pi^2|n|^2(\beta^{-1}+\widehat\g(n))t]$.
They have a positive diffusive gap and gain one derivative with an integrable $t^{-1/2}$ singularity. The Sobolev algebra estimate and the torus $L^1$ bound on the force give a quadratic bound $C\norm w_{H^\ell}^2$ for the flux. The resulting contraction in the exponentially weighted trajectory norm yields, for sufficiently small $\norm{w_0}_{H^\ell}$,
\begin{equation}\label{eq:periodic-background-decay}
 \sup_{t\ge0}e^{\omega t}\norm{w_t}_{H^\ell}
 \le C\norm{w_0}_{H^\ell}.
\end{equation}
In particular, $1/2\le\mu_t\le3/2$ after decreasing the initial radius. Since
$u_t=-\nabla\g*w_t-\beta^{-1}\nabla\log(1+w_t)$, Sobolev embedding and the force bound give
\begin{equation}\label{eq:periodic-velocity-budget}
 \frac\beta\kappa\int_0^\infty\norm{\nabla u_t}_{L^\infty}\dd t
 \le\frac{\beta C}{\kappa}\int_0^\infty\norm{w_t}_{H^\ell}\dd t
 \le\frac{\beta C}{\kappa\omega}\norm{\mu_0-1}_{H^\ell}.
\end{equation}
Choose $\kappa$ for the common density bound. The right side is smaller than the bound in \eqref{eq:periodic-expenditure} when the initial radius is sufficiently small. The Fourier and force estimates are uniform for the periodic heat cutoffs, so the same choice works before passing to the singular kernel.

For the periodic version of \cref{lem:weighted-commutator}, use
\begin{equation}\label{eq:periodic-gradient-translation}
 \norm{\nabla\psi(\cdot+z)-\nabla\psi}_{L^2(\T^\d)}
 \le 2^{1-\alpha}|z|^\alpha
       \norm{\nabla\psi}_{L^2(\T^\d)}^{1-\alpha}
       \norm{\nabla^2\psi}_{L^2(\T^\d)}^\alpha,
 \qquad 0<\alpha<1,
\end{equation}
where $z\in[-1/2,1/2]^\d$ and $\psi\in H^2(\T^\d)$. This follows by interpolating the bounds $2\norm{\nabla\psi}_2$ and $|z|\norm{\nabla^2\psi}_2$ for the left side. For $\max\{0,\s+1-\d/2\}<\alpha<1$, the periodic heat cutoffs satisfy
\begin{equation}\label{eq:periodic-force-square}
 \sup_{0\le\varepsilon\le1}
 \int_{\T^\d}\operatorname{dist}(z,0)^{2\alpha}
                  |\nabla\g_\varepsilon(z)|^2\dd z<\infty.
\end{equation}
The upper and positive lower bounds on $\mu_t$ allow comparison of the weighted and unweighted norms. Applying Cauchy--Schwarz in the two-point pairing with \eqref{eq:periodic-gradient-translation} and \eqref{eq:periodic-force-square} gives the required periodic mixed gradient--Hessian estimate. Thus, the arguments from the full moment through the first marginal and product-space projection apply on the torus. For arbitrary smooth positive initial data, the local contraction argument gives the same background bounds on an initial interval; shortening that interval gives \eqref{eq:periodic-expenditure}.
\end{remark}

\par\endgroup

\begingroup

\section{Unconfined consequences in self-similar variables}\label{sec:selfsimilar}

We apply the self-similar transformation of Serfaty and the third author \cite[Section~2]{RSselfsimilar} to the unconfined system, obtained from \eqref{eq:sde} by setting $V=0$. The transformed equation has confinement $|y|^2/4$. For the logarithmic interaction, \cref{thm:uniform} applies directly. For positive Riesz exponents, the interaction coefficient decreases in time; \cref{lem:decaying-coupling} supplies the Gaussian background estimates needed to apply \cref{cor:full-moment,prop:uniform-mark}.

An initial length scale $L>0$ will determine the strength of the transformed Riesz interaction. Set
\begin{equation}\label{eq:selfsimilar}
 \tau:=\log(1+t/L^2),\qquad
 y_i(\tau):=\frac{x_i(t)}{L e^{\tau/2}}.
\end{equation}
Write $D_\tau(y):=Le^{\tau/2}y$, so that $t=L^2(e^\tau-1)$ and $\mathrm dt/\mathrm d\tau=L^2e^\tau$. Define the transformed measures by
\begin{equation}\label{eq:selfsimilar-measures}
 \bar\mu_\tau(y):=(Le^{\tau/2})^\d\mu_{L^2(e^\tau-1)}(Le^{\tau/2}y),\qquad
 \bar P_{N,\tau}:=(D_\tau^{-1})^{\otimes N}_{\#}P_{N,L^2(e^\tau-1)}.
\end{equation}
Brownian scaling and the identity $\nabla\g(ay)=a^{-\s-1}\nabla\g(y)$ give
\begin{equation}\label{eq:transformed-sde}
 \dd y_i=\left[-\frac12y_i-
       \frac{c(\tau)}N\sum_{j\ne i}\nabla\g(y_i-y_j)\right]\dd\tau
                  +\sqrt{2/\beta}\dd\widetilde B_i,
\end{equation}
where the $\widetilde B_i$ are independent Brownian motions and
\begin{equation}\label{eq:selfsimilar-coupling}
 c(\tau):=\begin{cases}1,&\s=0,\\ L^{-\s}e^{-\s\tau/2},&\s>0.\end{cases}
\end{equation}
The force coefficient is $(Le^{\tau/2})^{-\s}$, while $\mathrm dt/\mathrm d\tau=(Le^{\tau/2})^2$ leaves the noise coefficient unchanged. This is the calculation of \cite[Lemma~2.1]{RSselfsimilar}, with the initial dilation $L$ included. The transformed mean-field equation, its drift, and its continuity velocity are respectively
\begin{align}
 \partial_\tau\bar\mu_\tau
 &=\beta^{-1}\Delta\bar\mu_\tau+
       \diver\left(\bar\mu_\tau\left(\frac y2+c(\tau)\nabla\g*\bar\mu_\tau\right)\right),
       \label{eq:selfsimilar-mf}\\
 \bar B_\tau&:=-\frac y2-c(\tau)\nabla\g*\bar\mu_\tau,\qquad
 \bar u_\tau:=\bar B_\tau-\beta^{-1}\nabla\log\bar\mu_\tau.
       \label{eq:selfsimilar-drift-velocity}
\end{align}
The bars distinguish these objects from the measures and fields in physical coordinates.

\begin{thm}[Explicit unconfined classes]\label{thm:unconfined}
Fix $\d\ge2$, $\beta>0$, an initial inverse temperature $\beta_0>\beta/2$, and an integer $m\ge3$. In both cases below, take the initial particle law to be $P_{N,0}=\Q_{N,\beta_0}(\mu_0)$.
\begin{enumerate}[(i)]
\item Suppose $\s=0$, and let $\pi$ denote the density of the thermal equilibrium measure \eqref{eq:thermal} for $V(y)=|y|^2/4$. There are $\delta>0$ and $C<\infty$, depending only on $\d,\beta,\beta_0,m$, such that for every $L>0$ and every $z_0$ of the form \eqref{eq:explicit-class} with $\norm{z_0}_{\Xspace^m}\le\delta$, the initial density
\begin{equation}\label{eq:unconfined-log-data}
 \mu_0(x)=L^{-\d}\frac{\pi(x/L)e^{z_0(x/L)}}{\int\pi e^{z_0}}
\end{equation}
satisfies both \eqref{eq:main-rates} and \eqref{eq:intro-uniform-TV} for the unconfined dynamics. The constant $C$ is independent of $N,k,t,L$ and of the higher derivatives and support of the perturbation.
\item Suppose $0<\s<\d/2$, and set
\begin{equation}\label{eq:selfsimilar-gaussian}
 \gamma(y):=\left(\frac\beta{4\pi}\right)^{\d/2}e^{-\beta|y|^2/4}.
\end{equation}
There are $L_*<\infty$, $\delta>0$, and $C<\infty$, depending only on $\d,\s,\beta,\beta_0,m$, such that for every $L\ge L_*$ and $z_0\in\Xspace^m$ with $\norm{z_0}_{\Xspace^m}\le\delta$, the initial density
\begin{equation}\label{eq:unconfined-riesz-data}
 \mu_0(x)=L^{-\d}\frac{\gamma(x/L)e^{z_0(x/L)}}{\int\gamma e^{z_0}}
\end{equation}
satisfies both \eqref{eq:main-rates} and \eqref{eq:intro-uniform-TV}. The constant $C$ is independent of $N,k,t,L$ within this class.
\end{enumerate}
\end{thm}

The passage from smooth to finite-order initial data in part \textup{(ii)} is justified at the end of the proof, after \eqref{eq:selfsimilar-velocity-integral}.

\begin{remark}[The one-dimensional cases]\label{rem:selfsimilar-one-dimensional}
The change of variables \eqref{eq:selfsimilar} is valid in every dimension. The restriction $\d\ge2$ in \cref{thm:unconfined,lem:decaying-coupling} comes from the locally integrable force estimates in \cref{lem:heat-force-estimates}: in the stated range, $\s+1<\d$. These estimates enter both the background construction below and the first-marginal argument. For $\d=1$ and $\s>0$, we expect that replacing the absolute-force estimates by more refined singular-integral estimates should allow one to treat the one-dimensional Riesz cases in the present Hilbert--Schmidt range $0<\s<1/2$.

The particle-law passage also uses the collision-free solution from \cref{prop:particle-wellposedness}. For $\d=1$ and $\s=0$, the collision-free construction is not available at fixed $\beta$ uniformly in $N$. We expect that an analogue should be accessible by working directly with an appropriate one-dimensional logarithmic particle dynamics through collisions, but we do not pursue this additional issue here.
\end{remark}

For the Riesz case, we construct the background as a perturbation of the Gaussian \eqref{eq:selfsimilar-gaussian}, which is invariant for the noninteracting transformed equation. Large $L$ makes the initial coupling small, and \eqref{eq:selfsimilar-coupling} makes the resulting forcing decay. The following lemma verifies the background hypotheses uniformly in the heat cutoff.

\begin{lemma}[A decaying weak interaction around the Gaussian]\label{lem:decaying-coupling}
Fix $\d\ge2$, $0<\s<\d/2$, $\beta>0$, and an integer $m\ge3$. There exist $\delta_*>0$, $C<\infty$, and $0<\omega<\min\{1/2,\s/2\}$, depending only on these parameters, with the following property. Let
\begin{equation}\label{eq:gaussian-coupling-data}
 c(\tau)=c_0e^{-\s\tau/2},\qquad c_0\ge0,\qquad
 z_0\in\Xspace^m,\qquad \norm{z_0}_{\Xspace^m}+c_0\le\delta_*.
\end{equation}
For $0\le\varepsilon\le1$, with $\g_0:=\g$, equation \eqref{eq:selfsimilar-mf} with $\g$ replaced by $\g_\varepsilon$ and initial density $\gamma e^{z_0}/\int\gamma e^{z_0}$ has a global positive solution $\bar\mu_\tau^\varepsilon$. Its continuity velocity $\bar u_\tau^\varepsilon$, defined by \eqref{eq:selfsimilar-drift-velocity} with the same replacement, satisfies
\begin{equation}\label{eq:gaussian-log-decay}
 \norm{\log(\bar\mu_\tau^\varepsilon/\gamma)}_{\Xspace^m}
   +\norm{\bar u_\tau^\varepsilon}_{W^{1,\infty}}
       \le C(\norm{z_0}_{\Xspace^m}+c_0)e^{-\omega\tau},\qquad \tau\ge0.
\end{equation}
For some $a>0$, the background bounds are
\begin{gather}
 \bar\mu_\tau^\varepsilon(y)\le Ce^{-a|y|^2},\label{eq:gaussian-background-density}\\
 \sup_{\substack{\tau\ge0\\0\le\varepsilon\le1}}
 \left(\norm{\nabla^2\log\bar\mu_\tau^\varepsilon}_{L^\infty}
 +\norm{\nabla\bar B_\tau^\varepsilon}_{L^\infty}
 +\norm{\nabla^2\bar B_\tau^\varepsilon}_{L^\infty}\right)\le C,
 \label{eq:gaussian-background-derivatives}\\
 \Var_{\bar\mu_\tau^\varepsilon}(\varphi)\le\frac4\beta
       \int|\nabla\varphi|^2\dd\bar\mu_\tau^\varepsilon,
       \qquad \varphi\in C_c^\infty(\R^\d).
 \label{eq:gaussian-background-poincare}
\end{gather}
The constants are independent of $\tau,\varepsilon,c_0,z_0$ under \eqref{eq:gaussian-coupling-data}. For every $T<\infty$,
\begin{equation}\label{eq:gaussian-cutoff-convergence}
 \sup_{0\le\tau\le T}\left(
 \norm{\log(\bar\mu_\tau^\varepsilon/\bar\mu_\tau^0)}_{\Xspace^m}
 +\norm{\bar u_\tau^\varepsilon-\bar u_\tau^0}_{W^{1,\infty}}\right)
 \xrightarrow[\varepsilon\downarrow0]{}0.
\end{equation}
If $z_0$ has finite higher norms, they remain finite on compact time intervals; the solutions are smooth for smooth initial data.
\end{lemma}
\begin{proof}
\parheading{The logarithmic equation.}
Fix $0\le\varepsilon\le1$, and write $\mu_z:=\gamma e^z/\int\gamma e^z$. We solve for $z$ modulo constants in the space \eqref{eq:X-norm}. The logarithmic equation is
\begin{equation}\label{eq:gaussian-log-PDE}
 \partial_\tau z=(\beta^{-1}\Delta-\tfrac12y\cdot\nabla)z
       +\beta^{-1}|\nabla z|^2+c(\tau)\mathcal F_\varepsilon(z),
\end{equation}
where
\begin{equation}\label{eq:coupling-nonlinearity}
 \mathcal F_\varepsilon(z):=\diver(\nabla\g_\varepsilon*\mu_z)
            -\frac\beta2 y\cdot(\nabla\g_\varepsilon*\mu_z)
            +(\nabla\g_\varepsilon*\mu_z)\cdot\nabla z.
\end{equation}
The Gaussian versions of \eqref{eq:normalized-first-remainder} and \eqref{eq:normalized-density-directional-derivative}, together with \eqref{eq:x-convolution}, give on a fixed small ball
\begin{equation}\label{eq:gaussian-forcing-bounds}
 \sup_{0\le\varepsilon\le1}\norm{\mathcal F_\varepsilon(0)}_{\Xspace^{m-1}}\le C,
 \qquad
 \sup_{0\le\varepsilon\le1}\norm{\mathcal F_\varepsilon(z)-\mathcal F_\varepsilon(w)}_{\Xspace^{m-1}}
 \le C\norm{z-w}_{\Xspace^m}.
\end{equation}
For the divergence and the term containing $y$, the identities used in \eqref{eq:gaussian-forcing-bounds} are
\begin{align}
 \diver(\nabla\g_\varepsilon*\mu_z)
 &=\sum_{i=1}^{\d}(\partial_i\g_\varepsilon)*(\partial_i\mu_z),
 \label{eq:gaussian-divergence-convolution}\\
 y\cdot(\nabla\g_\varepsilon*\mu_z)(y)
 &=\sum_{i=1}^{\d}\left[((x_i\partial_i\g_\varepsilon)*\mu_z)(y)
       +\bigl(\partial_i\g_\varepsilon*(x_i\mu_z)\bigr)(y)\right].
 \label{eq:gaussian-coordinate-convolution}
\end{align}
The first identity transfers the divergence to the density, so derivatives through order $m-1$ require derivatives of $\mu_z$ only through order $m$. The second identity places the factor $y$ either on the kernel or on the density, as in \eqref{eq:x-convolution}. Thus, \cref{lem:heat-force-estimates} gives the one-derivative loss in \eqref{eq:gaussian-forcing-bounds}. Membership in the completion $\Xspace^{m-1}$ also requires the decay
\begin{equation}\label{eq:gaussian-forcing-vanishing}
 \lim_{R\to\infty}\sup_{\substack{0\le\varepsilon\le1\\ \norm z_{\Xspace^m}\le\delta_*}}
       \sup_{|y|\ge R}\sum_{j=1}^{m-1}|\nabla^j\mathcal F_\varepsilon(z)(y)|=0.
\end{equation}
Indeed, derivatives of $\mu_z$ through order $m$, and their first spatial moments, have uniform Gaussian upper bounds on this ball. The kernels $\nabla\g_\varepsilon(x)$ and $x_i\partial_i\g_\varepsilon(x)$ are locally integrable and tend to zero as $|x|\to\infty$, uniformly in $\varepsilon$, since $\s>0$. Splitting each integral in the convolution variable $x$ into $|x|\le|y|/2$ and its complement proves \eqref{eq:gaussian-forcing-vanishing}.

Normalization is checked by the identities
\begin{align}
 \mu_z\mathcal F_\varepsilon(z)
 &=\diver\bigl(\mu_z\nabla\g_\varepsilon*\mu_z\bigr),
 \label{eq:gaussian-forcing-density-identity}\\
 \mu_z\left[(\beta^{-1}\Delta-\tfrac12y\cdot\nabla)z+
             \beta^{-1}|\nabla z|^2\right]
 &=\beta^{-1}\Delta\mu_z+\tfrac12\diver(y\mu_z).
 \label{eq:gaussian-OU-density-identity}
\end{align}
Both right-hand sides have integral zero. Choosing the representative centered with respect to $\gamma$ and differentiating $\log\int e^z\dd\gamma$ therefore cancels the centering constant, exactly as in \eqref{eq:centered-log-equation}--\eqref{eq:log-normalizer-derivative}. Hence a solution of \eqref{eq:gaussian-log-PDE} in the quotient gives the normalized solution of \eqref{eq:selfsimilar-mf}.

\parheading{Decay and cutoff convergence.}
Choose $0<\omega<\min\{1/2,\s/2\}$, and equip trajectories with the norm
\begin{equation}\label{eq:gaussian-trajectory-norm}
 \norm z_{\omega,m}:=\sup_{\tau\ge0}e^{\omega\tau}\norm{z_\tau}_{\Xspace^m}.
\end{equation}
In this proof, write $P_\tau$ for the Ornstein--Uhlenbeck semigroup \eqref{eq:OU-formula} with $A=\Id/2$, whose generator is $\beta^{-1}\Delta-\tfrac12y\cdot\nabla$. The Duhamel map is
\begin{equation}\label{eq:gaussian-Duhamel}
 (\mathcal T_\varepsilon z)_\tau
 :=P_\tau z_0+
   \int_0^\tau P_{\tau-r}
       \left[\beta^{-1}|\nabla z_r|^2+c_0e^{-\s r/2}\mathcal F_\varepsilon(z_r)\right]\dd r.
\end{equation}
The Ornstein--Uhlenbeck estimates \eqref{eq:OU-estimates} and \eqref{eq:gaussian-forcing-bounds} imply, for $\norm z_{\omega,m},\norm w_{\omega,m}\le R$,
\begin{align}
 \norm{\mathcal T_\varepsilon z}_{\omega,m}
 &\le C\norm{z_0}_{\Xspace^m}+C(c_0+R^2+c_0R),
       \label{eq:gaussian-Duhamel-ball}\\
 \norm{\mathcal T_\varepsilon z-\mathcal T_\varepsilon w}_{\omega,m}
 &\le C(R+c_0)\norm{z-w}_{\omega,m}.
       \label{eq:gaussian-Duhamel-contraction}
\end{align}
The inhomogeneous term $c(\tau)\mathcal F_\varepsilon(0)$ is controlled by
\begin{equation}\label{eq:gaussian-time-integrability}
 \sup_{\tau\ge0}e^{\omega\tau}\int_0^\tau(1+(\tau-r)^{-1/2})e^{-(\tau-r)/2-\s r/2}\dd r
 \le\int_0^\infty(1+r^{-1/2})e^{-(1/2-\omega)r}\dd r<\infty.
\end{equation}
Here, $\omega<\s/2$ bounds the remaining factor $e^{-(\s/2-\omega)r}$ by one. The quadratic and Lipschitz terms in \eqref{eq:gaussian-Duhamel-ball}--\eqref{eq:gaussian-Duhamel-contraction} have the same integrable time kernel.
The constant is independent of the heat cutoff. Choose $R=2C(\norm{z_0}_{\Xspace^m}+c_0)$ and then $\delta_*$ small enough that \eqref{eq:gaussian-Duhamel-ball} maps this ball into itself and the coefficient in \eqref{eq:gaussian-Duhamel-contraction} is at most $1/2$. The fixed point proves global existence and the logarithmic-density bound in \eqref{eq:gaussian-log-decay}.

For every integer $\ell>m$, the differentiated product and convolution estimates used in \eqref{eq:higher-tame} give
\begin{equation}\label{eq:gaussian-higher-tame}
 \norm{\beta^{-1}|\nabla z|^2+c(\tau)\mathcal F_\varepsilon(z)}_{\Xspace^{\ell-1}}
 \le C_\ell\left[(\norm z_{\Xspace^m}+c(\tau))\norm z_{\Xspace^\ell}+c(\tau)\right].
\end{equation}
The short-interval argument \eqref{eq:high-norm-continuation}, now with the displayed inhomogeneous term, propagates finite higher norms on every finite interval. Their size does not enter the low-norm contraction.

The force-difference estimates \eqref{eq:logmaps-cutoff-force-L1}--\eqref{eq:logmaps-cutoff-force-weighted} and \eqref{eq:x-convolution} give
\begin{equation}\label{eq:gaussian-forcing-cutoff}
 \sup_{\norm z_{\Xspace^m}\le R}
 \norm{\mathcal F_\varepsilon(z)-\mathcal F_0(z)}_{\Xspace^{m-1}}
 \le C_R\left(\varepsilon^{(\d-\s-1)/2}+\varepsilon^{(\d-\s)/2}\right)
 \xrightarrow[\varepsilon\downarrow0]{}0.
\end{equation}
Let $z^\varepsilon$ and $z^0$ be the trajectories constructed with $\g_\varepsilon$ and $\g$, respectively, from the same initial datum $z_0$. Thus,
\begin{equation}\label{eq:gaussian-fixed-point-pair}
 z^\varepsilon=\mathcal T_\varepsilon z^\varepsilon,\qquad
 z^0=\mathcal T_0z^0,\qquad
 z^\varepsilon|_{\tau=0}=z^0|_{\tau=0}=z_0.
\end{equation}
Subtracting these identities and using the contraction coefficient $1/2$ in \eqref{eq:gaussian-Duhamel-contraction} gives
\begin{equation}\label{eq:gaussian-fixed-point-comparison}
 \norm{z^\varepsilon-z^0}_{\omega,m}
 \le2\norm{\mathcal T_\varepsilon z^0-\mathcal T_0z^0}_{\omega,m}
 \le Cc_0\left(\varepsilon^{(\d-\s-1)/2}+\varepsilon^{(\d-\s)/2}\right)
 \xrightarrow[\varepsilon\downarrow0]{}0.
\end{equation}
The last inequality follows from \eqref{eq:gaussian-forcing-cutoff} and the integrable kernel in \eqref{eq:gaussian-time-integrability}. This proves convergence in \eqref{eq:gaussian-cutoff-convergence} for the logarithmic densities. Approximation in $\Xspace^m$ gives the same construction for finite-order initial data; for smooth data, the higher-norm comparison gives smooth solutions.

\parheading{The background bounds.}
Choose the representative satisfying $\int z_\tau^\varepsilon\dd\gamma=0$. By \eqref{eq:linear-growth-representative}, its absolute value is bounded by $C\norm{z_\tau^\varepsilon}_{\Xspace^m}(1+|y|)$. Jensen's inequality bounds the normalizing integral from below by one. Completing the square therefore gives \eqref{eq:gaussian-background-density}; differentiating the normalized density gives common Gaussian bounds for its derivatives through order $m$.

The force estimate \eqref{eq:thermal-force-split}, with derivatives placed on the density, now bounds $\nabla\g_\varepsilon*\bar\mu_\tau^\varepsilon$ and its first two derivatives uniformly in $\tau,\varepsilon$. The continuity velocity is
\begin{equation}\label{eq:gaussian-velocity}
 \bar u_\tau^\varepsilon=-\beta^{-1}\nabla z_\tau^\varepsilon
        -c(\tau)\nabla\g_\varepsilon*\bar\mu_\tau^\varepsilon.
\end{equation}
The logarithmic-density bound already proved and $c(\tau)=c_0e^{-\s\tau/2}$ give the velocity estimate in \eqref{eq:gaussian-log-decay}. For the convergence of the velocities, subtract \eqref{eq:gaussian-velocity} at cutoffs $\varepsilon$ and zero. The force-difference bound \eqref{eq:logmaps-cutoff-force-L1} and the Lipschitz estimate for the normalized density give, for every $T<\infty$,
\begin{equation}\label{eq:gaussian-velocity-cutoff-comparison}
 \sup_{\tau\in[0,T]}\norm{\bar u_\tau^\varepsilon-\bar u_\tau^0}_{W^{1,\infty}}
 \le C\left((1+c_0)\sup_{\tau\in[0,T]}\norm{z_\tau^\varepsilon-z_\tau^0}_{\Xspace^m}
       +c_0\varepsilon^{(\d-\s-1)/2}\right)\xrightarrow[\varepsilon\downarrow0]{}0.
\end{equation}
The limit follows from \eqref{eq:gaussian-fixed-point-comparison}. This completes \eqref{eq:gaussian-cutoff-convergence}. Moreover, decreasing $\delta_*$ if necessary,
\begin{equation}\label{eq:gaussian-curvature}
 \nabla^2(-\log\bar\mu_\tau^\varepsilon)
 =\frac\beta2\Id-\nabla^2z_\tau^\varepsilon
 \ge\frac\beta4\Id,\qquad \tau\ge0,\quad0\le\varepsilon\le1.
\end{equation}
The curvature criterion \cite[Section~4.8]{BGL} gives the Poincar\'e inequality \eqref{eq:gaussian-background-poincare}. Finally, \eqref{eq:selfsimilar-drift-velocity}, \eqref{eq:gaussian-curvature}, and the convolution bounds give \eqref{eq:gaussian-background-derivatives}.
\end{proof}

\begin{proof}[Proof of \cref{thm:unconfined}]
\parheading{The initial Gibbs measures.}
At $\tau=0$, write $X_N=LY_N$ and $\mu_0=(D_0)_\#\bar\mu_0$. Expanding the three terms in \eqref{eq:modulated} gives the modulated-energy transformation identity (cf.\ \cite[Lemma~2.3]{RSselfsimilar})
\begin{equation}\label{eq:selfsimilar-energy-scaling}
 N\Fr_N(LY_N,\mu_0)=
 \begin{cases}
 N\Fr_N(Y_N,\bar\mu_0)+\tfrac12\log L,&\s=0,\\
 L^{-\s}N\Fr_N(Y_N,\bar\mu_0),&\s>0.
 \end{cases}
\end{equation}
Thus,
\begin{equation}\label{eq:selfsimilar-partition-scaling}
 \K_{N,\beta_0}(\mu_0)=
 \begin{cases}
 L^{-\beta_0/2}\K_{N,\beta_0}(\bar\mu_0),&\s=0,\\
 \K_{N,\beta_0L^{-\s}}(\bar\mu_0),&\s>0.
 \end{cases}
\end{equation}
The deterministic logarithmic factor cancels between the Gibbs weight and its normalizing constant. In both cases,
\begin{equation}\label{eq:selfsimilar-preparation}
 \bar P_{N,0}=\Q_{N,\beta_0c(0)}(\bar\mu_0).
\end{equation}
For $\s=0$, the transformed dynamics have coupling one and confinement $|y|^2/4$, so \cref{thm:uniform} gives both marginal bounds for every $\tau\ge0$. Neither the transformed equation nor its initial density depends on $L$.

\parheading{The Riesz interaction.}
Set $c_0=L^{-\s}$. Choose $L_*\ge1$ and $\delta>0$ initially so that $\delta+L_*^{-\s}\le\delta_*$ in \cref{lem:decaying-coupling}. Its background bounds and \cref{prop:static-inputs} give static constants and a commutator coefficient $\kappa>0$ uniform over $\tau\ge0$, $0\le\varepsilon\le1$, $0\le c_0\le L_*^{-\s}$, and $\norm{z_0}_{\Xspace^m}\le\delta$. Fix $b>0$ as in \eqref{eq:lambda-coupling}, and write $\dot c:=\mathrm dc/\mathrm d\tau$. Since
\begin{equation}\label{eq:selfsimilar-coupling-variation}
 \dot c(\tau)=-\frac\s2c(\tau),\qquad
 \int_0^\infty|\dot c(\tau)|\dd\tau=c_0,
\end{equation}
the coefficient in \eqref{eq:lambda-coupling} has the common integrable bound
\begin{equation}\label{eq:selfsimilar-expenditure-bound}
 \int_0^\infty\left(\frac\beta\kappa\sup_{0\le\varepsilon\le1}\norm{\nabla\bar u_\tau^\varepsilon}_{L^\infty}
       +\frac\beta b|\dot c(\tau)|\right)\dd\tau
 \le\frac{\beta C}{\kappa\omega}(\delta+L_*^{-\s})+\frac\beta bL_*^{-\s}.
\end{equation}
Here, we used \eqref{eq:gaussian-log-decay}. The second term comes from the derivative of the coupling in the moving-reference defect \eqref{eq:defect-coupling}.

Choose a finite $p_0>2$ such that
\begin{equation}\label{eq:selfsimilar-initial-exponent}
 1-\frac1{p_0}<\min\left\{1,\frac{\beta_0}{\beta}\right\}.
\end{equation}
This is possible because $\beta_0>\beta/2$. Increase $L_*$ and decrease $\delta$ so that the right-hand side of \eqref{eq:selfsimilar-expenditure-bound} is strictly less than
\begin{equation}\label{eq:selfsimilar-expenditure-condition}
 \log\left(2\left(1-\frac1{p_0}\right)\right).
\end{equation}
These choices depend only on the fixed parameters in the theorem. The initial comparison in \cref{lem:initial-temperature}, with $b_0=\beta_0c_0$ and $b_1=\beta c_0$, uses the nonnegative coupling
\begin{equation}\label{eq:selfsimilar-comparison-coupling}
 p_0b_0-(p_0-1)b_1=c_0\bigl(p_0\beta_0-(p_0-1)\beta\bigr)\ge0.
\end{equation}
The comparison is uniform for $0\le c_0\le L_*^{-\s}$, since all three coupling parameters in \cref{lem:initial-temperature} lie in a fixed bounded nonnegative interval. The strict margin chosen in \eqref{eq:selfsimilar-expenditure-bound}--\eqref{eq:selfsimilar-expenditure-condition} is common to the same initial-data class. Consequently, \cref{cor:full-moment} gives a full second moment uniform in $N$, transformed time, the heat cutoff, and $L\ge L_*$.

Apply \cref{prop:initial-mark} at the initial coupling $\beta_0c_0$. Also, \eqref{eq:gaussian-log-decay} gives
\begin{equation}\label{eq:selfsimilar-velocity-integral}
 \sup_{0\le\varepsilon\le1}\int_0^\infty
      \norm{\nabla\bar u_\tau^\varepsilon}_{L^\infty}\dd\tau
 \le\frac C\omega(\norm{z_0}_{\Xspace^m}+c_0).
\end{equation}
Together with the Poincar\'e and derivative bounds of \cref{lem:decaying-coupling}, these estimates verify the hypotheses of \cref{prop:uniform-mark}. The term containing $\dot c$ in \eqref{eq:homogeneous-energy-identity} is nonpositive, so the homogeneous contraction estimate applies without an additional error. \Cref{lem:product-projection} then gives the marginal bounds.

For each fixed $N$, $L$, and finite transformed-time interval, case \textup{(iii)} of \cref{lem:approximation} passes these bounds to the singular particle system. That system is the image under \eqref{eq:selfsimilar} of the unconfined strong solution in \cref{prop:particle-wellposedness}. The localization is performed at fixed $N,L$; the marginal estimates have constants independent of both parameters. For finite-order initial data, approximate $z_0$ in $\Xspace^m$ by the smooth defining class and use the density, initial-measure, and path-localization arguments of \cref{lem:approximation}. The constants depend only on the common low-norm bound.

\parheading{Return to physical coordinates.}
For $t=L^2(e^\tau-1)$, \eqref{eq:selfsimilar-measures} gives
\begin{equation}\label{eq:selfsimilar-RN-identity}
 \frac{\mathrm d\bar P_{N,\tau}^{(k)}}{\mathrm d\bar\mu_\tau^{\otimes k}}(Y_k)
 =\frac{\mathrm dP_{N,t}^{(k)}}{\mathrm d\mu_t^{\otimes k}}
       \bigl(D_\tau^{\otimes k}Y_k\bigr).
\end{equation}
Changing variables in the defining integrals therefore yields
\begin{align}
 \chi^2(P_{N,t}^{(k)}\mid\mu_t^{\otimes k})
 &=\chi^2(\bar P_{N,\tau}^{(k)}\mid\bar\mu_\tau^{\otimes k}),
 \label{eq:selfsimilar-chi2-invariance}\\
 \Hr(P_{N,t}^{(k)}\mid\mu_t^{\otimes k})
 &=\Hr(\bar P_{N,\tau}^{(k)}\mid\bar\mu_\tau^{\otimes k}),
 \label{eq:selfsimilar-entropy-invariance}\\
 \TV(P_{N,t}^{(k)},\mu_t^{\otimes k})
 &=\TV(\bar P_{N,\tau}^{(k)},\bar\mu_\tau^{\otimes k}).
 \label{eq:selfsimilar-TV-invariance}
\end{align}
For relative entropy, this is \cite[Lemma~2.2]{RSselfsimilar}; the same Radon--Nikodym identity gives the other two equalities. For each fixed $L>0$, the time change $t=L^2(e^\tau-1)$ is a bijection of $[0,\infty)$ onto itself. Hence \eqref{eq:selfsimilar-chi2-invariance} gives
\begin{equation}\label{eq:selfsimilar-time-supremum}
 \sup_{t\ge0}\chi^2(P_{N,t}^{(k)}\mid\mu_t^{\otimes k})
 =\sup_{\tau\ge0}\chi^2(\bar P_{N,\tau}^{(k)}\mid\bar\mu_\tau^{\otimes k})
 \le C\frac{k^2}{N^2}.
\end{equation}
The entropy and total-variation identities give the corresponding time-uniform bounds. Thus, both assertions hold with the same constants as in the transformed estimates, independently of $L$ in the respective classes.
\end{proof}

\par\endgroup

\begingroup

\section{Failure of unrestricted full-particle chi-square control}\label{sec:periodic-counterexample-proof}
We prove \cref{thm:periodic-counterexample}. The construction has two parts. First, a mean-field density depending on one coordinate allows an explicit comparison between the product variance of a terminal test and the Brownian contribution to its fluctuation variance; see \cref{ce:prop:noisecross}. Second, an entropy estimate controls the nonlinear remainder in the empirical fluctuation equation, and \cref{ce:prop:noiseobstruction} converts the variance comparison into divergence of the full-particle chi-square. The singular particle law and the identities used in this argument are justified in \cref{ce:sec:singular}.

For the calculations in this section, it is convenient to use the angular variables $x=2\pi z$ and $\tau=4t$ for the periodic particle and mean-field equations specified in \cref{sec:intro-counterexample}. We reuse $\g$, $\mu$, $P_N$, $\Fr_N$, $\K_{N,b}$, and $\Q_{N,b}$ for the transformed kernel, measures, energy, partition function, and Gibbs law. These local conventions end with this section. The partition and static fluctuation estimates used below are those of \cite{DGRpartition}.

\subsection{An exact invariant stripe reduction}\label{ce:sec:stripe}

Write $\T:=\R/(2\pi\Z)$, and denote normalized Haar measure by $m$, with the dimension determined by the domain of integration. In these variables,
\begin{equation}\label{ce:eq:angkernel}
 \widehat{\g}(k):=\frac{\mathbf 1_{\{k\ne0\}}}{2|k|^4},\qquad D:=\frac{\pi^2}{\beta},\qquad
 \dd m(x):=(2\pi)^{-4}\dd x\quad\text{on }\T^4.
\end{equation}
The kernel is $\pi^2$ times the pullback of the kernel in \eqref{eq:periodic-ce-kernel}. The particle equation and generator are
\begin{align}
 \dd x_i(\tau)&=-\frac1N\sum_{j\ne i}\nabla\g(x_i(\tau)-x_j(\tau))\dd\tau
 +\sqrt{2D}\,\dd B_i(\tau),\label{ce:eq:angular-sde}\\
 L_N&:=D\sum_{i=1}^N\Delta_i-\frac1N\sum_{i\ne j}\nabla\g(x_i-x_j)\cdot\nabla_i.\label{ce:eq:angular-generator}
\end{align}
We use \eqref{eq:background}, \eqref{eq:modulated} and \eqref{eq:gibbs} without introducing a second set of definitions. The coupling in the initial Gibbs law becomes $D^{-1}$: multiplication of the modulated energy by $\pi^2$ is canceled by $D^{-1}=\beta/\pi^2$. Thus the transformed initial law is $\Q_{N,D^{-1}}(\mu_0)$.

On the circle, $\dd m(x):=\dd x/(2\pi)$. We call a density of the form $\rho(x)=\rho_1(x_1)$ a stripe density: it is constant along the three-dimensional tori $\{x_1=\text{constant}\}$. Its convolution with $\g$ also depends only on $x_1$. Integrating $\g$ in the other three coordinates keeps precisely its Fourier modes $(n,0,0,0)$, giving the circle kernel
\begin{equation}\label{ce:eq:W}
 W(x)=\sum_{n=1}^\infty\frac{\cos(nx)}{n^4}.
\end{equation}
The stripe equation is exactly
\begin{equation}\label{ce:eq:circleMF}
 \partial_\tau\rho=D\rho''+\partial_x(\rho W'*\rho),\qquad
 b_\tau=-W'*\rho_\tau,
\end{equation}
where convolution and densities are relative to $m$. For terminal data depending only on $x_1$, the solution of the backward linearized adjoint equation also depends only on $x_1$. The reduced operator is
\begin{equation}\label{ce:eq:circleadj}
 \mathcal L_\tau f(x)
 =Df''(x)-\int W'(x-y)\bigl(f'(x)-f'(y)\bigr)\rho_\tau(y)m(\dd y).
\end{equation}
The same coordinate integration gives \eqref{ce:eq:circleadj} from the full linearized adjoint. Fubini applies because $\g$ and $\nabla\g$ are locally integrable. We use the same symbol $\rho$ for the circle density and its extension $x\mapsto\rho(x_1)$ to $\T^4$; the domain of integration specifies the dimension of $m$.

On $[0,2\pi]$,
\begin{equation}\label{ce:eq:Wpoly}
 W(x)=\frac{\pi^4}{90}-\frac{\pi^2x^2}{12}
 +\frac{\pi x^3}{12}-\frac{x^4}{48}.
\end{equation}
Thus, $W\in C^{2,1}$ on the circle, $W$ is even, and
\begin{equation}\label{ce:eq:curvatures}
 W'(0)=W'(\pi)=0,\qquad
 W''(0)=-a,\quad W''(\pi)=b,
 \qquad a=\frac{\pi^2}{6},\quad b=\frac{\pi^2}{12}.
\end{equation}
The negative curvature at zero expands an isolated packet; the positive curvature at the antipode can confine a packet through the mass in the other packet.

\subsection{Two Gaussian packets with unequal masses}\label{ce:sec:packet}

Let $k_\delta$ be the periodized centered Gaussian of variance $\delta^2$, written as a density relative to $m$:
\begin{equation}\label{ce:eq:bump}
 k_\delta(x)=\frac{\sqrt{2\pi}}{\delta}
 \sum_{j\in\mathbb Z}\exp\!\left(-\frac{(x+2\pi j)^2}{2\delta^2}\right).
\end{equation}
It is smooth, even, strictly positive, and has integral one. Set
\begin{equation}\label{ce:eq:rhoinitial}
 p=\frac16,\qquad q=\frac56,\qquad
 \rho^\delta_0(x)=pk_\delta(x)+qk_\delta(x-\pi),
 \qquad D_\delta=\kappa\delta^2.
\end{equation}
The constant $\kappa$ is given in \eqref{ce:eq:kappa} below. Write $\rho^\delta_\tau$ for its solution of \eqref{ce:eq:circleMF}. For every fixed $\delta>0$, this is a smooth strictly positive trajectory on each compact time interval. Indeed, the drift is uniformly Lipschitz with constant $\norm{W''}_\infty$, independent of the density. The nonlinear diffusion is globally well posed by coupling, positivity follows from the maximum principle, and positive diffusion and convolution bootstrap the smooth initial density. All smooth norms needed below are finite for fixed $\delta$; they need not be uniform as $\delta\downarrow0$.

The initial law is invariant under $x\mapsto-x$. Uniqueness preserves this symmetry, and periodicity also makes it even about $\pi$. Therefore,
\begin{equation}\label{ce:eq:centersfixed}
 b^\delta_\tau(0)=b^\delta_\tau(\pi)=0.
\end{equation}
The two components start near $0$ and $\pi$. The drift generated by $p\delta_0+q\delta_\pi$ has derivatives $-\kappa$ at $0$ and $\kappa_b$ at $\pi$, where
\begin{equation}\label{ce:eq:kappa}
 \kappa:=qb-pa=\frac{\pi^2}{24}>0,
 \qquad \kappa_b:=qa-pb=\frac{\pi^2}{8}>0.
\end{equation}
Condition on which component of the initial mixture was chosen. Let $Y_a^\delta$ be the real-valued displacement from $0$ for the component of mass $p$, and let $Y_b^\delta$ be the displacement from $\pi$ for the component of mass $q$. Lift the periodic drift to $\R$ and solve
\begin{align}
 \dd Y_a^\delta(\tau)&=b_\tau^\delta(Y_a^\delta(\tau))\dd\tau+\sqrt{2\kappa}\,\delta\dd B_a(\tau),\label{ce:eq:minority-lift}\\
 \dd Y_b^\delta(\tau)&=b_\tau^\delta(\pi+Y_b^\delta(\tau))\dd\tau+\sqrt{2\kappa}\,\delta\dd B_b(\tau),\label{ce:eq:majority-lift}
\end{align}
with centered Gaussian initial displacements of variance $\delta^2$, independent of the Brownian motions. Modulo $2\pi$, their laws, with weights $p$ and $q$, recover $\rho^\delta_\tau m$. 

\begin{lemma}[Narrow-packet mean-field limit]\label{ce:lem:packet}
For every fixed $S<\infty$, as $\delta\downarrow0$, the measures $\rho^\delta_\tau m$ converge uniformly for $0\le\tau\le S$ in $\mathsf W_1$ to
\begin{equation}\label{ce:eq:atomic}
 \rho_*=p\delta_0+q\delta_\pi.
\end{equation}
The functions $b^\delta$ and $(b^\delta)'$ converge uniformly on $[0,S]\times\T$ to $b_*:=-W'*\rho_*$ and $b_*'$. With a common initial Gaussian and Brownian motion, $Y_a^\delta/\delta$ converges in $L^2(\Omega;C([0,S]))$ to the solution of
\begin{equation}\label{ce:eq:rareOU}
 \dd Z_\tau=-\kappa Z_\tau\dd\tau+\sqrt{2\kappa}\dd B_\tau,
 \qquad Z_0\sim\mathcal N(0,1).
\end{equation}
In particular,
\begin{equation}\label{ce:eq:v}
 \E[Z_\tau^2]=1.
\end{equation}
\end{lemma}
\begin{proof}
By \eqref{ce:eq:centersfixed},
\begin{equation}\label{ce:eq:lift-drift-bound}
 |b^\delta_\tau(y)|+|b^\delta_\tau(\pi+y)|\le2\norm{W''}_{L^\infty}|y|,
 \qquad y\in\R.
\end{equation}
Apply Doob's maximal inequality (see, e.g., \cite[Chapter~II, Section~1]{RevuzYor}) and Gronwall to \eqref{ce:eq:minority-lift} and \eqref{ce:eq:majority-lift}. For every finite $r\ge2$, this gives
\begin{equation}\label{ce:eq:moments}
 \E\!\left[\sup_{\tau\le S}|Y_a^\delta(\tau)|^r
 +\sup_{\tau\le S}|Y_b^\delta(\tau)|^r\right]
 \le C_{S,r}\delta^r.
\end{equation}
Here, $C_{S,r}$ is independent of $0<\delta\le1$. Consequently,
\begin{equation}\label{ce:eq:W1}
 \sup_{\tau\le S}\mathsf W_1(\rho^\delta_\tau m,\rho_*)\le C_S\delta.
\end{equation}
Because $W'$ and $W''$ are Lipschitz, this implies uniform convergence of $b^\delta$ and $(b^\delta)'$, with error $O_S(\delta)$.

At zero, $b_*'(0)=pa-qb=-\kappa$. We have that
\begin{equation}\label{ce:eq:drifterror}
 \left|\delta^{-1}b^\delta_\tau(\delta z)+\kappa z\right|
 \le C_S\delta|z|+C\delta|z|^2,\qquad \forall\,0\le\tau\le S,\ z\in\R.
\end{equation}
For large $z$, the same estimate follows from the Lipschitz bound after increasing $C$. Construct $Y_a^\delta$ and $Z$ with $Y_a^\delta(0)=\delta Z_0$ and $B_a=B$ in \eqref{ce:eq:minority-lift} and \eqref{ce:eq:rareOU}. The Brownian terms then cancel in the equation for $\delta^{-1}Y_a^\delta-Z$. Using \eqref{ce:eq:moments} with $r=4$ in the resulting integral equation yields
\begin{equation}\label{ce:eq:minority-strong-limit}
 \E\left[\sup_{0\le\tau\le S}|\delta^{-1}Y_a^\delta(\tau)-Z_\tau|^2\right]
 \le C_S\delta^2.
\end{equation}
The Ornstein--Uhlenbeck variance is
\begin{equation}\label{ce:eq:OU-variance-calculation}
 \E[Z_\tau^2]=e^{-2\kappa\tau}+2\kappa\int_0^\tau e^{-2\kappa(\tau-r)}\dd r=1,
\end{equation}
which proves \eqref{ce:eq:v}.
\end{proof}

Choose an even $\chi\in C_c^\infty((-\pi/4,\pi/4))$ equal to one on $(-\pi/8,\pi/8)$, and extend $\psi(x):=x\chi(x)$ periodically. Then $\psi$ is smooth and odd, equals $x$ near zero, and vanishes near $\pi$; in particular,
\begin{equation}\label{ce:eq:terminaltest}
 \psi'(0)=1,\qquad \psi'(\pi)=0,\qquad \rho^\delta_\tau(\psi)=0.
\end{equation}
The next variance is the variance of $N^{-1/2}\sum_i\psi(x_i)$ for independent $x_i$ with law $\rho_S^\delta m$. It is therefore the comparison variance for the product reference in \cref{ce:prop:noiseobstruction}.

\begin{cor}[Product variance of the terminal test]\label{ce:cor:productvariance}
For every fixed $S<\infty$, as $\delta\downarrow0$,
\begin{equation}\label{ce:eq:prodvar}
 \Var_{\rho^\delta_Sm}(\psi)
 =p\delta^2+o(\delta^2).
\end{equation}
\end{cor}
\begin{proof}
For the component of mass $p$ centered at zero, $\psi(Y_a^\delta)=Y_a^\delta$ until the displacement leaves a fixed neighborhood. The fourth-moment estimate in \eqref{ce:eq:moments} makes the error in the second moment $O_S(\delta^4)$. The same bound makes the contribution of the component centered at $\pi$ $O_S(\delta^4)$ because $\psi$ vanishes near $\pi$. Symmetry makes the mean zero. The $L^2$ convergence in \cref{ce:lem:packet} and $\E[Z_S^2]=1$ therefore give \eqref{ce:eq:prodvar}.
\end{proof}

\subsection{The backward test and the Brownian covariance}\label{ce:sec:adjoint}

Let $f^\delta_r$ solve
\begin{equation}\label{ce:eq:backward}
 \partial_r f^\delta_r+\mathcal L^\delta_r f^\delta_r=0,
 \qquad f^\delta_S=\psi.
\end{equation}
For each fixed $\delta>0$, the solution is classical and smooth. Set $g^\delta_r:=(f^\delta_r)'$. Differentiating \eqref{ce:eq:circleadj} gives
\begin{equation}\label{ce:eq:gradback}
 -\partial_r g^\delta_r
 =\kappa\delta^2(g^\delta_r)''+b^\delta_r(g^\delta_r)'
 +(b^\delta_r)'g^\delta_r
 +\int W''(x-y)g^\delta_r(y)\rho^\delta_r(y)m(\dd y).
\end{equation}


\begin{lemma}[Convergence of the adjoint gradients]\label{ce:lem:adjointlimit}
The functions $g^\delta$ converge uniformly on $[0,S]\times\T$ to the bounded Lipschitz mild solution of
\begin{equation}\label{ce:eq:limiting-gradient-adjoint}
 -\partial_r g^*_r(x)=b_*(x)\partial_xg^*_r(x)+b_*'(x)g^*_r(x)
 +\int_\T W''(x-y)g^*_r(y)\dd\rho_*(y),\qquad g^*_S=\psi'.
\end{equation}
Here, the limit is taken as $\delta\downarrow0$, so the diffusion coefficient $D_\delta=\kappa\delta^2$ vanishes. At $0$ and $\pi$, 
\begin{equation}\label{ce:eq:gradcenters}
 g^*_r(0)=p+q e^{-b(S-r)},\qquad
 g^*_r(\pi)=p(1-e^{-b(S-r)}).
\end{equation}
\end{lemma}
\begin{proof}
Reverse time in \eqref{ce:eq:gradback}. The Markov representation for the transport-diffusion equation gives contraction on bounded functions; see, e.g., \cite[Sections~7.1 and~8.1]{Oksendal}. Synchronous coupling gives the Lipschitz bound $e^{L(t-s)}$, where $L=\norm{W''}_\infty$. The other terms in \eqref{ce:eq:gradback} are bounded operators on bounded functions. Multiplication by $(b^\delta)'$ has a Lipschitz bound independent of $\delta$, because $W''$ is Lipschitz; the integral operator maps bounded functions to functions with Lipschitz norm bounded by $\norm{W'''}_{L^\infty}\norm g_\infty$.

Duhamel's formula and Gronwall therefore give
\begin{equation}\label{ce:eq:adjoint-uniform-bound}
 \sup_{0<\delta\le1}\sup_{0\le r\le S}
 \left(\norm{g^\delta_r}_{L^\infty}+\norm{\partial_xg^\delta_r}_{L^\infty}\right)\le C_S.
\end{equation} The same argument constructs $g^*$ along the deterministic flow of $b_*$. The transport-diffusion semigroup converges on Lipschitz functions to the deterministic semigroup, with error $O_S(\delta)$: couple the stochastic and deterministic characteristics, use the uniform $O(\delta)$ drift difference from \eqref{ce:eq:W1}, and then Gronwall. For a uniformly Lipschitz $g$, the difference of the integral operators is also $O_S(\delta)$, since $y\mapsto W''(x-y)g(y)$ has a Lipschitz bound uniform in $x$. Multiplication coefficients converge uniformly. Subtracting the two mild equations and applying Gronwall gives
\begin{equation}\label{ce:eq:adjoint-convergence-bound}
 \sup_{0\le r\le S}\norm{g^\delta_r-g^*_r}_{L^\infty}\le C_S\delta.
\end{equation}

Since $b_*(0)=b_*(\pi)=0$, the solutions of $\dot x=b_*(x)$ starting at either point remain there. The transport term in \eqref{ce:eq:limiting-gradient-adjoint} therefore contributes nothing when the mild equation is evaluated at $0$ or $\pi$, and the resulting two equations form a closed system. With $A_r=g^*_r(0)$ and $B_r=g^*_r(\pi)$,
\begin{equation}\label{ce:eq:adjoint2}
 -\partial_r\binom{A_r}{B_r}
 =b\begin{pmatrix}-q&q\\p&-p\end{pmatrix}\binom{A_r}{B_r},
 \qquad (A_S,B_S)=(1,0).
\end{equation}
The terms involving $W''(0)$ cancel: at zero, for example,
$(-\kappa)A_r-paA_r+qbB_r=-qbA_r+qbB_r$.
The vector $(p,p)$ is constant under this system and the other mode decays with rate $b$. Solving gives \eqref{ce:eq:gradcenters}.
\end{proof}

The limiting values in \eqref{ce:eq:gradcenters} determine the Brownian variance in \cref{ce:prop:noiseobstruction}. For the family just constructed, this variance is
\begin{equation}\label{ce:eq:noise}
 \mathcal V_\delta=2\kappa\delta^2\int_0^S
 \int|g^\delta_r(x)|^2\rho^\delta_r(x)m(\dd x)\dd r.
\end{equation}

\begin{prop}[Variance comparison]\label{ce:prop:noisecross}
For the fixed angular horizon $S=20$,
\begin{equation}\label{ce:eq:noiselimit}
 \lim_{\delta\downarrow0}\frac{\mathcal V_\delta}{\delta^2}
 =2\kappa p^2S+\frac{\kappa pq}{b}(1-e^{-2bS}),
\end{equation}
\begin{equation}\label{ce:eq:ratio}
 \lim_{\delta\downarrow0}\frac{\mathcal V_\delta}{\Var_{\rho^\delta_Sm}(\psi)}
 =2\kappa pS+\frac{\kappa q}{b}(1-e^{-2bS})
 =\frac{5\pi^2}{18}+\frac5{12}(1-e^{-10\pi^2/3})>2.
\end{equation}
Consequently, $\mathcal V_\delta>2\Var_{\rho^\delta_Sm}(\psi)$ for all sufficiently small $\delta>0$.
\end{prop}
\begin{proof}
By \cref{ce:lem:adjointlimit,ce:lem:packet}, the integrand in \eqref{ce:eq:noise}, after division by $2\kappa\delta^2$, converges uniformly in $r$ to
\begin{equation}\label{ce:eq:noisesquare}
 p\bigl(p+qe^{-b(S-r)}\bigr)^2+q\bigl(p(1-e^{-b(S-r)})\bigr)^2
 =p^2+pqe^{-2b(S-r)}.
\end{equation}
Integrating this identity over $0\le r\le S$ proves \eqref{ce:eq:noiselimit}. Dividing by the variance asymptotic in \eqref{ce:eq:prodvar} then gives \eqref{ce:eq:ratio}. For $S=20$, its first term is $5\pi^2/18>5/2$, so the limiting ratio is strictly larger than two.
\end{proof}

The choice $p=1/6$ and $q=5/6$ gives $\kappa>0$. With $D_\delta=\kappa\delta^2$, the rescaled displacement of the packet centered at zero converges to a nondegenerate stationary Ornstein--Uhlenbeck process. Consequently, its contribution to the product-reference variance remains $p\delta^2+o(\delta^2)$, while the common-translation Brownian variance accumulates linearly in $S$. Strict positivity of $\kappa$ is needed for this simultaneous presence of positive noise and stationary packet variance, not merely for localization at scale $\delta$ on a fixed time interval.

\subsection{Exponential moments of bounded centered statistics}\label{ce:sec:bounded-statistics}

We next establish the exponential moments used in \cref{ce:sec:entropy}. They control the reference defect \eqref{ce:eq:R} and the nonlinear empirical remainder \eqref{ce:eq:martremainders}. We use the bounded-kernel case of the degenerate $U$-statistic estimates recalled in \cite[Propositions~3.1 and~3.3]{DGRpartition}; the corresponding exponential estimate is also due to Jabin--Wang \cite[Theorem~4]{JabinWang2018}.

\begin{lemma}[Uniform small exponential moment]\label{ce:lem:Ustat}
Let $H$ be a bounded symmetric kernel and $\mu$ a probability measure such that
\begin{equation}\label{ce:eq:U-centering}
 \int H(x,y)\mu(\dd y)=\int H(x,y)\mu(\dd x)=0
 \quad\text{for $\mu$-a.e. remaining variable}.
\end{equation}
For independent $x_1,\ldots,x_N$ with common law $\mu$, set
\begin{equation}\label{ce:eq:UN}
 U_N:=\frac1N\sum_{i\ne j}H(x_i,x_j).
\end{equation}
There are universal constants $a,C>0$ such that, if $\norm H_\infty>0$,
\begin{equation}\label{ce:eq:Uexp}
 \sup_N\E_{\mu^{\otimes N}}\!\left[e^{a|U_N|/\norm H_\infty}\right]\le C.
\end{equation}
\end{lemma}
\begin{proof}
This is the bounded-kernel consequence of the cited $U$-statistic bounds. In the density setting of \cite[Theorem~4]{JabinWang2018}, apply that theorem to both signs of a sufficiently small multiple of $H/\norm H_\infty$, and use $e^{|x|}\le e^x+e^{-x}$. The theorem includes the diagonal sum; removing it changes the exponent by at most a universal constant. The decoupling formulation in \cite[Propositions~3.1 and~3.3]{DGRpartition} gives the same conclusion for an arbitrary probability measure $\mu$.
\end{proof}

For a bounded symmetric $c$ with its diagonal set to zero, define
\begin{equation}\label{ce:eq:Qdefinition}
 Q_{N,\mu}[c]:=N\iint c(x,y)(\mu_N-\mu)(\dd x)(\mu_N-\mu)(\dd y).
\end{equation}
If $\widetilde c$ is the Hoeffding centering of $c$, then
\begin{equation}\label{ce:eq:Qcompletion}
 Q_{N,\mu}[c]=\frac1N\sum_{i\ne j}\widetilde c(x_i,x_j)
 -\frac2N\sum_i c_\mu(x_i)+\overline c_\mu,
\end{equation}
where $c_\mu(x):=\int c(x,y)\mu(\dd y)$, $\overline c_\mu:=\iint c\dd\mu^{\otimes2}$, and $\widetilde c(x,y):=c(x,y)-c_\mu(x)-c_\mu(y)+\overline c_\mu$. Since
\begin{equation}\label{ce:eq:Q-correction-bound}
 \left|-\frac2N\sum_i c_\mu(x_i)+\overline c_\mu\right|\le3\norm c_\infty,
\end{equation}
\cref{ce:lem:Ustat} gives a small exponential moment of $|Q_{N,\mu}[c]|$, uniformly in $N$.

\subsection{Global entropy control relative to the moving Gibbs law}\label{ce:sec:entropy}

In this and the next two subsections, $t$ denotes angular time. Fix $D>0$, $S<\infty$, and a smooth probability density $\rho_t$ on $\T^4$, and set $\dd\mu_t:=\rho_t\dd m$. Assume that
\begin{equation}\label{ce:eq:four-dimensional-background}
 \partial_t\rho_t=D\Delta\rho_t+\diver(\rho_t\nabla(\g*\mu_t)),
 \qquad \inf_{[0,S]\times\T^4}\rho>0.
\end{equation}
Set $\mu_t(v):=\int_{\T^4}v\dd\mu_t$, and $\Gamma_{N,t}:=\Q_{N,D^{-1}}(\mu_t)$. The particle law $P_{N,t}$ solves \eqref{ce:eq:angular-sde} with initial law $\Gamma_{N,0}$, independent of the Brownian motions; its construction is given in \cref{ce:sec:singular}. Write
\begin{equation}\label{ce:eq:G}
 G_{N,t}:=\frac{\dd\Gamma_{N,t}}{\dd\mu_t^{\otimes N}}.
\end{equation}
By \eqref{eq:gibbs}, $G_{N,t}=\K_{N,D^{-1}}(\mu_t)^{-1}e^{-D^{-1}N\Fr_N(\cdot,\mu_t)}$, and hence
\begin{equation}\label{ce:eq:G2partition}
 \norm{G_{N,t}}_{L^2(\mu_t^{\otimes N})}^2
 =\frac{\K_{N,2D^{-1}}(\mu_t)}
 {\K_{N,D^{-1}}(\mu_t)^2}.
\end{equation}
The static partition theorem in \cite[Theorem 2.2(i), Remark 2.1]{DGRpartition}, applied at $D^{-1}$ and $2D^{-1}$, bounds the numerator. Jensen's inequality and the definition in \eqref{eq:modulated} give
$\K_{N,D^{-1}}(\mu)\ge e^{D^{-1}I_\mu}$, uniformly in $N$. Hence
\begin{equation}\label{ce:eq:staticG2}
 \sup_N\sup_{t\le S}\norm{G_{N,t}}_{L^2(\mu_t^{\otimes N})}<\infty.
\end{equation}
To check that the partition theorem applies uniformly for $0\le t\le S$, boundedness of $\rho$ on $[0,S]\times\T^4$ and $\g\in L^1(m)\cap L^2(m)$ give
\begin{equation}\label{ce:eq:static-trajectory-controls}
 \sup_{0\le t\le S}\left(\norm{\rho_t}_{L^\infty}+\norm{\g*\mu_t}_{L^\infty}
 +|I_{\mu_t}|+\norm{\widetilde{\g}_{\mu_t}}_{L^2(\mu_t^{\otimes2})}\right)<\infty,
\end{equation}
where $\widetilde{\g}_{\mu_t}$ is the kernel centered in each variable under $\mu_t$. The bounded convolution also bounds the one-body logarithmic exponential moments. These are the scalar controls in \cite[Appendix~A]{DGRpartition}. For the trajectory $\rho^\delta_t(x_1)$ constructed in \cref{ce:sec:packet}, these bounds may depend on the fixed $\delta$.

For a Lipschitz vector field $v:\T^4\to\R^4$, define
\begin{equation}\label{ce:eq:cvector}
 c_v(x,y):=\mathbf1_{\{x\ne y\}}(v(x)-v(y))\cdot\nabla\g(x-y),
\end{equation}
Then
\begin{equation}\label{ce:eq:commutator-sup-bound}
 \norm{c_v}_{L^\infty}\le C\norm{\nabla v}_{L^\infty}.
\end{equation}
 The mean-field continuity velocity is
\begin{equation}\label{ce:eq:angular-continuity-velocity}
 u_t:=-\nabla(\g*\mu_t)-D\nabla\log\rho_t,
 \qquad \partial_t\rho_t+\diver(\rho_tu_t)=0.
\end{equation}

\begin{lemma}[Exact reference defect and entropy estimate]\label{ce:lem:globalentropy}
The moving Gibbs defect is
\begin{equation}\label{ce:eq:R}
 \partial_t\Gamma_{N,t}-L_N^*\Gamma_{N,t}=R_{N,t}\Gamma_{N,t},\qquad
 R_{N,t}=-\frac1{2D}\left(Q_{N,\mu_t}[c_{u_t}]
 -\E_{\Gamma_{N,t}}[Q_{N,\mu_t}[c_{u_t}]]\right).
\end{equation}
For the actual particle law $P_{N,t}$ started from $\Gamma_{N,0}$,
\begin{equation}\label{ce:eq:Hbounds}
 \sup_N\sup_{t\le S}\Hr(P_{N,t}\mid\Gamma_{N,t})<\infty,
 \qquad
 \sup_N\sup_{t\le S}\Hr(P_{N,t}\mid\mu_t^{\otimes N})<\infty.
\end{equation}
\end{lemma}
\begin{proof}
Off the diagonal, the Gibbs score satisfies
\begin{equation}\label{ce:eq:score}
 D\nabla_i\log\Gamma_{N,t}=b_i^N-u_t(x_i).
\end{equation}
Consequently, $L_N^*\Gamma_{N,t}=-\sum_i\diver_i(u_t(x_i)\Gamma_{N,t})$. The material derivative of \eqref{eq:modulated} and the derivative of its partition function are
\begin{align}
 \left(\partial_t+\sum_i u_t(x_i)\cdot\nabla_i\right)N\Fr_N(\cdot,\mu_t)
 &=\frac12Q_{N,\mu_t}[c_{u_t}],\label{ce:eq:angular-material-derivative}\\
 \partial_t\log\K_{N,D^{-1}}(\mu_t)
 &=-\frac1{2D}\E_{\Gamma_{N,t}}[Q_{N,\mu_t}[c_{u_t}]].\label{ce:eq:angular-normalizer-derivative}
\end{align}
These are the angular versions of the computation in \cref{sec:static}; substitution in the normalized density gives \eqref{ce:eq:R}.

The kernel $c_{u_t}$ is uniformly bounded on the fixed time interval. By \cref{ce:lem:Ustat,ce:eq:Qcompletion,ce:eq:staticG2} and Cauchy--Schwarz, there are $a>0$ and $K<\infty$, independent of $N$ and $t$, such that
\begin{equation}\label{ce:eq:Rexp}
 \E_{\Gamma_{N,t}}[e^{a|R_{N,t}|}]\le K.
\end{equation}
Indeed, if $a_0>0$ is small enough for the product exponential moment, then
\begin{equation}\label{ce:eq:commutator-Gibbs-transfer}
 \E_{\Gamma_{N,t}}[e^{a_0|Q_{N,\mu_t}[c_{u_t}]|}]
 \le\norm{G_{N,t}}_{L^2(\mu_t^{\otimes N})}
 \left(\E_{\mu_t^{\otimes N}}[e^{2a_0|Q_{N,\mu_t}[c_{u_t}]|}]\right)^{1/2}\le C.
\end{equation}
Jensen then bounds $|\E_{\Gamma_{N,t}}[Q_{N,\mu_t}[c_{u_t}]]|$ by $a_0^{-1}\log C$. Choosing $a$ in terms of $a_0$ and $D$ proves \eqref{ce:eq:Rexp}. Set $h_N(t):=\Hr(P_{N,t}\mid\Gamma_{N,t})$. The exact entropy identity is
\begin{equation}\label{ce:eq:entropyidentity}
 h_N'(t)=-D\int\left|\nabla\log\frac{\dd P_{N,t}}{\dd\Gamma_{N,t}}\right|^2\dd P_{N,t}
 -\E_{P_{N,t}}[R_{N,t}].
\end{equation}
Its singular-law justification is given in \cref{ce:sec:singular}. The Donsker--Varadhan lemma relative to $\Gamma_{N,t}$ and \eqref{ce:eq:Rexp} imply
\begin{equation}\label{ce:eq:entropygronwall}
 h_N'(t)\le a^{-1}(h_N(t)+\log K),\qquad h_N(0)=0.
\end{equation}
Gronwall proves the first bound in \eqref{ce:eq:Hbounds} for every finite $S$. Finally,
\begin{equation}\label{ce:eq:Hproduct}
 \Hr(P_{N,t}\mid\mu_t^{\otimes N})
 =h_N(t)+\E_{P_{N,t}}[\log G_{N,t}]
 \le2h_N(t)+\log\int G_{N,t}^2\dd\mu_t^{\otimes N}.
\end{equation}
Here, the Donsker--Varadhan lemma is applied relative to $\Gamma_{N,t}$ with observable $\log G_{N,t}$, and $\int G_{N,t}\dd\Gamma_{N,t}=\int G_{N,t}^2\dd\mu_t^{\otimes N}$. This proves the second bound.
\end{proof}

\begin{lemma}[Two consequences for the actual laws]\label{ce:lem:transfer}
If $(c_t)_{0\le t\le S}$ is a family of symmetric commutator kernels with $\sup_{t\le S}\norm{c_t}_{L^\infty}<\infty$, then
\begin{equation}\label{ce:eq:Qactual}
 \sup_N\sup_{t\le S}\E_{P_{N,t}}[|Q_{N,\mu_t}[c_t]|]<\infty.
\end{equation}
If $(v_t)_{0\le t\le S}$ is a family of real-valued functions with $\sup_{t\le S}\norm{v_t}_{L^\infty}<\infty$, then
\begin{equation}\label{ce:eq:linearactual}
 \sup_N\sup_{t\le S}
 \E_{P_{N,t}}\!\left[\left|N^{-1/2}\sum_i(v_t(x_i)-\mu_t(v_t))\right|^2\right]<\infty.
\end{equation}
\end{lemma}
\begin{proof}
For the first assertion, apply \eqref{ce:eq:commutator-Gibbs-transfer} with $c_t$ in place of $c_{u_t}$ and then the Donsker--Varadhan lemma with \eqref{ce:eq:Hbounds}. For the second assertion, Hoeffding's lemma gives
\begin{equation}\label{ce:eq:linear-product-mgf}
 \E_{\mu_t^{\otimes N}}\left[e^{\theta N^{-1/2}\sum_i(v_t(x_i)-\mu_t(v_t))}\right]
 \le e^{\theta^2\operatorname{osc}(v_t)^2/8},\qquad\theta\in\R;
\end{equation}
see also \cite[Section~8.2]{DGRpartition}. Integrating \eqref{ce:eq:linear-product-mgf} with $\theta=\sqrt{2a}\,Z$, where $Z$ is an independent standard Gaussian, yields
\begin{equation}\label{ce:eq:linear-product-square-mgf}
 \E_{\mu_t^{\otimes N}}\left[e^{a|N^{-1/2}\sum_i(v_t(x_i)-\mu_t(v_t))|^2}\right]
 \le\left(1-\frac a2\operatorname{osc}(v_t)^2\right)^{-1/2}.
\end{equation}
Choose $a>0$ so that $a\sup_{t\le S}\operatorname{osc}(v_t)^2<1$. Cauchy--Schwarz transfers the moment with parameter $a/2$ to $\Gamma_{N,t}$, and the Donsker--Varadhan lemma transfers its first moment to $P_{N,t}$. This proves \eqref{ce:eq:linearactual}.
\end{proof}

\subsection{Brownian noise as a terminal chi-square obstruction}\label{ce:sec:obstruction}

We now use \cref{ce:lem:transfer} to control the remainder and predictable quadratic variation in the empirical fluctuation equation \eqref{ce:eq:empirical-adjoint-differential}. The variance $A$ in \eqref{ce:eq:noisecriterion} is the Brownian variance computed in \eqref{ce:eq:noise} for the stripe construction. It suffices that the initial statistics have a tight family of laws.

\begin{prop}[Noise-variance obstruction]\label{ce:prop:noiseobstruction}
Let $\rho_t$ satisfy \eqref{ce:eq:four-dimensional-background}, and let $P_{N,0}=\Gamma_{N,0}$. For $\psi\in C^\infty(\T^4;\R)$, let $f^r$ solve
\begin{align}
 \partial_r f^r+\mathcal L_r f^r&=0,\qquad f^S=\psi,\label{ce:eq:full-adjoint}\\
 \mathcal L_r f(x)&:=D\Delta f(x)-\int\nabla\g(x-y)\cdot(\nabla f(x)-\nabla f(y))\dd\mu_r(y).\label{ce:eq:full-adjoint-generator}
\end{align}
Then $f$ is the unique classical solution, with
\begin{equation}\label{ce:eq:full-adjoint-bound}
 \sup_{0\le r\le S}\bigl(\norm{f^r}_{L^\infty}+\norm{\nabla f^r}_{L^\infty}+\norm{\nabla^2 f^r}_{L^\infty}\bigr)<\infty.
\end{equation}
Set
\begin{equation}\label{ce:eq:noisecriterion}
 V:=\Var_{\mu_S}(\psi)>0,\qquad
 A:=2D\int_0^S\mu_r(|\nabla f^r|^2)\dd r.
\end{equation}
If $A>2V$, then
\begin{equation}\label{ce:eq:chiobstruction}
 \lim_{N\to\infty}\chi^2(P_{N,S}\mid\mu_S^{\otimes N})=+\infty.
\end{equation}
\end{prop}
\begin{proof}
\parheading{The backward solution and empirical identity.}
The local diffusion with drift $-\nabla\g*\mu_r$ has smooth coefficients. The remaining term $\nabla\g*(\rho_r\nabla f)$ maps $C^m$ to $C^{m-1}$, with norm bounded by $C_{m,\rho}\|\nabla\g\|_{L^1}$. The usual one-derivative parabolic smoothing estimate therefore makes its Duhamel kernel integrable. Successive short-interval contraction and differentiation give existence, uniqueness, and \eqref{ce:eq:full-adjoint-bound}; see, e.g., \cite[Chapter~5]{Lunardi} for the local parabolic estimates. In the stripe case, this is the solution already constructed in \eqref{ce:eq:backward}.

Set $\mu_{N,r}:=N^{-1}\sum_i\delta_{x_i(r)}$ and $\xi_{N,r}:=\sqrt N(\mu_{N,r}-\mu_r)$. The finite-particle adjoint identity of \cite[Proposition~8.2]{DGRpartition}, in the angular normalization, is
\begin{equation}\label{ce:eq:exactmart}
 \xi_{N,S}(\psi)=Y_N+M_N(S)+R_N,\qquad Y_N:=\xi_{N,0}(f^0),
\end{equation}
\begin{align}
 R_N&:=-\frac1{2\sqrt N}\int_0^S Q_{N,\mu_r}[c_{\nabla f^r}]\dd r,\label{ce:eq:martremainders}\\
 M_N(t)&:=\sqrt{\frac{2D}{N}}\sum_i\int_0^t\nabla f^r(x_i(r))\cdot\dd B_i(r),\qquad 0\le t\le S.\label{ce:eq:martingale-process}
\end{align}
The process $M_N$ is a continuous square-integrable martingale. Here, $c_{\nabla f}$ is defined by \eqref{ce:eq:cvector}, with diagonal zero. The coefficient can also be checked directly: substituting $\mu_{N,r}=\mu_r+N^{-1/2}\xi_{N,r}$ in the symmetrized empirical drift gives
\begin{equation}\label{ce:eq:empirical-adjoint-differential}
 \dd\xi_{N,r}(f^r)=\xi_{N,r}(\partial_r f^r+\mathcal L_rf^r)\dd r
 -\frac{Q_{N,\mu_r}[c_{\nabla f^r}]}{2\sqrt N}\dd r+\dd M_N(r).
\end{equation}
The first term vanishes by \eqref{ce:eq:full-adjoint}. All pair sums retain the normalization $1/N$.

The Hessian bound makes $c_{\nabla f^r}$ uniformly bounded, so \eqref{ce:eq:Qactual} gives $\E[|R_N|]=O(N^{-1/2})$. The predictable quadratic variation of $M_N$ is
\begin{equation}\label{ce:eq:martbracket}
 \langle M_N\rangle_S=2D\int_0^S\mu_{N,r}(|\nabla f^r|^2)\dd r.
\end{equation}
Consequently,
\begin{equation}\label{ce:eq:bracketdifference}
 \langle M_N\rangle_S-A=\frac{2D}{\sqrt N}\int_0^S\xi_{N,r}(|\nabla f^r|^2)\dd r.
\end{equation}
Thus, \eqref{ce:eq:linearactual} and Cauchy--Schwarz give
$\E[|\langle M_N\rangle_S-A|]\le C_{D,S,f}N^{-1/2}\to0$. We also have the deterministic bound $\langle M_N\rangle_S\le 2DS\sup_{0\le r\le S}\norm{\nabla f^r}_{L^\infty}^2$.

\parheading{The initial statistic and the independent Gaussian limit.}
The initial variables $Y_N$ are tight. Indeed, Cauchy--Schwarz, \eqref{ce:eq:staticG2}, and the bounded fourth moment of the independent centered sum give
\begin{equation}\label{ce:eq:initial-statistic-tightness}
 \E[Y_N^2]\le\norm{G_{N,0}}_{L^2(\mu_0^{\otimes N})}
 \left(\E_{\mu_0^{\otimes N}}[|\xi_{N,0}(f^0)|^4]\right)^{1/2}\le C_{\rho_0,D,f^0}.
\end{equation} Along any subsequence one can therefore choose a further subsequence on which $Y_N$ converges in law, say to a probability law $\nu$.

To prove that the limiting Gaussian variable from $M_N(S)$ is independent of the limit of $Y_N$, set, for $\theta\in\R$,
\begin{equation}\label{ce:eq:complexmart}
 Z_N(r):=\exp\!\left(i\theta M_N(r)+\frac{\theta^2}{2}\langle M_N\rangle_r\right).
\end{equation}
Its modulus is bounded by the deterministic bound on $\langle M_N\rangle_S$, so $Z_N$ is a true martingale and
$\E[Z_N(S)\mid\mathcal F_0]=1$. Since $Y_N$ is $\mathcal F_0$-measurable, write $e^{i\theta M_N(S)}=Z_N(S)e^{-\theta^2\langle M_N\rangle_S/2}$ and use the conditional expectation above. For every $\alpha,\theta\in\R$, the mean-value estimate for the last exponential gives
\begin{equation}\label{ce:eq:factorizationestimate}
 \left|\E[e^{i\alpha Y_N+i\theta M_N(S)}]-e^{-\theta^2A/2}\E[e^{i\alpha Y_N}]\right|
 \le C_\theta\E[|\langle M_N\rangle_S-A|]\longrightarrow0.
\end{equation}
Thus, the joint characteristic functions converge to those of an independent pair with laws $\nu$ and $\mathcal N(0,A)$. Hence, along the further subsequence,
\begin{equation}\label{ce:eq:convolutionlimit}
 \xi_{N,S}(\psi)\ \longrightarrow\ \nu*\mathcal N(0,A)\quad\text{in law}.
\end{equation}
Under the product law $\mu_S^{\otimes N}$, the same function $X_N\mapsto N^{-1/2}\sum_i(\psi(x_i)-\mu_S(\psi))$ converges in law to $\mathcal N(0,V)$ by the usual central limit theorem.

\parheading{The limiting chi-square divergence.}
For any probability measure $\nu$, the convolution $\nu*\mathcal N(0,A)$ in \eqref{ce:eq:convolutionlimit} has infinite chi-square relative to $\mathcal N(0,V)$ when $A>2V$. To see this without an assumption on $\nu$, choose $R<\infty$ with $\nu([-R,R])>0$. The Lebesgue density $p$ of $\nu*\mathcal N(0,A)$ satisfies
\begin{equation}\label{ce:eq:convtaillower}
 p(x)\ge\frac{\nu([-R,R])}{\sqrt{2\pi A}}
 \exp\!\left(-\frac{(|x|+R)^2}{2A}\right).
\end{equation}
If $\gamma_V$ is the density of $\mathcal N(0,V)$, then
\begin{equation}\label{ce:eq:Gaussian-quotient-lower}
 \frac{p(x)^2}{\gamma_V(x)}
 \ge\frac{\nu([-R,R])^2\sqrt{2\pi V}}{2\pi A}
 \exp\left[\left(\frac1{2V}-\frac1A\right)x^2-\frac{2R}{A}|x|-\frac{R^2}{A}\right].
\end{equation}
The positive quadratic coefficient makes this lower bound nonintegrable.

Suppose, contrary to \eqref{ce:eq:chiobstruction}, that there are $N_j\uparrow\infty$ such that
\begin{equation}\label{ce:eq:bounded-chi-subsequence}
 \sup_j\chi^2(P_{N_j,S}\mid\mu_S^{\otimes N_j})<\infty.
\end{equation}
Apply data processing to the map $X_N\mapsto N^{-1/2}\sum_i(\psi(x_i)-\mu_S(\psi))$ under these two laws; see \cite[Section~2]{NishiyamaSason2020}. Along the further subsequence selected above, its laws converge to $\nu*\mathcal N(0,A)$ and $\mathcal N(0,V)$, respectively. Joint lower semicontinuity under weak convergence would then give
\begin{equation}\label{ce:eq:limiting-chi-contradiction}
 \chi^2\bigl(\nu*\mathcal N(0,A)\mid\mathcal N(0,V)\bigr)
 \le\liminf_j\chi^2(P_{N_j,S}\mid\mu_S^{\otimes N_j})<\infty,
\end{equation}
contrary to \eqref{ce:eq:Gaussian-quotient-lower}. Here, lower semicontinuity follows, for example, from \cite[Theorem~19 and~(6)]{vanErvenHarremoes2014}, applied to the R\'enyi divergence of order two. Every bounded subsequence is excluded, proving \eqref{ce:eq:chiobstruction}.
\end{proof}

The static central limit theorem of \cite[Proposition~2.11]{DGRpartition} additionally identifies $\nu$ as a centered Gaussian law. We use this identification only for the central limit statement in \cref{thm:periodic-counterexample}; the chi-square divergence proved above requires only tightness of $Y_N$.

\subsection{Actual singular laws and the entropy identity}\label{ce:sec:singular}

We justify the entropy and empirical identities for the singular diffusion \eqref{ce:eq:angular-sde}. The collision-free construction in \cref{app:particle-wp} has the same local form on the torus: the periodic kernel is $-\pi^2\log|x|$ plus a smooth function near zero, and the internal contribution of the smooth part to the cluster-dispersion equation is bounded by a constant times the dispersion, as is the exterior-field remainder. Here, we instead use the square-integrable force to give a short approximation argument that also identifies the reversible stationary law needed for entropy differentiation. Throughout this subsection, $N$ and $D>0$ are fixed.

Let
\begin{equation}\label{ce:eq:stationary}
 \dd\Pi_N:=Z_N^{-1}\exp\!\left(-\frac1{DN}\sum_{i<j}\g(x_i-x_j)\right)\dd m^{\otimes N}.
\end{equation}

\parheading{Stationary domination and convergence of the particle laws.}
Let $\g_\varepsilon:=e^{\varepsilon\Delta}\g$, and let $\Pi_N^\varepsilon$ and $\Gamma_{N,0}^\varepsilon$ be the corresponding stationary and initial laws. The lower bound on $\g$ and the bound on $\|\nabla\g\|_{L^2(\T^4)}$ are preserved by heat smoothing. Dominated convergence of the stationary densities, including their normalizing constants, therefore gives
\begin{equation}\label{ce:eq:stationary-uniform-bounds}
 \sup_{0<\varepsilon\le1}\left(
 \left\|\frac{\dd\Pi_N^\varepsilon}{\dd m^{\otimes N}}\right\|_{L^\infty}
 +\int|b_N^\varepsilon|^2\dd\Pi_N^\varepsilon\right)<\infty.
\end{equation}
Moreover,
\begin{equation}\label{ce:eq:preparedstationaryratio}
 \frac{\dd\Gamma_{N,0}^\varepsilon}{\dd\Pi_N^\varepsilon}(X)
 =C_{N,\varepsilon}\prod_{i=1}^N
 \rho_0(x_i)\exp\!\left(D^{-1}(\g_\varepsilon*\mu_0)(x_i)\right).
\end{equation}
Since $\rho_0$ is bounded above and below and $\sup_\varepsilon\|\g_\varepsilon*\mu_0\|_\infty<\infty$, normalization yields
\begin{equation}\label{ce:eq:initial-stationary-domination}
 0<c_N\le\frac{\dd\Gamma_{N,0}^\varepsilon}{\dd\Pi_N^\varepsilon}\le C_N<\infty.
\end{equation}
This domination also holds on path space. Indeed, for every nonnegative measurable path functional $\Phi$, disintegrating with respect to the initial configuration gives
\begin{equation}\label{ce:eq:path-law-domination}
 \E_{\Gamma_{N,0}^\varepsilon}[\Phi(X)]
 =\int \E_x[\Phi(X)]\frac{\dd\Gamma_{N,0}^\varepsilon}{\dd\Pi_N^\varepsilon}(x)\dd\Pi_N^\varepsilon(x)
 \le C_N\E_{\Pi_N^\varepsilon}[\Phi(X)].
\end{equation}
Here, $\E_x$ refers to the same cutoff diffusion with deterministic initial configuration $x$.
Stationarity, \eqref{ce:eq:stationary-uniform-bounds}, and \eqref{ce:eq:path-law-domination} imply
\begin{equation}\label{ce:eq:driftmodulus}
 \E\!\left[\int_s^t|b_N^\varepsilon(X_r)|^2\dd r\right]\le C_N|t-s|,
 \qquad
 \E\!\left[\left|\int_s^t b_N^\varepsilon(X_r)\dd r\right|^2\right]\le C_N|t-s|^2.
\end{equation}
Kolmogorov's tightness criterion (see, e.g., \cite[Chapter~I, Theorem~2.1]{RevuzYor}) applies to the drift integrals, and the Brownian terms are tight in continuous paths.

Fix two particle indices $i\ne j$, and choose a smooth function $0\le F_r\le1$ of their relative position that equals one when $\operatorname{dist}(x_i,x_j)\le r$ and vanishes when this distance is at least $2r$. For sufficiently small $r>0$, the four-dimensional relative coordinate gives
\begin{equation}\label{ce:eq:collision-cutoff-norms}
 \|F_r\|_{L^1(m^{\otimes N})}+\|F_r\|_{L^2(m^{\otimes N})}^2\le C_Nr^4,
 \quad \|\Delta F_r\|_{L^1(m^{\otimes N})}\le C_Nr^2,
 \quad \|\nabla F_r\|_{L^2(m^{\otimes N})}\le C_Nr.
\end{equation}
For the stationary cutoff diffusion, It\^o's formula gives
\begin{equation}\label{ce:eq:collision-ito}
 F_r(X_t)=F_r(X_0)+\int_0^t(D\Delta F_r+b_N^\varepsilon\cdot\nabla F_r)(X_u)\dd u
 +\sqrt{2D}\int_0^t\nabla F_r(X_u)\cdot\dd B_u.
\end{equation}
Taking the supremum and expectation, and applying Doob's inequality and the It\^o isometry to the last term (see, e.g., \cite[Chapters~II and~IV]{RevuzYor}), yields
\begin{align}
 \E_{\Pi_N^\varepsilon}\!\left[\sup_{t\le S}F_r(X_t)\right]
 &\le\|F_r\|_{L^1(\Pi_N^\varepsilon)}+DS\|\Delta F_r\|_{L^1(\Pi_N^\varepsilon)}\notag\\
 &\quad+S\|b_N^\varepsilon\|_{L^2(\Pi_N^\varepsilon)}\|\nabla F_r\|_{L^2(\Pi_N^\varepsilon)}
 +\sqrt{8DS}\|\nabla F_r\|_{L^2(\Pi_N^\varepsilon)}\notag\\
 &\le C_{N,D,S}(r^4+r^2+r)\le C_{N,D,S}r.\label{ce:eq:capacitybound}
\end{align}
By \eqref{ce:eq:path-law-domination}, the same estimate holds up to a constant for the initial law $\Gamma_{N,0}^\varepsilon$. Summing over pairs gives
\begin{equation}\label{ce:eq:collision-probability}
 \Pb\left(\min_{0\le t\le S}\min_{i<j}\operatorname{dist}(x_i^\varepsilon(t),x_j^\varepsilon(t))<r\right)
 \le C_{N,D,S}r.
\end{equation}
The event in \eqref{ce:eq:collision-probability} is open in the space of continuous paths. Portmanteau therefore gives the same bound for any path limit, and letting $r\downarrow0$ excludes collisions. On each collision-free compact configuration set, the drifts converge in every smooth norm. The stopped martingale problem consequently passes to the limit. Pathwise uniqueness away from collisions identifies the limit with the singular strong solution and removes the subsequence. These arguments apply both to stationary and to Gibbs initial laws. Stationarity and reversibility pass to the limit, and \eqref{ce:eq:path-law-domination} identifies the latter path law as a bounded change of its initial distribution.

\parheading{Entropy differentiation.}
Write $P_{N,t}=p_t\Pi_N$ and $\Gamma_{N,t}=e^{A_t}\Pi_N$. The function $A_t$ is smooth with bounded derivatives on each fixed time interval, because the singular pair factors cancel in this ratio. The bounds $0<c_N\le p_0\le C_N$ pass to $p_t$ under the reversible Markov semigroup.

We first justify integration by parts across the collision set. Choose $\chi_r$ vanishing near all pair collisions and converging to one elsewhere, with $\int|\nabla\chi_r|^2\dd\Pi_N\to0$, as supplied by \eqref{ce:eq:collision-cutoff-norms}. For a smooth test $\varphi$,
\begin{equation}\label{ce:eq:generator-boundary-term}
 \left|\int\varphi\nabla A_t\cdot\nabla\chi_r\dd\Pi_N\right|
 \le\|\varphi\|_{L^\infty}\|\nabla A_t\|_{L^\infty}
 \left(\int|\nabla\chi_r|^2\dd\Pi_N\right)^{1/2}\longrightarrow0.
\end{equation}
Furthermore, $L_NA_t=D\Delta A_t+b_N\cdot\nabla A_t\in L^2(\Pi_N)$ by \eqref{ce:eq:stationary-uniform-bounds}. Thus $A_t$ belongs to the generator domain, and the reversible gradient form is the closure of $D\int|\nabla\varphi|^2\dd\Pi_N$ without a collision-boundary term. The entropy chain rule for this form, applied to the bounded positive semigroup density $p_t$, gives the relative de Bruijn identity (cf.\ \cite[Chapter~5]{BGL})
\begin{equation}\label{ce:eq:statentropy}
 \frac\dd{\dd t}\int p_t\log p_t\dd\Pi_N
 =-D\int|\nabla\log p_t|^2p_t\dd\Pi_N.
\end{equation}
The identity holds for a.e.\ $t$ and in integrated form starting at zero; form approximation and the bounds on $p_t$ justify the logarithmic test. Differentiating the smooth cross term now gives
\begin{equation}\label{ce:eq:crossentropy}
 \frac\dd{\dd t}\Hr(P_{N,t}\mid\Gamma_{N,t})
 =-D\int|\nabla\log p_t|^2p_t\dd\Pi_N
 -\int p_t(\partial_t A_t+L_NA_t)\dd\Pi_N.
\end{equation}
Since $L_Ne^{A_t}=e^{A_t}(L_NA_t+D|\nabla A_t|^2)$, the reference defect is
\begin{equation}\label{ce:eq:RviaA}
 R_{N,t}=\partial_t A_t-L_NA_t-D|\nabla A_t|^2.
\end{equation}
Reversibility gives $\int p_tL_NA_t\dd\Pi_N=-D\int\nabla p_t\cdot\nabla A_t\dd\Pi_N$. Substitution into \eqref{ce:eq:crossentropy} therefore yields
\begin{equation}\label{ce:eq:entropy-square-completion}
 \frac{\dd}{\dd t}\Hr(P_{N,t}\mid\Gamma_{N,t})
 =-D\int|\nabla\log p_t-\nabla A_t|^2p_t\dd\Pi_N
 -\int R_{N,t}p_t\dd\Pi_N,
\end{equation}
which is \eqref{ce:eq:entropyidentity}. The empirical It\^o identities follow by stopping away from the collision set and then removing the stopping, using compact-time positive pair separation and the bounded symmetrized commutator in \eqref{ce:eq:empirical-adjoint-differential}.

\parheading{Finiteness of chi-square at fixed particle number.}
The law $P_{N,t}$ is bounded above by a fixed-$N$ multiple of $\Pi_N$. The latter has a bounded density relative to $m^{\otimes N}$, and $\rho_t$ is bounded below on $[0,S]\times\T^4$. Consequently, $\chi^2(P_{N,t}\mid\mu_t^{\otimes N})$ is finite for every fixed $N$ and $t\le S$.

\subsection{Completion of the theorem and the fluctuation covariance}\label{ce:sec:completion}

\begin{proof}[Proof of \cref{thm:periodic-counterexample}]
Choose the family in \eqref{ce:eq:rhoinitial} and the terminal test in \eqref{ce:eq:terminaltest}. By \cref{ce:prop:noisecross}, there is $\delta_0>0$ such that $\mathcal V_\delta>2\Var_{\rho^\delta_Sm}(\psi)$ for all $0<\delta<\delta_0$. Set $\beta_*=24/\delta_0^2$. Given any $\beta>\beta_*$, fix $\delta=\sqrt{24/\beta}<\delta_0$.

The angular four-dimensional trajectory $\rho^\delta_\tau(x_1)$ is smooth and strictly positive on $[0,20]$. The actual full singular particle laws exist by \cref{ce:sec:singular}. Static moments, the global moving-reference entropy estimate, and the backward-test estimates satisfy every hypothesis of \cref{ce:prop:noiseobstruction}. Writing $\dd\mu_{20}^\delta(x):=\rho_{20}^\delta(x_1)\dd m(x)$ on $\T^4$, we obtain
\begin{equation}\label{ce:eq:angular-final-divergence}
 \lim_{N\to\infty}\chi^2(P_{N,20}\mid(\mu_{20}^\delta)^{\otimes N})=+\infty.
\end{equation} Under \eqref{ce:eq:angkernel}, this is exactly \cref{eq:periodic-ce-failure} at physical time $5$ with $\beta=\pi^2/(\kappa\delta^2)=24/\delta^2$ and
\begin{equation}\label{ce:eq:physicalprofile}
 \mu_0(z)=\rho^\delta_0(2\pi z_1).
\end{equation}
The initial uniform chi-square bound is \eqref{ce:eq:staticG2}. The full entropy bound is \eqref{ce:eq:Hbounds}; finiteness at each fixed $N$ was established in \cref{ce:sec:singular}. Finally, $k=N$ in \cref{eq:periodic-ce-unrestricted-target} would give a bounded sequence in \cref{eq:periodic-ce-failure}, a contradiction.

Moreover, the static central limit theorem of \cite[Proposition 2.11]{DGRpartition} identifies the initial covariance as a positive screened covariance. Hence the terminal covariance on the test $\psi$ satisfies
\begin{equation}\label{ce:eq:prepared-covariance}
 \Sigma_\delta(\psi,\psi)
 =\Sigma_{\delta,0}(f^\delta_0,f^\delta_0)+\mathcal V_\delta
 \ge\mathcal V_\delta>2\Var_{\rho^\delta_Sm}(\psi).
\end{equation}
The initial Gaussian fluctuation and the asymptotically independent Brownian Gaussian fluctuation therefore give the central limit theorem \eqref{eq:intro-periodic-clt}, and \eqref{ce:eq:prepared-covariance} gives its strict variance inequality.
\end{proof}


\par\endgroup

\appendix
\begingroup

\section{Thermal equilibrium measures and logarithmic mean-field stability}\label{app:background}

We prove \cref{prop:background-stability} in four steps. The thermal equilibrium estimates in \cref{app:thermal-profiles} give the density bounds and one-body Poincar\'e inequality. \Cref{app:logarithmic-maps} establishes the convolution estimates needed to construct the linearized evolution in \cref{app:linear-semigroup}. Finally, \cref{app:nonlinear-stability} uses the one-derivative smoothing estimate to solve the nonlinear logarithmic equation.

The kernel and cutoff are normalized as in \eqref{eq:kernel} and \eqref{eq:heat-cutoff}. A superscript $\varepsilon$ denotes replacement of $\g$ by $\g_\varepsilon$, with $\varepsilon=0$ corresponding to the singular kernel. Bounds stated uniformly in the cutoff hold for $\varepsilon\in[0,\varepsilon_*]$, for a sufficiently small fixed $\varepsilon_*>0$. The time intervals and any uniformity in time are specified separately.

\subsection{Thermal equilibrium measures}\label{app:thermal-profiles}
For the Coulomb interaction, existence, uniqueness, and the Euler--Lagrange characterization of the thermal equilibrium measure are proved in \cite[Lemmas~2.1--2.2]{ArmstrongSerfatyThermal}. We give the argument for the present log/Riesz range, including the derivative estimates and convergence under heat truncation. Convexity of $-\log\pi$ outside a ball will give the one-body Poincar\'e inequality.

\begin{lemma}[Thermal regularity and a one-body spectral gap]\label{lem:thermal}
For every confinement satisfying \eqref{eq:confinement-class}--\eqref{eq:confinement-decay}, every $\beta>0$, and kernel \eqref{eq:kernel}, the thermal equation \eqref{eq:thermal} has a unique probability density $\pi$. It is smooth and positive. Set $U_\pi:=-\log\pi$. There are constants $a,M,C_P>0$ such that
\begin{equation}\label{eq:thermal-bounds}
 \pi(x)\le Me^{-a|x|^2},\qquad
 \norm{\nabla^j\pi}_{L^1}+\sup_x e^{a|x|^2/2}|\nabla^j\pi(x)|<\infty,\quad(j\ge0),
\end{equation}
\begin{equation}\label{eq:thermal-score}
 \nabla U_\pi=\beta\nabla V+\beta\nabla\g*\pi,
 \qquad \norm{\nabla^2U_\pi}_\infty<\infty,\qquad
 \nabla^2U_\pi\ge\frac{\beta\lambda_A}{2}\Id\quad\text{outside a ball},
\end{equation}
and the probability measure $\pi(x)\dd x$ satisfies \eqref{eq:one-body-PI}. The corresponding bounds hold uniformly for $\pi^\varepsilon$. For every integer $j\ge0$ and $0<a'<a/2$,
\begin{equation}\label{eq:thermal-cutoff-convergence}
 \norm{\nabla^j(\pi^\varepsilon-\pi)}_{L^1}
 +\sup_x e^{a'|x|^2}|\nabla^j(\pi^\varepsilon-\pi)(x)|\xrightarrow[\varepsilon\downarrow0]{}0.
\end{equation}
\end{lemma}
\begin{proof}
\parheading{Existence and uniqueness.}
Consider the free energy
\begin{equation}\label{eq:meanfield-freeenergy}
 \mathcal F(\rho)=\beta^{-1}\int\rho\log\rho\dd x
              +\int V(x)\rho(x)\dd x
              +\frac12\iint\g(x-y)\rho(x)\rho(y)\dd x\dd y.
\end{equation}
A Gaussian trial density has finite free energy. For $\s>0$, the interaction is nonnegative. For $\s=0$, use
\begin{equation}\label{eq:thermal-log-lower}
 -\log|x-y|\ge-\log(1+|x|)-\log(1+|y|),
 \qquad \log(1+r)\le\epsilon r^2+C_\epsilon.
\end{equation}
Combining part of the quadratic confinement with entropy relative to a broader Gaussian gives coercivity in both cases. Consequently, a minimizing sequence has uniformly bounded second moment and relative entropy with respect to such a Gaussian. Tightness and lower semicontinuity give a minimizer. To see the latter point directly for the logarithm, its negative tail is uniformly integrable by the second-moment bound, whereas its positive singular part is lower semicontinuous after truncation. The entropy plus a quadratic term is a Gaussian relative entropy up to a constant and is lower semicontinuous. The same argument applies to the cutoffs uniformly.

The entropy term excludes a set of positive Lebesgue measure on which a minimizer vanishes. Indeed, adding a small bounded density on such a set, while subtracting the same mass on a set where the minimizer is bounded above and below, produces a negative entropy variation of order $\epsilon|\log\epsilon|$. Restricting first to sets where the potentials are bounded makes the other variations of order $\epsilon$. Exhaustion gives the claimed positivity almost everywhere. Bounded variations where the density and its reciprocal are bounded then yield \eqref{eq:thermal}, again by exhaustion.

The heat representation shows that the interaction quadratic form is nonnegative on zero-mass densities. Strict convexity of entropy therefore gives uniqueness of the minimizer. The same convexity shows that any solution of the Euler--Lagrange relation is the minimizer.

\parheading{Uniform density and derivative bounds.}
Use the normalizations $b=0$ and $V(0)=0$ following \eqref{eq:confinement-decay}, and write
\begin{equation}\label{eq:sublinear-remainder}
 \widetilde V(x)=V(x)-\frac12x\cdot Ax.
\end{equation}
Given $\delta>0$, choose $R_\delta$ so that $|\nabla\widetilde V(x)|\le\delta$ for $|x|\ge R_\delta$. Integrate $\nabla\widetilde V$ along the segment from $0$ to $x$, using its global bound inside $B_{R_\delta}$ and the bound $\delta$ outside. This gives
\begin{equation}\label{eq:sublinear-growth}
 |\widetilde V(x)|\le\delta|x|+C_\delta.
\end{equation}
After decreasing $\delta$ and completing the square,
\begin{equation}\label{eq:quadratic-lower-bound}
 V(x)\ge\frac{\lambda_A}{4}|x|^2-C.
\end{equation}
Let $\gamma_A$ be the Gaussian proportional to $e^{-\beta x\cdot Ax/2}$. The function $\widetilde V$ is integrable under $\gamma_A$, and in both interaction cases $\sup_y(\g_\varepsilon*\gamma_A)(y)\le C$ uniformly. Jensen's inequality under $\gamma_A$ gives a lower bound on the thermal normalizer:
\begin{equation}
 Z_{\pi^\varepsilon} =Z_{\gamma_A}\int e^{-\beta(\widetilde V+h_{\pi^\varepsilon}^\varepsilon)}\dd\gamma_A \ge Z_{\gamma_A} \exp\left[-\beta\int\widetilde V\dd\gamma_A -\beta\int(\g_\varepsilon*\gamma_A)\dd\pi^\varepsilon\right] \ge c_{V,A,\beta}>0. \label{eq:thermal-Zlower}
\end{equation}
For a Riesz kernel, $h_{\pi^\varepsilon}^\varepsilon\ge0$, and \eqref{eq:quadratic-lower-bound} immediately gives a Gaussian upper bound. For the logarithm, the uniform second moment and \eqref{eq:thermal-log-lower} give
\begin{equation}\label{eq:thermal-log-density}
 h_{\pi^\varepsilon}^\varepsilon(x)\ge-\log(1+|x|)-C,
 \qquad
 \pi^\varepsilon(x)\le C(1+|x|)^\beta e^{-c|x|^2}.
\end{equation}
The second inequality follows from \eqref{eq:quadratic-lower-bound} and \eqref{eq:thermal-Zlower}. It implies \eqref{eq:thermal-bounds} for $j=0$ with a weaker Gaussian exponent. The singular minimizer satisfies the same bound.

The force estimate in \cref{lem:heat-force-estimates} gives
\begin{equation}\label{eq:thermal-force-split}
 \norm{\nabla\g_\varepsilon*f}_{L^\infty}
 \le \norm{\nabla\g_\varepsilon}_{L^1(B_1)}\norm f_{L^\infty}
    +\norm{\nabla\g_\varepsilon}_{L^\infty(B_1^c)}\norm f_{L^1}
 \le C(\norm f_{L^1}+\norm f_{L^\infty}).
\end{equation}
The convolution $\nabla\g*\pi$ is continuous: its local part is controlled by translation continuity in $L^1$, and its far-field part by dominated convergence. Differentiating the Euler--Lagrange equation gives
\begin{equation}\label{eq:thermal-gradient-identity}
 \nabla\pi=-\beta(\nabla V+\nabla\g*\pi)\pi,
 \qquad
 \nabla^j(\nabla\g*\pi)=\nabla\g*\nabla^j\pi.
\end{equation}
The first identity gives a polynomial times Gaussian bound for $\nabla\pi$. The second identity and \eqref{eq:thermal-force-split} then permit induction in the derivative order. This proves smoothness and \eqref{eq:thermal-bounds}, also for the cutoff equilibrium densities. Each polynomial factor is absorbed by the weaker Gaussian weight in that display.

The derivatives of $\nabla\widetilde V$, and the force and its derivatives convolved with a Gaussian-decaying density, tend to zero at infinity. For the derivatives this follows by splitting $|y|\le|x|/2$ and its complement, differentiating the smooth far-field kernel on the first region, and using the Gaussian derivative bounds on the second. The same splitting is uniform in the cutoff. Since $\nabla^2V=A+\nabla^2\widetilde V$ and $\nabla^2\widetilde V\to0$, this proves \eqref{eq:thermal-score} after enlarging the exterior ball.

\parheading{The one-body Poincar\'e inequality.}
By \eqref{eq:thermal-score}, $U_\pi$ has bounded Hessian and is uniformly convex outside a ball, with constants independent of $\varepsilon$ for $U_{\pi^\varepsilon}$. A bounded smooth radial modification $\chi$ may therefore be chosen, independently of $\varepsilon$, so that $\nabla^2(U_{\pi^\varepsilon}+\chi)\ge\rho_1\Id$ for some $\rho_1>0$. Indeed, choose $\chi'(r)=Kr$ on the fixed interior ball, then decrease $\chi'$ to zero on a sufficiently wide annulus with its negative radial second derivative smaller than half the exterior convexity constant. Keep $\chi'\ge0$ on that annulus. Applying the curvature criterion to the measure proportional to $e^{-U_{\pi^\varepsilon}-\chi}$, and then the bounded-perturbation principle (see, e.g., \cite[Sections 4.2 and 4.8]{BGL}), gives
\begin{equation}\label{eq:thermal-Poincare-transfer}
 \operatorname{Var}_{\pi^\varepsilon}(\psi)
 \le e^{\operatorname{osc}\chi}\rho_1^{-1}
       \int|\nabla\psi|^2\dd\pi^\varepsilon,
 \qquad(\psi\in H^1(\pi^\varepsilon),\ 0\le\varepsilon\le\varepsilon_*).
\end{equation}
Thus, the Poincar\'e constant is independent of the heat cutoff.

\parheading{Convergence of the thermal equilibrium measures.}
Tightness and the preceding bounds give subsequential convergence of the cutoff minimizers $\pi^\varepsilon$ as $\varepsilon\downarrow0$. Lower semicontinuity and testing against the singular minimizer $\pi$ identify every limit with $\pi$; alternatively pass the Euler--Lagrange equation by the same near-/far-field splitting and use uniqueness. The force satisfies
\begin{equation}\label{eq:thermal-force-convergence-explicit}
 \norm{\nabla\g_\varepsilon-\nabla\g}_{L^1(\R^\d)}\longrightarrow0,
 \qquad
 \sup_{r\le|x|\le R}|\nabla^j(\nabla\g_\varepsilon-\nabla\g)(x)|\longrightarrow0
\end{equation}
as $\varepsilon\downarrow0$, for every $0<r<R<\infty$ and $j\ge0$; moreover, $\sup_{0\le\varepsilon\le\varepsilon_*}\norm{\nabla\g_\varepsilon}_{L^\infty(B_1^c)}<\infty$, by \cref{lem:heat-force-estimates}. The differentiated Euler--Lagrange equations and the uniform Gaussian bounds give local convergence of every derivative by induction. To obtain \eqref{eq:thermal-cutoff-convergence}, first use local derivative convergence on $B_R$. On $B_R^c$, the Gaussian derivative bound, uniform in $\varepsilon$, gives, for $a'<a/2$,
\begin{equation}\label{eq:thermal-tail-convergence}
 \sup_{0\le\varepsilon\le\varepsilon_*}\sup_{|x|\ge R}
 e^{a'|x|^2}|\nabla^j(\pi^\varepsilon-\pi)(x)|
 \le C_j e^{-(a/2-a')R^2}.
\end{equation}
The corresponding $L^1$ tail is uniformly small as well. Combining these tail bounds with convergence on $B_R$, then letting $\varepsilon\downarrow0$ and $R\to\infty$, proves \eqref{eq:thermal-cutoff-convergence}.

\end{proof}

\subsection{Elementary bounds in the logarithmic derivative spaces}\label{app:logarithmic-maps}
Recall that constants are identified with zero in $\Xspace^j$. When a pointwise representative is needed we impose $\pi(z)=0$. Then
\begin{equation}\label{eq:linear-growth-representative}
 |z(x)|\le\norm{\nabla z}_\infty\left(|x|+\int|y|\dd\pi(y)\right),\qquad
 \norm z_{L^2(\pi)}\le\sqrt{C_P}\norm z_{\Xspace^j}.
\end{equation}
In particular, bounded sets in $\Xspace^j$ have a common linear-growth bound.

For an integrable function $f$ with Gaussian-decaying derivatives, the interaction term in the logarithmic equation is
\begin{equation}\label{eq:Rf}
 \mathcal R_\pi f:=\pi^{-1}\diver\bigl(\pi(\nabla\g*f)\bigr)
 =\diver(\nabla\g*f)-\nabla U_\pi\cdot(\nabla\g*f).
\end{equation}
The next lemma separates the linear perturbation of the Ornstein--Uhlenbeck generator from the quadratic remainder. Both map $\Xspace^m$ into $\Xspace^{m-1}$; the loss of one derivative is compensated by the smoothing estimate in \cref{app:linear-semigroup}.

\begin{lemma}[Convolution and nonlinear estimates]\label{lem:X-maps}
For $m\ge3$, the operator
\begin{equation}\label{eq:D-operator}
 Dz:=-(\nabla\widetilde V+\nabla\g*\pi)\cdot\nabla z+
          \mathcal R_\pi\bigl(\pi(z-\pi(z))\bigr)
\end{equation}
maps $\Xspace^m$ boundedly into $\Xspace^{m-1}$. For the representative $Dz$ defined by \eqref{eq:D-operator},
\begin{equation}\label{eq:D-uniform-decay}
 \lim_{R\to\infty}\sup_{\norm z_{\Xspace^m}\le1}
       \sup_{|x|\ge R}\sum_{j=0}^{m-1}|\nabla^jDz(x)|=0.
\end{equation}
Let $P_t^A$ be the Ornstein--Uhlenbeck semigroup given by \eqref{eq:OU-formula}. For every $0<a<b<\infty$, the set
\begin{equation}\label{eq:OU-D-compact-family}
 \{P_t^ADz:\ t\in[a,b],\ \norm z_{\Xspace^m}\le1\}
\end{equation}
has compact closure in $\Xspace^m$. In particular, $P_t^AD:\Xspace^m\to\Xspace^m$ is compact for every $t>0$.

Let $\mu_z=\pi e^z/\int\pi e^z$, take $\pi(z)=0$, and set
\begin{equation}\label{eq:nonlinear-remainder}
 \mathcal N(z):=\beta^{-1}|\nabla z|^2+
       (\nabla\g*(\mu_z-\pi))\cdot\nabla z+
       \mathcal R_\pi(\mu_z-\pi-\pi z).
\end{equation}
For every sufficiently small fixed ball in $\Xspace^m$,
\begin{equation}\label{eq:N-bounds}
 \norm{\mathcal N(z)}_{\Xspace^{m-1}}\le C\norm z_{\Xspace^m}^2,
 \qquad
 \norm{\mathcal N(z)-\mathcal N(w)}_{\Xspace^{m-1}}
       \le C(\norm z_{\Xspace^m}+\norm w_{\Xspace^m})\norm{z-w}_{\Xspace^m}.
\end{equation}
The corresponding cutoff maps converge to these maps in operator norm for $D$ and locally uniformly, with their Lipschitz bounds, for $\mathcal N$.
\end{lemma}
\begin{proof}
\parheading{Convolution bounds and the derivative count.}
We first estimate $\mathcal R_\pi f$ and the terms defining $D$. If $z$ is centered, \eqref{eq:linear-growth-representative} and \cref{lem:thermal} give common Gaussian bounds with polynomial prefactors for all derivatives of $\pi z$ through order $m$. Every convolution with the force is bounded by splitting its local $L^1$ and bounded far-field parts. The coefficient $\nabla U_\pi=\beta Ax+\beta\nabla\widetilde V+\beta\nabla\g*\pi$ requires one additional observation. The coefficient $\nabla\widetilde V$ and its derivatives are bounded and tend to zero at infinity. For the linear part of $\nabla U_\pi$, use
\begin{equation}\label{eq:x-convolution}
 x_i(\partial_j\g*f)(x)=((z_i\partial_j\g)*f)(x)
                       +(\partial_j\g*(y_i f))(x).
\end{equation}
Both kernels on the right are locally integrable and bounded on $B_1^c$. The second term is estimated by \eqref{eq:thermal-force-split} with $y_i f$ in place of $f$. In the divergence term of $\mathcal R_\pi$, write
\begin{equation}\label{eq:Rpi-derivative-count}
 \diver(\nabla\g*f)=\sum_{i=1}^{\d}(\partial_i\g)*(\partial_i f).
\end{equation}
Apply \eqref{eq:x-convolution} after transferring further derivatives to $f$. Derivatives of $\mathcal R_\pi f$ through order $m-1$ therefore use derivatives of $f$ through order $m$, and no higher derivative of the singular kernel is needed. More explicitly,
\begin{equation}\label{eq:Rpi-weighted-density-bound}
 \sum_{j=0}^{m-1}\norm{\nabla^j\mathcal R_\pi f}_{L^\infty}
 \le C\sum_{j=0}^m\left(
     \norm{(1+|x|)\nabla^j f}_{L^1}
     +\norm{(1+|x|)\nabla^j f}_{L^\infty}\right).
\end{equation}
For $f=\pi(z-\pi(z))$, \eqref{eq:linear-growth-representative} and Leibniz' rule give
\begin{equation}\label{eq:pi-z-weighted-bound}
 \sum_{j=0}^m\left(
     \norm{(1+|x|)\nabla^j[\pi(z-\pi(z))]}_{L^1}
     +\norm{(1+|x|)\nabla^j[\pi(z-\pi(z))]}_{L^\infty}\right)
 \le C\norm z_{\Xspace^m}.
\end{equation}
These estimates and the bounded derivatives of the coefficient of $\nabla z$ in \eqref{eq:D-operator} prove the boundedness of $D$.

\parheading{Decay at infinity.}
For positive Riesz exponents, the first term on the right of \eqref{eq:x-convolution} tends to zero, because $z_i\partial_j\g(z)=O(|z|^{-\s})$; the remaining terms also decay. For the logarithm, the same conclusion uses $\int f=0$. The kernel $z\mapsto z_i\partial_j\g(z)$ is bounded on $\R^\d\setminus\{0\}$, and its gradient is bounded by $C|z|^{-1}$. Its value at the origin does not affect the convolution. Thus, subtracting its value at $x$ gives
\begin{equation}\label{eq:log-convolution-cancellation}
 \left|((z_i\partial_j\g)*f)(x)\right|
 \le\frac C{|x|}\int|y|\,|f(y)|\dd y+C\int_{|y|>|x|/2}|f(y)|\dd y
 \xrightarrow[|x|\to\infty]{}0.
\end{equation}
The moments and tails on the right are common on the bounded input sets under consideration. Derivatives of $f$ have zero integral as well. This proves uniform decay for the centered $f=\pi(z-\pi(z))$ and for every derivative needed above. The coefficients $\nabla\widetilde V$ and $\nabla\g*\pi$, together with all derivatives required here, tend to zero at infinity. Their products with $\nabla z$ therefore tend to zero at infinity, uniformly on the bounded input sets. This proves \eqref{eq:D-uniform-decay}. It also shows that these images lie in the stated completion, by smooth spatial cutoff.

\parheading{Compactness after Ornstein--Uhlenbeck smoothing.}
Fix $0<a<b<\infty$. For $t\in[a,b]$, the Gaussian formula \eqref{eq:OU-formula} for $P_t^A$ gives two additional bounded derivatives of the image after convolution. For $t\in[a,b]$ and $\norm z_{\Xspace^m}\le1$, the derivatives of $P_t^ADz$ vanish uniformly at infinity, since the singular values of $e^{-At}$ stay away from zero and the Gaussian convolution kernels have uniform integrable tails. Hence the image is equicontinuous through order $m$, uniformly vanishes at infinity, and is bounded. Arzel\`a--Ascoli on a large ball, followed by the uniform tail bound, proves compactness in the full bounded-derivative norm and hence in $\Xspace^m$. This argument is uniform for $t$ in a compact positive interval.

\parheading{The nonlinear estimates.}
Take $\pi(z)=0$. Jensen gives $\int e^z\dd\pi\ge1$, and \eqref{eq:linear-growth-representative} controls $e^{|z(x)|}$ by an exponential of $C\norm z_{\Xspace^m}(1+|x|)$. Taylor's formula gives
\begin{equation}\label{eq:normalized-exponential-remainder}
 e^z-1-z=z^2\int_0^1(1-r)e^{rz}\dd r,
 \qquad
 \int e^z\dd\pi-1=\int(e^z-1-z)\dd\pi.
\end{equation}
In particular, the exact quadratic remainder of the normalized density is
\begin{equation}\label{eq:normalized-density-remainder-formula}
 \mu_z-\pi-\pi z
 =\frac{\pi}{\int e^z\dd\pi}
   \left[e^z-1-z-(1+z)\int(e^z-1-z)\dd\pi\right].
\end{equation}
Differentiate this identity through order $m$. Each term in its right-hand side contains at least two factors of $z$ or its derivatives; undifferentiated factors of $z$ have at most linear growth. The Gaussian derivative bounds for $\pi$ absorb these polynomial factors and $e^{C\norm z(1+|x|)}$ on a fixed $\Xspace^m$ ball. The analogous first-order identity gives, for every integer $q\ge0$,
\begin{align}
 \sum_{j=0}^m\left(\norm{(1+|x|)^q\nabla^j(\mu_z-\pi)}_{L^1}
       +\norm{(1+|x|)^q\nabla^j(\mu_z-\pi)}_{L^\infty}\right)
       &\le C_q\norm z_{\Xspace^m},\label{eq:normalized-first-remainder}\\
 \sum_{j=0}^m\left(\norm{(1+|x|)^q\nabla^j(\mu_z-\pi-\pi z)}_{L^1}
       +\norm{(1+|x|)^q\nabla^j(\mu_z-\pi-\pi z)}_{L^\infty}\right)
       &\le C_q\norm z_{\Xspace^m}^2.\label{eq:normalized-second-remainder}
\end{align}
For the difference estimates, differentiate the normalized density along the segment joining $w$ to $z$:
\begin{equation}\label{eq:normalized-density-directional-derivative}
 \frac{\mathrm d}{\mathrm d\theta}\mu_{w+\theta(z-w)}
 =\mu_{w+\theta(z-w)}
   \left[z-w-\int(z-w)\dd\mu_{w+\theta(z-w)}\right].
\end{equation}
The same weighted bounds, integrated in $\theta$, yield
\begin{align}
 &\sum_{j=0}^m\left(
 \norm{(1+|x|)\nabla^j[\mu_z-\mu_w-\pi(z-w)]}_{L^1}
 +\norm{(1+|x|)\nabla^j[\mu_z-\mu_w-\pi(z-w)]}_{L^\infty}\right)\notag\\
 &\hspace{2em}\le C(\norm z_{\Xspace^m}+\norm w_{\Xspace^m})
                      \norm{z-w}_{\Xspace^m}.
 \label{eq:normalized-remainder-Lipschitz}
\end{align}
The same sum without the linear subtraction is bounded by $C\norm{z-w}_{\Xspace^m}$. Apply \eqref{eq:Rpi-weighted-density-bound} to the quadratic remainder and its difference. The two remaining terms in \eqref{eq:nonlinear-remainder} satisfy
\begin{equation}\label{eq:nonlinear-local-products}
 \norm{|\nabla z|^2}_{\Xspace^{m-1}}
 +\norm{[\nabla\g*(\mu_z-\pi)]\cdot\nabla z}_{\Xspace^{m-1}}
 \le C\norm z_{\Xspace^m}^2,
\end{equation}
by Leibniz' rule, \eqref{eq:thermal-force-split}, and \eqref{eq:normalized-first-remainder}; their differences satisfy the corresponding bilinear bound. These are precisely \eqref{eq:N-bounds}. The density remainders have integral zero, so the decay argument above applies to their convolution terms. The remaining term $|\nabla z|^2$ tends to the constant $|a_z|^2$ at infinity, and its positive-order derivatives vanish there, as required by the quotient space.

\parheading{Cutoff convergence.}
The preceding estimates hold on every fixed bounded $\Xspace^m$ ball, with constants depending on its radius. The heat representation gives the following bounds on the force difference $\nabla\g_\varepsilon-\nabla\g$:
\begin{align}
 \norm{\nabla\g_\varepsilon-\nabla\g}_{L^1}
 &\le C\varepsilon^{(\d-\s-1)/2},\label{eq:logmaps-cutoff-force-L1}\\
 \norm{|x|(\nabla\g_\varepsilon-\nabla\g)}_{L^1}
 &\le C\varepsilon^{(\d-\s)/2}.\label{eq:logmaps-cutoff-force-weighted}
\end{align}
Indeed, these follow by integrating respectively $\norm{\nabla p_r}_{L^1}=C r^{-1/2}$ and $\norm{|x|\nabla p_r}_{L^1}=C$ over $0<r<\varepsilon$ with the weight in \eqref{eq:heat-cutoff}. Both powers are positive in the stated range.

The operators $D^\varepsilon$ and $D$ use the centered densities $\pi^\varepsilon(z-\pi^\varepsilon(z))$ and $\pi(z-\pi(z))$, respectively. Choose the representative satisfying $\pi(z)=0$. By \eqref{eq:linear-growth-representative},
\begin{equation}\label{eq:cutoff-centering-convergence}
 \sup_{\substack{\norm z_{\Xspace^m}\le1\\ \pi(z)=0}}|\pi^\varepsilon(z)|
 \le C\norm{(1+|x|)(\pi^\varepsilon-\pi)}_{L^1}
 \xrightarrow[\varepsilon\downarrow0]{}0.
\end{equation}
Here, \eqref{eq:thermal-cutoff-convergence} implies convergence of the weighted norm on the right. For the density difference, the exact identity is
\begin{equation}\label{eq:cutoff-centered-density-difference}
 \pi^\varepsilon(z-\pi^\varepsilon(z))-\pi z
 =(\pi^\varepsilon-\pi)z-\pi^\varepsilon\pi^\varepsilon(z),
 \qquad \pi(z)=0.
\end{equation}
Leibniz' rule, $|z(x)|\le C(1+|x|)$, and the unit bound on the positive-order derivatives of $z$ give, for $p\in\{1,\infty\}$,
\begin{align}
 &\sup_{\substack{\norm z_{\Xspace^m}\le1\\ \pi(z)=0}}
   \sum_{j=0}^m\left\|(1+|x|)\nabla^j
       [\pi^\varepsilon(z-\pi^\varepsilon(z))-\pi z]\right\|_{L^p}\notag\\
 &\quad\le C\sum_{j=0}^m
       \left\|(1+|x|)^2\nabla^j(\pi^\varepsilon-\pi)\right\|_{L^p}
       +C\norm{(1+|x|)(\pi^\varepsilon-\pi)}_{L^1}
 \xrightarrow[\varepsilon\downarrow0]{}0.
 \label{eq:cutoff-centered-density-norms}
\end{align}
The term containing $\pi^\varepsilon(z)$ is bounded using \eqref{eq:cutoff-centering-convergence} and the cutoff-uniform Gaussian derivative bounds for $\pi^\varepsilon$ in \cref{lem:thermal}. The remaining terms tend to zero by \eqref{eq:thermal-cutoff-convergence}. Combining this with \eqref{eq:logmaps-cutoff-force-L1}--\eqref{eq:logmaps-cutoff-force-weighted} in \eqref{eq:x-convolution} proves
\begin{equation}\label{eq:cutoff-D-convergence}
 \norm{D^\varepsilon-D}_{\Xspace^m\to\Xspace^{m-1}}
 \xrightarrow[\varepsilon\downarrow0]{}0.
\end{equation}
For the normalized densities, use \eqref{eq:normalized-density-remainder-formula} and \eqref{eq:normalized-density-directional-derivative} with $\pi^\varepsilon$ in place of $\pi$; subtracting a constant from $z$ does not change either normalized density. The same estimates give, for each fixed $R>0$,
\begin{equation}\label{eq:cutoff-N-convergence-explicit}
 \sup_{\norm z_{\Xspace^m}\le R}
 \norm{\mathcal N^\varepsilon(z)-\mathcal N(z)}_{\Xspace^{m-1}}
 \xrightarrow[\varepsilon\downarrow0]{}0.
\end{equation}
The Lipschitz bounds on that ball are independent of $\varepsilon\in[0,\varepsilon_*]$.
\end{proof}

\subsection{The linearized semigroup}\label{app:linear-semigroup}
We first construct the linearized evolution in the logarithmic derivative norm, then prove its exponential decay. The one-body Poincar\'e inequality gives decay in weighted $L^2$. To obtain decay in $\Xspace^m$, we also use compactness of the difference between the linearized and Ornstein--Uhlenbeck semigroups.

The linearized operator is
\begin{equation}\label{eq:linearized-operator}
 \mathcal L_\pi=-\beta^{-1}\Aop_\pi(\Id+\beta T_\pi),
\end{equation}
where $\Aop_\pi$ is the weighted Laplacian in \cref{lem:weighted-elliptic} and $T_\pi$ is the centered interaction operator in \cref{lem:screened-inverse}.

The anisotropic Ornstein--Uhlenbeck semigroup with generator $\beta^{-1}\Delta-Ax\cdot\nabla$ acts on the quotient spaces by
\begin{equation}\label{eq:OU-formula}
 P_t^Az(x)=\E\left[z\left(e^{-At}x+\Sigma_t^{1/2}Z\right)\right],
 \qquad \Sigma_t=\beta^{-1}A^{-1}(\Id-e^{-2At}),\qquad Z\sim\mathcal N(0,\Id).
\end{equation}
Here, $\lambda_A>0$ denotes the smallest eigenvalue of $A$. For a smooth representative of an element of $\Xspace^j$,
\begin{equation}\label{eq:OU-derivative-identity}
 \nabla^jP_t^Az=(e^{-At})^{\otimes j}P_t^A(\nabla^jz).
\end{equation}
Strong continuity holds first on affine-plus-compact smooth functions. The boundedness supplied by \eqref{eq:OU-derivative-identity} and density then extend it to $\Xspace^j$. Differentiating the Gaussian formula \eqref{eq:OU-formula} also gives
\begin{equation}\label{eq:OU-estimates}
 \norm{P_t^A}_{\Xspace^j\to\Xspace^j}\le C_Ae^{-\lambda_A t},\qquad
 \norm{P_t^A}_{\Xspace^{j-1}\to\Xspace^j}\le C_{A,j}(1+t^{-1/2})e^{-\lambda_A t},
                    \quad j\ge2,\ t>0.
\end{equation}
Throughout this subsection, constants may depend on the fixed parameters $\beta,A,j$.

\begin{lemma}[A stable linearized semigroup in logarithmic derivative norms]\label{lem:linear-X}
The operator $\mathcal L_\pi$ in \eqref{eq:linearized-operator} generates on $\Xspace^m$ a strongly continuous semigroup $S(t)$. There exist $C>0$ and an exponential decay rate $\omega_0>0$ such that
\begin{equation}\label{eq:S-X}
 \norm{S(t)}_{\Xspace^m\to\Xspace^m}\le C e^{-\omega_0t},\qquad
 \norm{S(t)}_{\Xspace^{m-1}\to\Xspace^m}
       \le C(1+t^{-1/2})e^{-\omega_0t}.
\end{equation}
The constants can be chosen uniformly for sufficiently small heat cutoffs.
\end{lemma}
\begin{proof}
\parheading{Construction and local smoothing.}
On affine-plus-compact smooth representatives, the thermal equation \eqref{eq:thermal} gives
\begin{equation}\label{eq:linear-OU-decomposition}
 \mathcal L_\pi=(\beta^{-1}\Delta-Ax\cdot\nabla)+D,
\end{equation}
with $D$ defined by \eqref{eq:D-operator}. We construct $S$ through the Volterra equation
\begin{equation}\label{eq:S-volterra}
 S(t)z=P_t^Az+\int_0^tP_{t-r}^A D S(r)z\dd r.
\end{equation}
By \cref{lem:X-maps} and \eqref{eq:OU-estimates}, the integral operator on $C([0,T];\Xspace^m)$ has norm at most $C(T+\sqrt T)$. A contraction on a short interval, followed by continuation, gives a unique solution and bounds on every finite interval. Uniqueness gives the semigroup property, and the integral equation gives strong continuity at zero.

For data in $\Xspace^{m-1}$, use continuity in that space together with
\begin{equation}\label{eq:local-smoothing-norm}
 \sup_{0<t\le T}t^{1/2}\norm{S(t)z}_{\Xspace^m}<\infty.
\end{equation}
The only additional convolution is
\begin{equation}\label{eq:smoothing-convolution}
 t^{1/2}\int_0^t(t-r)^{-1/2}r^{-1/2}\dd r=\pi t^{1/2}.
\end{equation}
Thus, the same contraction applies when $T$ is small and gives the local one-derivative smoothing estimate. The integral term is continuous at zero in $\Xspace^{m-1}$ because its norm there is bounded by a constant times $\int_0^t r^{-1/2}\dd r$.

\parheading{Compactness of the semigroup difference.}
Fix $t>0$ and $0<\eta<t$. For the part of \eqref{eq:S-volterra} where $t-r<\eta$, the local bounds give
\begin{equation}\label{eq:compact-short-integral}
 \left\|\int_{t-\eta}^tP_{t-r}^ADS(r)\dd r\right\|_{\Xspace^m\to\Xspace^m}
 \le C_t(\eta+\sqrt\eta).
\end{equation}
For $0\le r\le t-\eta$, set $M_t:=\sup_{0\le r\le t}\norm{S(r)}_{\Xspace^m\to\Xspace^m}$. By \eqref{eq:OU-D-compact-family},
\begin{equation}\label{eq:compact-integral-image-set}
 \mathcal C_{t,\eta}:=
 \overline{\{P_s^ADz:\ \eta\le s\le t,\ \norm z_{\Xspace^m}\le M_t\}}
 \quad\text{is compact in }\Xspace^m.
\end{equation}
For every $z$ with $\norm z_{\Xspace^m}\le1$,
\begin{equation}\label{eq:compact-integral-hull}
 \frac1{t-\eta}\int_0^{t-\eta}P_{t-r}^ADS(r)z\dd r
 \in\overline{\operatorname{conv}}(\mathcal C_{t,\eta}).
\end{equation}
The closed convex hull of a compact subset of a Banach space is compact, so this part of the integral is a compact operator. Taking $\eta\downarrow0$ in \eqref{eq:compact-short-integral}, and using that compact operators are norm closed, gives
\begin{equation}\label{eq:compact-semigroup-difference}
 S(t)-P_t^A:\Xspace^m\longrightarrow\Xspace^m
 \quad\text{compact for every }t>0.
\end{equation}

\parheading{Identification with the Hilbert-space semigroup.}
On $L^2_0(\pi)$, set $B:=\Id+\beta T_\pi$. This operator is bounded and self-adjoint, and $B\ge\Id$ by positivity of the centered interaction. By the representation theorem for densely defined closed nonnegative forms (see, e.g., \cite[Chapter VI, Section 2]{Kato}), there is a unique nonnegative self-adjoint operator $H$ associated with the quadratic form
\begin{equation}\label{eq:Hilbert-form-domain}
 a\longmapsto\norm{\nabla(B^{1/2}a)}_{L^2(\pi)}^2,
 \qquad
 \{a\in L^2_0(\pi):B^{1/2}a\in H^1(\pi)\}.
\end{equation}
The form is closed and densely defined because $B^{1/2}$ is bounded and invertible. The Poincar\'e inequality gives
\begin{equation}\label{eq:Hilbert-form-gap}
 \norm{\nabla(B^{1/2}a)}_{L^2(\pi)}^2
 \ge C_P^{-1}\norm{B^{1/2}a}_{L^2(\pi)}^2
 \ge C_P^{-1}\norm a_{L^2(\pi)}^2.
\end{equation}
Consequently, the semigroup realizing $-\beta^{-1}\Aop_\pi B$ is
\begin{equation}\label{eq:L2-semigroup}
 S_2(t)=B^{-1/2}e^{-tH/\beta}B^{1/2},\qquad
 \norm{S_2(t)}_{L^2_0(\pi)\to L^2_0(\pi)}
       \le\sqrt{\norm B}\,e^{-t/(\beta C_P)}.
\end{equation}

We check that this is the same evolution as \eqref{eq:S-volterra} after choosing centered representatives. Schur's test, applied as in \eqref{eq:force-schur}, gives
\begin{equation}\label{eq:B-form-domain}
 \norm{\nabla T_\pi a}_{L^2(\pi)}
  =\norm{\nabla\g*(\pi a)}_{L^2(\pi)}\le C\norm a_{L^2(\pi)}.
\end{equation}
Thus, $B$ maps $H^1(\pi)$ continuously to itself. Its inverse does also: for $a=B^{-1}f$, the identity $a=f-\beta T_\pi a$ and \eqref{eq:B-form-domain} control $\nabla a$ by $\norm f_{H^1(\pi)}$. The Hilbert-space evolution satisfies
\begin{equation}\label{eq:linear-weak-equation}
 \langle\partial_ta,\psi\rangle_{H^1(\pi)^*,H^1(\pi)}
 =-\beta^{-1}\int\nabla(Ba)\cdot\nabla\psi\dd\pi
 \qquad(\psi\in H^1(\pi)).
\end{equation}
For a difference of two weak solutions in $L^2_{\mathrm{loc}}(H^1(\pi))$ with time derivative in $L^2_{\mathrm{loc}}(H^1(\pi)^*)$, testing by $Ba$ yields
\begin{equation}\label{eq:linear-weak-uniqueness}
 \frac{\mathrm d}{\mathrm dt}\langle a,Ba\rangle_{L^2(\pi)}
 =-2\beta^{-1}\norm{\nabla(Ba)}_{L^2(\pi)}^2.
\end{equation}
Time mollification justifies the identity, since $B$ is self-adjoint and preserves the form domain. It proves uniqueness for zero initial difference.

For affine-plus-compact smooth initial data, the centered representative of \eqref{eq:S-volterra} satisfies \eqref{eq:linear-weak-equation}. Writing $a_t:=S(t)z-\pi(S(t)z)$ for this representative, \eqref{eq:linear-growth-representative} and \eqref{eq:B-form-domain} give, for every $T<\infty$,
\begin{equation}\label{eq:X-semigroup-energy-class}
 a\in C([0,T];L^2_0(\pi))\cap L^2(0,T;H^1(\pi)),
 \qquad \partial_ta\in L^2(0,T;H^1(\pi)^*).
\end{equation}
Thus, the centered solution lies in the class for which weak uniqueness was just proved. The spatial cutoff and operator-core argument of \cref{lem:weighted-elliptic} justify the weak equation and the energy test. Uniqueness, and then density in $\Xspace^m$, give
\begin{equation}\label{eq:semigroup-realization-identity}
 S(t)z-\pi(S(t)z)=S_2(t)(z-\pi(z)).
\end{equation}
The expression on the left is independent of the representative chosen for the quotient element $S(t)z$.

\parheading{Exponential decay in $\Xspace^m$.}
Choose $t_0$ so that $\norm{P_{t_0}^A}_{\Xspace^m\to\Xspace^m}<1/4$. By \eqref{eq:compact-semigroup-difference}, $S(t_0)$ is a compact perturbation of $P_{t_0}^A$. For $|\zeta|>1/4$,
\begin{equation}\label{eq:Fredholm-factorization}
 \zeta\Id-S(t_0)
 =(\zeta\Id-P_{t_0}^A)
 \left[\Id-(\zeta\Id-P_{t_0}^A)^{-1}(S(t_0)-P_{t_0}^A)\right].
\end{equation}
Here, the spectral argument is on the complexification $\Xspace^m_{\mathbb C}$ of the real Banach space $\Xspace^m$, with $P_{t_0}^A$ and $S(t_0)$ extended complex linearly. The Fredholm alternative shows that every spectral value $\zeta\in\sigma(S(t_0)|_{\Xspace^m_{\mathbb C}})$ with $|\zeta|>1/4$ is an eigenvalue of $S(t_0)$ on this space; see, e.g., \cite[Chapters IV--V]{EngelNagel}.

Let $z\ne0$ be a corresponding eigenvector in $\Xspace^m_{\mathbb C}$. Its centered representative $a=z-\pi(z)$ is nonzero in $L^2_0(\pi;\mathbb C)$. By \eqref{eq:semigroup-realization-identity} and the semigroup property,
\begin{equation}\label{eq:eigenvector-iteration}
 |\zeta|^n\norm a_{L^2(\pi)}
 =\norm{S_2(nt_0)a}_{L^2(\pi)}
 \le\sqrt{\norm B}\,e^{-nt_0/(\beta C_P)}\norm a_{L^2(\pi)}.
\end{equation}
Taking $n$th roots gives
\begin{equation}\label{eq:eigenvalue-bound}
 |\zeta|\le e^{-t_0/(\beta C_P)}<1.
\end{equation}
All other spectral values lie in the disk of radius $1/4$. Hence $S(t_0)$ has spectral radius strictly below one. The spectral-radius formula gives $\norm{S(nt_0)}\le Cq^n$ for some $q<1$. Boundedness for $0\le t\le t_0$ then proves the first estimate in \eqref{eq:S-X}. For $t\ge1$, the semigroup property and the local smoothing estimate give
\begin{equation}\label{eq:longtime-one-derivative-smoothing}
 \norm{S(t)}_{\Xspace^{m-1}\to\Xspace^m}
 \le\norm{S(t-1)}_{\Xspace^m\to\Xspace^m}
      \norm{S(1)}_{\Xspace^{m-1}\to\Xspace^m}
 \le C e^{-\omega_0t}.
\end{equation}
Together with the local estimate for $0<t\le1$, this proves the second bound in \eqref{eq:S-X}.

\parheading{Estimates uniform in the heat cutoff.}
The convergence in \eqref{eq:cutoff-D-convergence}, the bounds on $S^\varepsilon$ on $[0,T]$, independent of $\varepsilon$, and \eqref{eq:S-volterra} imply, for every $T<\infty$,
\begin{equation}\label{eq:cutoff-semigroup-convergence}
 \sup_{0\le t\le T}\norm{S^\varepsilon(t)-S(t)}_{\Xspace^m\to\Xspace^m}
 \le C_T\norm{D^\varepsilon-D}_{\Xspace^m\to\Xspace^{m-1}}
 \xrightarrow[\varepsilon\downarrow0]{}0.
\end{equation}
Choose $t_1$ with $\norm{S(t_1)}<1/4$. For sufficiently small cutoffs,
\begin{equation}\label{eq:cutoff-common-contraction}
 \norm{S^\varepsilon(t_1)}_{\Xspace^m\to\Xspace^m}<1/2.
\end{equation}
Iteration and the bounds on $[0,t_1]$, uniform for $0\le\varepsilon\le\varepsilon_*$, give an exponential decay rate independent of $\varepsilon$. The local smoothing constants are also independent of $\varepsilon$. All these operators act on the same quotient space $\Xspace^m$; the centering used in \eqref{eq:semigroup-realization-identity} does not change that space when $\pi$ is replaced by $\pi^\varepsilon$.
\end{proof}

\subsection{The nonlinear equation and its stability}\label{app:nonlinear-stability}
\begin{proof}[Proof of \cref{prop:background-stability}]
\parheading{The normalized logarithmic equation.}
For $\mu_z=\pi e^z/\int\pi e^z$, expansion of \eqref{eq:mf} gives, modulo spatial constants,
\begin{equation}\label{eq:log-PDE}
 \partial_tz=\mathcal L_\pi z+\mathcal N(z).
\end{equation}
To verify that normalization does not alter the density equation, first observe that the definitions imply
\begin{equation}\label{eq:logarithmic-equation-expansion}
 \mathcal L_\pi z+\mathcal N(z)
 =\mu_z^{-1}\left[\beta^{-1}\Delta\mu_z+
      \diver\bigl(\mu_z(\nabla V+\nabla\g*\mu_z)\bigr)\right],
 \qquad
 \int(\mathcal L_\pi z+\mathcal N(z))\dd\mu_z=0.
\end{equation}
The integral identity follows by integration by parts using the Gaussian decay. Choose the representative with $\pi(z)=0$. Its evolution and the derivative of its normalizer are
\begin{align}
 \partial_tz&=\mathcal L_\pi z+\mathcal N(z)
       -\pi(\mathcal L_\pi z+\mathcal N(z)),\label{eq:centered-log-equation}\\
 \partial_t\log\int e^z\dd\pi
      &=\int\partial_tz\dd\mu_z
       =-\pi(\mathcal L_\pi z+\mathcal N(z)).\label{eq:log-normalizer-derivative}
\end{align}
Subtracting the second equality from the first yields the original equation for $\partial_t\log\mu_z$. Thus, solving the quotient equation \eqref{eq:log-PDE} produces the normalized mean-field solution.

\parheading{Small-data stability and local existence.}
The integral equation in $\Xspace^m$ is
\begin{equation}\label{eq:nonlinear-Duhamel}
 z_t=S(t)z_0+\int_0^tS(t-r)\mathcal N(z_r)\dd r.
\end{equation}
Fix $0<\omega<\omega_0$. By \cref{lem:X-maps,lem:linear-X},
\begin{equation}\label{eq:nonlinear-weighted-bound}
 \sup_{t\ge0}e^{\omega t}
 \left\|\int_0^tS(t-r)\mathcal N(z_r)\dd r\right\|_{\Xspace^m}
 \le C\left(\sup_{r\ge0}e^{\omega r}\norm{z_r}_{\Xspace^m}\right)^2
       \int_0^\infty(1+r^{-1/2})e^{-(\omega_0-\omega)r}\dd r.
\end{equation}
The integral satisfies
\begin{equation}\label{eq:nonlinear-integrability}
 \int_0^\infty (1+r^{-1/2})e^{-(\omega_0-\omega)r}\dd r<\infty.
\end{equation}
Applying the difference estimate in \eqref{eq:N-bounds} in the same way shows that \eqref{eq:nonlinear-Duhamel} is a contraction on a ball of radius $2C\norm{z_0}_{\Xspace^m}$ in the exponentially weighted trajectory norm, when the initial norm is sufficiently small. This proves global existence, uniqueness in that class, and \eqref{eq:log-decay}. On any fixed bounded initial ball, the local Lipschitz estimate instead gives a contraction on a positive interval because $\int_0^T(1+r^{-1/2})\dd r=T+2\sqrt T$ tends to zero. This proves local existence for the data in part (iii).

\parheading{Higher regularity.}
For each integer $\ell>m$, the same product and convolution estimates give, on a bounded $\Xspace^m$ ball,
\begin{equation}\label{eq:higher-tame}
 \norm{\mathcal N(z)}_{\Xspace^{\ell-1}}
       \le C_\ell\norm z_{\Xspace^m}\norm z_{\Xspace^\ell}.
\end{equation}
Indeed, for $1\le j\le\ell$, bounded-derivative interpolation gives
\begin{equation}\label{eq:bounded-derivative-interpolation}
 \norm{\nabla^jz}_{L^\infty}
 \le C_\ell\norm{\nabla z}_{L^\infty}^{1-(j-1)/(\ell-1)}
             \norm{\nabla^\ell z}_{L^\infty}^{(j-1)/(\ell-1)}.
\end{equation}
Apply this inequality to the products in Leibniz' rule and to the differentiated Taylor remainder \eqref{eq:normalized-density-remainder-formula}. On a fixed $\Xspace^m$ ball, one obtains
\begin{align}
 \norm{|\nabla z|^2}_{\Xspace^{\ell-1}}
 &\le C_\ell\norm z_{\Xspace^m}\norm z_{\Xspace^\ell},\label{eq:tame-gradient-product}\\
 \norm{[\nabla\g*(\mu_z-\pi)]\cdot\nabla z}_{\Xspace^{\ell-1}}
 &\le C_\ell\norm z_{\Xspace^m}\norm z_{\Xspace^\ell},\label{eq:tame-nonlocal-product}\\
 \sum_{j=0}^\ell\left(\norm{(1+|x|)\nabla^j(\mu_z-\pi-\pi z)}_{L^1}
 +\norm{(1+|x|)\nabla^j(\mu_z-\pi-\pi z)}_{L^\infty}\right)
 &\le C_\ell\norm z_{\Xspace^m}\norm z_{\Xspace^\ell}.
 \label{eq:tame-density-remainder}
\end{align}
For the last estimate, every differentiated Taylor term contains at least two factors of $z$ or its derivatives; \eqref{eq:bounded-derivative-interpolation} places at most one factor in the highest norm, and \eqref{eq:linear-growth-representative} and the Gaussian bounds control undifferentiated factors. Applying \eqref{eq:Rpi-weighted-density-bound} with $m=\ell$ proves the same bound for $\mathcal R_\pi(\mu_z-\pi-\pi z)$. The three terms in \eqref{eq:nonlinear-remainder} therefore give \eqref{eq:higher-tame}.

The one-derivative smoothing estimate for $S$ at order $\ell$ follows from \eqref{eq:S-volterra}--\eqref{eq:smoothing-convolution} with $m$ replaced by $\ell$. On an interval of length $h$, the high-norm contribution of the integral in \eqref{eq:nonlinear-Duhamel} is at most
\begin{equation}\label{eq:high-norm-continuation}
 C_\ell(h+\sqrt h)
 \sup_{t\le r\le t+h}\norm{z_r}_{\Xspace^m}
 \sup_{t\le r\le t+h}\norm{z_r}_{\Xspace^\ell}.
\end{equation}
Choose $h$ so that the coefficient of the last supremum is below $1/2$, and continue on successive intervals. Every finite higher norm remains finite on compact time intervals. The differentiated equation then gives time regularity. Hence the smooth initial data in the proposition remain smooth, while the low-norm constants in \eqref{eq:log-decay} are independent of their higher derivatives.

\parheading{Convergence of the cutoff solutions.}
By \cref{lem:X-maps,lem:linear-X}, the solutions lie in a common bounded $\Xspace^m$ ball on each fixed finite interval of existence. To compare them, use the same Ornstein--Uhlenbeck semigroup for every cutoff. The decomposition \eqref{eq:linear-OU-decomposition} gives
\begin{equation}\label{eq:nonlinear-common-OU}
 z_t^\varepsilon=P_t^Az_0+
 \int_0^tP_{t-r}^A\bigl(D^\varepsilon z_r^\varepsilon+
                    \mathcal N^\varepsilon(z_r^\varepsilon)\bigr)\dd r.
\end{equation}
This identity also holds for $\varepsilon=0$. Fix a finite interval $[0,T]$ and a radius $R$ containing both trajectories. On $[0,h]$, with $0<h\le\min\{1,T\}$, subtract the two identities and use \eqref{eq:OU-estimates} and the uniform local Lipschitz bounds. We obtain
\begin{align}
 \sup_{0\le t\le h}\norm{z_t^\varepsilon-z_t}_{\Xspace^m}
 &\le C_R(h+\sqrt h)
       \sup_{0\le t\le h}\norm{z_t^\varepsilon-z_t}_{\Xspace^m}\notag\\
 &\quad+C_R(h+\sqrt h)\left(
       \norm{D^\varepsilon-D}_{\Xspace^m\to\Xspace^{m-1}}
       +\sup_{\norm w_{\Xspace^m}\le R}
        \norm{\mathcal N^\varepsilon(w)-\mathcal N(w)}_{\Xspace^{m-1}}
                       \right).
 \label{eq:nonlinear-cutoff-comparison}
\end{align}
Choose $h$, independently of $\varepsilon$, so that $C_R(h+\sqrt h)<1/2$ and absorb the first term. The remaining terms tend to zero by \eqref{eq:cutoff-D-convergence} and \eqref{eq:cutoff-N-convergence-explicit}. On the next interval of length at most $h$, the same estimate has the additional initial term $C\norm{z_h^\varepsilon-z_h}_{\Xspace^m}$, which already tends to zero. Iterating over finitely many intervals proves
\begin{equation}\label{eq:logarithmic-cutoff-convergence}
 \sup_{0\le t\le T}\norm{z_t^\varepsilon-z_t}_{\Xspace^m}\xrightarrow[\varepsilon\downarrow0]{}0.
\end{equation}
This argument applies both to the small global solutions and to the larger local solutions. The higher-regularity bounds above are uniform in the cutoff, so repeating the comparison at each finite higher order gives convergence in that order for smooth data. Representatives may be centered with respect to the fixed measure $\pi$, since the equation and norm are defined modulo constants.

\parheading{Density estimates and convergence.}
For the representative $\pi(z_t)=0$, Jensen's inequality and \eqref{eq:linear-growth-representative} give
\begin{equation}\label{eq:perturbed-Gaussian-bound}
 \int e^{z_t}\dd\pi\ge1,
 \qquad
 \mu_t(x)\le M\exp\!\left[-a|x|^2+
      \norm{\nabla z_t}_{L^\infty}\left(|x|+\int|y|\dd\pi(y)\right)\right].
\end{equation}
If $\sup_{t\in I}\norm{z_t}_{\Xspace^m}<\infty$, with $I=[0,\infty)$ for small data and $I=[0,T_0]$ for the local solution, \eqref{eq:perturbed-Gaussian-bound} gives a Gaussian upper bound with constants independent of $t\in I$. The same constants may be chosen for $\mu_t^\varepsilon$ when the $\Xspace^m$ bound is uniform in $\varepsilon\in[0,\varepsilon_*]$. Differentiating the normalized exponential gives the corresponding density-derivative bounds. In particular, \eqref{eq:thermal-cutoff-convergence} and \eqref{eq:logarithmic-cutoff-convergence} imply
\begin{equation}\label{eq:density-cutoff-L1-explicit}
 \sup_{0\le t\le T}\norm{\mu_t^\varepsilon-\mu_t}_{L^1(\R^\d)}
 \xrightarrow[\varepsilon\downarrow0]{}0,
 \qquad [0,T]\subset I.
\end{equation}
For smooth initial data, the convergence also holds for every spatial derivative uniformly on $[0,T]\times B_R$, for each $R<\infty$.

\parheading{The logarithmic Hessian and Poincar\'e inequality.}
The logarithmic Hessian is
\begin{equation}\label{eq:perturbed-log-Hessian}
 \nabla^2(-\log\mu_t)=\nabla^2U_\pi-\nabla^2z_t.
\end{equation}
For the small global solution, decrease $\delta_0$ so that $\sup_{t\ge0}\norm{\nabla^2z_t}_{L^\infty}$ is less than one quarter of $\beta\lambda_A$. Then \eqref{eq:thermal-score} and \eqref{eq:perturbed-log-Hessian} imply uniform convexity of $-\log\mu_t$ outside a fixed ball, independently of $t\ge0$. Applying the bounded modification and Poincar\'e perturbation argument of \cref{lem:thermal} gives a Poincar\'e constant independent of time.

For a larger local solution, continuity of $t\mapsto z_t$ in $\Xspace^m$ and \eqref{eq:X-tail-properties} give
\begin{equation}\label{eq:local-trajectory-uniform-Hessian-tail}
 \lim_{R\to\infty}\sup_{0\le t\le T_0}\sup_{|x|\ge R}
            |\nabla^2z_t(x)|=0.
\end{equation}
Indeed, the image of $[0,T_0]$ is compact in $\Xspace^m$, so a finite cover reduces this assertion to \eqref{eq:X-tail-properties}. Hence $-\log\mu_t$ is uniformly convex outside a ball with radius independent of $t\in[0,T_0]$. By \eqref{eq:logarithmic-cutoff-convergence} and the thermal Hessian bounds, the same radius, enlarged if necessary, and convexity constant work for $\mu_t^\varepsilon$, $0\le t\le T_0$, $0\le\varepsilon\le\varepsilon_*$. The Poincar\'e constant is therefore independent of both time and cutoff on this interval.

\parheading{Drift and continuity velocity.}
The first two drift derivatives can be estimated without differentiating the singular force:
\begin{equation}\label{eq:background-drift-derivatives}
 \nabla^jB_t=-\nabla^{j+1}V-\nabla\g*\nabla^j\mu_t
 \qquad(j=1,2).
\end{equation}
The confinement derivatives are bounded, and \eqref{eq:thermal-force-split} applies to the Gaussian-decaying density derivatives. This proves \eqref{eq:background-derivative-bound}.

Finally, the stationary thermal equation cancels the confinement in the continuity velocity:
\begin{equation}\label{eq:velocity-from-log}
 u_t=-\beta^{-1}\nabla z_t-\nabla\g*(\mu_t-\pi).
\end{equation}
The first-order version of \eqref{eq:normalized-first-remainder}, together with \eqref{eq:thermal-force-split}, yields
\begin{equation}\label{eq:velocity-X-control}
 \norm{u_t}_{W^{1,\infty}}\le C\norm{z_t}_{\Xspace^m}.
\end{equation}
This proves \eqref{eq:background-velocity-decay}. The same estimates with constants independent of $\varepsilon\in[0,\varepsilon_*]$, the force convergence from \cref{lem:heat-force-estimates}, and \eqref{eq:logarithmic-cutoff-convergence} give
\begin{equation}\label{eq:velocity-cutoff-convergence-explicit}
 \sup_{0\le t\le T}\norm{u_t^\varepsilon-u_t}_{W^{1,\infty}}
 \xrightarrow[\varepsilon\downarrow0]{}0,
 \qquad[0,T]\subset I.
\end{equation} Together with \eqref{eq:logarithmic-cutoff-convergence} and \eqref{eq:density-cutoff-L1-explicit}, this proves \eqref{eq:background-cutoff-convergence} and completes the proof of \cref{prop:background-stability}.
\end{proof}

\par\endgroup

\begingroup

\section{Finite-particle well-posedness and noncollision}\label{app:particle-wp}

This appendix constructs the singular particle law used in \cref{lem:approximation}. The localization and noncollision arguments are standard (cf.\ \cite[\S4]{RSdiffusive} for singular Riesz-type particle systems and \cite[Lemma~15]{FournierJourdain2015} for the cluster-dispersion argument in an attractive planar system). We include the details for the interactions and initial configurations considered here. For strong solutions of infinite systems with singular interactions, see also \cite{OsadaTanemura2020}.

We first construct a maximal strong solution by truncating the force, then exclude collapse of a cluster by descending induction on its cardinality. Here, the quadratic dispersion has constant Hessian, so no term involving $\Delta\g$ occurs, while the repulsive internal force has a nonnegative contribution.

\begin{prop}[Global strong well-posedness and noncollision]\label{prop:particle-wellposedness}
Let $N\ge2$, $\d\ge2$, $0\le\s<\d/2$, and $\beta>0$. Suppose $V\in C_{\mathrm{loc}}^2(\R^\d)$ and $\norm{\nabla^2V}_{L^\infty}<\infty$. If the initial configuration $X_N(0)$ is almost surely pairwise distinct and independent of the Brownian motions $(B_i)_{i\in[N]}$, then \eqref{eq:sde} admits a pathwise unique global strong solution. Moreover, for every $T<\infty$,
\begin{equation}\label{eq:particle-uniform-separation}
 \Pb\left(\min_{0\le t\le T}\min_{1\le i<j\le N}|x_i(t)-x_j(t)|>0\right)=1.
\end{equation}
\end{prop}

We first take the initial configuration $X_N(0)$ to be deterministic and pairwise distinct. Set
\begin{equation}\label{eq:particle-collision-set}
 F:=-\nabla\g,\qquad
 \mathfrak D_N:=\{X_N\in(\R^\d)^N:x_i=x_j\text{ for some }i\ne j\text{ in }[N]\}.
\end{equation}
The particle generator on the collision complement $(\R^\d)^N\setminus\mathfrak D_N$ is
\begin{equation}\label{eq:particle-generator-noncollision}
 \mathcal L_N:=\frac1\beta\sum_{i=1}^N\Delta_{x_i}
 +\sum_{i=1}^N\left(\frac1N\sum_{j\ne i}F(x_i-x_j)-\nabla V(x_i)\right)\cdot\nabla_{x_i}.
\end{equation}

\begin{lemma}[Maximal strong solution]\label{lem:particle-maximal}
There are a stopping time $\tau_\star\in(0,\infty]$ and a pathwise unique strong solution of \eqref{eq:sde} on $[0,\tau_\star)$ such that, on $\{\tau_\star<\infty\}$,
\begin{equation}\label{eq:particle-collision-alternative}
 \liminf_{t\uparrow\tau_\star}\min_{1\le i<j\le N}|x_i(t)-x_j(t)|=0.
\end{equation}
\end{lemma}
\begin{proof}
Choose smooth odd, globally Lipschitz vector fields $F_\varepsilon$ on $\R^\d$ satisfying $F_\varepsilon=F$ on $\{|x|\ge\varepsilon\}$. These force truncations are used only in this construction; they need not be the heat truncations from \cref{eq:heat-cutoff}. Since $\nabla V$ is globally Lipschitz, the regularized system has a pathwise unique global strong solution (see, e.g., \cite[\S5.2]{Oksendal}). Set $\varepsilon_n:=2^{-n}$ and let $X_N^{(n)}$ be the corresponding solution, with the same initial configuration and Brownian motions for every $n$. For all sufficiently large $n$, the initial pair distances exceed $\varepsilon_n$. For $m\ge n$, set
\begin{equation}\label{eq:particle-coarse-entrance}
 \tau_{m,n}:=\inf\{t\ge0:\min_{i<j}|x_i^{(m)}(t)-x_j^{(m)}(t)|\le\varepsilon_n\},
 \qquad \tau_n:=\tau_{n,n}.
\end{equation}
For $m>n$, both processes solve the same locally Lipschitz equation on $(\R^\d)^N\setminus\mathfrak D_N$ up to $\tau_n\wedge\tau_{m,n}$. Stopped pathwise uniqueness and continuity therefore give
\begin{equation}\label{eq:particle-pasting-consistency}
 \tau_{m,n}=\tau_n,
 \qquad
 X_N^{(m)}=X_N^{(n)}\ \text{on }[0,\tau_n],
 \qquad
 \tau_m\ge\tau_n.
\end{equation}
By \eqref{eq:particle-pasting-consistency}, the process defined by $X_N(t):=X_N^{(n)}(t)$ for $t<\tau_n$ is independent of the choice of $n$ and solves \eqref{eq:sde} on $[0,\tau_\star)$, where $\tau_\star:=\lim_n\tau_n$. The drift is locally Lipschitz on $(\R^\d)^N\setminus\mathfrak D_N$, so the same localization proves pathwise uniqueness there. If $\tau_\star<\infty$, then at times $\tau_n\uparrow\tau_\star$ some pair has distance $\varepsilon_n$; after passing to a subsequence, one fixed pair does so. This proves \eqref{eq:particle-collision-alternative}.
\end{proof}

For $I\subset[N]$ with cardinality $k\ge2$, define
\begin{equation}\label{eq:cluster-dispersion}
 \overline x_I:=\frac1k\sum_{i\in I}x_i,
 \qquad
 Q_I(X_N):=\sum_{i\in I}|x_i-\overline x_I|^2
 =\frac1k\sum_{\substack{i,j\in I\\i<j}}|x_i-x_j|^2.
\end{equation}
For $j\notin I$, one has
\begin{equation}\label{eq:cluster-enlargement}
 Q_{I\cup\{j\}}=Q_I+\frac{k}{k+1}|x_j-\overline x_I|^2.
\end{equation}

\begin{lemma}[Cluster dispersion equation]\label{lem:cluster-ito}
Before $\tau_\star$,
\begin{equation}\label{eq:cluster-ito}
 \dd Q_I=\left(\frac{2\d(k-1)}\beta+\mathcal A_I+\mathcal R_I\right)\dd t+\dd\mathcal M_I,
 \qquad
 \dd\langle\mathcal M_I\rangle_t=\frac8\beta Q_I\dd t,
\end{equation}
where
\begin{equation}\label{eq:cluster-internal-term}
 \mathcal A_I=-\frac2N\sum_{\substack{i,j\in I\\i<j}}
 (x_i-x_j)\cdot\nabla\g(x_i-x_j)\ge0.
\end{equation}
For every $\delta>0$, there are $q_{I,\delta}>0$ and $C_{I,\delta}<\infty$ such that
\begin{equation}\label{eq:cluster-remainder-bound}
 \mathcal R_I\ge-C_{I,\delta}Q_I
\end{equation}
whenever $Q_I<q_{I,\delta}$ and
\begin{equation}\label{eq:larger-cluster-separated}
 Q_{I\cup\{j\}}\ge\delta\qquad\text{for every }j\notin I.
\end{equation}
For the full cluster, \eqref{eq:larger-cluster-separated} is vacuous and the constants do not depend on $\delta$.
\end{lemma}
\begin{proof}
The elementary identities
\begin{equation}\label{eq:cluster-derivatives}
 \nabla_{x_i}Q_I=2(x_i-\overline x_I),
 \qquad
 \sum_{i\in I}\Delta_{x_i}Q_I=2\d(k-1),
 \qquad
 \sum_{i\in I}|\nabla_{x_i}Q_I|^2=4Q_I
\end{equation}
give the diffusion drift in \eqref{eq:cluster-ito} by It\^o's formula. The martingale is
\begin{equation}\label{eq:cluster-martingale}
 \mathcal M_I(t):=2\sqrt{\frac2\beta}
 \sum_{i\in I}\int_0^t(x_i(r)-\overline x_I(r))\cdot\dd B_i(r),
 \qquad t<\tau_\star.
\end{equation}
Independence of the Brownian motions and \eqref{eq:cluster-derivatives} give its quadratic variation. Pairing the two orientations of each internal interaction gives \eqref{eq:cluster-internal-term}. For $\s>0$,
\begin{equation}\label{eq:riesz-radial-sign}
 z\cdot\nabla\g(z)=-|z|^{-\s},
\end{equation}
whereas for $\s=0$ the same quantity equals $-1$. Thus, $\mathcal A_I\ge0$.

For an exterior index $j\notin I$, subtract the common value $F(\overline x_I-x_j)$ and use $\sum_{i\in I}(x_i-\overline x_I)=0$. Its contribution is
\begin{equation}\label{eq:cluster-exterior-contribution}
 \frac2N\sum_{i\in I}(x_i-\overline x_I)\cdot
 \bigl(F(x_i-x_j)-F(\overline x_I-x_j)\bigr).
\end{equation}
If \eqref{eq:larger-cluster-separated} holds and $Q_I\le\delta/16$, then \eqref{eq:cluster-enlargement} gives
$|\overline x_I-x_j|\ge\sqrt{\delta-Q_I}$. Thus, the entire segment used to compare the two force values stays away from the origin:
\begin{equation}\label{eq:cluster-segment-separation}
 |\overline x_I-x_j+\theta(x_i-\overline x_I)|
 \ge\sqrt{\delta-Q_I}-\sqrt{Q_I}\ge\frac{\sqrt\delta}{2},
 \qquad i\in I,\quad 0\le\theta\le1.
\end{equation}
Consequently,
\begin{equation}\label{eq:cluster-exterior-bound}
 \left|\frac2N\sum_{i\in I}(x_i-\overline x_I)\cdot
 \bigl(F(x_i-x_j)-F(\overline x_I-x_j)\bigr)\right|
 \le\frac2N\sup_{|z|\ge\sqrt\delta/2}|\nabla F(z)|\,Q_I.
\end{equation} Subtracting $\nabla V(\overline x_I)$ in the confinement contribution similarly gives
\begin{equation}\label{eq:cluster-confinement-contribution}
 -2\sum_{i\in I}(x_i-\overline x_I)\cdot
 \bigl(\nabla V(x_i)-\nabla V(\overline x_I)\bigr)
 \ge-2\norm{\nabla^2V}_{L^\infty}Q_I.
\end{equation}
Summing the exterior and confinement contributions proves \eqref{eq:cluster-remainder-bound}.
\end{proof}

\begin{lemma}[Noncollapse of a cluster]\label{lem:cluster-barrier}
Let $I$ have cardinality $k\ge2$ and suppose $\d(k-1)\ge2$. Assume that every one-point strict enlargement $J\supsetneq I$ satisfies
\begin{equation}\label{eq:larger-cluster-induction-hypothesis}
 \Pb\left(\inf_{0\le t<T\wedge\tau_\star}Q_J(X_N(t))=0\right)=0
 \qquad\text{for every }T<\infty.
\end{equation}
Then the same assertion holds for $I$.
\end{lemma}
\begin{proof}
Fix $\delta>0$ and stop at
\begin{equation}\label{eq:cluster-exterior-stopping-time}
 \eta_{I,\delta}:=\inf\left\{t<\tau_\star:
 \min_{j\notin I}Q_{I\cup\{j\}}(X_N(t))\le\delta\right\},
\end{equation}
with $\eta_{I,\delta}=\infty$ for the full cluster. Choose $0<q_0<\min\{q_{I,\delta},1\}$ as in \cref{lem:cluster-ito}. Let $\Phi\in C_{\mathrm{loc}}^2((0,\infty))$ be nonnegative and nonincreasing, equal to $-\log q$ on $(0,q_0/4]$, and constant on $[q_0/2,\infty)$. On $\{Q_I\le q_0/4\}$, It\^o's formula and \cref{lem:cluster-ito} give
\begin{equation}\label{eq:cluster-log-generator}
 \mathcal L_N\Phi(Q_I)=\frac{4-2\d(k-1)}{\beta Q_I}-\frac{\mathcal A_I+\mathcal R_I}{Q_I}
 \le\frac{4-2\d(k-1)}{\beta Q_I}+C_{I,\delta}\le C_{I,\delta}.
\end{equation}
The coefficient of $Q_I^{-1}$ is nonpositive also when $\d(k-1)=2$. In the transition region $q_0/4\le Q_I\le q_0/2$, the inequality $\Phi'\le0$ and \eqref{eq:cluster-remainder-bound} give
$\Phi'\mathcal R_I\le C_{I,\delta}Q_I|\Phi'|$, while $\Phi'\mathcal A_I\le0$ and $Q_I|\Phi''|$ is bounded. Hence \eqref{eq:cluster-log-generator}, with a larger constant, holds before $\eta_{I,\delta}$.

For $0<a<\min\{Q_I(X_N(0)),q_0/4,e^{-1}\}$, set
\begin{equation}\label{eq:cluster-hitting-time}
 \sigma_{I,a}:=\inf\{t<\tau_\star:Q_I(X_N(t))\le a\}.
\end{equation}
Apply It\^o's formula up to
$T\wedge\eta_{I,\delta}\wedge\sigma_{I,a}\wedge\tau_n$.
The martingale in this formula is square integrable, since
\begin{equation}\label{eq:cluster-stopped-martingale-bound}
 \frac8\beta\int_0^{T\wedge\eta_{I,\delta}\wedge\sigma_{I,a}\wedge\tau_n}
 Q_I(X_N(r))\bigl|\Phi'(Q_I(X_N(r)))\bigr|^2\dd r
 \le\frac{8T}\beta\sup_{q\ge a}q|\Phi'(q)|^2<\infty.
\end{equation}
Taking expectations and using \eqref{eq:cluster-log-generator} gives
\begin{equation}\label{eq:cluster-barrier-expectation}
 \E\!\left[\Phi\!\left(Q_I\bigl(X_N(T\wedge\eta_{I,\delta}\wedge\sigma_{I,a}\wedge\tau_n)\bigr)\right)\right]
 \le \Phi(Q_I(X_N(0)))+C_{I,\delta}T.
\end{equation}
On $\{\sigma_{I,a}\le T\wedge\eta_{I,\delta}\wedge\tau_n\}$, continuity gives
$Q_I(X_N(\sigma_{I,a}))=a$. Since $\Phi\ge0$, \eqref{eq:cluster-barrier-expectation} implies
\begin{equation}\label{eq:cluster-stopped-hitting-bound}
 \Pb(\sigma_{I,a}\le T\wedge\eta_{I,\delta}\wedge\tau_n)
 \le\frac{\Phi(Q_I(X_N(0)))+C_{I,\delta}T}{|\log a|}.
\end{equation}
The events on the left increase to
$\{\sigma_{I,a}\le T\wedge\eta_{I,\delta},\ \sigma_{I,a}<\tau_\star\}$.
Therefore,
\begin{equation}\label{eq:cluster-hitting-probability}
 \Pb(\sigma_{I,a}\le T\wedge\eta_{I,\delta},\ \sigma_{I,a}<\tau_\star)
 \le\frac{\Phi(Q_I(X_N(0)))+C_{I,\delta}T}{|\log a|}.
\end{equation}
Letting $a\downarrow0$ excludes approach of $Q_I$ to zero on
$[0,T\wedge\tau_\star)$ unless $\eta_{I,\delta}<T\wedge\tau_\star$.
By the induction hypothesis,
\begin{equation}\label{eq:cluster-exterior-stopping-vanishes}
 \Pb(\eta_{I,\delta}<T\wedge\tau_\star)
 \le\sum_{j\notin I}
 \Pb\!\left(\inf_{0\le t<T\wedge\tau_\star}
 Q_{I\cup\{j\}}(X_N(t))\le\delta\right)
 \xrightarrow[\delta\downarrow0]{}0.
\end{equation}
Letting $\delta\downarrow0$ completes the proof.
\end{proof}

\begin{proof}[Proof of \cref{prop:particle-wellposedness}]
For the full cluster, \cref{lem:cluster-barrier} applies without an induction hypothesis. Descend in the cardinality of $I$. Since $\d\ge2$, every cluster with $|I|\ge2$ satisfies $\d(|I|-1)\ge2$, and therefore
\begin{equation}\label{eq:no-cluster-collapse}
 \Pb\left(\inf_{0\le t<T\wedge\tau_\star}Q_I(X_N(t))=0\right)=0
 \qquad\text{for every }T<\infty.
\end{equation}
If $\tau_\star$ were finite, \eqref{eq:particle-collision-alternative} and the finiteness of the set of pairs would contradict \eqref{eq:no-cluster-collapse} for a two-particle cluster. Hence $\tau_\star=\infty$ almost surely.

Since $\tau_n\uparrow\infty$ almost surely, for every finite $T$ and almost every sample there is an $n$ such that $T<\tau_n$. By the construction in \cref{lem:particle-maximal}, $X_N(t)=X_N^{(n)}(t)$ for all $t\in[0,T]$. The regularized solution is global, so $X_N$ cannot explode in space. Continuity and the absence of collisions then give \eqref{eq:particle-uniform-separation}. Local pathwise uniqueness on $(\R^\d)^N\setminus\mathfrak D_N$, as in \cref{lem:particle-maximal}, gives global pathwise uniqueness.

For a random pairwise-distinct initial configuration $X_N(0)$ independent of the Brownian motions $(B_i)_{i\in[N]}$, apply the deterministic construction conditionally on $X_N(0)$. Equivalently, use the measurable solution maps of the regularized equations and their pathwise stabilization in \cref{lem:particle-maximal}. This gives the asserted global strong solution and completes the proof.
\end{proof}

\par\endgroup

\begingroup
\section{Proof of the iid counterexample}\label{app:iid-counterexample}
We prove \cref{prop:iid-counterexample} in five steps. After constructing the stationary measure and its confinement, we establish entropy and fixed-particle density bounds for the singular dynamics. These bounds justify an adjoint fluctuation formula from independent initial data. We then estimate the amplification of a particular test and compare the resulting Gaussian variance with its product-reference variance. The final subsection deduces the chi-square divergence and the marginal entropy bound.

\subsection{The stationary measure and confinement}\label{iid:subsec:construction}
Throughout this appendix, $\d=4$ and $0\le\s<1$, and $\g$ is the whole-space kernel \eqref{eq:kernel}.
Let $\gamma_a^{(m)}$ denote the centered Gaussian density on $\R^m$ with covariance $a^2\Id$, and write $x=(r,z)\in\R\times\R^3$. For $0<\delta<1/2$, set
\begin{equation}\label{iid:eq:parameters}
 \eta_\delta:=e^{-\delta^{-2}},\qquad \beta_\delta^{-1}:=\delta^6,\qquad T_\delta:=\sqrt\delta,
\end{equation}
\begin{equation}\label{iid:eq:mu}
 \mu_\delta(r,z):=\frac{1-\eta_\delta}{2}\sum_{\sigma=\pm1}
 \gamma_\delta^{(1)}(r-\sigma)\gamma_1^{(3)}(z)
 +\eta_\delta\gamma_2^{(4)}(r,z).
\end{equation}
The first two components of $\mu_\delta$ concentrate near the parallel hyperplanes $r=\pm1$. The third component, although negligible as $\delta\downarrow0$, gives a Gaussian lower bound at every fixed $\delta$ and determines the confinement at infinity. As elsewhere, we use the same symbol for a density and its probability measure. Define
\begin{equation}\label{iid:eq:V}
 h_\delta:=\g*\mu_\delta,\qquad
 V_\delta:=-\beta_\delta^{-1}\log\mu_\delta-h_\delta.
\end{equation}
The definition of $V_\delta$ makes $\mu_\delta$ stationary; an additive constant in $V_\delta$ has no effect on the dynamics. We first obtain the entropy and fluctuation estimates at fixed $\delta$, then let $\delta\downarrow0$ to select the counterexample. Constants in this latter limit may depend on $\s$ and the fixed test cutoffs, but not on $\delta$. Once $\delta$ has been chosen, it remains fixed as $N\to\infty$.

\begin{lemma}[Stationary measure and confinement]\label{iid:lem:stationary-measure}
The density in \eqref{iid:eq:mu} is smooth, positive, and even in $r$. For every fixed $\delta$, it has a Gaussian upper bound, a bounded logarithmic Hessian, a one-body Poincar\'e inequality, and finite logarithmic or Riesz energy. Moreover,
\begin{equation}\label{iid:eq:score-hessian}
 \norm{\nabla^2\log\mu_\delta}_{\infty}\le C\delta^{-4},
\end{equation}
and
\begin{equation}\label{iid:eq:score-first}
 \sup_{\substack{|r|\le R\\z\in\R^3}}|\partial_r\log\mu_\delta(r,z)|
 \le C_R\delta^{-2},\qquad R>0.
\end{equation}
For $j=1,2$,
\begin{equation}\label{iid:eq:kernel-row}
 \sup_{0<\delta<1/2}\sup_x\int |\nabla^j\g(x-y)|\,\mu_\delta(\dd y)<\infty.
\end{equation}
The confinement in \eqref{iid:eq:V} satisfies \eqref{iid:eq:force-class}, and $\mu_\delta$ is its thermal equilibrium measure at inverse temperature $\beta_\delta$. The stationary trajectory satisfies the intrinsic one-body and static hypotheses of \cref{thm:finite}.
\end{lemma}
\begin{proof}
We write
\begin{equation}\label{iid:eq:profile-ratios}
 \mu_\delta=\eta_\delta\gamma_2^{(4)}(1+R_++R_-),\qquad
 R_\pm(r,z):=\frac{8(1-\eta_\delta)}{\eta_\delta\delta}
 \exp\left(-\frac{(r\mp1)^2}{2\delta^2}+\frac{r^2}{8}
 -\frac{3|z|^2}{8}\right).
\end{equation}
For fixed $\delta$, each $R_\pm$ and all its derivatives decay rapidly at infinity. Consequently, $\log(1+R_++R_-)$ and all its derivatives have the same property. In particular, the logarithmic derivatives of order at least two are bounded.

There are constants $0<c_\delta\le C_\delta<\infty$ such that $c_\delta\gamma_2^{(4)}\le\mu_\delta\le C_\delta\gamma_2^{(4)}$. This gives the Gaussian upper bound and, by comparison with the Gaussian Poincar\'e inequality,
\begin{equation}\label{iid:eq:poincare}
 \Var_{\mu_\delta}(f)\le 4(C_\delta/c_\delta)
 \int|\nabla f|^2\,\dd\mu_\delta.
\end{equation}
The bounded density and Gaussian upper bound imply every static integrability condition used below, including square integrability of the centered potential and finiteness of the logarithmic energy.

\parheading{Logarithmic derivative bounds.}
For the estimate uniform in $\delta$, write $w_0,w_+,w_-$ for the three mixture summands divided by $\mu_\delta$, and $\ell_0,\ell_+,\ell_-$ for their logarithms before division. Direct differentiation gives
\begin{equation}\label{iid:eq:mixture-hessian}
 \nabla^2\log\mu_\delta
 =\sum_iw_i\nabla^2\ell_i
 +\frac12\sum_{i,j}w_iw_j
 (\nabla\ell_i-\nabla\ell_j)\otimes
 (\nabla\ell_i-\nabla\ell_j).
\end{equation}
The Gaussian Hessians have size at most $C\delta^{-2}$. The difference of the two narrow-component scores is the constant vector $(2\delta^{-2},0)$. For a narrow and the broad component, the score difference is $\nabla\log R_\pm$.

Completing the square yields $R_\pm=e^{a_\delta-Q_\pm}$, where $0<a_\delta\le C\delta^{-2}$ and $Q_\pm$ is a nonnegative quadratic form centered at a fixed point of bounded size, with largest quadratic coefficient at most $C\delta^{-2}$. Therefore,
\begin{equation}\label{iid:eq:score-ratio-bound}
 |\nabla\log R_\pm|^2\le C\delta^{-2}Q_\pm,
 \qquad
 w_0w_\pm\le\min\{1,R_\pm\}.
\end{equation}
The elementary inequality
\begin{equation}\label{iid:eq:quadratic-mixture-bound}
 \sup_{u\ge0}u\min\{1,e^{a-u}\}\le a+1\qquad(a\ge0)
\end{equation}
proves \eqref{iid:eq:score-hessian} from \eqref{iid:eq:mixture-hessian}. The first-coordinate score is a convex combination of $-(r\mp1)/\delta^2$ and $-r/4$, which proves \eqref{iid:eq:score-first}.

\parheading{Uniform convolution bounds.}
The three transverse variables provide the integrability needed uniformly as the first-coordinate variance vanishes. Indeed, $|\nabla^j\g(x)|\le C_\s|x|^{-\s-j}$ for $j=1,2$, and $\s+2<3$. Thus, for $j=1,2$ and $\sigma\in\{1,2\}$,
\begin{equation}\label{iid:eq:transverse-integrability}
 \sup_{b,z}\int_{\R^3}(b^2+|z-y|^2)^{-(\s+j)/2}
 \gamma_\sigma^{(3)}(y)\,\dd y
 \le \norm{\gamma_\sigma^{(3)}}_\infty
 \int_{|y|<1}|y|^{-\s-j}\,\dd y+1<\infty.
\end{equation}
Integrating the remaining one-dimensional probability density proves \eqref{iid:eq:kernel-row}.
\parheading{Confinement and stationarity.}

For fixed $\delta$, every derivative of $\mu_\delta$ is a Schwartz function. The force $\nabla\g$ is locally integrable and bounded outside every neighborhood of the origin, with limit zero at infinity. Differentiating its convolution by putting the additional derivatives on $\mu_\delta$ shows that $\nabla h_\delta$ is smooth and every derivative is bounded and tends to zero at infinity. By \eqref{iid:eq:profile-ratios},
\begin{equation}\label{iid:eq:force-explicit}
 \nabla V_\delta(x)=\frac{\beta_\delta^{-1}}{4}x
 -\beta_\delta^{-1}\nabla\log(1+R_+(x)+R_-(x))-\nabla h_\delta(x).
\end{equation}
This proves \eqref{iid:eq:force-class}; in particular, the second and third derivatives of $V_\delta$ are bounded.

The defining identity for the confinement gives
\begin{equation}\label{iid:eq:stationary-drift}
 B_t=-\nabla V_\delta-\nabla h_\delta
 =\beta_\delta^{-1}\nabla\log\mu_\delta,
 \qquad
 u_t=B_t-\beta_\delta^{-1}\nabla\log\mu_\delta=0.
\end{equation}
Hence $\beta_\delta^{-1}\Delta\mu_\delta-\diver(B_t\mu_\delta)=0$. The density comparison above, \eqref{iid:eq:score-hessian}, and \eqref{iid:eq:poincare} supply the one-body bounds. The first two derivatives of $B_t$ are bounded for fixed $\delta$ by \eqref{iid:eq:profile-ratios}. The partition-function bounds and exponential transport-commutator estimates follow from \cite[Theorems~2.2 and~2.8]{DGRpartition}, respectively.

The identity \eqref{iid:eq:V} is precisely \eqref{eq:thermal}. Since \eqref{iid:eq:force-explicit} verifies \eqref{eq:confinement-class}--\eqref{eq:confinement-decay}, \cref{lem:thermal} identifies $\mu_\delta$ as the unique thermal equilibrium measure. Its logarithmic perturbation relative to that equilibrium measure is zero.
\end{proof}

\subsection{Particle laws and entropy}\label{iid:subsec:particle-laws}
Fix $\delta\in(0,1/2)$ throughout this subsection and suppress its subscript on $\beta,V,\mu$, and $h$. Set
\begin{equation}\label{iid:eq:Gibbs}
 \Gamma_N:=\Q_{N,\beta}(\mu).
\end{equation}
By \eqref{iid:eq:V} and \eqref{eq:modulated}, $\Gamma_N$ is the invariant canonical Gibbs measure of \eqref{eq:sde} with confinement $V$. The particle measure started from $\mu^{\otimes N}$ is generally nonstationary. Since $\nabla^2V$ is bounded and the initial configuration is almost surely pairwise distinct, \cref{prop:particle-wellposedness} gives the global noncolliding strong solution and the compact-time separation \eqref{eq:particle-uniform-separation}.

For the entropy and density estimates, use convolution with the heat kernel $p_\varepsilon$ of covariance $2\varepsilon\Id$, and set
\begin{equation}\label{iid:eq:regularization}
 \g_{(\varepsilon)}:=\g*p_\varepsilon,\qquad
 h_{(\varepsilon)}:=\g_{(\varepsilon)}*\mu,\qquad
 V_{(\varepsilon)}:=-\beta^{-1}\log\mu-h_{(\varepsilon)}.
\end{equation}
This convolution mollification is distinct from the heat-parameter truncation $\g_\varepsilon$ in \eqref{eq:heat-cutoff}. The symbols $\Fr_{N,(\varepsilon)}$, $\K_{N,b,(\varepsilon)}$, and $\Q_{N,b,(\varepsilon)}$ denote \eqref{eq:modulated}--\eqref{eq:gibbs} with $\g$ replaced by $\g_{(\varepsilon)}$. Write $\Gamma_{N,(\varepsilon)}:=\Q_{N,\beta,(\varepsilon)}(\mu)$ and $P_{N,t}^{(\varepsilon)}$ for the law of the smooth particle system with $V_{(\varepsilon)}$, $\g_{(\varepsilon)}$, and the same iid initial law. The target confinement remains the fixed $V$; $V_{(\varepsilon)}$ is used only in this auxiliary construction. At every cutoff,
\begin{equation}\label{iid:eq:compensated-drift}
 -\nabla V_{(\varepsilon)}-\nabla h_{(\varepsilon)}
 =\beta^{-1}\nabla\log\mu.
\end{equation}
The mollified particle drift is globally Lipschitz at fixed $N,\varepsilon$; see, e.g., \cite[Section 5.2]{Oksendal}. Since $\nabla\g$ is locally integrable and the derivatives of $\mu$ decay at Gaussian rates, $\nabla h_{(\varepsilon)}\to\nabla h$ locally with its first derivatives. The pair force and its first derivatives converge uniformly on compact sets away from the diagonal. Coupling the initial configurations identically and using the same Brownian motions, the localization argument \eqref{eq:approximation-path-Gronwall}--\eqref{eq:approximation-path-probability} therefore applies. At each fixed $N,T$, it gives
\begin{equation}\label{iid:eq:path-convergence}
 P_{N,\cdot}^{(\varepsilon)}\ \longrightarrow\ P_{N,\cdot}
 \quad\text{as path laws on }C([0,T];(\R^4)^N),\qquad\varepsilon\downarrow0.
\end{equation}
Here, the positive minimum pair separation is supplied by \cref{prop:particle-wellposedness}. This argument uses neither a Gibbs initial law nor a uniform-in-$N$ path-convergence estimate.


\begin{lemma}[Uniform relative entropy]\label{iid:lem:entropy}
For the initial law \eqref{iid:eq:initial-law},
\begin{equation}\label{iid:eq:entropy-both}
 \sup_{N\ge2}\sup_{t\ge0}
 \left\{\Hr(P_{N,t}\mid\Gamma_N)+\Hr(P_{N,t}\mid\mu^{\otimes N})\right\}<\infty.
\end{equation}
\end{lemma}
\begin{proof}
We first prove estimates uniform in the mollification parameter $\varepsilon$, then pass to the singular dynamics. The auxiliary diffusion has invariant measure $\Gamma_{N,(\varepsilon)}$, and its initial entropy is
\begin{equation}\label{iid:eq:initial-entropy}
 \Hr(\mu^{\otimes N}\mid\Gamma_{N,(\varepsilon)})
 =\log\K_{N,\beta,(\varepsilon)}(\mu)-\beta I_{(\varepsilon)},
 \qquad I_{(\varepsilon)}:=\frac12\iint\g_{(\varepsilon)}\,\dd\mu^{\otimes2}.
\end{equation}

\parheading{Partition bounds uniform in the mollification parameter.}
Set $\mu_{(\varepsilon)}:=\mu*p_{\varepsilon/2}$. Let $y_1,\ldots,y_N$ have density $p_{\varepsilon/2}$ independently of one another and of $X_N\sim\mu^{\otimes N}$, and set $Y_N:=(y_1,\ldots,y_N)$. Conditional expectation gives
\begin{equation}\label{iid:eq:energy-jensen}
 N\Fr_{N,(\varepsilon)}(X_N,\mu)
 =\E_Y\bigl[N\Fr_N(X_N+Y_N,\mu_{(\varepsilon)})\mid X_N\bigr].
\end{equation}
The pair term uses the difference of two independent heat increments; the one-body term uses one increment and the convolution defining $\mu_{(\varepsilon)}$. Their variances add to $2\varepsilon$ in both cases. Thus
\begin{equation}\label{iid:eq:partition-jensen}
 \K_{N,b,(\varepsilon)}(\mu)\le\K_{N,b}(\mu_{(\varepsilon)}),\qquad b>0.
\end{equation}
For $0\le\varepsilon\le1$, the mollified background is the Gaussian mixture
\begin{equation}\label{iid:eq:mollified-mixture}
 \mu_{(\varepsilon)}(r,z)=\frac{1-\eta_\delta}{2}\sum_{\sigma=\pm1}
 \gamma_{\sqrt{\delta^2+\varepsilon}}^{(1)}(r-\sigma)
 \gamma_{\sqrt{1+\varepsilon}}^{(3)}(z)
 +\eta_\delta\gamma_{\sqrt{4+\varepsilon}}^{(4)}(r,z).
\end{equation}
For fixed $\delta$, the means are fixed and the covariance matrices in \eqref{iid:eq:mollified-mixture} range over a compact subset of the positive-definite matrices. Consequently, the family has common $L^\infty$ and Gaussian-moment bounds. Its Riesz-potential bounds and, when $\s=0$, its absolute logarithmic-energy bounds are uniform as well: each relevant integral is a finite continuous function of the covariance parameters on this compact set. For $0<\s<1$, nonnegativity of $\g$ and the uniform potential bound also control the exponential negative-potential term in the partition estimate. For $\s=0$,
\begin{equation}\label{iid:eq:mollified-log-potential}
 h_{\mu_{(\varepsilon)}}(x)\ge-\log(1+|x|)-C_\delta,
 \qquad 0\le\varepsilon\le1,
\end{equation}
so the uniform Gaussian moments control that exponential term as well. In particular, for each fixed $b>0$,
\begin{equation}\label{iid:eq:mollified-static-controls}
 \sup_{0\le\varepsilon\le1}\left\{
 \norm{\widetilde\g_{\mu_{(\varepsilon)}}}_{L^2((\mu_{(\varepsilon)})^{\otimes2})}
 +\norm{h_{\mu_{(\varepsilon)}}-2I_{\mu_{(\varepsilon)}}}_{L^2(\mu_{(\varepsilon)})}
 +\int e^{-b(h_{\mu_{(\varepsilon)}}-2I_{\mu_{(\varepsilon)}})}\dd\mu_{(\varepsilon)}
 \right\}<\infty.
\end{equation}
Here, $\widetilde\g_{\mu_{(\varepsilon)}}$ is the kernel centered with respect to $\mu_{(\varepsilon)}$ as in \eqref{eq:static-centered-kernel}. Tracking the constants in \cite[Appendix~A]{DGRpartition}, the large-$N$ threshold and the partition bound in \cite[Theorem~2.2(i)]{DGRpartition} may therefore be chosen uniformly for $0\le\varepsilon\le1$, both for $b=\beta$ and $b=2\beta$. For the finitely many particle numbers below that threshold, continuity of $\varepsilon\mapsto\K_{N,b}(\mu_{(\varepsilon)})$ on $[0,1]$, obtained by dominated convergence, gives the same uniform conclusion. Also, $I_{(\varepsilon)}\to I_\mu$, and Jensen gives $\K_{N,b,(\varepsilon)}\ge e^{bI_{(\varepsilon)}}$. Combining these facts with \eqref{iid:eq:partition-jensen}, we obtain, for every fixed $b>0$,
\begin{equation}\label{iid:eq:uniform-mollified-partition}
 1\le e^{-bI_{(\varepsilon)}}\K_{N,b,(\varepsilon)}(\mu)\le C_{\delta,b},
 \qquad N\ge2,\quad 0<\varepsilon\le1.
\end{equation}

\parheading{Entropy contraction and change of reference.}
Entropy contraction under the smooth Markov semigroup, as in \cite[Section~5.2]{BGL}, gives
\begin{equation}\label{iid:eq:entropy-contraction}
 \Hr(P_{N,t}^{(\varepsilon)}\mid\Gamma_{N,(\varepsilon)})
 \le\log\K_{N,\beta,(\varepsilon)}-\beta I_{(\varepsilon)}.
\end{equation}
For $G_{N,(\varepsilon)}:=\dd\Gamma_{N,(\varepsilon)}/\dd\mu^{\otimes N}$, the partition-function identity is
\begin{equation}\label{iid:eq:Gibbs-second-moment}
 \int G_{N,(\varepsilon)}^2\dd\mu^{\otimes N}
 =\frac{\K_{N,2\beta,(\varepsilon)}(\mu)}{\K_{N,\beta,(\varepsilon)}(\mu)^2}.
\end{equation}
Applying the Donsker--Varadhan lemma to $\log G_{N,(\varepsilon)}$ with reference $\Gamma_{N,(\varepsilon)}$, and using \eqref{iid:eq:entropy-contraction}--\eqref{iid:eq:Gibbs-second-moment}, gives
\begin{equation}\label{iid:eq:entropy-transfer}
 \Hr(P_{N,t}^{(\varepsilon)}\mid\mu^{\otimes N})
 \le\log\K_{N,2\beta,(\varepsilon)}(\mu)-2\beta I_{(\varepsilon)}.
\end{equation}
For fixed $\delta$, these bounds are uniform in $t\ge0$, $N\ge2$, and $0<\varepsilon\le1$.

\parheading{Passage to the singular dynamics.}
At fixed $N$, the normalized Gibbs densities converge in total variation. For $\s>0$, nonnegativity of the mollified pair potential and boundedness of the one-body convolution give a common multiple of $\mu^{\otimes N}$. For $\s=0$,
\begin{equation}\label{iid:eq:log-majorant}
 \g_{(\varepsilon)}(x)\ge-C-\log(1+|x|),\qquad 0<\varepsilon\le1,
\end{equation}
by Jensen applied to $\log|x+Y|$. This gives a fixed polynomial times a Gaussian as a majorant of the unnormalized density. Off-diagonal convergence, the nullity of the collision set under the product measure, and the positive lower normalizer bound prove the claim. Together with \eqref{iid:eq:path-convergence}, joint weak lower semicontinuity of relative entropy \cite[Theorem~19]{vanErvenHarremoes2014} gives, for every fixed $N$ and $t$,
\begin{equation}\label{iid:eq:singular-entropy-transfer}
 \Hr(P_{N,t}\mid\Gamma_N)+\Hr(P_{N,t}\mid\mu^{\otimes N})
 \le\liminf_{\varepsilon\downarrow0}
 \left(\Hr(P_{N,t}^{(\varepsilon)}\mid\Gamma_{N,(\varepsilon)})
 +\Hr(P_{N,t}^{(\varepsilon)}\mid\mu^{\otimes N})\right)\le C_\delta.
\end{equation}
The constant is independent of $N$ and $t$, so taking their suprema proves \eqref{iid:eq:entropy-both}.
\end{proof}

\begin{lemma}[Finiteness at fixed particle number]\label{iid:lem:fixed-N-finite}
For fixed $\delta,N$ and $T<\infty$, there is a finite constant $C_{\delta,N,T}$ such that
\begin{equation}\label{iid:eq:Gaussian-domination}
 P_{N,t}\le C_{\delta,N,T}\,(\gamma_2^{(4)})^{\otimes N},\qquad 0\le t\le T.
\end{equation}
Consequently, the fixed-$N$ chi-square divergences relative to $\mu_\delta^{\otimes N}$ are finite.
\end{lemma}
\begin{proof}
For $0\le \s<1$ and $\d=4$,
\begin{equation}\label{iid:eq:superharmonic}
 \Delta\g(x)=(\s-2)|x|^{-\s-2}\le0\quad(x\ne0).
\end{equation}
This identity also holds distributionally. Use the approximation \eqref{iid:eq:regularization} and set
\begin{equation}\label{iid:eq:mollified-confinement-remainder}
 \widetilde V_{(\varepsilon)}(x):=V_{(\varepsilon)}(x)-\frac{|x|^2}{8\beta}.
\end{equation} For fixed $\delta$, uniformly in $0<\varepsilon\le1$,
\begin{equation}\label{iid:eq:remainder-tail-bound}
 \norm{\Delta\widetilde V_{(\varepsilon)}}_\infty+
 \norm{x\cdot \nabla\widetilde V_{(\varepsilon)}(x)}_\infty\le C_\delta.
\end{equation}
For the logarithmic-density term this follows from the rapid decay in \eqref{iid:eq:profile-ratios}. For the convolution term, bounded derivatives follow from local integrability of $\nabla^2\g$ and the common Gaussian bounds on $\mu*p_\varepsilon$. For the factor of $x$, use
\begin{equation}\label{iid:eq:virial-convolution}
 x\cdot(\nabla\g*\rho)(x)
 =\bigl((y\cdot\nabla\g(y))*\rho\bigr)(x)
 +\int y\cdot\nabla\g(x-y)\rho(y)\,\dd y.
\end{equation}
The first convolution is $-|x|^{-\s}*\rho$, or $-1$ when $\s=0$. The second is bounded by splitting the force into its local $L^1$ part and bounded far part and using the uniform Gaussian bounds on $|y|\rho(y)$.

\parheading{A Gaussian supersolution.}
Set $q:=\gamma_2^{(4)}$ and compare the smooth particle density with $q^{\otimes N}$. Applying the forward generator to this reference density and dividing by $q^{\otimes N}$ gives the scalar source
\begin{align}\label{iid:eq:Gaussian-source}
 S_{N,(\varepsilon)}(X_N):={}&\sum_i\left(\Delta\widetilde V_{(\varepsilon)}(x_i)
 -\frac14x_i\cdot \nabla\widetilde V_{(\varepsilon)}(x_i)\right)\notag\\
 &+\frac1N\sum_{i<j}\left(2\Delta\g_{(\varepsilon)}(x_i-x_j)
 -\frac14(x_i-x_j)\cdot\nabla\g_{(\varepsilon)}(x_i-x_j)\right).
\end{align}
The Ornstein--Uhlenbeck part cancels exactly because $q$ has covariance $4\Id$ and the quadratic part of $V_{(\varepsilon)}$ is $|x|^2/(8\beta)$.

There is a finite constant $C_\s$ such that, off zero,
\begin{equation}\label{iid:eq:radial-source-bound}
 2\Delta\g(x)-\tfrac14x\cdot\nabla\g(x)
 =2(\s-2)|x|^{-\s-2}+\tfrac14|x|^{-\s}\le C_\s.
\end{equation}
When $\s=0$, the last factor $|x|^{-\s}$ equals one. Gaussian integration by parts gives
\begin{equation}\label{iid:eq:heat-virial}
 x\cdot\nabla\g_{(\varepsilon)}
 =(x\cdot\nabla\g)*p_\varepsilon-2\varepsilon\Delta\g_{(\varepsilon)}.
\end{equation}
Since $\Delta\g\le0$ distributionally, \eqref{iid:eq:heat-virial} gives
\begin{equation}\label{iid:eq:mollified-radial-source}
 2\Delta\g_{(\varepsilon)}-\tfrac14x\cdot\nabla\g_{(\varepsilon)}
 =\left(2\Delta\g-\tfrac14x\cdot\nabla\g\right)*p_\varepsilon
   +\tfrac\varepsilon2\Delta\g_{(\varepsilon)}\le C_\s.
\end{equation}
Thus, \eqref{iid:eq:remainder-tail-bound} implies $S_{N,(\varepsilon)}\le C_\delta N$, uniformly in the mollification parameter.

At a fixed cutoff the quotient equation has a globally Lipschitz auxiliary drift, at most linear growth, and a bounded scalar source. Its initial datum is at most $(\sup\mu/q)^N$. The Feynman--Kac formula (see, e.g., \cite[Section 8.2]{Oksendal}) gives
\begin{equation}\label{iid:eq:Gaussian-cutoff-domination}
 P_{N,t}^{(\varepsilon)}\le (\sup\mu/q)^N e^{C_\delta Nt}\,q^{\otimes N}.
\end{equation}
The Feynman--Kac solution is bounded with bounded first spatial derivative at each fixed cutoff: the initial quotient and the source have bounded first derivatives, and the auxiliary drift has bounded derivative. Multiplying it by $q^{\otimes N}$ and integrating the forward equation against expanding spatial cutoffs proves conservation of mass; the boundary terms vanish by Gaussian decay. Uniqueness for the smooth particle forward equation therefore identifies this solution with the actual quotient. The bound passes to the singular law by weak convergence and nonnegative continuous compactly supported tests. Finally $\mu\ge\eta_\delta q$, so the density ratio $P_{N,t}/\mu^{\otimes N}$ is essentially bounded at every fixed $N,T$. This proves the claim. No uniform bound in $N$ is asserted by this lemma.
\end{proof}

\subsection{Fluctuations from independent initialization}\label{iid:subsec:fluctuations}
Continue to fix $\delta$ and suppress it from the notation in this subsection. Set
\begin{equation}\label{iid:eq:linear-operators}
 Lf:=\beta^{-1}\bigl(\Delta f+\nabla\log\mu\cdot\nabla f\bigr),\qquad
 \mathcal Bf(x):=\int\nabla\g(x-y)\cdot\nabla f(y)\,\mu(\dd y),\qquad
 \mathcal A:=L+\mathcal B.
\end{equation}
The operator $\mathcal A$ is the adjoint of the linearized mean-field equation at $\mu$. We denote its evolution by $U_t$, so that $\partial_tU_t\phi=\mathcal A U_t\phi$ and $U_0\phi=\phi$. The next proposition separates the fluctuations of the independent initial configuration from those generated by the Brownian motions. The required adjoint regularity will be verified for the test used in the counterexample in \cref{iid:lem:adjoint-estimate}.

\begin{prop}[Fluctuations from iid initial data]\label{iid:prop:clt}
Suppose $f_r:=U_{T-r}\phi$, $0\le r\le T$, is bounded with bounded first and second spatial derivatives, jointly continuous in time and space, and solves the adjoint equation classically locally. At-most-linear growth of $\partial_rf_r$ suffices. Then the initial law \eqref{iid:eq:initial-law} gives
\begin{equation}\label{iid:eq:clt}
 \frac1{\sqrt N}\sum_{i=1}^N\bigl(\phi(x_i(T))-\mu(\phi)\bigr)
 \longrightarrow\mathcal N(0,W_T),
\end{equation}
where
\begin{equation}\label{iid:eq:covariance}
 W_T:=\Var_\mu(U_T\phi)+2\beta^{-1}\int_0^T\int|\nabla U_r\phi|^2\,\dd\mu\dd r.
\end{equation}
\end{prop}
\begin{proof}
For a bounded test $f$, set $S_N^t(f):=\sqrt N(\mu_N^t(f)-\mu(f))$. The exact finite-particle calculation, including the confinement, gives
\begin{equation}\label{iid:eq:fluctuation-duality}
 S_N^T(\phi)=S_N^0(f_0)+M_N(T)-\frac1{2\sqrt N}
 \int_0^T N\mathsf A_1(\nabla f_r,X_N(r),\mu)\,\dd r,
\end{equation}
\begin{equation}\label{iid:eq:martingale}
 M_N(t):=\sqrt{\frac{2}{\beta N}}\sum_i\int_0^t\nabla f_r(x_i(r))\cdot\dd B_i(r),
\end{equation}
where $\mu_N^t:=N^{-1}\sum_{i=1}^N\delta_{x_i(t)}$ and $\mathsf A_1$ is the transport commutator in \eqref{eq:A1}. In particular, $N\mathsf A_1/2$ has the normalization $C_N$ used in \cref{sec:static}. This is the finite-particle adjoint computation of \cite[Proposition~8.2]{DGRpartition}, with the additional one-body drift $-\nabla V\cdot\nabla f$. The computation is valid for the independent initialization used here.

By the exponential commutator estimate \cite[Theorem~2.8]{DGRpartition}, the bounded Hessians of $f_r$, and \cref{iid:lem:entropy},
\begin{equation}\label{iid:eq:remainder}
 \sup_{N,r\le T}\E_{P_{N,r}}[N|\mathsf A_1(\nabla f_r,X_N,\mu)|]<\infty.
\end{equation}
Here the estimate for both signs first bounds the exponential of the absolute value under $\Gamma_N$; the Donsker--Varadhan lemma transfers it to $P_{N,r}$. Consequently,
\begin{equation}\label{iid:eq:remainder-limit}
 \E\left[\left|\frac1{2\sqrt N}\int_0^T
 N\mathsf A_1(\nabla f_r,X_N(r),\mu)\dd r\right|\right]
 \le\frac{C_{\delta,T}}{\sqrt N}\xrightarrow[N\to\infty]{}0.
\end{equation}
To justify \eqref{iid:eq:fluctuation-duality}, first apply It\^o's formula up to exit from collision-free compact sets. The estimate \eqref{iid:eq:remainder} controls the time integral of the commutator in $L^1$. The at-most-linear terms are integrable by \cref{iid:lem:entropy} and the Gaussian exponential moments of $\mu$. Finally, boundedness of $\nabla f_r$ gives $L^2$ convergence of the stopped stochastic integrals. These estimates justify removing the stopping times.

\parheading{The martingale and the initial fluctuation.}
For a bounded test $\zeta$, the product-measure calculation \eqref{ce:eq:linear-product-mgf}--\eqref{ce:eq:linear-product-square-mgf} applies to the centered function $\zeta-\mu(\zeta)$ on $(\R^4,\mu)$. Combined with \cref{iid:lem:entropy} and the Donsker--Varadhan lemma, it yields
\begin{equation}\label{iid:eq:empirical-lln}
 \E_{P_{N,r}}[|\mu_N(\zeta)-\mu(\zeta)|]
 \le C_\delta\operatorname{osc}(\zeta)N^{-1/2},\qquad r\ge0.
\end{equation}
Set $v_T:=2\beta^{-1}\int_0^T\mu(|\nabla f_r|^2)\dd r$. Applying \eqref{iid:eq:empirical-lln} to $|\nabla f_r|^2$ in the predictable quadratic variation of \eqref{iid:eq:martingale} gives
\begin{equation}\label{iid:eq:bracket-limit}
 \E[|\langle M_N\rangle_T-v_T|]
 \le2\beta^{-1}\int_0^T\E[|\mu_N^r(|\nabla f_r|^2)-\mu(|\nabla f_r|^2)|]\dd r
 \le\frac{C_{\delta,T}}{\sqrt N}\xrightarrow[N\to\infty]{}0.
\end{equation}
We also have the deterministic bound $\langle M_N\rangle_T\le2\beta^{-1}T\sup_{0\le r\le T}\norm{\nabla f_r}_\infty^2$. Let $\mathcal F_0:=\sigma(X_N(0))$, and use the filtration generated by the initial configuration and the Brownian motions. For each real $\theta$, the exponential
$t\mapsto\exp(i\theta M_N(t)+\theta^2\langle M_N\rangle_t/2)$ is a bounded-in-modulus true complex martingale on $[0,T]$. Therefore,
\begin{equation}\label{iid:eq:conditional-exponential-martingale}
 \E\left[\left.
 \exp\left(i\theta M_N(T)+\frac{\theta^2}{2}\langle M_N\rangle_T\right)
 \right|\mathcal F_0\right]=1.
\end{equation}
For every bounded $\mathcal F_0$-measurable random variable $Z_N$ with $|Z_N|\le1$, \eqref{iid:eq:conditional-exponential-martingale} and \eqref{iid:eq:bracket-limit} imply
\begin{equation}\label{iid:eq:independence}
 \E[Z_Ne^{i\theta M_N(T)}]
 -e^{-\theta^2v_T/2}\E[Z_N]\longrightarrow0.
\end{equation}
The calculation in \eqref{ce:eq:factorizationestimate} applies with this $\mathcal F_0$-measurable multiplier and bounds the difference by $C\E[|\langle M_N\rangle_T-v_T|]$.

The independent-sample central limit theorem gives variance $\Var_\mu(f_0)$ for the initial statistic. Taking $Z_N=e^{i\theta S_N^0(f_0)}$ in \eqref{iid:eq:independence} factors the limiting characteristic function into the initial Gaussian factor and the Brownian Gaussian factor. Together with \eqref{iid:eq:remainder-limit}, this proves \eqref{iid:eq:clt}--\eqref{iid:eq:covariance}.
\end{proof}


\subsection{Amplification and variance comparison}\label{iid:subsec:amplification}
We now restore the dependence on $\delta$: $L_\delta$ and $\mathcal B_\delta$ are the operators in \eqref{iid:eq:linear-operators} with $\mu=\mu_\delta$ and $\beta=\beta_\delta$, and $U_t^\delta$ is the corresponding evolution. We let $\delta\downarrow0$. The test below has product variance of order one, whereas its image under $\mathcal B_\delta$ has $L^2(\mu_\delta)$ norm of order $\delta^{-1}$. This difference produces the variance comparison at time $T_\delta=\sqrt\delta$.

Choose an even smooth cutoff $\chi$ equal to one on $[-1/4,1/4]$ and supported in $(-1/3,1/3)$. Set
\begin{equation}\label{iid:eq:test}
 \psi(r):=(r-1)\chi(r-1)+(r+1)\chi(r+1),
 \qquad \phi_\delta(r,z):=\delta^{-1}\psi(r).
\end{equation}
The function $\psi$ is odd. In each narrow component, $\phi_\delta$ measures displacement within that component in units of its width; it does not count the distance between the two components. We have
\begin{equation}\label{iid:eq:variance-name}
 \Var_{\mu_\delta}(\phi_\delta)\longrightarrow1\qquad\text{as }\delta\downarrow0.
\end{equation}
The convergence follows from the Gaussian mixture \eqref{iid:eq:mu} and the Gaussian tail bound \cite[Section~5.4]{BGL} in the two neighborhoods where $\psi(r)=r\mp1$. The broad component contributes at most $C\eta_\delta/\delta^2\to0$. Symmetry gives zero mean.

\begin{lemma}[Interaction amplification]\label{iid:lem:interaction-amplification}
There is a constant $c_\s>0$ such that
\begin{equation}\label{iid:eq:interaction-limit}
 \delta^2\norm{\mathcal B_\delta\phi_\delta}_{L^2(\mu_\delta)}^2\longrightarrow c_\s^2,
\end{equation}
where
\begin{equation}\label{iid:eq:c-s}
 c_\s^2:=\int_{\R^3}\left(\int_{\R^3}
 \frac{\gamma_1^{(3)}(y)\,\dd y}{(4+|z-y|^2)^{(\s+2)/2}}\right)^2
 \gamma_1^{(3)}(z)\,\dd z>0.
\end{equation}
Moreover, uniformly in small $\delta$,
\begin{equation}\label{iid:eq:B-bounds}
 \norm{\mathcal B_\delta f}_\infty+\norm{\nabla \mathcal B_\delta f}_\infty
 \le C_\s\norm{\nabla f}_\infty.
\end{equation}
\end{lemma}
\begin{proof}
The bounds in \eqref{iid:eq:B-bounds} follow from \eqref{iid:eq:kernel-row}. Write
\begin{equation}\label{iid:eq:Hdelta}
 H_\delta(x):=\delta \mathcal B_\delta\phi_\delta(x)
 =\int\partial_1\g(x-y)\psi'(y_1)\,\mu_\delta(\dd y).
\end{equation}
These functions and their gradients are uniformly bounded by \eqref{iid:eq:kernel-row}. Define
\begin{equation}\label{iid:eq:Hlimit}
 H_0(r,z):=\frac12\int\bigl(\partial_1\g(r-1,z-y)
 +\partial_1\g(r+1,z-y)\bigr)\gamma_1^{(3)}(y)\,\dd y.
\end{equation}
The mean value theorem and \eqref{iid:eq:transverse-integrability} with $j=2$ control the error from replacing each narrow Gaussian by its center. Since $\psi'=1$ for $|r-\sigma|\le1/4$, $\sigma=\pm1$, the remaining errors are bounded by the Gaussian tails and the weight of the broad component. Thus, for some $c>0$,
\begin{align}\label{iid:eq:sheet-force-convergence}
 \norm{H_\delta-H_0}_\infty
 &\le C_\s\left[\sum_{\sigma=\pm1}\int_\R
 \bigl(|r-\sigma|+\mathbf1_{\{|r-\sigma|\ge1/4\}}\bigr)
 \gamma_\delta^{(1)}(r-\sigma)\dd r+\eta_\delta\right]\notag\\
 &\le C_\s\left(\delta+e^{-c/\delta^2}+\eta_\delta\right)
 \xrightarrow[\delta\downarrow0]{}0.
\end{align}
Also, $\mu_\delta\rightharpoonup\tfrac12(\delta_{-1}+\delta_1)\otimes\gamma_1^{(3)}$. The limit $H_0$ is bounded and continuous. On the two limiting hyperplanes,
\begin{equation}\label{iid:eq:sheet-values}
 H_0(1,z)=-\int(4+|z-y|^2)^{-(\s+2)/2}\gamma_1^{(3)}(y)\,\dd y,
 \qquad H_0(-1,z)=-H_0(1,z).
\end{equation}
Thus,
\begin{equation}\label{iid:eq:interaction-norm-limit}
 \int |H_\delta|^2\dd\mu_\delta
 \longrightarrow\frac12\int_{\R^3}
 \left(|H_0(-1,z)|^2+|H_0(1,z)|^2\right)\gamma_1^{(3)}(z)\dd z
 =c_\s^2>0.
\end{equation}
The uniform convergence in \eqref{iid:eq:sheet-force-convergence} justifies replacing $H_\delta$ by $H_0$ before taking the weak limit of the measures. Since $H_\delta=\delta\mathcal B_\delta\phi_\delta$, this proves \eqref{iid:eq:interaction-limit}.
\end{proof}

\begin{lemma}[Backward evolution]\label{iid:lem:adjoint-estimate}
For $0\le t\le1$, the adjoint evolution $U_t^\delta\phi_\delta$ exists in the regularity class of \cref{iid:prop:clt}. There is a constant independent of small $\delta$ such that
\begin{equation}\label{iid:eq:gradient-evolution}
 \norm{\nabla U_t^\delta f}_\infty\le C\norm{\nabla f}_\infty,
 \qquad
 \norm{\nabla(U_t^\delta\phi_\delta-\phi_\delta)}_\infty
 \le Ct/\delta,
\end{equation}
for the bounded smooth functions used below, and
\begin{equation}\label{iid:eq:adjoint-expansion}
 \norm{U_t^\delta\phi_\delta-\phi_\delta-t\mathcal B_\delta\phi_\delta}_{L^2(\mu_\delta)}
 \le C\left(t\delta^3+\frac{t^2}{\delta}
 +\delta^2t^{3/2}\right).
\end{equation}
\end{lemma}
\begin{proof}
Let $P_t^\delta$ be the Markov semigroup of $L_\delta$. By \eqref{iid:eq:score-hessian}, its drift has Lipschitz constant at most $C\beta_\delta^{-1}\delta^{-4}\le C$. Differentiating the stochastic flow gives
\begin{equation}\label{iid:eq:P-gradient}
 \norm{\nabla P_t^\delta f}_\infty\le e^{Ct}\norm{\nabla f}_\infty.
\end{equation}
The semigroup preserves $\mu_\delta$, is self-adjoint on $L^2(\mu_\delta)$, and has Dirichlet form $\beta_\delta^{-1}\int|\nabla f|^2\dd\mu_\delta$. Its spectral theorem gives
\begin{equation}\label{iid:eq:P-difference}
 \norm{(P_t^\delta-\Id)f}_{L^2(\mu_\delta)}
 \le\sqrt{\beta_\delta^{-1} t}\norm{\nabla f}_{L^2(\mu_\delta)}.
\end{equation}
Indeed, $(1-e^{-a})^2\le a$ for $a\ge0$.

\parheading{Construction and regularity.}
The Volterra equation
\begin{equation}\label{iid:eq:U-mild}
 U_t^\delta f=P_t^\delta f+
 \int_0^tP_{t-r}^\delta \mathcal B_\delta U_r^\delta f\,\dd r
\end{equation}
is solved on the concrete class needed below as follows. For fixed $\delta$, equip $C_b^1(\R^4)$ with
\begin{equation}\label{iid:eq:picard-norm}
 \norm{f}_{C_b^1}:=\norm{f}_\infty+\norm{\nabla f}_\infty.
\end{equation}
Starting from the bounded smooth test $\phi_\delta$ in \eqref{iid:eq:test}, define
\begin{equation}\label{iid:eq:picard-iterates}
 u^{(0)}_t:=P_t^\delta\phi_\delta,
 \qquad
 u^{(n+1)}_t:=P_t^\delta\phi_\delta
 +\int_0^tP_{t-r}^\delta \mathcal B_\delta u^{(n)}_r\,\dd r.
\end{equation}
\eqref{iid:eq:B-bounds} and \eqref{iid:eq:P-gradient} imply, for $0\le t\le\tau$,
\begin{equation}\label{iid:eq:picard-contraction}
 \sup_{0\le t\le\tau}
 \norm{u^{(n+1)}_t-u^{(n)}_t}_{C_b^1}
 \le C\tau
 \sup_{0\le t\le\tau}
 \norm{u^{(n)}_t-u^{(n-1)}_t}_{C_b^1}.
\end{equation}
Choosing $\tau$ so that $C\tau<1$ gives convergence in the bounded-function plus bounded-gradient norm. Repeating the construction from the final value on each interval gives a unique mild solution on $[0,1]$ and the first inequality in \eqref{iid:eq:gradient-evolution}.

The test $\phi_\delta$ is smooth, bounded, and compactly supported in the first coordinate. Stochastic-flow continuity gives $P_t^\delta\phi_\delta\to\phi_\delta$ locally uniformly, together with its first derivative, as $t\downarrow0$. The integral term in \eqref{iid:eq:U-mild} is $O(t)$ in the norm \eqref{iid:eq:picard-norm}; hence the mild solution has the required initial trace on this concrete class.

For fixed $\delta$, $\sup_{x\in\R^4}\int|\nabla^3\g(x-y)|\,\mu_\delta(\dd y)<\infty$, since $\s+3<4$ and $\mu_\delta$ is bounded. Thus, for the bounded smooth functions under consideration,
\begin{equation}\label{iid:eq:fixed-delta-second-derivatives}
 \norm{\nabla^2\mathcal B_\delta f}_\infty\le C_\delta\norm{\nabla f}_\infty,
 \qquad
 \norm{\nabla^2P_t^\delta f}_\infty
 \le C_\delta\left(\norm{\nabla f}_\infty+\norm{\nabla^2f}_\infty\right),
 \quad 0\le t\le1.
\end{equation}
The second estimate follows by differentiating the stochastic flow twice; the first two derivatives of its drift are bounded at fixed $\delta$. Applying these estimates to differences of successive Picard iterates gives
\begin{equation}\label{iid:eq:picard-second-derivatives}
 \sup_{0\le t\le\tau}\norm{\nabla^2(u^{(n+1)}_t-u^{(n)}_t)}_\infty
 \le C_\delta\tau\sup_{0\le t\le\tau}
       \norm{u^{(n)}_t-u^{(n-1)}_t}_{C_b^1}.
\end{equation}
The right-hand side is summable in $n$ by \eqref{iid:eq:picard-contraction}. Hence the second derivatives converge uniformly on each short interval. Iteration gives bounded second derivatives on $[0,1]$, with a bound allowed to depend on $\delta$.

Passing to the limit in \eqref{iid:eq:picard-iterates} and applying local parabolic regularity \cite[Chapter~5]{Lunardi} gives a classical solution on $(0,1]\times\R^4$. The differentiated flow is locally continuous at time zero, and the integral term is $O(t)$ also in the bounded-second-derivative norm at fixed $\delta$. Thus, the required spatial derivatives have their initial trace. The time derivative has at most linear growth by \eqref{iid:eq:linear-operators}, giving the regularity in \cref{iid:prop:clt}.

\parheading{Estimates uniform in $\delta$.}
Since $\psi$ has fixed compact support and bounded derivatives of every order, \eqref{iid:eq:score-first} and \eqref{iid:eq:score-hessian} imply
\begin{equation}\label{iid:eq:Lphi}
 \norm{L_\delta\phi_\delta}_\infty\le C\delta^3,
 \qquad
 \norm{\nabla L_\delta\phi_\delta}_\infty\le C\delta.
\end{equation}
The identity $P_t^\delta\phi_\delta-\phi_\delta
=\int_0^tP_r^\delta L_\delta\phi_\delta\dd r$ follows from It\^o; all terms on the right are bounded. Taking gradients and using \eqref{iid:eq:P-gradient} gives a bound $Ct\delta$. Subtract $\phi_\delta$ in \eqref{iid:eq:U-mild}, apply \eqref{iid:eq:B-bounds}, and use the already proved gradient estimate for $U_r^\delta\phi_\delta$. This proves the second inequality in \eqref{iid:eq:gradient-evolution}.

Finally, subtract $\phi_\delta+t\mathcal B_\delta\phi_\delta$ in \eqref{iid:eq:U-mild} to obtain
\begin{equation}\label{iid:eq:three-errors}
 U_t^\delta\phi_\delta-\phi_\delta-t\mathcal B_\delta\phi_\delta
 =(P_t^\delta-\Id)\phi_\delta
 +\int_0^t(P_{t-r}^\delta\mathcal B_\delta U_r^\delta-\mathcal B_\delta)\phi_\delta\,\dd r.
\end{equation}
Adding and subtracting $P_{t-r}^\delta\mathcal B_\delta\phi_\delta$ inside the integral separates the right-hand side into three terms. Their $L^2(\mu_\delta)$ norms are bounded, respectively, by $Ct\delta^3$, $Ct^2/\delta$, and $C\sqrt{\beta_\delta^{-1}}\,t^{3/2}/\delta$. For the last bound, use \eqref{iid:eq:P-difference} and \eqref{iid:eq:B-bounds}; for the second, use $L^2$ contraction, \eqref{iid:eq:B-bounds}, and \eqref{iid:eq:gradient-evolution}. Since $\sqrt{\beta_\delta^{-1}}=\delta^3$, this is \eqref{iid:eq:adjoint-expansion}.
\end{proof}

\begin{prop}[Variance comparison]\label{iid:prop:variance-crossing}
Let $W_\delta$ be the variance \eqref{iid:eq:covariance} at $T=T_\delta$ for $\phi=\phi_\delta$. Then
\begin{equation}\label{iid:eq:variance-asymptotic}
 \lim_{\delta\downarrow0}\Var_{\mu_\delta}(\phi_\delta)=1,\qquad
 \lim_{\delta\downarrow0}\delta W_\delta=c_\s^2>0.
\end{equation}
In particular, $W_\delta>2\Var_{\mu_\delta}(\phi_\delta)$ for every sufficiently small $\delta$.
\end{prop}
\begin{proof}
Reflection of the first coordinate preserves $\mu_\delta$, $L_\delta$, and $\mathcal B_\delta$, so $U_t^\delta\phi_\delta$ is odd and has zero $\mu_\delta$ mean. At $T_\delta=\sqrt\delta$, the right-hand side of \eqref{iid:eq:adjoint-expansion} is bounded independently of small $\delta$. Since $\sqrt\delta T_\delta\mathcal B_\delta\phi_\delta=H_\delta$,
\begin{equation}\label{iid:eq:leading-amplification}
 \norm{\sqrt\delta\,U_{T_\delta}^\delta\phi_\delta-H_\delta}_{L^2(\mu_\delta)}
 \le\sqrt\delta\,\norm{\phi_\delta}_{L^2(\mu_\delta)}+C\sqrt\delta
 \xrightarrow[\delta\downarrow0]{}0.
\end{equation}
Here, \eqref{iid:eq:variance-name} bounds the norm of $\phi_\delta$, and \cref{iid:lem:interaction-amplification} gives $\norm{H_\delta}_{L^2(\mu_\delta)}\to c_\s$. Consequently, $\delta\Var_{\mu_\delta}(U_{T_\delta}^\delta\phi_\delta)\to c_\s^2$. The noise part of \eqref{iid:eq:covariance} is bounded by
\begin{equation}\label{iid:eq:noise-small}
 2\beta_\delta^{-1}\int_0^{T_\delta}\norm{\nabla U_r^\delta\phi_\delta}_{L^2(\mu_\delta)}^2\dd r
 \le C \beta_\delta^{-1} T_\delta\delta^{-2}=C\delta^{9/2}.
\end{equation}
This proves \eqref{iid:eq:variance-asymptotic}.
\end{proof}


\subsection{Completion and marginal entropy}\label{iid:subsec:completion}
\begin{proof}[Proof of \cref{prop:iid-counterexample}]
The background and entropy assertions follow from \cref{iid:lem:stationary-measure,iid:lem:entropy} and \cref{prop:particle-wellposedness}, and fixed-particle chi-square finiteness follows from \cref{iid:lem:fixed-N-finite}. Fix $\delta$ small enough that $W_\delta>2\Var_{\mu_\delta}(\phi_\delta)$. All subsequent limits are in $N$, with $\s,\delta,\beta_\delta,V_\delta,\mu_\delta,T_\delta$, and the test $\phi_\delta$ fixed.

The product reference satisfies the independent-sample central limit theorem with variance $\Var_{\mu_\delta}(\phi_\delta)$. By \cref{iid:prop:clt,iid:lem:adjoint-estimate}, the particle law satisfies the central limit theorem with variance $W_\delta$. Apply the same statistic $X_N\mapsto N^{-1/2}\sum_i(\phi_\delta(x_i)-\mu_\delta(\phi_\delta))$ to both laws. Data processing and joint weak lower semicontinuity \cite[Theorem 19 and (6)]{vanErvenHarremoes2014} give
\begin{equation}\label{iid:eq:data-processing}
 \liminf_{N\to\infty}\chi^2(P_{N,T_\delta}\mid\mu_\delta^{\otimes N})
 \ge\chi^2\bigl(\mathcal N(0,W_\delta)\mid\mathcal N(0,\Var_{\mu_\delta}(\phi_\delta))\bigr).
\end{equation}
For $W,v>0$, direct substitution of the Gaussian densities gives
\begin{equation}\label{iid:eq:Gaussian-threshold}
 1+\chi^2\bigl(\mathcal N(0,W)\mid\mathcal N(0,v)\bigr)
 =\frac{\sqrt v}{\sqrt{2\pi}\,W}
 \int_\R\exp\left[\left(\frac1{2v}-\frac1W\right)x^2\right]\dd x.
\end{equation}
With $W=W_\delta$ and $v=\Var_{\mu_\delta}(\phi_\delta)$, the coefficient of $x^2$ is positive because $W_\delta>2v$. The integral, and hence the right-hand side of \eqref{iid:eq:data-processing}, is infinite. This proves divergence of the whole sequence in \eqref{iid:eq:chi-main}.
\end{proof}

\begin{remark}[Marginal entropy from independent initial data]\label{iid:rem:local}
The same example is locally chaotic in relative entropy, in the sense that
\begin{equation}\label{iid:eq:local-entropy}
 \sup_{t\ge0}\Hr(P_{N,t}^{(k)}\mid\mu_\delta^{\otimes k})
 \le C_\delta\frac{k}{N},\qquad 1\le k\le N.
\end{equation}
To see this, for fixed $t$, set $H_j:=\Hr(P_{N,t}^{(j)}\mid\mu_\delta^{\otimes j})$ and $H_0:=0$. The chain rule, data processing, and exchangeability give, for $2\le j<N$,
\begin{equation}\label{iid:eq:entropy-increments}
 H_{j+1}-H_j=H_1+\Hr\bigl(P_{N,t}^{(j+1)}\mid P_{N,t}^{(j)}\otimes P_{N,t}^{(1)}\bigr)\ge H_j-H_{j-1}.
\end{equation}
For $j=1$, the same inequality follows from $\Hr(P_{N,t}^{(2)}\mid(P_{N,t}^{(1)})^{\otimes2})\ge0$. Thus, averaging these nondecreasing increments gives
\begin{equation}\label{iid:eq:marginal-entropy-averages}
 \frac{H_k}{k}=\frac1k\sum_{j=1}^k(H_j-H_{j-1})
 \le\frac1N\sum_{j=1}^N(H_j-H_{j-1})=\frac{H_N}{N}.
\end{equation}
Therefore, \eqref{iid:eq:entropy-main} proves \eqref{iid:eq:local-entropy}. In particular, every sequence $k_N=o(N)$ satisfies
\begin{equation}\label{iid:eq:mesoscopic-entropy-limit}
 \sup_{t\ge0}\Hr(P_{N,t}^{(k_N)}\mid\mu_\delta^{\otimes k_N})
 \le C_\delta\frac{k_N}{N}\xrightarrow[N\to\infty]{}0.
\end{equation}
The example rules out a bound $\chi^2(P_{N,t}^{(k)}\mid\mu_\delta^{\otimes k})\le C_\delta k^2/N^2$ holding simultaneously for all $1\le k\le N$, since it fails at $k=N$ and $t=T_\delta$. It does not settle the sharp chi-square or relative-entropy rate $k^2/N^2$, or the total-variation rate $k/N$, for fixed marginals or for sublinear growing marginals.

\end{remark}


\par\endgroup

\providecommand{\bysame}{\leavevmode\hbox to3em{\hrulefill}\thinspace}
\providecommand{\MR}{\relax\ifhmode\unskip\space\fi MR }
% \MRhref is called by the amsart/book/proc definition of \MR.
\providecommand{\MRhref}[2]{%
  \href{http://www.ams.org/mathscinet-getitem?mr=#1}{#2}
}
\providecommand{\href}[2]{#2}
\begin{thebibliography}{CdCRS25}

\bibitem[BGL14]{BGL}
Dominique Bakry, Ivan Gentil, and Michel Ledoux, \emph{Analysis and Geometry of Markov Diffusion Operators}, Grundlehren der mathematischen Wissenschaften, vol.~348, Springer, Cham, 2014, \href{https://doi.org/10.1007/978-3-319-00227-9}{doi:10.1007/978-3-319-00227-9}.


\bibitem[AL16]{AngiuliLorenzi}
Luciana Angiuli and Luca Lorenzi, \emph{On the estimates of the derivatives of solutions to nonautonomous Kolmogorov equations and their consequences}, Riv. Mat. Univ. Parma (N.S.) \textbf{7} (2016), no.~2, 421--471, \href{https://arxiv.org/abs/1702.02428}{arXiv:1702.02428}.


\bibitem[BJS25]{BJS}
Didier Bresch, Pierre-Emmanuel Jabin, and Juan Soler, \emph{A new approach to
  the mean-field limit of {V}lasov--{F}okker--{P}lanck equations}, Anal. PDE
  \textbf{18} (2025), 1037--1064,
  \href{https://doi.org/10.2140/apde.2025.18.1037}{doi:10.2140/apde.2025.18.1037}.


\bibitem[BJW19]{BJW}
Didier Bresch, Pierre-Emmanuel Jabin, and Zhenfu Wang, \emph{Modulated free
  energy and mean field limit}, 2019,
  \href{https://arxiv.org/abs/1912.03152}{arXiv:1912.03152}.


\bibitem[CL95]{CarlenLoss}
Eric~A. Carlen and Michael Loss, \emph{Optimal smoothing and decay estimates
  for viscously damped conservation laws, with applications to the 2-{D}
  {N}avier--{S}tokes equation}, Duke Math. J. \textbf{81} (1995), no.~1,
  135--157,
  \href{https://doi.org/10.1215/S0012-7094-95-08110-1}{doi:10.1215/S0012-7094-95-08110-1}.


\bibitem[DGR26a]{DGRpartition}
Matias~G. Delgadino, Rishabh Gvalani, and Matthew Rosenzweig, \emph{Sharp
  mean-field estimates for diffusive log/{R}iesz gases in the
  {H}ilbert--{S}chmidt regime}, 2026,
  \href{https://arxiv.org/abs/2609.02743v1}{arXiv:2609.02743v1}.


\bibitem[DGR26b]{DGRchaos}
Matias~G. Delgadino, Rishabh Gvalani, and Matthew Rosenzweig, \emph{Sharp two-sided size of chaos estimates for diffusive gases in the
  {H}ilbert--{S}chmidt regime}, Manuscript, 2026.


\bibitem[DJ25]{DJBrownian}
Mitia Duerinckx and Pierre-Emmanuel Jabin, \emph{Correlation estimates for
  {B}rownian particles with singular interactions}, 2025,
  \href{https://arxiv.org/abs/2510.01507}{arXiv:2510.01507}.


\bibitem[DJ26]{DJFluctuations}
Mitia Duerinckx and Pierre-Emmanuel Jabin, \emph{Singular mean-field limits for fluctuations around equilibrium},
  2026, \href{https://arxiv.org/abs/2605.28979}{arXiv:2605.28979}.


\bibitem[Lun95]{Lunardi}
Alessandra Lunardi, \emph{Analytic Semigroups and Optimal Regularity in Parabolic Problems}, Progress in Nonlinear Differential Equations and their Applications, vol.~16, Birkh\"auser, Basel, 1995, \href{https://doi.org/10.1007/978-3-0348-9234-6}{doi:10.1007/978-3-0348-9234-6}.


\bibitem[EN00]{EngelNagel}
Klaus-Jochen Engel and Rainer Nagel, \emph{One-Parameter Semigroups for Linear Evolution Equations}, Graduate Texts in Mathematics, vol.~194, Springer, New York, 2000, \href{https://doi.org/10.1007/b97696}{doi:10.1007/b97696}.


\bibitem[Gro75]{Gross}
Leonard Gross, \emph{Logarithmic {S}obolev inequalities}, Amer. J. Math.
  \textbf{97} (1975), no.~4, 1061--1083,
  \href{https://doi.org/10.2307/2373688}{doi:10.2307/2373688}.


\bibitem[HCR25]{HCR}
Elias Hess-Childs and Keefer Rowan, \emph{Higher-order propagation of chaos in
  {$L^2$} for interacting diffusions}, Probab. Math. Phys. \textbf{6} (2025),
  no.~2, 581--646,
  \href{https://doi.org/10.2140/pmp.2025.6.581}{doi:10.2140/pmp.2025.6.581}.


\bibitem[KLL10]{KLL}
Markus Kunze, Luca Lorenzi, and Alessandra Lunardi, \emph{Nonautonomous Kolmogorov parabolic equations with unbounded coefficients}, Trans. Amer. Math. Soc. \textbf{362} (2010), no.~1, 169--198, \href{https://doi.org/10.1090/S0002-9947-09-04738-2}{doi:10.1090/S0002-9947-09-04738-2}.


\bibitem[Lac22]{LackerGibbs}
Daniel Lacker, \emph{Quantitative approximate independence for continuous mean field Gibbs measures}, Electron. J. Probab. \textbf{27} (2022), Paper No.~15, 1--21,
\href{https://doi.org/10.1214/22-EJP743}{doi:10.1214/22-EJP743}.

\bibitem[Lac23]{Lacker}
Daniel Lacker, \emph{Hierarchies, entropy, and quantitative propagation of
  chaos for mean field diffusions}, Probab. Math. Phys. \textbf{4} (2023),
  377--432,
  \href{https://doi.org/10.2140/pmp.2023.4.377}{doi:10.2140/pmp.2023.4.377}.


\bibitem[LLF23]{LackerLeFlem}
Daniel Lacker and Luc Le~Flem, \emph{Sharp uniform-in-time propagation of
  chaos}, Probab. Theory Related Fields \textbf{187} (2023), 443--480,
  \href{https://doi.org/10.1007/s00440-023-01192-x}{doi:10.1007/s00440-023-01192-x}.



\bibitem[NRS22]{NguyenRosenzweigSerfaty}
Quoc Hung Nguyen, Matthew Rosenzweig, and Sylvia Serfaty, \emph{Mean-field limits of Riesz-type singular flows}, Ars Inven. Anal. (2022), Paper No.~4, 45~pp.,
\href{https://doi.org/10.15781/nvv7-jy87}{doi:10.15781/nvv7-jy87}.

\bibitem[Ros25]{RosenzweigJEDP}
Matthew Rosenzweig, \emph{Commutators, mean-field, and supercritical mean-field limits for {Coulomb/Riesz} gases}, Journ\'ees \'equations aux d\'eriv\'ees partielles (2025), Talk No.~7, 32~pp.,
\href{https://doi.org/10.5802/jedp.698}{doi:10.5802/jedp.698}.

\bibitem[RS25]{RSlsi}
Matthew Rosenzweig and Sylvia Serfaty, \emph{Modulated logarithmic {S}obolev inequalities and generation of
  chaos}, Ann. Fac. Sci. Toulouse Math. (6) \textbf{34} (2025), no.~1,
  107--134, \href{https://doi.org/10.5802/afst.1807}{doi:10.5802/afst.1807}.



\bibitem[RS24]{RSselfsimilar}
Matthew Rosenzweig and Sylvia Serfaty, \emph{Relative entropy and modulated free energy without confinement via self-similar transformation}, arXiv:2402.13977v1 (2024), \url{https://arxiv.org/abs/2402.13977v1}.

\bibitem[Wan26]{WangSharpLocal}
Songbo Wang, \emph{Sharp local propagation of chaos for mean field particles with {$W^{-1,\infty}$} kernels}, J. Funct. Anal. \textbf{290} (2026), no.~3, 111240,
  \href{https://doi.org/10.1016/j.jfa.2025.111240}{doi:10.1016/j.jfa.2025.111240}.


\begingroup
\bibitem[Bre11]{Brezis}
Haim Brezis, \emph{Functional Analysis, Sobolev Spaces and Partial Differential Equations}, Universitext, Springer, New York, 2011, \href{https://doi.org/10.1007/978-0-387-70914-7}{doi:10.1007/978-0-387-70914-7}.

\par\endgroup

\begingroup
\bibitem[\O ks03]{Oksendal}
Bernt \O ksendal, \emph{Stochastic Differential Equations: An Introduction with Applications}, sixth ed., Universitext, Springer, Berlin, 2003, \href{https://doi.org/10.1007/978-3-642-14394-6}{doi:10.1007/978-3-642-14394-6}.
\par\endgroup


\begingroup
\bibitem[Lee90]{Lee1990}
A.~J. Lee, \emph{U-Statistics: Theory and Practice}, Statistics: Textbooks and Monographs, vol.~110, Marcel Dekker, New York, 1990.

\bibitem[NS20]{NishiyamaSason2020}
Tomohiro Nishiyama and Igal Sason, \emph{On relations between the relative entropy and $\chi^2$-divergence, generalizations and applications}, Entropy \textbf{22} (2020), no.~5, article~563, \href{https://doi.org/10.3390/e22050563}{doi:10.3390/e22050563}; equation numbers refer to \href{https://arxiv.org/abs/2004.11197v3}{arXiv:2004.11197v3}.
\par\endgroup


\begingroup
\bibitem[AS20]{ArmstrongSerfatyThermal}
Scott Armstrong and Sylvia Serfaty, \emph{Thermal approximation of the equilibrium measure and obstacle problem}, arXiv:1912.13018v5 (2020), \url{https://arxiv.org/abs/1912.13018v5}.

\bibitem[Kat95]{Kato}
Tosio Kato, \emph{Perturbation Theory for Linear Operators}, second ed., Classics in Mathematics, Springer, Berlin, 1995, \href{https://doi.org/10.1007/978-3-642-66282-9}{doi:10.1007/978-3-642-66282-9}.
\par\endgroup

\bibitem[JW18]{JabinWang2018}
P.-E. Jabin and Z. Wang,
\emph{Quantitative estimates of propagation of chaos for stochastic systems with $W^{-1,\infty}$ kernels},
Invent. Math. \textbf{214} (2018), no.~1, 523--591.

\bibitem[RY99]{RevuzYor}
D. Revuz and M. Yor,
\emph{Continuous Martingales and Brownian Motion}, third ed.,
Grundlehren der mathematischen Wissenschaften, vol.~293, Springer, Berlin, 1999.

\bibitem[vEH14]{vanErvenHarremoes2014}
T. van Erven and P. Harremo\"es,
\emph{R\'enyi divergence and Kullback--Leibler divergence},
IEEE Trans. Inform. Theory \textbf{60} (2014), no.~7, 3797--3820.


\bibitem[FJ15]{FournierJourdain2015}
Nicolas Fournier and Benjamin Jourdain, \emph{Stochastic particle approximation
of the {K}eller--{S}egel equation and two-dimensional generalization of
{B}essel processes}, 2015,
\href{https://arxiv.org/abs/1507.01087v1}{arXiv:1507.01087v1}.

\bibitem[RW25]{RenWangGibbs}
Zhenjie Ren and Songbo Wang, \emph{Size of chaos for Gibbs measures of mean field interacting diffusions}, Probab. Theory Related Fields (2025),
\href{https://doi.org/10.1007/s00440-025-01435-z}{doi:10.1007/s00440-025-01435-z}.

\bibitem[RS23]{RSdiffusive}
Matthew Rosenzweig and Sylvia Serfaty, \emph{Global-in-time mean-field convergence for singular {R}iesz-type diffusive flows}, Ann. Appl. Probab. \textbf{33} (2023), no.~2, 754--798, \href{https://doi.org/10.1214/22-AAP1833}{doi:10.1214/22-AAP1833}.

\bibitem[OT20]{OsadaTanemura2020}
Hirofumi Osada and Hideki Tanemura, \emph{Infinite-dimensional stochastic differential equations and tail $\sigma$-fields}, Probab. Theory Related Fields \textbf{177} (2020), 1137--1242, \href{https://doi.org/10.1007/s00440-020-00981-y}{doi:10.1007/s00440-020-00981-y}.


\begingroup
\bibitem[Ser26]{SerfatyBook}
Sylvia Serfaty, \emph{Lectures on Coulomb and Riesz Gases}, Colloquium Publications, vol.~70, American Mathematical Society, 2026. Preliminary version: \href{https://arxiv.org/abs/2407.21194}{arXiv:2407.21194}.

\bibitem[CdCRS25a]{CdCRSRepulsive}
Antonin Chodron de Courcel, Matthew Rosenzweig, and Sylvia Serfaty,
\emph{Sharp uniform-in-time mean-field convergence for singular periodic Riesz flows},
Ann. Inst. H. Poincar\'e Anal. Non Lin\'eaire \textbf{42} (2025), no.~2, 391--472,
\href{https://doi.org/10.4171/AIHPC/105}{doi:10.4171/AIHPC/105}.

\bibitem[CdCRS25b]{CdCRSAttractive}
Antonin Chodron de Courcel, Matthew Rosenzweig, and Sylvia Serfaty,
\emph{The attractive log gas: stability, uniqueness, and propagation of chaos},
Commun. Amer. Math. Soc. \textbf{5} (2025), 695--773,
\href{https://doi.org/10.1090/cams/55}{doi:10.1090/cams/55}.

\bibitem[BJW23]{BJWPKS}
Didier Bresch, Pierre-Emmanuel Jabin, and Zhenfu Wang,
\emph{Mean field limit and quantitative estimates with singular attractive kernels},
Duke Math. J. \textbf{172} (2023), no.~13, 2591--2641,
\href{https://doi.org/10.1215/00127094-2022-0088}{doi:10.1215/00127094-2022-0088}.

\bibitem[WZ26]{WZ26}
Zhenfu Wang and Xianliang Zhao, \emph{Uniform Partition-Function Estimates for Coulomb Modulated Energy at All Positive Temperatures}, 2026,
\href{https://arxiv.org/abs/2609.00867v1}{arXiv:2609.00867v1}.

\bibitem[FW26]{FW26}
Xuanrui Feng and Zhenfu Wang, \emph{Modulated Gibbs Measure and Mean-Field Limit for 3D Vlasov--Poisson--Fokker--Planck Equations}, 2026,
\href{https://arxiv.org/abs/2609.08640v1}{arXiv:2609.08640v1}.

\bibitem[DGR26c]{DGRAttractiveII}
Matias~G. Delgadino, Rishabh Gvalani, and Matthew Rosenzweig,
\emph{The attractive logarithmic gas II: mean-field dynamics}, In preparation, 2026.
\par\endgroup

\end{thebibliography}

\end{document}
